\documentclass[11pt, a4paper, BCOR=10mm, DIV=11]{scrbook}
\linespread{1.25}
\usepackage[utf8]{inputenc}
\usepackage{graphicx}
%\usepackage[a4paper, margin=2.5cm]{geometry}
\usepackage{hyperref}
\usepackage{amsmath}
\usepackage{enumitem}
%\usepackage{mdframed}
\usepackage{amsthm}
\usepackage{amssymb}
\usepackage{mathtools}
%\usepackage{fdsymbol}
\usepackage{cite}
\graphicspath{ {images/} }
\usepackage{xcolor}         % Extended colors
\usepackage{color}         % Color extended names
\usepackage{nomencl}
\usepackage{lipsum}
\usepackage[tickmarkheight=0.1cm, colorinlistoftodos]{todonotes}
\makenomenclature
\renewcommand{\nomname}{Notation Index}
\mathsurround=2pt
\usepackage[framemethod=TikZ]{mdframed}
\mdfsetup{%
   middlelinecolor=black,
   middlelinewidth=.5pt,
   backgroundcolor=gray!2,
   roundcorner=5pt}
\newsavebox{\selvestebox}
\newenvironment{takeaways}
  {
   \begin{lrbox}{\selvestebox}%
   \begin{minipage}{12.4cm}
     \textbf{Takeaways}
   }
  {\end{minipage}\end{lrbox}%
   \begin{center}
\setlength\fboxsep{.5cm}
   \colorbox[HTML]{F8E0E0}{\usebox{\selvestebox}}
   \end{center}}
% Theorem handle


\newtheorem{gtheorem}{Theorem}[chapter]
\newenvironment{theorem}
  {\begin{mdframed}\begin{gtheorem}}
  {\end{gtheorem}\end{mdframed}}
\newtheorem*{theorem*}{Theorem}
\newtheorem*{definition*}{Definition}
\newtheorem*{lemma*}{Lemma}
\newenvironment{ftheorem}
  {\begin{mdframed}\begin{gtheorem}}
  {\end{gtheorem}\end{mdframed}}
\newenvironment{ftheorem*}
  {\begin{mdframed}\begin{theorem*}}
  {\end{theorem*}\end{mdframed}}
\newtheorem{fassumption}{Assumption}
\newenvironment{assumption}
{\begin{mdframed}[
   backgroundcolor=red!2,
    ]
    \begin{fassumption}}
  {\end{fassumption}\end{mdframed}}

\newtheorem*{assumptions*}{Assumptions}
\newtheorem{corollary}{Corollary}[gtheorem]
\newtheorem{definition}{Definition}[chapter]

\theoremstyle{definition}

\newtheorem*{remark/}{Remark}
\newenvironment{remark}
  {\renewcommand{\qedsymbol}{$\diamondsuit$}%
   \pushQED{\qed}\begin{remark/}}
  {\popQED\end{remark/}}
\newtheorem{example/}{Example}[chapter]
\newenvironment{example}
  {\renewcommand{\qedsymbol}{$\diamondsuit$}%
   \pushQED{\qed}\begin{example/}}
  {\popQED\end{example/}}
\newtheorem*{example*/}{Example}
\newenvironment{example*}
  {\renewcommand{\qedsymbol}{$\diamondsuit$}%
   \pushQED{\qed}\begin{example*/}}
  {\popQED\end{example*/}}
\newtheorem*{reflection*/}{Reflection}
\newenvironment{reflection*}
  {\renewcommand{\qedsymbol}{$\spadesuit$}%
   \pushQED{\qed}\begin{reflection*/}}
  {\popQED\end{reflection*/}}
\newtheorem{problem}{Problem}
\newenvironment{fproblem}
  {\begin{mdframed}\begin{problem}}
  {\end{problem}\end{mdframed}}
  %%%% plain %%%%
\theoremstyle{plain}
\newtheorem{flemma}{Lemma}[chapter]
\newenvironment{lemma}
{\begin{mdframed}\begin{flemma}}
  {\end{flemma}\end{mdframed}}
\newtheorem{proposition}{Proposition}[chapter]
\newtheorem*{proposition*}{Proposition}
%\newtheorem{subassumption}{}[assumption]

\renewcommand{\proofname}{\textbf{Proof}}
% definition of constants

\newcommand{\CP}{C_\P}
\newcommand{\Ctau}{C_\tau}
\newcommand{\LearnRate}{\varepsilon_n}

%
\DeclarePairedDelimiterX{\inner}[2]{\langle}{\rangle}{#1, #2}
\newcommand{\C}{\mathbb{C}}
\newcommand{\G}{\mathbb{G}}
\newcommand{\E}{\mathbf{E}}
\renewcommand{\P}{\mathbf{P}}
\newcommand{\R}{\mathbb{R}}
\newcommand\norm[1]{\left\lVert#1\right\rVert}

\title{
  {
    Causal Inference in Observational Studies
  }
  \\
  {\large Universität Stuttgart}
  \\
  {\includegraphics{unistuttgart_logo_deutsch.jpg}}
}

\author{Ioan Scheffel}
\date{\today}
\setlength {\marginparwidth }{2cm}

\begin{document}

%\listoftodos

\maketitle

\tableofcontents 

%\chapter{Introduction}
%% SAY THAT RESEARCHERS ARE OFTEN LEFT WITH OBSERVATIONAL STUDYS TO ANSWER THEIR QUESTIONS

% SAY SOMETHING ABOUT PS METHODS AS TO DO WITH CAUSAL INFERENCE IN OS

% SAY THAT MODEL DEPENDENCY IS OFTEN A PROBLEM WITH PS

% MENTION DIFFERENT APPROACHES TO ALLEVIATE THIS PROBLEM

% HIGHLIGHT THE METHODS BEING DISCUSSED IN THE THESIS

% OUTLINE STRUCTURE 
%% - WEIGHTING APPROACH BINARY
%%  - MODIFICATIONS IN BINARY SETTING TO ESTIMATE ATE
%%  - EXTENSION TO CONTINUOUS TREATMENT
%% - MATCHING APPROACH BINARY 
%%  - POSSIBLE EXTENSION
%% - SIMULATION STUDY
%% - REAL DATA USECASES


Researchers are often left with observational studies to answer questions about causality. When confounders are present the task of infering causality can become arbitrarily complex. Propensity score methods \cite{Rosenbaum1983}, e.g. inverse probability weighting or matching, are popular methods to adjust for confounders. Usually these methods rely heavily on estimates of the true propensity score, which are known to suffer from model dependencies and misspecification\cite{Kang2007}. This issue becomes more pressing when moving from binary to continuous treatment\cite{Hirano2005}. Therefore methods have been developed to directly target imbalances in the data\cite{Fong2018}\cite{Hainmueller2012}\cite{Zubizarreta2015}.
We take a closer look at \cite{Wang2019} and extend the analysis to settings with continuous treatment 
\cite{Vegetabile2020}\cite{Tubbicke2020}.

%\chapter{Causal Inference}
%In this chapter we want to give a introduction to causal inference.
We particularly highlight the role of propensity score analysis and explain its importance in observational studies.


\section{The Rubin Causal Model}
from wiki:
The Rubin causal model (RCM), also known as the Neyman–Rubin causal model,[1] is an approach to the statistical analysis of cause and effect based on the framework of potential outcomes, named after Donald Rubin. The name "Rubin causal model" was first coined by Paul W. Holland.[2] The potential outcomes framework was first proposed by Jerzy Neyman in his 1923 Master's thesis,[3] though he discussed it only in the context of completely randomized experiments.[4] Rubin extended it into a general framework for thinking about causation in both observational and experimental studies.[1]
\section{Propensity Score Analysis}
from wiki:
In the statistical analysis of observational data, propensity score matching (PSM) is a statistical matching technique that attempts to estimate the effect of a treatment, policy, or other intervention by accounting for the covariates that predict receiving the treatment. PSM attempts to reduce the bias due to confounding variables that could be found in an estimate of the treatment effect obtained from simply comparing outcomes among units that received the treatment versus those that did not. Paul R. Rosenbaum and Donald Rubin introduced the technique in 1983.[1]

The possibility of bias arises because a difference in the treatment outcome (such as the average treatment effect) between treated and untreated groups may be caused by a factor that predicts treatment rather than the treatment itself. In randomized experiments, the randomization enables unbiased estimation of treatment effects; for each covariate, randomization implies that treatment-groups will be balanced on average, by the law of large numbers. Unfortunately, for observational studies, the assignment of treatments to research subjects is typically not random. Matching attempts to reduce the treatment assignment bias, and mimic randomization, by creating a sample of units that received the treatment that is comparable on all observed covariates to a sample of units that did not receive the treatment.

For example, one may be interested to know the consequences of smoking. An observational study is required since it is unethical to randomly assign people to the treatment 'smoking.' The treatment effect estimated by simply comparing those who smoked to those who did not smoke would be biased by any factors that predict smoking (e.g.: gender and age). PSM attempts to control for these biases by making the groups receiving treatment and not-treatment comparable with respect to the control variables. 

from a paper:
Propensity score weighting is one of the techniques used in controlling for selection biases in non-
experimental studies. Propensity scores can be used as weights to account for selection assignment
differences between treatment and comparison groups. One of the advantages of this approach is
that all the individuals in the study can be used for the outcomes evaluation


\section{Weighting beyond the PS}
from \cite{Wang2019}:
Conventionally, the weights are estimated by modeling the propensities of receiving treatment or exhibiting missingness and then inverting the predicted propensities. However, with this approach it can be difficult to properly adjust for or balance the observed covariates. The reason is that this approach only balances covariates in expectation, by the law of large numbers, but in any particular data set it can be difficult to balance covariates, especially if the data set is small or if the covariates are sparse (Zubizarreta et al., 2011). In addition, this approach can result in very unstable estimates when a few observations have very large weights (e.g., Kang and Schafer 2007). To address these problems, a number of methods have been proposed recently. Instead of explicitly modeling the propensities of treatment or missingness, these methods directly balance the covariates. Some of these methods also minimize a measure of dispersion of the weights.

Most of these weighting methods balance covariates exactly rather than approximately. This is a subtle but important difference because approximate balance can trade bias for variance whereas exact balance cannot. Also, exact balance may not admit a solution whereas approximate balance may do so. For a fixed sample size, approximate balance may balance more functions of the covariates than exact balance.

\chapter{Introduction}
How does action change an outcome?
How should I guide my actions towards a better outcome?
The first question is about causality, the second about ethics.

How do causality and ethics reflect on statistics?
If you have not spent much time thinking about study design, this is a good way to start: 
As an analyst, ask yourself “Who acted? Who assigned treatments?"
As researcher -- plan your study accurately. You can ask yourself “How do we act? How do we assign treatment? Can we act?"

Let's say, you gather a sample from a study population, assign treatment (but forget how you did it). Some units get the drug, others don't. Then the statistical analysis shows a strong correlation of treatment and outcome. You hurry to your supervisor. “How was treatment assigned", asks she. “I forgot", says you.
“How do you know your analysis is correct then?"
You show her the data and together find out, that all units that received treatment were significantly taller than the rest of the sample.
After all, is the drug or the height responsible for the change  in outcome?
You realise, that the data is worthless for answering this question.
But you are lucky: It is just grass and fertiliser you were studying.

You get a second chance. A new medication needs testing before it enters the market. 
A company shall recruit participants, but the board requires you to write an outline for the study.
You carefully  explain steps to minimize risks for participants. You include plans to meet other requirements of human research.
Then you have to decide how to assign treatment.
No hand waving this time. You talk to your supervisor.
“Last time, too many tall blades received fertiliser. The distribution of treatment was not really random..."
You decide to determine treatment status by the flip a fair coin.
You call the procedure 'randomization'.

Randomization does not work in important situations. Would you smoke if a coin tells you to? If you say yes - you are likely to be a smoker. The point is that forcing someone to smoke is unethical. But so is not investigating the risks of smoking.

Say, a professor is curious if the smoking habits of his students affect their grades. 
At a beginning break he observes the backyard through his field glasses.
His assistant gets to know his plans. He warns him. “Many students attend parties the night before exams. Maybe they are also more likely to smoke." “I shall see this for myself..." says the professor. He puts away the glasses. After a while, he visits the local club.
He talks to a few of his students. Some smoke, some don't. The chats are enjoyable. He thinks: “Some of best students celebrate \textbf{before} the exam."

I hope, by now it's clear that sometimes it's all about how treatment was assigned.
%You sit in a plain room. There are two chairs. You sit in one of them.
%A person enters the room. “Subject 3225057? The coin decided. You have to smoke!"
%On your way out, anguished cries hit you from another room.
%“Please! Don't make me smoke!"
%You feel mild anxiety and wake up. It was just a dream. But it recurs. You talk to a friend. “I dream about a randomized study about the effects of smoking. A cold, indifferent coin decides who smokes. The non-smokers suffer a lot." When you realise, that this is all non-sense, you feel silly. But your friend
%But you happen to be a professor at the university.
%One day, in your office, you lean out of the window.“Give me the field glasses please. I want to observe my students smoking in the backyard."“What for?" asks your assistant. “Let's see, if this affects their grades" you murmur and look through the eyepieces. The issue comes up again at lunch. “Did you consider that some of the students that smoke also go to a party the night before an exam."“Well, maybe I should. How about I see for myself?" The night before the next exam, you openly enter the place of the party. You find some of the students you deem smartest at the bar. You're having a conversation. But it's to loud in there and in the morning you have an headache.
%
%
%%randomized trials versus observational studies
%
%Is study design more important then statistical analysis?
%
%I think, they are at least equal. 
%
%But a bad analysis can be undone,
%whereas a bad design can not.
%
%You have to stick with the data.
%
%If you are not familiar with study design the distinction between randomized and observational study is helpful.
%
%If you read the literature and are unsure about the design of a study, ask for this terms.
%
%You are likely to find an answer.
%
%
%It is all about how we collect the data.
%
%Say, we want to test the effect of a drug in a study population.
%
%There usually are differences among the units of the study population.
%
%Some are more healthy than others.
%
%We form a treatment and control group, that is, one group takes the drug and the other doesn't.
%
%Then we compare the groups by their health. 
%
%Then a critcal review comes in. What do you mean by healthy.
%
%We mean this and that.
%
%It seems you did not consider this factor.
%
%Maybe the drug is not effective, but the effect we see in your analysis comes from something else.
%
%What do we answer to this?
%
%
%A good method to avoid this awkward situation is to randomize.
%
%For every unit of the population we toss a fair coin that decides if they get the drug.
%
%Now comes the critic.
%
%From the tables it seems there is an effect. But what about unknown influence?
%
%We answer: Does the coin now of them?
%
%It is not ideal, but this way you can prevent systematic damage to your analysis.
%
%
%What if we can't decide who gets treatment?
%
%Don't think treatment has to be something good, it should not carry any label of good or bad.
%
%But what about smoking?
%
%Would you smoke if a coin tells you to?
%
%So this is unethical.
%
%But it is also unethical not to investigate the effects of smoking on the health.
%
%Let's accept, that we sometimes (often?) can not control who gets treatment.
%
%Some smoke, some don't, and we mearly observe.
%
%This is typical example for an observational study.
%
%Honestly, this is an oversimplification, but I hope you get the point.
%
%Who still is insulted by the tone will maybe like\cite{Rubin2007}.
%
%%propensity score
%
%In \cite{Rubin2007} you will find the propensity score.
%
%The propensity score is the individual probability to receive treatment, that is,
%\begin{gather}
%  \P
%  [T=1|X]
%\end{gather}
%if $T$ is the random variable that decides abuot treatment and $X$ is the vector that carries your individual information.
%
%This concept goes back to \cite{Rosenbaum1983}.
%
%It is maybe worth to stop here and think about this definition and its connection to the two study designs.
%
%Discover it for yourself.
%
%\begin{reflection*}
%What is the propensity score in the above example.
%How does the propensity score behave in rs and os?
%\end{reflection*}
%
%
%
%
%
%

%\chapter{True Introduction}
%Starting point:
Propensitiy score analysis~\cite{Rosenbaum1983}.
Two major branches: PS-weighting and matching.
We study a weighting method.
Procedure: Estimate PS, invert, weight.
Problem: Extreme weights when PS close to 0.
Bias when estimation model is misspecified.
Solution: Balance some measure of dependence simultaneously, e.g.
Covariate balancing PS~\cite{Imai2014a}.
Other solution: Doubly robust estimators\cite{H2005}.
They incorporate treatment and outcome model.
Problem: bad results if either are (slightly) misspecified\cite{Kang2007}.
A third option is obtaining weights (seamingly) unrelated to PS.
Entropy Balancing \cite{Hainmueller2012}, ?balancing\cite{Zubizarreta2015}.
Problem: contstraint $\delta=0$ to strict. Bad convergence.
Solution: relax to $\delta\to 0$ for $N\to\infty$.
Paper with mathematical analysis\cite{Wang2019}. Surprising connection
to PS. Also doubly robust\cite{Zhao2017a}.
This attracted my attention. 

I choose different basis as in \cite{Gyorfi2002}. Does analysis work?
Consistency? Asymptotic Normality? Beyond that?

\cite{Wang2019}: Proofs are substandard. Many mistakes. Missing assumptions. Theorems have to be differently.

This thesis is no erratum of\cite{Wang2019}, but can be consulted for writing one.

We thank Wang and Zubi for discussions.

%\chapter{Introduction}
 % We consider a study population in which we want to test the effect of a treatment.
We introduce the \textbf{indicator of treatment} $T\in \left\{ 0,1 \right\}$.
For each treatment level there exist the \textbf{marginal potential outcomes}
$(Y(0),Y(1))$. We would like to estimate $\E[Y(1)]$. If we succeed the same technique
shall yield an estimate of $\E[Y(0)]$. We shall compare $\E[Y(1)]$ and $\E[Y(0)]$ and 
find out something about the effect of the treatment in the population.

The data we acquire is independent and identically distributed. But usually
\begin{gather}
  Y(1)|T=1 \nsim Y(1) 
  \,,
\end{gather}
that is, $T=1$ carries more information than observing the outcome under treatment.
We say that $Y(1)|T=1$ is \textbf{confounded}. To extract that plus of information from $T=1$ and put it where it belongs by collecting more data.
We gather it in $X\in \R^d$ and assume
\begin{gather}
  (Y(0),Y(1))
  \perp
  T
  \ 
  |
  \ 
  X
  \,,
\end{gather}
that is, \textbf{conditional unconfoundedness}.
Thus, we end up collecting $N\in \mathbb{N}$ independent and identically distributed copies of 
$(T,X,Y(T))$. For convenience, we assume that the first $n\in \mathbb{N}$ copies have $T=1$.

A natural estimator for $\E[Y(1)]$ is the weighted mean
\begin{gather}
  \frac{1}{n}
  \sum_{i=1}^{n} 
  w_i Y_i
  \,.
\end{gather}
The weights should satisfy (in a broader sense)
\begin{gather}
  w_i\cdot Y_i \to Y(1)
  \qquad 
  \text{for}\ 
  N
  \ 
  \to
  \ 
  \infty
  \,.
\end{gather}
One class of such weights has been recently analyzed in \cite{Wang2019}.
We take ideas and extend.

\subsection*{The algorithm}
\begin{fproblem}
  \label{bw:1:primal}
\begin{align*}
  %%%% objective %%%%
    &\underset{w_1, \ldots, w_n \in \R}
    {\text{minimize}}
    &&\qquad\qquad
    \sum_{i = 1}^{n} 
    f(w_i)
    &&&
    \\
    %%%% w_i T_i >= 0 %%%%
    &\text{subject to}
    &&\qquad\qquad
    w_i 
    \ge
    0
    &&&
    \qquad
    \text{for all}\ 
    i \in \left\{ 1, \ldots, n \right\}
    \,,
    \\
    %%%% 1/n sum w = 1 %%%%
    & 
    &&\qquad\qquad
    \frac{1}{N}
    \sum_{i=1}^{n} 
    w_i
    =1
    \\
    %%%% box constraints %%%%
    & 
    &&\qquad
    \left| 
      \frac{1}{N} 
      \left( 
      \sum_{i = 1}^{n} 
      w_i
      B_k(X_i)
      -
      \sum_{i=1}^{N} 
      B_k(X_i)
      \right)
    \right|
    \ 
    \le 
    \ 
    \delta_k
    &&&
    \qquad
    \text{for all}\ 
    k \in \left\{ 1, \ldots, N \right\}
    \,.
\end{align*}
\end{fproblem}
This is a (convex) optimization problem. We will talk about the \textbf{objective function} $f$ and
the \textbf{equality} and \textbf{inequality constraints}, especially about the 
\textbf{regression basis} $B$.

\section*{Objective Function}
Strictly speaking, we consider the sum
\begin{gather}
  [w_1,\ldots,w_n]^\top
  \ 
  \mapsto
  \ 
  \sum_{i=1}^{n} 
  f(w_i)
\end{gather}
as the objective function. It is natural to consider the dual formulation of the optimization problem. This involves the \textbf{convex conjugate}(cf.Definition~?) of the original objective function. We show in Example that for the sum this is
\begin{gather}
  [\lambda_1,\ldots,\lambda_n]^\top
  \ 
  \mapsto
  \ 
  \sum_{i=1}^{n} 
  f^*(\lambda_i)
\end{gather}
where $f^*$ is the Legendre transformation of $f$.

In the sequel we need $f$ to be strictly convex and its convex conjugate (or Legendre transformation) to be continuously differentiable and strictly non-decreasing.
Two popular choices of $f$ are the \textbf{negative entropy} and the \textbf{sample variance}.
\subsection*{Negative Entropy}
We define the negative entropy to be
\begin{gather}
  f
  \colon
  [0,\infty)
  \to
  \R
  ,\quad
  w
  \mapsto
  \begin{cases}
    0\quad\text{if}\ w=0,\\
    w\log w\quad
    \text{else}.
  \end{cases}
\end{gather}
It is strictly convex. To compute its Legendre transformation we note, that
\begin{gather}
  (f^{'})^{-1}
  =
  \lambda\mapsto
  e^{\lambda-1}
\end{gather}
Thus
  \begin{align*}
  f^*
  (\lambda)
  &
  \ 
  =
  \ 
  \lambda
    \cdot
    (f^{'})^{-1}(\lambda)
  \ 
    -
  \ 
    f
    \left( 
      (f^{'})^{-1}(\lambda)
    \right)
    \\
  &
  \ 
  =
  \ 
  \lambda
    \cdot
  e^{\lambda-1}
  \ 
    -
  \ 
  e^{\lambda-1}
  \log
  \left( 
  e^{\lambda-1}
  \right)
  \\
  &
  \ 
  =
  \ 
  e^{\lambda-1}
  \,.
  \end{align*}
  Thus $f^*$ is smooth and strictly non-decreasing.


  \subsection*{Sample Variance}
We define the sample variance to be
\begin{gather}
  f
  \colon
  \R
  \to
  \R
  ,\quad
  w
  \mapsto
  (w-1/n)^2
\end{gather}
It is strictly convex. To compute its Legendre transformation we note, that
\begin{gather}
  (f^{'})^{-1}
  =
  \lambda\mapsto
  \frac{\lambda}{2}
  +
  \frac{1}{n}
\end{gather}
Thus
  \begin{align*}
  f^*
  (\lambda)
  &
  \ 
  =
  \ 
  \lambda
    \cdot
    \left( 
  \frac{\lambda}{2}
  +
  \frac{1}{n}
    \right)
  \ 
    -
  \ 
    \left( 
    \left( 
  \frac{\lambda}{2}
  +
  \frac{1}{n}
    \right)
    -
    \frac{1}{n}
    \right)
    ^2
    \\
  &
  \ 
  =
  \ 
  \frac{\lambda^2}{4}
  +
  \frac{\lambda}{n}
  \,.
  \end{align*}
  Thus $f^*$ is smooth.
  To eliminate some variables in the optimization problem,
  we need $f^*$ also to be
  strictly non-decreasing. But the sample variance violates this assumption.

\section*{Constraints}
Let's turn our attention to the constraints.
The first constraint makes sure we do not extrapolate from the poputation.
The second constraint norms the weights. 
The third constraint controls the bias of the resulting estimator.
\section*{Regression Basis}
We adopt ideas from \cite{Gyorfi2002}. Another angle would be sieve estimates\cite{Newey1997a} where the number of basis functions can grow slower than $N$.
Their notion of (weak) consistency~\cite[Definitien~1.1]{Gyorfi2002} for noiseless estimands
is
\begin{gather}
  \E
  \left[ 
    \int_\mathcal{X}
    \left| 
    \sum_{k=1}^{N} 
    B_k(x)
    \cdot
    m(X_k)
    -
    m(x)
    \right|
    ^2
    \P_X(dx)
  \right]
  \to
  0
  \qquad
  \text{as}
  \ 
  n\to \infty
  \,.
\end{gather}
Universal consistency in this sense holds, if this is true for all distributions with
$\E[m(X)^2]<\infty$(cf.\cite[Definition~1.3]{Gyorfi2002}).

We adopt a slightly different notion of consistency. The next theorem dose the translation work. 
\begin{theorem}
  Assume
  $\E[m(X)^2]<\infty$
  and the basis function are 
(weak) universal consistency in the sense of 
\cite[Definitien~1.3]{Gyorfi2002}.
Then it holds for all $\varepsilon>0$
\begin{gather}
  \P
  \left[ 
    \left| 
    \sum_{k=1}^{N} 
    B_k(X)
    \cdot
    m(X_k)
    -
    m(X)
    \right|
    \ge
    \varepsilon
  \right]
  \to
  0
  \qquad
  \text{as}
  \ 
  n\to \infty
  \,.
\end{gather}
\end{theorem}
\begin{proof}
  By Markov's inequality it holds
  \begin{align*}
    &
  \P
  \left[ 
    \left| 
    \sum_{k=1}^{N} 
    B_k(X)
    \cdot
    m(X_k)
    -
    m(X)
    \right|
    \ge
    \varepsilon
  \right]
  \\
  &
  \ 
  \le
  \ 
  \frac
  {
  \E
  \left[ 
    \left| 
    \sum_{k=1}^{N} 
    B_k(X)
    \cdot
    m(X_k)
    -
    m(X)
    \right|
    ^2
  \right]
  }
  {\varepsilon^2}
  \\
  &
  \ 
  =
  \ 
  \frac
  {
  \E
  \left[ 
    \E
    \left[ 
    \left| 
    \sum_{k=1}^{N} 
    B_k(X)
    \cdot
    m(X_k)
    -
    m(X)
    \right|
    ^2
    |
    X_1,
    \ldots,
    X_N
    \right]
  \right]
  }
  {\varepsilon^2}
  \\
  &
  \ 
  =
  \ 
  \frac
  {
  \E
  \left[ 
    \int_\mathcal{X}
    \left| 
    \sum_{k=1}^{N} 
    B_k(x)
    \cdot
    m(X_k)
    -
    m(x)
    \right|
    ^2
    \P_X(dx)
  \right]
  }
  {\varepsilon^2}
  \,.
  \end{align*}
  The last equality is due to \cite[(1.2)]{Gyorfi2002}.
  By the weak universal consistency of $B$
  the last expression goes to $0$ as $N\to \infty$.
\end{proof}

Classical choices of the basis functions are \textbf{partitioning estimates} and 
\textbf{kernel estimates}(cf.\cite[§4,§5]{Gyorfi2002}).

\subsection*{Partitioning Estimates}
We consider a partition
$
  \mathcal{P}_N
  =
  \left\{ 
    A_{N,1}
    ,
    A_{N,2}
    ,
    \ldots
  \right\}
$
of $ \R^d $
and define
$ A_N(x) $ to be the cell of $ \mathcal{P}_N $ containing $x$.
We define $N$ basis functions $B_k$ of the covariates by
\begin{gather*}
  B_k(x)
  :=
  \frac{
  \mathbf{1}_{X_k \in A_N(x)}
  }{
  \sum_{j=1}^{N} 
  \mathbf{1}_{X_j \in A_N(x)}
  }
  \,,
  \qquad
  k=
  1,\ldots,N
  \,.
\end{gather*}
The euclidian norm of the basis functions is bounded above by $1$.
\begin{gather*}
  \norm{B(x)}^2
  =
  \sum_{k=1}^{n} 
  \left( 
  \frac{
  \mathbf{1}_{X_k \in A_n(x)}
  }{
  \sum_{j=1}^{n} 
  \mathbf{1}_{X_j \in A_n(x)}
  }
  \right)
  ^2
  \le
  \sum_{k=1}^{n} 
  \frac{
  \mathbf{1}_{X_k \in A_n(x)}
  }{
  \sum_{j=1}^{n} 
  \mathbf{1}_{X_j \in A_n(x)}
  }
  =1
  \,.
\end{gather*}
Under mild conditions, the basis functions are universally consistent.
\begin{theorem}
  \label{bw:i:bf:pe:th:c}
  If for each sphere $S$ centered at the origin 
  \begin{gather}
    \max
    _
    {
      j\colon
      A_{N,j} 
      \cap
      S
      \neq
      \emptyset
    }
    \mathrm{diam}
    \ 
      A_{N,j} 
      \ 
      \to
      \ 
      0
      \qquad
      \text{for}\ 
      N\to \infty 
  \end{gather}
  and
  \begin{gather}
    \frac
    {
    \#
    \left\{  
      j\colon
      A_{N,j} 
      \cap
      S
      \neq
      \emptyset
    \right\}
    }
    {N}
      \ 
      \to
      \ 
      0
      \qquad
      \text{for}\ 
      N\to \infty 
  \end{gather}
  then the partitioning regression function estimate 
  (definition)
  is
  universally consistent (definition).
\end{theorem}
\begin{proof}
  \cite[Theorem~4.2.]{Gyorfi2002}
\end{proof}
\begin{corollary}
Assume
  $\E[m(X)^2]<\infty$
  and the basis functions $B$ belong to a partitioning estimate.
  Furthermore assume that the conditions of 
  Theorem~\ref{bw:i:bf:pe:th:c} are met.
Then it holds for all $\varepsilon>0$
\begin{gather}
  \P
  \left[ 
    \left| 
    \sum_{k=1}^{N} 
    B_k(X)
    \cdot
    m(X_k)
    -
    m(X)
    \right|
    \ge
    \varepsilon
  \right]
  \to
  0
  \qquad
  \text{as}
  \ 
  n\to \infty
  \,.
\end{gather}
\end{corollary}
\subsection*{Kernel Estimates}
Let $K\colon \R^d\to [0,1]$ (bounded kernel) and $h_n>0$ (bandwith).
For examples see \cite[§5.1.]{Gyorfi2002}.
We define
\begin{gather}
  B_k(x)
  :=
  \frac
  {
    K \left( \frac{x-X_k}{h_n} \right)
  }
  {
    \sum_{i=1}^{N} 
    K \left( \frac{x-X_i}{h_n} \right)
  }
  \,.
\end{gather}
By the boundedness of the kernel it follows
$\norm{B(x)}\le 1$.
\begin{theorem}
  \label{bw:i:bf:ke:th:c}
  Assume that there are 
  balls
  $S_{0,r}$ of radius $r$ 
  and
  balls
  $S_{0,R}$ of radius $R$ 
  centered at the origin with $0<r\le R$, and a constant $b>0$ such that
  \begin{gather}
    \mathbf{1}_{\left\{ x\in S_{0,R} \right\}}
    \ge
    K(x)
    \ge
    b
    \cdot
    \mathbf{1}_{\left\{ x\in S_{0,r} \right\}}
  \end{gather}
\end{theorem}
(boxed kernel). Then for bandwiths with $h_n\to0$ and
$n\cdot h_n^d\to\infty$ as $n\to \infty$ the kernel estimate is weakly universally consistent.
\begin{corollary}
Assume
  $\E[m(X)^2]<\infty$
  and the basis functions $B$ belong to a kernel estimate.
  Furthermore assume that the conditions of 
  Theorem~\ref{bw:i:bf:ke:th:c} are met.
Then it holds for all $\varepsilon>0$
\begin{gather}
  \P
  \left[ 
    \left| 
    \sum_{k=1}^{N} 
    B_k(X)
    \cdot
    m(X_k)
    -
    m(X)
    \right|
    \ge
    \varepsilon
  \right]
  \to
  0
  \qquad
  \text{as}
  \ 
  n\to \infty
  \,.
\end{gather}
\end{corollary}

In the sequel we mainly work with the dual problem.

\section*{Dual Problem}
\begin{ftheorem*}
  The dual of 
  Problem~\ref{bw:1:primal}
  is the unconstrained optimization problem 
\begin{gather*}
  \underset
  {\lambda_0,\ldots,\lambda_{N}\in \R}
  {
    \mathrm{minimize}
  }
  \quad
  \frac{1}{N}
\sum_{i=1} 
  ^N
  \left[ 
    \,
  T_i
  \cdot
  f^*
  \!
  \left( 
\lambda_0
+
\inner
{B(X_i)}
{\lambda}
  \right)
  \ 
-
\ 
  \left( 
\lambda_0
+
\inner
{B(X_i)}
{\lambda}
  \right)
  \,
  \right]
  \ 
+
\ 
\inner
{\delta}
{|\lambda|}
  \,.
\end{gather*}
  where
  \begin{gather*}
  f^*
  \,
  \colon
  \, 
  \R
  \ 
  \to
  \ 
  \R
  \,
  ,
  \qquad 
  x^*
  \ 
  \mapsto
  \ 
    x^*
    \!
    \cdot
    (f^{'})^{-1}(x^*)
  \ 
    -
  \ 
    f
    \left( 
      (f^{'})^{-1}(x^*)
    \right)
  \end{gather*}
  is the Legendre transformation of $f$,
  the vector
  $
    B(X_i)
    =
    \left[ 
      B_1(X_i)
      ,
      \ldots
      ,
      B_n(X_i)
    \right]
    ^\top
  $
  denotes 
  the $N$ basis functions of the covariates 
  of unit $i\in \left\{ 1, \ldots, N \right\}$
  and
  $
    \left| \lambda \right|
    =
    \left[ 
      \left| \lambda_1 \right|
      ,
      \ldots
      ,
      \left| \lambda_N \right|
    \right]
    ^\top
    ,
  $
  where $\left| \,\cdot\, \right|$
  is the absolute value of a real-valued scalar.
  Moreover,
  if $\lambda^\dagger$
  is an optimal solution of the above problem 
  then the optimal solution to problem
  Problem~\ref{bw:1:primal}
  is given by
  \begin{gather*}
    w_i^\dagger
    \ 
    =
    \ 
    (f^{'})^{-1}
    \left( 
      \inner
      {B(X_i)}
      {\lambda^\dagger}
      +
      \lambda_0^\dagger
    \right)
    \qquad
    \text{for}\ 
    i\in \left\{ 1\ldots,n \right\}
    \,.
  \end{gather*}
\end{ftheorem*}
\subsection*{Plan of proof}
  We want to apply Theorem~\ref{cv:ts:th}.
  To this end, we find the suitable \textbf{matrix notation}.
  (\cite[p.20-22]{Wang2019} fail to do so. The problem is, that they divide by 0 in the second display on p.21).
  Theorem~\ref{cv:ts:th} covers only parts of the constraints, 
  so we apply the argument in \cite[p.19-20]{Wang2019} to eliminate the remaining \textbf{non-negativity constraints}. 
%%%%%%%%%%%%%%%  
%%%% PROOF %%%%
%%%%%%%%%%%%%%%
\begin{proof}
  \subsubsection*{Matrix notation}
  We consider the vector of basis functions of the covariates 
  of unit $i\in \left\{ 1, \ldots, n \right\}$, that is,
\begin{align*}
    B(X_i)
    &
    \ 
    :=
    \ 
    \left[ 
      B_1(X_i)
      ,
      \ldots
      ,
      B_N(X_i)
    \right]
    ^\top
    \,,
    \intertext{the constraints vectors}
d
    &
    \ 
    :=
    \ 
    \begin{bmatrix}
      0_n
      \\
      -N\cdot\delta 
      \pm\,
      \sum_{i = 1}^{N} B_k(X_i)
    \end{bmatrix}
    \,,
    \\
  a
&
  \ 
    :=
    \ 
    N
    \intertext{
  the matrix of the basis functions of the treated
  }
    \mathbf{B}(\mathbf{X})
    &
    \ 
    :=
    \ 
    \begin{bmatrix}
      B(X_1), \ldots, B(X_n)
    \end{bmatrix}
    \intertext{
      and
  the constraint matrices
    }
    \mathbf{U}
    &
    \ 
    :=
    \ 
    \begin{bmatrix}
      \mathbf{I}_n
      \\
      \pm\,\mathbf{B}(\mathbf{X})
    \end{bmatrix}
        \,.
        \\
    \mathbf{A}
    &
    \ 
    :=
    \ 
      \mathrm{1}_n
  \end{align*}
  By Example~\ref{cv:cc:ex}
the convex conjugate of the objective function 
of 
  Problem~\ref{bw:1:primal} is
\begin{gather*}
  [x^*_1,\ldots,x^*_n]^\top
  \ 
  \mapsto
  \ 
  \sum_{i=1}^n f^*(x^*_i)
  \,,
\end{gather*}
Before we apply Theorem~\ref{cv:ts:th}
we eliminate the non-negativity constraints.
To this end, we consider the objective function $G$ of the dual
problem and update it until we reach its final form.
We write
\begin{gather}
  \lambda_d
  =:
  \begin{bmatrix}
    \rho\\
    \lambda^+\\
    \lambda^-
  \end{bmatrix}
\end{gather}
\begin{align*}
G(\lambda_d,\lambda_0)
&
\ 
=
\ 
G
\left( 
\rho,
  \lambda^+,
  \lambda^-,
\lambda_0
\right)
\\
&
\ 
:=
\ 
    \sum_{i = 1}^{N} 
    -
    f^*
    \left( 
    \rho_i
    +
    \lambda_0
    +
    \inner{B(X_i)}
    {\lambda^+-\lambda^-}
    \right)
    \ 
    +
    \ 
    \left( 
    \lambda_0
    +
    \inner{B(X_i)}
    {\lambda^+-\lambda^-}
    \right)
\\
&
    \qquad 
    -
    \ 
    N
    \cdot
    \inner{\delta}
    {\lambda^++\lambda^-}
\end{align*}
Since we maximize $G$ and $f^*$ is strictly non-decreasing, 
$\rho=0$ is optimal. We update $G$.
\begin{align*}
G
\left( 
  \lambda^+,
  \lambda^-,
\lambda_0
\right)
&
\ 
=
\ 
    \sum_{i = 1}^{N} 
    -
    f^*
    \left( 
    \lambda_0
    +
    \inner{B(X_i)}
    {\lambda^+-\lambda^-}
    \right)
    \ 
    +
    \ 
    \left( 
    \lambda_0
    +
    \inner{B(X_i)}
    {\lambda^+-\lambda^-}
    \right)
\\
&
    \qquad 
    -
    \ 
    N
    \cdot
    \inner{\delta}
    {\lambda^++\lambda^-}
\end{align*}

\subsubsection*{Non-negativity constraints}

  Next we want to remove the non-negativity constraints on $\lambda^\pm$.
  We show 
  for all $i \in \left\{ 1,\ldots,N \right\}$
\begin{alignat*}{2}
  \text{either}
  &
  &&
  \qquad
      \lambda_i^+ > 0
  \\
  \text{or}
  &
  &&
  \qquad
      \lambda_i^- > 0
  \,.
\end{alignat*}
Assume towards a contradiction that 
there exists
$i \in \left\{ 1,\ldots,N \right\}$
such that
$
      \lambda_i^+ > 0
$
and
$
      \lambda_i^- > 0
$ 
and that 
$\lambda^\pm$ is optimal.
Consider
  \begin{gather}
    \tilde{\lambda}
    \ 
    :=
    \ 
    \begin{bmatrix}
      \ 
      \lambda_1^+,
      \ldots,
      \ 
      \lambda_i^+
      \!
      -
      (
      \lambda_i^+
      \!
      \land
      \lambda_i^-
      )\,,
      \ 
      \ldots,
      \lambda_N^+,
      \ 
      \lambda_1^-,
      \ldots,
      \lambda_i^-
      \!
      -
      (
      \lambda_i^+
      \!
      \land
      \lambda_i^-
      )\,,
      \ 
      \ldots,
      \lambda_N^-
      \,,
      \lambda_0
    \end{bmatrix}
    ^\top
    \,.
  \end{gather}
  Since 
  $
      \lambda_i^\pm
      -
      (
      \lambda_i^+
      \!
      \land
      \lambda_i^-
      )
      \ge 
      0
  $,
  the perturbed vector $\tilde{\lambda}$ is in the domain of the 
  optimization problem.
  But 
  \begin{align}
  G(\tilde{\lambda})
  -
  G(\lambda)
  \ 
  =
  \ 
  2
  N
  \cdot
  \delta_i
  \cdot
      (
      \lambda_i^+
      \!
      \land
      \lambda_i^-
      )
  \ 
  >
  \ 
  0
  \,,
  \end{align}
  which contradicts the optimality of $\lambda$.
But then 
$
\lambda^\pm_i
\ge 0
$
collapses to
$
\lambda_i\in \R
$ 
for all
$i\in \left\{ 0,\ldots,N \right\}$, that is,
$ \lambda_i=\lambda_i^+\!-\lambda_i^- $.
Note that
$ |\lambda_i|=\lambda_i^+\!+\lambda_i^- $.

We update the objective function one more time.
Multiplying with $-1/N$ and introducing $T$ we get

\begin{gather*}
  \underset
  {\lambda_0,\ldots,\lambda_{N}\in \R}
  {
    \mathrm{minimize}
  }
  \quad
  \frac{1}{N}
\sum_{i=1} 
  ^N
  \left[ 
    \,
  T_i
  \cdot
  f^*
  \!
  \left( 
\lambda_0
+
\inner
{B(X_i)}
{\lambda}
  \right)
  \ 
-
\ 
  \left( 
\lambda_0
+
\inner
{B(X_i)}
{\lambda}
  \right)
  \,
  \right]
  \ 
+
\ 
\inner
{\delta}
{|\lambda|}
  \,.
\end{gather*}
We apply Theorem~\ref{cv:ts:th} to finish the proof.
\end{proof}
We have gathered all the tools to tackle consistency of the weighted mean.

\chapter{The Optimization Problem behind the Weights}
Let $(T_1,X_1),\ldots,(T_N,X_N)$ be independent and identically-distributed copies of $T$ and $X$ (see Introduction page). 
We gather them in the (random) data set 
\begin{gather*}
D_N:=\left\{ (T_i,X_i)\colon i\in \left\{ 1,\ldots,N \right\} \right\}
\,.
\end{gather*}
Furthermore, let
\begin{gather*}
  n
  \ 
  :=
  \ 
  \# 
  \left\{ 
    i\in \left\{ 1,\ldots,N \right\}
    \colon
    T_i=1
  \right\}
\end{gather*}
be the number of treated units. This is a random variable. We assume the order $T_i=1$ for all $i\le n$.

Let
$\# A\in\mathbb{N}_0$ be the number of elements of a finite set $A$, let
$\overline{\R}:=\R\cup \left\{ \infty \right\}$ 
be the extended real numbers, 
$\Sigma_\mathcal{X}$ be a $\sigma$-algebra on the covariate space $\mathcal{X}$,
and let $\mathcal{B}(\R)$ denote the Borel-$\sigma$-algebra on the real numbers.

For a convex function $\varphi\colon \R\to\overline{\R}$,
a finite set of measurable (basis-)functions on the covariate space
\begin{gather*}
  \mathfrak{B}
  :
  =
  \left\{ 
B_k\colon (\mathcal{X},\Sigma_\mathcal{X})\to(\R,\mathcal{B}(\R))
\ 
|
\ 
k\in \left\{ 1,\ldots,\# \mathfrak{B} \right\}
  \right\}
\end{gather*}
and a constraints vector $\delta=[\delta_1,\ldots,\delta_{\# \mathfrak{B}}]^\top>0$
we consider the following (random) convex optimization problem.
\newpage
\begin{fproblem}
  \label{bw:1:primal}
\begin{align*}
  %%%% objective %%%%
    &\underset{w_1, \ldots, w_n \in \R}
    {\text{minimize}}
    &&\qquad\qquad
    \sum_{i = 1}^{n} 
    \varphi(w_i)
    &&&
    \\
    %%%% w_i T_i >= 0 %%%%
    &\text{subject to}
    &&\qquad\qquad
    w_i 
    \ge
    0
    &&&
    \qquad
    \text{for all}\ 
    i \in \left\{ 1, \ldots, n \right\}
    \,,
    \\
    %%%% 1/n sum w = 1 %%%%
    & 
    &&\qquad\qquad
    \frac{1}{N}
    \sum_{i=1}^{n} 
    w_i
    =1
    \\
    %%%% box constraints %%%%
    & 
    &&\qquad
    \left| 
      \frac{1}{N} 
      \left( 
      \sum_{i = 1}^{n} 
      w_i
      B_k(X_i)
      -
      \sum_{i=1}^{N} 
      B_k(X_i)
      \right)
    \right|
    \ 
    \le 
    \ 
    \delta_k
    &&&
    \qquad
    \text{for all}\ 
    k \in \left\{ 1, \ldots, \# \mathfrak{B} \right\}
    \,.
\end{align*}
\end{fproblem}
What is random in Problem~\ref{bw:1:primal}?
First, the dimension of the search space $(w\in\R^n)$ depends on the random variable $n$. 
Thus, we only compute weights for the treated units (the ones with $T_i=1$).
Next consider the \textbf{objective function}
\begin{gather*}
    \sum_{i = 1}^{n} 
    \varphi(w_i)
    \,.
\end{gather*}
The number of summands is random (again $n$). We sometimes use the equivalent notation
\begin{gather*}
    \sum_{i = 1}^{N} 
    T_i
    \cdot
    \varphi(w_i)
    \,,
\end{gather*}
where we set the weights of the untreated (the ones with $T_i=0$) to some arbitrary value in the domain of $\varphi$.
Let's consider the constraints. There is no randomness in the first two constraints.
\begin{gather*}
    w_i 
    \ge
    0
    \qquad
    \text{for all}\ 
    i \in \left\{ 1, \ldots, n \right\}
    \quad
    \text{and}
    \quad
    \frac{1}{N}
    \sum_{i=1}^{n} 
    w_i
    =1
    \,.
\end{gather*}
They only make sure, that the weights (divided by $N$) form a convex combination.
If, for example, the outcome space $\mathcal{Y}$ is convex we make sure that a weighted-mean-estimate of $\E[Y(1)]$ satisfies
\begin{gather*}
  \widehat{Y}(1) 
  \ 
  :=
  \ 
  \frac{1}{N}
  \sum_{i=1}^{n} 
  w_i\cdot Y_i
  \ 
  \in
  \ 
  \mathcal{Y}
\end{gather*}
or that a weighted-mean-estimate of the distribution function of $Y(1)$ satisfies
\begin{gather*}
  \widehat{F}_{Y(1)} 
  \ 
  :=
  \ 
  \frac{1}{N}
  \sum_{i=1}^{n} 
  w_i\cdot \mathbf{1}\left\{ Y_i\le z \right\}
  \ 
  \in
  \ 
  [0,1]
  \,.
\end{gather*}
There remain the box constraints
\begin{gather*}
    \left| 
      \frac{1}{N} 
      \left( 
      \sum_{i = 1}^{n} 
      w_i
      B_k(X_i)
      -
      \sum_{i=1}^{N} 
      B_k(X_i)
      \right)
    \right|
    \ 
    \le 
    \ 
    \delta_k
    \qquad
    \text{for all}\ 
    k \in \left\{ 1, \ldots, \# \mathfrak{B} \right\}
    \,.
\end{gather*}
Here the number of summands in
\begin{gather*}
      \sum_{i = 1}^{n} 
      w_i
      B_k(X_i)
\end{gather*}
is again random, and we sometimes use the equivalent notation
\begin{gather*}
      \sum_{i = 1}^{N} 
      T_i
      \cdot
      w_i
      \cdot
      B_k(X_i)
      \,.
\end{gather*}
The (basis-)functions in $\mathfrak{B}$ can be random functions, for example, if they depend on the data $D_N$.
Also the constraint vector $\delta$ can depend on the data (see \cite[Algorithm~1 on page 11]{Wang2019}).


Next, we formulate and discuss assumptions on (the core of) the objective function $\varphi$.
\subsection{Objective Function}
\begin{assumption}
  \label{asu:objective_f}
  The objective function $\varphi\colon \R\to\overline{\R}$ of Problem~\ref{bw:1:primal} 
  satisfies the following conditions
  \begin{itemize}
    \item
      $\varphi(0)=0$ and $\varphi(x)=\infty$ for all $x<0$
    \item
      $\varphi$ is strictly convex and continuously differentiable on $(0,\infty)$ with derivative $\varphi^{'}$
    \item
      $\varphi^{'}(0,\infty)=\R$
    \item
      the inverse of the derivative $(\varphi^{'})^{-1}$ is Lipschitz continuous and continuously differentiable on $\R$.
  \end{itemize}
\end{assumption}
The next lemma provides a link to the assumptions on the objective function in Theorem~\ref{cv:ts:th}.
\newpage
\begin{lemma}
  \label{1065}
  Let Assumption\ref{asu:objective_f} hold true. Then the convex conjugate of $\varphi$ (see \eqref{def:convex_conjugate}) is
  \begin{gather*}
    \varphi^*
    \colon
    \R
    \ 
    \to
    \ 
    \R
    \,,
    \quad
    x^*
    \ 
    \mapsto
    \ 
    x^*
    \!
    \cdot
    (
    \varphi^{'}
    )^{-1}
    (x^*)
    \ 
    -
    \ 
    \varphi
    \left( 
      (
    \varphi^{'}
    )^{-1}
    (x^*)
    \right)
    \,.
  \end{gather*}
  Furthermore, $\varphi^*$ is strictly convex and continuously differentiable on $\R$.
\end{lemma}
\begin{proof}
We define
\begin{gather*}
 \phi
 \ 
 \colon
 [0,\infty)
 \times
 \R
 \ 
 \to
 \ 
 \R
 \,,
 \quad
 (x,x^*)
 \ 
 \mapsto
 \ 
 x\cdot x^*
 \ 
 -
 \ 
 \varphi(x)
 \,.
\end{gather*}
Let $x^*\in\R$.
Since
(by assumption)
      $\varphi$ is continuously differentiable on $(0,\infty)$ with derivative $\varphi^{'}$,
      so is $\phi(\cdot,x^*)$ with derivative
      satisfying 
      \begin{gather*}
        \frac{\partial}{\partial x}
        \phi(x,x^*)
        \ 
        =
 \ 
        x^*
        -
        \varphi^{'}(x)
        \qquad
        \text{for all}\ 
        x\in(0,\infty)
        \,.
      \end{gather*}
  It follows that 
  \begin{gather*}
    z
    \ 
    :=
    \ 
      (
    \varphi^{'}
    )^{-1}
    (x^*)
  \end{gather*}
  is an extreme value of $\phi(\cdot,x^*)$.
  Note, that $x^*\in\R$ is in the domain of 
  $
      (
    \varphi^{'}
    )^{-1}
  $
  by the assumption $\varphi^{'}(0,\infty)=\R$.
  Since $\varphi$ is strictly convex, $\phi(\cdot,x^*)$ is strictly concave. 
  Thus,
  $z>0$ is the unique maximum in $(0,1)$.
  By the continuity of $\phi(\cdot,x^*)$ on $[0,\infty)$ it follows, that $z$ is the unique maximum on $[0,\infty)$.
  Thus
  \begin{align*}
    \varphi^*(x^*)
    &
    \ 
    =
    \ 
    \sup_{x\in\R}
    x\cdot x^* - \varphi(x)
    \ 
    =
    \ 
    \sup_{x\in [0,\infty)}
    x\cdot x^* - \varphi(x)
    \ 
    =
    \ 
    \sup_{x\in [0,\infty)}
    \phi(x,x^*)
    \\
    &
    \ 
    =
    \ 
    \phi(z,x^*)
    \\
    &
    \ 
    =
    \ 
    x^*
    \!
    \cdot
    (
    \varphi^{'}
    )^{-1}
    (x^*)
    \ 
    -
    \ 
    \varphi
    \left( 
      (
    \varphi^{'}
    )^{-1}
    (x^*)
    \right)
    \qquad 
    \text{for all}\ 
    x^*\in\R
    \,.
  \end{align*}
  Now we proof the second statement.
  Since
  $
      (
    \varphi^{'}
    )^{-1}
  $
  is (by assumption) continuously differentiable, it holds
  \begin{align}
    \label{0098}
    \begin{split}
    \frac{\partial}{\partial x^*}
     \varphi^*(x^*)
    &
    \ 
    =
    \ 
    (
    \varphi^{'}
    )^{-1}
    (x^*)
    \ 
    +
    \ 
    x^*
    \!
    \cdot
    \frac{\partial}{\partial x^*}
    (
    \varphi^{'}
    )^{-1}
    (x^*)
    \ 
    -
    \ 
    \varphi^{'}
    \left( 
      (
    \varphi^{'}
    )^{-1}
    (x^*)
    \right)
    \cdot
    \frac{\partial}{\partial x^*}
    (
    \varphi^{'}
    )^{-1}
    (x^*)
    \\
    &
    \ 
    =
    \ 
    (
    \varphi^{'}
    )^{-1}
    (x^*)
    \ 
    +
    \ 
    x^*
    \!
    \cdot
    \frac{\partial}{\partial x^*}
    (
    \varphi^{'}
    )^{-1}
    (x^*)
    \ 
    -
    \ 
    x^*
    \cdot
    \frac{\partial}{\partial x^*}
    (
    \varphi^{'}
    )^{-1}
    (x^*)
    \\
    &
    \ 
    =
    \ 
    (
    \varphi^{'}
    )^{-1}
    (x^*)
    \qquad
    \text{for all}\ 
    x^*\in\R
    \,.
    \end{split}
  \end{align}
  Since $\varphi$ is strictly convex and continuously differentiable, 
  $\varphi^{'}$ is continuous and strictly non-decreasing.
  Thus 
  $
    (
    \varphi^{'}
    )^{-1}
  $
  is continuous and strictly non-decreasing.
  It follows from \eqref{0098} that $\varphi^*$ is strictly convex and continuously differentiable.
\end{proof}
The next lemma completes the link.
\begin{lemma}
  Let Assumption~\ref{asu:objective_f} hold true. Then 
\begin{gather*}
  \Phi
  \ 
  :
  \ 
  \R^n
  \to
  \ 
  \overline{\R}
  \,
  ,
  \qquad
  [w_1,\ldots,w_n]^\top
  \ 
  \mapsto
  \ 
  \sum_{i=1}^n \varphi(w_i)
  \,,
\end{gather*}
satisfies Assumption~\ref{cv:ts:asu}.
\end{lemma}
\begin{proof}
  By Example~\ref{cv:cc:ex}
  the convex conjugate of $\Phi$ is 
\begin{gather*}
  \Phi^*
  \ 
  :
  \ 
  \R^n
  \to
  \ 
  \overline{\R}
  \,
  ,
  \qquad
  [\lambda_1,\ldots,\lambda_n]^\top
  \ 
  \mapsto
  \ 
  \sum_{i=1}^n \varphi^*(\lambda_i)
  \,,
\end{gather*}
where $\varphi^*$ is the convex conjugate of $\varphi$.
By Assumption~\ref{asu:objective_f} $\varphi$ is strictly convex. Thus,
$\Phi$ is strictly convex. By Lemma~\ref{1065}, $\varphi^*$ continuously differentiable on $\R$. Thus,
$\Phi$ is continuously differentiable on $\R^n$.
It follows the statement of Assumption~\ref{cv:ts:asu} for $\Phi$.
\end{proof}


Next we discuss some concrete choices of $\varphi$

\begin{example}
  For two discrete distributions
  \begin{gather*}
    p:=[p_1,\ldots,p_N]
    \qquad
    \text{and}
    \qquad
    q:=[q_1,\ldots,p_N]
  \end{gather*}
  we consider the following distance measure
  \begin{gather*}
   D(p|q)
   \ 
   :=
   \ 
   \sum_{i=1}^{N} 
   p_i
   \cdot
   \log
   \left( 
     \frac{p_i}{q_i}
   \right)
   \,.
  \end{gather*}
  This is known as the Kullback-Leibler-Entropy.
  In \cite[§3.1]{Hainmueller2012} the author connects this concept to a convex optimization problem.
  The idea is, to optimize the Kullback-Leibler-Entropy of the distribution induced by the weights and some base weights.
  For example, if we choose
  \begin{gather*}
    w:=
    \frac{1}{N}[w_1\cdot T_1,\ldots,w_N\cdot T_N]
    \qquad
    \text{and}
    \qquad
    q:=\frac{1}{N}[1,\ldots,1]
  \end{gather*}
  we get
  \begin{align}
    \label{8926}
   D(w|q)
   \ 
   =
   \ 
   \frac{1}{N}
   \sum_{i=1}^{N} 
w_i\cdot T_i
   \cdot
   \log
   \left( 
     \frac{w_i\cdot T_i}{N}\cdot N
   \right)
   \ 
   =
   \ 
   \frac{1}{N}
   \sum_{i=1}^{n} 
w_i
   \cdot
   \log
   \left( 
     w_i
   \right)
   \,,
  \end{align}
  where we set $"0\cdot \log(0)"=0$. Thus, the optimization problem
  \begin{gather*}
    \underset{w_1, \ldots, w_n \in \R}
    {\text{minimize}}
    \qquad\qquad
    \sum_{i = 1}^{n} 
w_i
   \cdot
   \log
   \left( 
     w_i
   \right)
  \end{gather*}
  produces the same optimal solutions as minimizing the Kullback-Leibler-Entropy \eqref{8926} with respect to $w$.
Thus, we consider 
\begin{gather*}
  \varphi
  \ 
  \colon
  \R
  \ 
  \to
  \ 
  \overline{\R}
  \,,
  \qquad
  x
  \ 
  \mapsto
  \ 
  \begin{cases}
    x\cdot \log x\ &\text{if}\ x>0\,, \\
    0 \ &\text{if}\, x=0\,,\\
    \infty\ &\text{if}\ x<0\,.
  \end{cases}
\end{gather*}
We show, that this choices satisfies Assumption~\ref{asu:objective_f}.
By definition it holds \textit{(i)}.
For \textit{(ii)} note, that the second derivative on $(0,\infty)$ is $x\mapsto 1/x$. Thus (by the second derivative test), $\varphi$ is strictly convex.
Clearly it is also twice continuously differentiable.
For the continuity in $0$ note, that $\lim_{x\to 0} x\log x=0$. 
For \textit{(iii)} note, that $x\mapsto \log x$ is continuous and strictly non-decreasing on $\R$. 
Since 
\begin{align*}
  \lim_{x\to 0}\log x\ =\ -\infty
  \qquad
  \text{and}
  \qquad
  \lim_{x\to \infty}\log x\ =\ \infty
  \,,
\end{align*}
and $\varphi^{'}=(x\mapsto \log x +1)$ on $(0,\infty)$, it follows \textit{(iii)}. 
Finally, it holds $(\varphi^{'})^{-1}=(x\mapsto \exp(x-1))$ on $\R$. 
Thus, it follows \textit{(iv)}.
\end{example}
\begin{example}
  In a similar setting, the authors of \cite{Zubizarreta2015} choose the sample variance of the weights as objective function, that is,
\begin{gather*}
  \varphi
  \ 
  \colon
  \R
  \ 
  \to
  \ 
  \overline{\R}
  \,,
  \qquad
  x
  \ 
  \mapsto
  \ 
  \begin{cases}
    \left(
      x-\frac
      {1}
      {n}
    \right)^2\ &\text{if}\ x\ge0\,, \\
    \infty\ &\text{if}\ x<0\,.
  \end{cases}
\end{gather*}
\end{example}



Next, we derive the dual formulation of Problem~\ref{bw:1:primal}.
\subsection{Dual Problem}
\begin{lemma}
  \label{matrix_notation}
  A matrix formulation of Problem~\ref{bw:1:primal} is 
\begin{align}
  \label{cv:ts:primal}
  %%%% objective %%%%
    &\underset{w \in \R^n}
    {\mathrm{minimize}}
    &&\qquad\qquad
    \Phi(w)
    &&&
    \\
    %%%% Ax >= b %%%%
    \nonumber
    &\mathrm{subject}\ \mathrm{to} 
    &&\qquad\qquad
    \mathbf{U}w
    \ 
    \ge
    \ 
    d
    \,,
    \\
    \nonumber
    &
    &&\qquad\qquad
    \mathbf{A}w
    \ 
    =
    \ 
    a
    \,,
\end{align}
with objective function
\begin{gather*}
  \Phi
  \ 
  :
  \ 
  \R^n
  \to
  \ 
  \overline{\R}
  \,
  ,
  \qquad
  [w_1,\ldots,w_n]^\top
  \ 
  \mapsto
  \ 
  \sum_{i=1}^n \varphi(w_i)
  \,,
\end{gather*}
inequality matrix and vector
\begin{alignat*}{2}
    \mathbf{U}
    &
    \ 
    :=
    \ 
    \begin{bmatrix}
      \mathbf{I}_n
      \\
      \pm\,\mathbf{B}(\mathbf{X})
    \end{bmatrix}
    \in
    \R^{(n+  2 N)\times n}
        \qquad
    &&
d
    \ 
    :=
    \ 
    \begin{bmatrix}
      0_n
      \\
      -N\cdot\delta 
      \ 
      \pm\ 
      \sum_{i = 1}^{N} B(X_i)
    \end{bmatrix}
    \in
    \R^{n+  2 N}
    \,,
    \intertext{and equality matrix and vector}
    \mathbf{A}
    &
    \ 
    :=
    \ 
      \mathrm{1}_n
      ^\top
      \in\R^{1\times n}
      \qquad
    &&
    a
  \ 
    :=
    \ 
    N
    \in\mathbb{N}
    \,.
\end{alignat*}
\end{lemma}

\begin{proof}
  Recall that the box constraints of Problem~\ref{bw:1:primal} are
  \begin{gather*}
        \left| 
      \frac{1}{N} 
      \left( 
      \sum_{i = 1}^{n} 
      w_i
      B_k(X_i)
      -
      \sum_{i=1}^{N} 
      B_k(X_i)
      \right)
    \right|
    \ 
    \le 
    \ 
    \delta_k
    \qquad
    \text{for all}\ 
    k\in \left\{ 1,\ldots, N \right\}
    \,.
  \end{gather*}
  Put differently, it holds both
  \begin{align*}
    -
      \sum_{i = 1}^{n} 
      w_i
      B_k(X_i)
    \ge 
    -
    N
    \delta_k
      -
      \sum_{i=1}^{N} 
      B_k(X_i)
      \quad 
    \text{and}
      \quad
      \sum_{i = 1}^{n} 
      w_i
      B_k(X_i)
    \ge 
    -
    N
    \delta_k
      +
      \sum_{i=1}^{N} 
      B_k(X_i)
  \end{align*}
  for all 
  $
    k\in \left\{ 1,\ldots, N \right\}
  $. In matrix notation this is 
  \begin{gather*}
    \pm\mathbf{B}(\mathbf{X})w
    \ 
    \ge
    \ 
    [d_{n+1},\ldots, d_{n+  2 N}]^\top
    \,.
  \end{gather*}
  Proving the rest of the statements is straightforward. We omit the details.
\end{proof}
\begin{remark}
  The inequality constraints of
  Lemma~\ref{matrix_notation} differ from its counterpart
  \cite[Proof of Lemma~1]{Wang2019}.
  We don't transform the variable $w$, but shift to $d$ what prevents us from keeping $w$.
  Note, that the choice of
  \cite[Proof of Lemma~1]{Wang2019} leads to a mistake on page 21.
  The mistake is most obvious in the second display, where the first implication follows from dividing by 0.
  I discussed this with the authors and proposed a version of Lemma\ref{matrix_notation} to solve the problem. I think it's best not to transform variables, because the mistake comes from (wrongly) calculating the convex conjugate of the (more complicated) transformed version of the objective function. The subsequent analysis even simplifies with my version.

  I was surprised to find the (exact) same mistake in the earlier paper 
  \cite[page 35 second display]{Chan2016}. 
  There is no reference in
  \cite[Proof of Lemma~1]{Wang2019} 
  to
  \cite{Chan2016}. Yet the formulation and the mistake are the same.
  Did the authors of \cite{Wang2019} (inadvertently?) plagiarize
  the mathematical analysis of 
  \cite{Chan2016}
  ?
\end{remark}

\begin{lemma}
  Consider the optimization problem
\begin{align}
  \label{9993}
  \begin{split}
  \underset
  {\begin{smallmatrix}
\rho\,,\, \lambda^+,\,\lambda^-\ge 0 \\
\lambda_0\in\R
  \end{smallmatrix}}
  {
    \mathrm{maximize}
  }
  \quad
  &
  -
\sum_{i=1} 
  ^n
    \,
  \varphi^*
  \!
  \left( 
    \rho_i
    +
\lambda_0
+
\inner
{B(X_i)}
{
\lambda^+
-
\lambda^-
}
  \right)
  \\
  &
+
\ 
\sum_{i=1}^{N} 
  \left( 
\lambda_0
+
\inner
{B(X_i)}
{
\lambda^+
-
\lambda^-
}
  \right)
  \,
  \ 
-
\ 
\inner
{\delta}
{
\lambda^+
+
\lambda^-
}
  \,.
  \end{split}
\end{align}
If Assumption~\ref{asu:objective_f} holds true 
and there exists an optimal solution 
$
(\rho^\dagger,\lambda_0^\dagger,\lambda^{+,\dagger},\lambda^{-,\dagger})
$
then the unique optimal solutions to Problem~\ref{bw:1:primal} are 
\begin{gather*}
  w^\dagger_i
  \ 
  :=
  \ 
  (
  \varphi^{'}
  )^{-1}
  \left(
    \rho^\dagger_i
  \ 
    +
  \ 
\lambda_0^\dagger
  \ 
+
  \ 
\inner
{B(X_i)}
{
  \lambda^{+,\dagger}
-
\lambda^{-,\dagger}
}
  \right)
  \qquad
  \text{for all}\ 
  i\in
  \left\{ 1,\ldots,n \right\}
  \,.
\end{gather*}
\end{lemma}
\begin{proof}
  By Lemma~\ref{matrix_notation},
  Problem~\ref{bw:1:primal} has the form required in Theorem~\ref{cv:ts:th}.
  By Assumption~\ref{asu:objective_f} and Lemma~\ref{9991} the objective function $\Phi$ of Problem~\ref{bw:1:primal}
  satisfies Assumption~\ref{cv:ts:asu}.
  Thus we can apply
  Theorem~\ref{cv:ts:th} to Problem~\ref{bw:1:primal}.
  Calculations yield the result.
\end{proof}
The next Theorem aims at simplifying this result. 
\newpage
\begin{ftheorem}
  \label{dual_solution_th}
  Consider the optimization problem
\begin{align}
  \label{dual}
  \begin{split}
  \underset
  {\begin{smallmatrix}
      \rho&\in&&\R^N 
      \\
      \lambda_0 & \in&&\R
      \\
      \lambda&\in&&\R^{N}
  \end{smallmatrix}}
  {
    \mathrm{minimize}
  }
  \quad
  \frac{1}{N}
\sum_{i=1} 
  ^N
  &
  \Big[
  T_i
  \cdot
  \varphi^*
  \!
  \left( 
    \rho_i
    +
\lambda_0
+
\inner
{B(X_i)}
{
\lambda
}
  \right)
  \ 
  -
  \ 
\lambda_0
-
\inner
{B(X_i)}
{
\lambda
}
\Big]
  \ 
+
\ 
\inner
{\delta}
{
  |\lambda|
}
  \,,
  \\
  \mathrm{subject}\ \mathrm{to}
  \quad
  \qquad
  &
  \rho_i \ge 0 
  \quad 
  \mathrm{for}\ \mathrm{all}\ i\le n
  \qquad 
  \mathrm{and}
  \qquad
  \rho_i=0
  \quad 
  \mathrm{for}\ \mathrm{all}\ i>n
  \,.
\end{split}
\end{align}
If Assumption~\ref{asu:objective_f} holds true 
and there exists an optimal solution 
$
(\rho^\dagger,\lambda_0^\dagger,\lambda^\dagger)
$
then the unique optimal solutions to Problem~\ref{bw:1:primal} are 
\begin{gather*}
  w^\dagger_i
  \ 
  :=
  \ 
  (
  \varphi^{'}
  )^{-1}
  \left(
    \rho^\dagger_i
  \ 
    +
  \ 
\lambda_0^\dagger
  \ 
+
  \ 
\inner
{B(X_i)}
{
\lambda^{\dagger}
}
  \right)
  \qquad
  \text{for all}\ 
  i\in
  \left\{ 1,\ldots,n \right\}
  \,.
\end{gather*}
\end{ftheorem}

\begin{proof}
  Assume that
$
  (\rho^\dagger,\lambda_0^\dagger,\lambda^{+,\dagger},\lambda^{-,\dagger})
$
is an optimal solution to Problem~\ref{9993}.
We write
\begin{align*}
  G
  (\rho,\lambda_0,\lambda^+,\lambda^-)
  &
  \ 
  :=
  \ 
 -
\sum_{i=1} 
  ^n
    \,
  \varphi^*
  \!
  \left( 
    \rho_i
    +
\lambda_0
+
\inner
{B(X_i)}
{
\lambda^+
-
\lambda^-
}
  \right)
  \\
  &
  \qquad
+
\ 
\sum_{i=1}^{N} 
  \left( 
\lambda_0
+
\inner
{B(X_i)}
{
\lambda^+
-
\lambda^-
}
  \right)
  \,
  \ 
-
\ 
\inner
{\delta}
{
\lambda^+
+
\lambda^-
}
  \,.
\end{align*}
 To eliminate the remaining constraints, 
  we paraphrase \cite[pages~19-20]{Wang2019}.
  We show 
  for all $i \in \left\{ 1,\ldots,N \right\}$
\begin{alignat}{2}
  \notag
  \text{either}
  &
  &&
  \qquad
  \lambda_i^{+,\dagger} > 0
  \\
  \label{9992}
  \text{or}
  &
  &&
  \qquad
  \lambda_i^{-,\dagger} > 0
  \,.
\end{alignat}
Assume towards a contradiction that 
\begin{gather}
  \label{1232}
  \text{
there exists
  } 
  \ 
i \in \left\{ 1,\ldots,N \right\}
\ 
\text{such that}
\qquad
  \lambda_i^{+,\dagger} > 0
  \qquad 
  \text{and}
  \qquad
  \lambda_i^{-,\dagger} > 0
  \,.
\end{gather}
Consider
  \begin{align*}
    \tilde{\lambda}^{+,\dagger}
    &
    \ 
    :=
    \ 
    \begin{bmatrix}
      \ 
      \lambda_1^{+,\dagger}
      \ldots,
      \ 
      \lambda_i^{+,\dagger}
      \!
      \ 
      -
      \ 
      (
      \lambda_i^{+,\dagger}
      \!
      \land
      \lambda_i^{-,\dagger}
      )\,,
      \ 
      \ldots,
      \lambda_{N}^{+,\dagger}
    \end{bmatrix}
    ^\top
    \intertext{and}
    \tilde{\lambda}^{-,\dagger}
    &
    \ 
    :=
    \ 
    \begin{bmatrix}
      \ 
      \lambda_1^{-,\dagger}
      \ldots,
      \ 
      \lambda_i^{-,\dagger}
      \!
      \ 
      -
      \ 
      (
      \lambda_i^{+,\dagger}
      \!
      \land
      \lambda_i^{-,\dagger}
      )\,,
      \ 
      \ldots,
      \lambda_{N}^{-,\dagger}
    \end{bmatrix}
    ^\top
    \,.
  \end{align*}
  Since
  \begin{gather*}
      \lambda_i^{\pm,\dagger}
      \!
      \ 
      -
      \ 
      (
      \lambda_i^{+,\dagger}
      \!
      \land
      \lambda_i^{-,\dagger}
      )
      \ 
      \ge 
      \ 
      0
      \,,
  \end{gather*}
  the perturbed vectors $\tilde{\lambda}^{\pm,\dagger}$ are  in the domain of the 
  optimization problem.
  By Assumption~\eqref{1232} and $\delta>0$ it follows
  \begin{align*}
  G
  \left( 
  \rho^\dagger,\lambda_0^\dagger,\tilde{\lambda}^{+,\dagger},\tilde{\lambda}^{-,\dagger}
  \right)
  \ 
  -
  \ 
  G
  \left( 
  \rho^\dagger,\lambda_0^\dagger,\lambda^{+,\dagger},\lambda^{-,\dagger}
  \right)
  \ 
  =
  \ 
  2
  \cdot
  \delta_i
  \cdot
      (
      \lambda_i^{+,\dagger}
      \!
      \land
      \lambda_i^{-,\dagger}
      )
  \ 
  >
  \ 
  0
  \,,
  \end{align*}
  which contradicts the optimality of
$
  (\rho^\dagger,\lambda^{+,\dagger},\lambda^{-,\dagger},\lambda_0^\dagger)
$
(it is supposed to be a maximum in the domain of the optimization problem)
.
It follows \eqref{9992}.
But then 
$
\lambda^{\pm,\dagger}_i
\ge 0
$
collapses to
$
\lambda_i^\dagger\in \R
$ 
for all
$i\in \left\{ 0,\ldots,N \right\}$, that is, we set
\begin{gather*}
 \lambda_i^\dagger
 \ 
 =
 \ 
 \lambda_i^{+,\dagger}
 \ 
 -
 \ 
 \lambda_i^{-,\dagger}
 \qquad
 \text{and}
 \qquad
|\lambda_i^\dagger|
\ 
=
\ 
\lambda_i^{+,\dagger}
\ 
+
\ 
\lambda_i^{-,\dagger}
\,.
\end{gather*}
Thus, we can extend the domain of Problem~\ref{9993} to $\lambda\in\R^{N}$ and update the objective function in the following way
(without changing the optimal solution).
\begin{align*}
  G
  (\rho,\lambda_0,\lambda)
  &
  \ 
  :=
  \ 
 -
\sum_{i=1} 
  ^n
    \,
  \varphi^*
  \!
  \left( 
    \rho_i
    +
\lambda_0
+
\inner
{B(X_i)}
{
\lambda
}
  \right)
  \\
  &
  \qquad
+
\ 
\sum_{i=1}^{N} 
  \left( 
\lambda_0
+
\inner
{B(X_i)}
{
\lambda
}
  \right)
  \,
  \ 
-
\ 
\inner
{\delta}
{
  |\lambda|
}
  \,.
\end{align*}
Multiplying $G$ with $-1/N$ doesn't change the solution either
(if we search instead for a minimum).
To finish the proof, we choose the notation with $T_i$ instead of $n$. This extends the domain of $\rho$ to $\R^N_{\ge 0}$, but the 
new $\rho_i$ are not effective because of $T_i=0$ for all $i>n$. 
Thus we may set them to 0.
\end{proof}


%Bring optimal solutions together. with argmax meas. and Proposition for aSsumption and assumption.

%It is useful too define a weights function
%New chapter?
In the formulation of Theorem~\ref{dual} we encounter "If (...) there exists the optimal solution $(\rho^\dagger,\lambda_0^\dagger,\lambda)$ ... " .
To be able to study asymptotic properties of the solutions we have to become independent of this assumption.
For this we need some tools from functional analysis.
\subsection{Argmax Measurability Theorem}
We follow \cite{Aliprantis2007}.
A \textbf{correspondence} $\psi$ from a set $S_1$ to a set $S_2$ assigns to each $s_1\in S_1$ a subset $\psi(s_1)\subset S_2$.
To clarify that we map $s_1$ to a set, we use the double arrow, that is,
$
  \psi
  \colon
  S_1
  \twoheadrightarrow
  S_2
$.
  Let 
  $(\mathcal{Z},\Sigma_{\mathcal{Z}})$ be a measurable space and $\mathcal{S}$  a topological space.
  We say, that a correspondence 
  $
  \psi
  \colon
  \mathcal{Z}
  \twoheadrightarrow
  \mathcal{S}
  $
  is 
  \textbf{
  weakly measurable
  },
  if
  \begin{gather*}
    \left\{ 
      z\in \mathcal{Z}
      \ 
      |
      \ 
      \psi(z)
      \cap
      O
      \neq
      \emptyset
    \right\}
    \in
    \Sigma_{\mathcal{Z}}
    \qquad
    \text{for all open subsets}
    \ 
    O\subset \mathcal{S}
    \,.
  \end{gather*}
  A \textbf{selector} from a correspondence $\psi\colon \mathcal{Z}\twoheadrightarrow \mathcal{S}$ is a function $s\colon \mathcal{Z}\to \mathcal{S}$ that satisfies 
  \begin{align*}
s(z)\in\psi(z)
\qquad
\text{for all}\ 
z\in\mathcal{Z}
\,.
  \end{align*}
  

\begin{definition}
  Let 
  $(\mathcal{Z},\Sigma_{\mathcal{Z}})$ be a measurable space, and let $\mathcal{S}_1$ and $\mathcal{S}_2$  be topological space.
  A function 
  $f\colon \mathcal{Z}\times \mathcal{S}_1 \to \mathcal{S}_2$
  is a \textbf{Caratheodory function} if
  \begin{align*}
    f(\cdot,s_1)
    &
    \colon
    \mathcal{Z}\to \mathcal{S}_2
    \qquad
    \text{is}\ 
    (\Sigma_{\mathcal{Z}},\mathcal{B}(\mathcal{S}_2))-measurable
    \ 
    \text{for all}
    \ 
    s_1\in \mathcal{S}_1
    \,,
    \intertext{and}
    f(z,\cdot)
    &
    \colon
    \mathcal{Z}\to \mathcal{S}_2
    \qquad
    \text{is continuous for all}\ 
    z\in \mathcal{Z}
    \,.
  \end{align*}
\end{definition}
\begin{theorem}
  \label{th:argmax}
  Let $\mathcal{S}$ be a separable metrizable space and
  $
  (\mathcal{Z},\Sigma_{\mathcal{Z}})
  $
  a measurable space.
  Let $\psi\colon \mathcal{Z} \twoheadrightarrow \mathcal{S}$ be a weakly measurable correspondence with non-empty compact values, and suppose
  $f\colon \mathcal{Z}\times \mathcal{S} \to \R$
  is a Caratheodory function. Define the value function 
  $m\colon \mathcal{Z}\to \R$ by
  \begin{gather*}
    m(z):=\max_{s\in\psi(z)}f(z,s)
    \,,
  \end{gather*}
  and the correspondence 
  $\mu\colon \mathcal{Z}\twoheadrightarrow \mathcal{S}$ of maximizers by
  \begin{gather*}
    \mu(z):= \left\{ 
      s\in \psi(z)
      |
      f(z,s)=m(z)
    \right\}
    \,.
  \end{gather*}
  Then the value function $m$ is measurable, 
  the argmax correspondence $\mu$ has non-empty and compact values,
  is measurable and admits a measurable selector.
\end{theorem}
\begin{proof}
  \cite[Theorem~18.19]{Aliprantis2007}
\end{proof}


%\subsection{Propensity Score Function}

%\begin{definition}
  We define the \textbf{propensity score function} by
  \begin{gather*}
    \pi
    \ 
    \colon
    \R^d
    \ 
    \to
    \ 
    [0,1]
    \,,
    \qquad
    x
    \ 
    \mapsto
    \ 
    \P
    [
    T=1
    |
    X=x
    ]
    \,,
  \end{gather*}
  that is, 
  as the conditional probability of treatment given individual characteristics $x\in\R^d$.
\end{definition}
\begin{remark}
  We define $\pi$ on the whole $\R^d$, although $\mathcal{X}\subset \R^d$ may be a much smaller (possibly finite or countable) subset.
  The reason is that we want to assume continuity.
\end{remark}
\begin{assumption}
  \label{asu:ps}
  The propensity score function $\pi$ satisfies
  \begin{itemize}
    \item
      $\pi(x)\in (0,1)$ for all $x\in\R^d$
    \item
      $\pi$ is continuously differentiable on $\R^d$
  \end{itemize}
\end{assumption}

Next we give a standard example for a propensity score model

\begin{example}
  Logistic regression
\end{example}


%Next we formulate and discuss assumptions on the basis functions.
%\subsection{Basis Functions}
%In our setting basis functions of the covariates are stochastic processes indexed over the covariate space $\mathcal{X}$.
We consider a set $\mathfrak{B}$ of stochastic processes 
\begin{gather*}
  \left\{ 
      B_k(x)
      \ 
      \colon
      (
      \Omega,
      \mathcal{A}
      ,
      \P
      )
      \ 
      \to
      \ 
      (
      \R
      ,
      \mathcal{B}(\R)
      )
      \ 
      |
      \ 
      x\in\mathcal{X}
  \right\}
  \qquad
  \text{for}\ 
k\in \left\{ 1,\ldots,\# \mathfrak{B} \right\}
\end{gather*}
and
denote the vector of the basis functions (also a stochastic processes indexed over $\mathcal{X}$)
as
\begin{align*}
B(x)
  \ 
  :=
  \ 
  [B_1(x),\ldots,B_{\# \mathfrak{B}}(x)]^\top
  \ 
  \colon
  \ 
      (
      \Omega,
      \mathcal{A}
      ,
      \P
      )
      \to
      (
      \R^{\#\mathfrak{B}}
      ,
      \mathcal{B}
      (
      \R^{\#\mathfrak{B}}
      )
      )
    \qquad
  \text{for}\ 
  x\in\mathcal{X}
  \,.
\end{align*}
\begin{gather*}
\end{gather*}
\begin{assumption}
  \label{asu:basis}
  The set of basis functions 
  $
  \mathfrak{B}
  $
  satisfies
  \begin{enumerate}[label=(\roman*)]
    \item
      $\norm{B(x)}_2\lesssim 1$ for all $x\in \mathcal{X}$.
    \item
      There exist random vectors
      $
      \lambda^*_{\varphi^{'}\circ\, 1/\pi}
      $
      and
      $
      \lambda^*_{F_{Y(1)(z|\cdot)}}
      $,
      for all $z\in\R$,
      with values in $\R^{\#\mathfrak{B}}$
      such that
      \begin{align}
        \label{asu:basis:ii:1}
        \frac{1}{N}
        \sum_{i=1}^{N} 
        \left| 
        \inner{B(X_i)}
        {
      \lambda^*_{\varphi^{'}\circ\, 1/\pi}
      }
        -
        \varphi^{'}\left( \frac{1}{\pi(X_i)} \right)
        \right|
        &
        \ 
        \overset{\P}
        {
        \to 
        }
        \ 
        0
        \qquad
        \text{for}
        \ 
        N\to\infty
        \,,
        \intertext{and}
        \label{asu:basis:ii:2}
        \sqrt{N}
        \max_{i\in \left\{ 1,\ldots,N \right\}}
        \sup_{z\in\R}
        \left| 
        \inner{B(X_i)}
        {
      \lambda^*_{F_{Y(1)}(z|\cdot)}
      }
        -
        F_{Y(1)}(z|X_i)
        \right|
        &
        \ 
        \overset{\P}
        {
        \to 
        }
        \ 
        0
        \qquad
        \text{for}
        \ 
        N\to\infty
        \,.
      \end{align}
    \end{enumerate}
\end{assumption}

Next, we give some examples. 
We begin with a histogram basis \cite[§4]{Gyorfi2002}
\begin{example}
We consider a sequence of partitions
$
\left( 
  \mathcal{P}_N
  =
  \left\{ 
    A_{N,1}
    ,
    A_{N,2}
    ,
    \ldots
  \right\}
\right)
$
of $ \R^d $
and define
$ A_N(x) $ to be the cell of $ \mathcal{P}_N $ containing $x$.
We also assume
uniform partition width that decreases to 0, that is,
\begin{gather}
  \label{8881}
  \text{
for all $j\in\mathbb{N}$
it holds
  }
  \qquad
  \lambda^d(A_{N,j})
  \ 
  =
  \ 
  h_N^d
  \ 
  \to
  \ 
  0
  \qquad
  \text{for}
  \ 
  N\to\infty
  \,.
\end{gather}
We define $N$ basis functions $B_k$ of the covariates by
\begin{gather*}
  B_k(x)
  \ 
  :=
  \ 
  \frac{
  \mathbf{1}{\left\{ X_k \in A_N(x) \right\}}
  }{
  \sum_{j=1}^{N} 
  \mathbf{1}{\left\{ X_j \in A_N(x) \right\}}
  }
  \qquad
  \text{for}
  \,
  k\in
  \left\{ 
  1,\ldots,N
  \right\}
  \,,
\end{gather*}
where we keep to the convention $"0/0=0"$.
If at least one $B_k(x)>0$, the basis functions sum to 1. If $B_k(x)=0$ for all basis functions, the sum is 0.
Thus
\begin{align}
  \label{8882}
  \sum_{k=1}^{N}
  B_k(x)
  &
  \ 
  \in
  \ 
  \left\{ 0,1 \right\}
  \qquad
  \text{for all}
  \ 
  x\in\R^d
  \,.
  \notag
  \intertext{
Since $B_i(X_i)>0$, it holds
  }
  \sum_{k=1}^{N}
  B_k(X_i)
  &
  \ 
  =
 \  
  1
  \qquad
  \text{for all}\ 
  i\in
  \left\{ 1,\ldots,N \right\}
  \,.
\end{align}
We check \textit{(i)} in Assumption~\ref{asu:basis}. 
Since 
\begin{gather*}
B_k(x)
\ 
\in
\ 
[0,1]
\qquad
\text{for all}\ 
x\in\R^d
\ 
\text{and for all}\ 
k\in \left\{ 1,\ldots,N \right\}
\,,
\end{gather*}
it holds
\begin{gather*}
  \label{basis_l2_bdd}
  \norm{B(x)}_2^2
  \ 
  =
  \ 
  \sum_{k=1}^{N} 
  B_k(x)
  ^2
  \ 
  \le
  \ 
  \sum_{k=1}^{N} 
  B_k(x)
  \ 
  \in
  \ 
  \left\{ 0,1 \right\}
  \quad
  \text{
    for all
  }
x\in\R^d
\,
\,.
\end{gather*}
Next we check \eqref{asu:basis:ii:1} in \textit{(ii)} in Assumption~\ref{asu:basis}.
To this end we consider
\begin{gather*}
  \lambda^*_{\varphi^{'}\circ\,1/\pi}
  \ 
  :=
  \ 
  \left[ 
    \varphi^{'}
    \left( 
      \frac{1}{\pi(X_1)}
    \right)
    ,
    \ldots
    ,
    \varphi^{'}
    \left( 
      \frac{1}{\pi(X_N)}
    \right)
  \right]
  ^\top
  \,.
\end{gather*}
Since $f$ is continuous and the $X_k$ are random vectors,
$\lambda^*_f$ is also a random vector.
It follows 
\begin{align*}
  &
  \frac{1}{N}
  \sum_{i=1}^{N} 
        \left| 
        \inner{B(X_i)}
        {
  \lambda^*_{\varphi^{'}\circ\,1/\pi}
      }
        -
    \varphi^{'}
    \left( 
      \frac{1}{\pi(X_i)}
    \right)
        \right|
        \\
  &
  \ 
        =
  \ 
  \frac{1}{N}
  \sum_{i=1}^{N} 
        \left| 
        \sum_{k=1}^{N} 
        B_k(X_i)
        \cdot
    \varphi^{'}
    \left( 
      \frac{1}{\pi(X_k)}
    \right)
        -
    \varphi^{'}
    \left( 
      \frac{1}{\pi(X_i)}
    \right)
        \right|
        \\
  &
  \ 
        =
  \ 
  \frac{1}{N}
  \sum_{i=1}^{N} 
        \left| 
        \sum_{k=1}^{N} 
        B_k(X_i)
        \cdot
        \left( 
    \varphi^{'}
    \left( 
      \frac{1}{\pi(X_k)}
    \right)
        -
    \varphi^{'}
    \left( 
      \frac{1}{\pi(X_i)}
    \right)
        \right)
        \right|
\\
  &
  \ 
        =
  \ 
  \frac{1}{N}
  \sum_{i=1}^{N} 
        \left| 
        \sum_{k=1}^{N} 
        B_k(X_i)
        \cdot
        \mathbf{1}
        \left\{ X_k\in A_N(X_i) \right\}
        \cdot
        \left( 
    \varphi^{'}
    \left( 
      \frac{1}{\pi(X_k)}
    \right)
        -
    \varphi^{'}
    \left( 
      \frac{1}{\pi(X_i)}
    \right)
    \right)
        \right|
        \\
  &
  \ 
        \le
  \ 
  \frac{1}{N}
  \sum_{i=1}^{N} 
        \sum_{k=1}^{N} 
        B_k(X_i)
        \cdot
        \omega
        ( 
        \varphi^{'} 
        \circ
        (x\mapsto 1/x)
        \circ
        \pi
        ,
        \lambda^d(A_N(X_i))
        )
        \\
  &
  \ 
        =
  \ 
        \omega
        ( 
        \varphi^{'} 
        \circ
        (x\mapsto 1/x)
        \circ
        \pi
        ,
        h_N^d
        )
        \ 
        \to
        \ 
        0
        \qquad
        \text{for}
        \ 
        N\to\infty
        \,,
\end{align*}
where $\omega$ is the modulus of continuity.
The second equality is due to \eqref{8881}, the third equality follows from the definition of the basis functions, and 
the inequality follows from \eqref{8881}, the convexity of the absolute value and the continuity of $f$.
The convergence in due to \eqref{8882}.

Next, we check \eqref{asu:basis:ii:2} in \textit{(ii)} in Assumption~\ref{asu:basis}.
To this end, we consider
\begin{gather*}
  \lambda^*_{F_{Y(1)}(z|\cdot)}
  \ 
  :=
  \ 
  \left[ 
    F_{Y(1)}(z|X_1)
    ,
    \ldots
    ,
    F_{Y(1)}(z|X_N)
  \right]
  ^\top
  \qquad
  \text{for}\ 
  z\in\R
  \,.
\end{gather*}
To obtain the desired convergence, we need to assume some regularity for 
$
    F_{Y(1)}(z|\cdot)
$.
Let it be continuity for the moment.
With similar arguments as before we get
    \begin{align*}
      &
         \sqrt{N}
        \max_{i\in \left\{ 1,\ldots,N \right\}}
        \sup_{z\in\R}
        \left| 
        \inner{B(X_i)}
        {
      \lambda^*_{F_{Y(1)}(z|\cdot)}
      }
        -
        F_{Y(1)}(z|X_i)
        \right|
        \\
        &
        \ 
        =
        \ 
         \sqrt{N}
        \sup_{z\in\R}
        \omega
        \left( 
        F_{Y(1)}(z|\cdot)
        ,
        h_N^d
        \right)
    \end{align*}
    Thus, if the interplay of $F_{Y(1)}$ and $h_N$ is sufficiently good we get convergence.
    The most abstract assumption would be
    \begin{gather}
      \label{6674}
         \sqrt{N}
        \sup_{z\in\R}
        \omega
        \left( 
        F_{Y(1)}(z|\cdot)
        ,
        h_N^d
        \right)
        \to
        0
        \,.
    \end{gather}
    To be more concrete, if, for example,
$
    F_{Y(1)}(z|\cdot)
    $
    is $\alpha$-Hölder continuous for all $z\in\R$ with 
    $\alpha\in(0,1)$ and $\sqrt{N}h_N^{d\cdot \alpha}\to 0$ it holds
    \eqref{6674}.

\end{example}

%We continue with kernel bases \cite[§5]{Gyorfi2002}.
%\begin{example}
%Let
%$
%  K
%  \colon
%  \R^d
%  \to
%  [0,\infty)
%$
%be a measurable function with $K(0)>0$. 
%We call $K$ a kernel function.
%We assume that there exists
%$R>0$ such that 
%\begin{gather}
%  \label{8884}
%  K(x)
%  \ 
%  \le
%  \ 
%  \mathbf{1}\left\{ \norm{x}_2\le R\right\}
%  \,,
%\end{gather}
%that is, the kernel function has compact support.
%For a decreasing sequence (the bandwith of the kernel) $(h_N)\in (0,\infty)$ with $h_N\to 0$ for $N\to\infty$,
%we define $N$ basis functions $B_k$ of the covariates by
%\begin{gather*}
%  B_k(x)
%  \ 
%  :=
%  \ 
%  \frac{
%    K \left( \frac{\norm{x-X_k}_2}{h_N} \right)
%  }{
%  \sum_{j=1}^{N} 
%    K \left( \frac{\norm{x-X_j}_2}{h_N} \right)
%  }
%  \qquad
%  \text{for}
%  \,
%  k\in
%  \left\{ 
%  1,\ldots,N
%  \right\}
%  \,,
%\end{gather*}
%where (again) we keep to the convention $"0/0=0"$.
%Since $K(0)>0$, as in the previous example, it follows
%\begin{align}
%  \label{8883}
%  \sum_{k=1}^{N}
%  B_k(x)
%  &
%  \ 
%  \in
%  \ 
%  \left\{ 0,1 \right\}
%  \qquad
%  \text{for all}
%  \ 
%  x\in\R^d
%  \,,
%  \notag
%  \\
%  \sum_{k=1}^{N}
%  B_k(X_i)
%  &
%  \ 
%  =
% \  
%  1
%  \qquad
%  \text{for all}\ 
%  i\in
%  \left\{ 1,\ldots,N \right\}
%  \intertext{and}
%  \notag
%  \norm{B(x)}_2^2
%  &
%  \ 
%  \in
%  \ 
%  \left\{ 0,1 \right\}
%  \qquad
%  \text{for all}
%  \ 
%  x\in\R^d
%  \,.
%\end{align}
%This verifies the first condition of Assumption~\ref{asu:basis}.
%To verify the second condition, note, that
%\begin{align}
%  \label{8885}
%  \notag
%  K(x)
%  &
%  \ 
%  =
%  \ 
%  \mathbf{1}\left\{ \norm{x}_2\le R \right\} 
%  K(x)
%  \qquad
%  \text{for all}
%  \ 
%  x\in\R^d
%  \intertext{and thus}
%  B_k(x)
%  &
%  \ 
%  =
%  \ 
%  B_k(x)
%  \cdot
%  \mathbf{1}\left\{ \norm{x-X_k}_2\le R\cdot h_N \right\} 
%  \qquad
%  \text{for all}
%  \ 
%  x\in\R^d
%  \ 
%  \text{and for all}
%  \ 
%  k\in \left\{ 1,\ldots,N \right\}
%  \,.
%\end{align}
%Let 
%$f$ be a continuous function
%and consider (as in the previous example)
%\begin{gather*}
%  \lambda^*_f
%  \ 
%  :=
%  \ 
%  [f(X_1),\ldots,f(X_N)]^\top
%  \,.
%\end{gather*}
%It follows
%\begin{align*}
%  &
%        \left| 
%        \inner{B(X_i)}{\lambda^*_f}
%        -
%        f(X_i)
%        \right|
%        \\
%  &
%  \ 
%        =
%  \ 
%        \left| 
%        \sum_{k=1}^{N} 
%        B_k(X_i)
%        \cdot
%        f(X_k)
%        -
%        f(X_i)
%        \right|
%        \ 
%        =
%        \ 
%        \left| 
%        \sum_{k=1}^{N} 
%        B_k(X_i)
%        \cdot
%        \left( 
%        f(X_k)
%        -
%        f(X_i)
%        \right)
%        \right|
%\\
%  &
%  \ 
%        =
%  \ 
%        \left| 
%        \sum_{k=1}^{N} 
%        B_k(X_i)
%        \cdot
%  \mathbf{1}\left\{ \norm{x-X_k}_2\le R\cdot h_N \right\} 
%        \cdot
%        \left( 
%        f(X_k)
%        -
%        f(X_i)
%        \right)
%        \right|
%        \\
%  &
%  \ 
%        \le
%  \ 
%        \sum_{k=1}^{N} 
%        B_k(X_i)
%        \cdot
%        \omega
%        ( 
%        f,
%        R\cdot h_N
%        )
%        \\
%  &
%  \ 
%        =
%  \ 
%        \omega
%        ( 
%        f,
%        R\cdot h_N
%        )
%        \to
%        0
%        \qquad
%        \text{for}
%        \ 
%        N\to\infty
%        \,,
%\end{align*}
%The convergence in due to $h_N\to 0$ for $N\to \infty$. 
%\end{example}
We now have a large class of possible basis function. We may proceed.

\subsection{Pseudo Solutions}
There are two observations that help achieving our goal.
First, if we solve Problem~\ref{dual} on a compact search space, there always exists a measurable solution (if the objective function is Caratheodory).
Second (with hindsight), we want optimal solutions to estimate an oracle parameter $\lambda^*$ well. 
To this end, we will define a (random) compact parameter space that always contains the oracle parameter $\lambda^*$ and restrict Problem~\ref{dual} to it.
Then Theorem~\ref{th:argmax} gives us (under weak measurability conditions) a measurable solution of the restricted Problem~\ref{dual}.
In some cases, this is only a local solution. It becomes global if it is in the interior of the compact parameter space.
Thus we call it the restricted, or the pseudo solution.

Consider 
the random vector (the oracle parameter)
\begin{align*}
  \lambda^*
  \ 
  :=
  \ 
  \left(
  0_{N},0,
  \left[
  \varphi^{'}
  \left(
  \frac
  {1}
  {\pi(X_i)}
  \right)
\right]_{i\in \left\{
  1,\ldots,N
\right\}}
  \right)
\end{align*}
with values in 
$\R^N_{\ge 0}\times \R\times \R^N$.
Clearly,  
\begin{align*}
 \lambda^*
 \qquad
 \text{is}
 \qquad
 \left(
( 
\sigma
\left( 
  (T_i,X_i)_{i\in 
\left\{ 1,\ldots,N \right\}
  } 
\right)
,
\mathcal{B}
\left( 
\R^N_{\ge 0}\times \R\times \R^N
\right)
\right)
-\text{measurable}
\,.
\end{align*}
Next we define for $N\in\mathbb{N}$ the correspondence  
\begin{align*}
  \Theta_N\colon
  &
  \left( 
  \Omega,
\sigma
\left( 
  (T_i,X_i)_{i\in 
\left\{ 1,\ldots,N \right\}
  } 
\right)
  \right)
  \ 
\to
  \ 
\R^N_{\ge 0}\times \R\times \R^N
\\
&
\omega
  \ 
\mapsto
  \ 
\left\{ 
  \left( 
  \rho,\lambda_0,\lambda
  \right)
  \ 
  \in
  \ 
\R^N_{\ge 0}\times \R\times \R^N
  \ 
\colon
  \ 
\norm{
  \left( 
  \rho,\lambda_0,\lambda
  \right)
-\lambda^*}_2\le 1
\right\}
\,.
\end{align*}

\begin{lemma}
  \label{Theta_maes}
  For all $N\in\mathbb{N}$ it holds that 
  $\Theta_N$ is a weakly-measurable correspondence  with non-empty compact values.
\end{lemma}
\begin{proof}
  Let $N\in\mathbb{N}$. 
  Since 
  \begin{align*}
  \Theta_N(\omega)
  \ 
  \text{
  is the closed unit ball
  in
  }
  \ 
\R^N_{\ge 0}\times \R\times \R^N
\quad 
\text{around}
\quad 
  \lambda^*(\omega) 
  \,,
  \end{align*}
  it is non-empty and compact for all $\omega\in\Omega$.
  To prove that $\Theta_N$ is weakly-measurable,
  let 
  $
  O\subset
\R^N_{\ge 0}\times \R\times \R^N
  $
  be an open set and note, that
  \begin{align*}
    \left\{ 
      \Theta_N\cap O \neq \emptyset
    \right\}
    &
    \ 
    =
    \ 
    \left\{ 
      \exists
  \left( 
  \rho,\lambda_0,\lambda
  \right)
  \ 
  \in
  \ 
  O
\ 
\colon
\ 
\norm{
  \left( 
  \rho,\lambda_0,\lambda
  \right)
-\lambda^*}_2\le 1
    \right\}
    \\
    &
    \ 
    =
    \ 
    \left\{ 
      \exists
  \left( 
  \rho,\lambda_0,\lambda
  \right)
  \ 
  \in
  \ 
  O
  \cap
  \left( 
  \mathbb{Q}^N_{\ge 0}\times \mathbb{Q}\times \mathbb{Q}^N
  \right)
\ 
\colon
\ 
\norm{
  \left( 
  \rho,\lambda_0,\lambda
  \right)
-\lambda^*}_2\le 1
    \right\}
    \\
&
    \ 
    =
    \ 
    \bigcup_{
  \left( 
  \rho,\lambda_0,\lambda
  \right)
  \ 
  \in
  \ 
  O
  \cap
  \left( 
  \mathbb{Q}^N_{\ge 0}\times \mathbb{Q}\times \mathbb{Q}^N
  \right)
    }
    \left\{ 
\norm{
  \left( 
  \rho,\lambda_0,\lambda
  \right)
-\lambda^*}_2\le 1
    \right\}
    \,.
  \end{align*}
  Since
  \begin{align*}
    \left\{ 
\norm{
  \left( 
  \rho,\lambda_0,\lambda
  \right)
-\lambda^*}_2\le 1
    \right\}
    \in
\sigma
\left( 
  (T_i,X_i)_{i\in 
\left\{ 1,\ldots,N \right\}
  } 
\right)
\quad
\text{for all}\quad
  \left( 
  \rho,\lambda_0,\lambda
  \right)
  \in
\R^N_{\ge 0}\times \R\times \R^N
\,,
  \end{align*}
  and the union is countable, it follows
  $
    \left\{ 
      \Theta_N\cap O \neq \emptyset
    \right\}
    \in
\sigma
\left( 
  (T_i,X_i)_{i\in 
\left\{ 1,\ldots,N \right\}
  } 
\right)
  $.
  Thus $\Theta_N$ is weakly-measurable.
\end{proof}
We are ready to construct the measurable (pseudo) solution.
\begin{lemma}
  \label{lem:pseud_sol}
  For all $N\in\mathbb{N}$ there exists a (pseudo) solution
  \begin{align*}
    s_N
    \ 
    \colon
   \  
    \Omega
    \ 
    \to
    \ 
  \R^N_{\ge 0}
  \times
  \R
  \times
  \R^{N}
  \end{align*}
  to
  Problem~\ref{dual} restricted to $\Theta_N$ 
  that is
  \begin{align*}
  \left(
    \sigma \left( \left( T_i,X_i \right)_{i\in \left\{ 1,\ldots,N \right\}} \right),\mathcal{B}
  \left(
  \R^N_{\ge 0}
  \times
  \R
  \times
  \R^{N}
  \right)
  \right)
  -\text{measurable}
  \,.
  \end{align*}
  Furthermore,
if $s_N\in \mathrm{int}\, \Theta_N$, then it is the global solution
  $
  (\rho^\dagger,\lambda_0^\dagger,\lambda^\dagger)
  $.
\end{lemma}
\begin{proof}
  By Lemma~\ref{Theta_maes}
  the correspondence $\Theta_N$ satisfies the conditions of
  Theorem~\ref{th:argmax}.
Next, we consider the (random) objective function of (the maximize version of) Problem~\ref{dual}, that is,
  \begin{align*}
    &
  G\colon
  \left(
  \Omega,
\sigma
\left( 
  (T_i,X_i)_{i\in 
\left\{ 1,\ldots,N \right\}
  } 
  \right)
  \right)
  \times
  \left(
  \R^N_{\ge 0}
  \times
  \R
  \times
  \R^{N}
  \right)
  \to
  \R\cup \left\{
    -\infty
  \right\}
  \intertext{with}
    &
  G(\omega,(\rho,\lambda_0,\lambda))
  \ 
  =
  \ 
  -\infty
  \qquad 
  \text{if}\ 
  \rho_i\neq 0\  \text{for some}\ i>n
  \,,
  \\
  \intertext{and}\qquad
  &
  G(\omega,(\rho,\lambda_0,\lambda))
  \\
  &
  \ 
  =
  \ 
  \frac{1}{N}
\sum_{i=1} 
  ^N
  \Big[
    -
  T_i(\omega)
  \cdot
  \varphi^*
  \!
  \left( 
    \rho_i
    +
\lambda_0
+
\inner
{B(X_i)(\omega)}
{
\lambda
}
  \right)
  \ 
  +
  \ 
\lambda_0
+
\inner
{B(X_i)(\omega)}
{
\lambda
}
\Big]
  \ 
-
\ 
\inner
{\delta(\omega)}
{
  |\lambda|
}
\,.
  \end{align*}
  Clearly, the (random) objective function $G$  is a Caratheodory function
  . This follows from the (assumed) continuity of $\varphi^*$ and the measurability 
  of all random variables included.
  Since $G$ is also strictly concave, it has a unique argmax $s_N$  in $\Theta_N$. 
  By Theorem~\ref{th:argmax} this is $
  \left(
    \sigma
    \left( \left( T_i,X_i \right)_{i\in \left\{ 1,\ldots,N \right\}}\right)
    ,\mathcal{B}
  \left(
  \R^N_{\ge 0}
  \times
  \R
  \times
  \R^{N}
  \right)
  \right)
$-measurable.
Furthermore, by the strict concavity of $G$,
if $s_N\in \mathrm{int}\,\Theta_N$, then it's the global optimal solution. 
\end{proof}
We will prove later, that with probability going to 1 the latter will be the case.
Furthermore, we will prove, that the (measurable) solution converges to the random variable in probability.


Before we define the weights function we specify the basis $\mathfrak{B}$.
\subsection{basis}
Let $
\left(
\mathcal{P}_N
\right)
$
denote a sequence of countable, $\mathcal{B}$-measurable partitions 
\begin{align*}
\mathcal{P}_N= \left\{
  A_{N,1},
  A_{N,2},
  \ldots
\right\}
\subset \mathcal{B}(\R^d)
\end{align*}
of $\R^d$, that is, 
\begin{align*}
  A_{N,i}\cap A_{N,j}=\emptyset
  \qquad
  \text{if}\ i\neq j
  \qquad
  \text{and}
  \qquad 
  \bigcap_{i\in\mathbb{N}}A_{N,i}
  \ 
  =
  \ 
  \R^d
  \,.
\end{align*}
We define
$ A_N(x) $ to be the cell of $ \mathcal{P}_N $ containing $x$, that is,
\begin{align*}
  A_N
  \colon
  \R^d 
  \ 
  \twoheadrightarrow 
  \ 
  \R^d  
  \,,\qquad
  x
  \ 
  \mapsto
  \ 
  A_N(x)
  \,,
\end{align*}
where $A_N(x)$ is the only cell containing $x$. 

Next, we define the (empirical) basis vector
\begin{align*}
  B\colon
  \R^d\times \R^{d\cdot N}
  \to
  \R
  \,,
  \qquad
  (x,(x_1,\ldots,x_N))\mapsto
  \frac
  {
    \left[
    \mathbf{1}
    _{
      A_N(x)
    }
    (x_k)
    \right]
    _{k\in \left\{
        1,\ldots,N
    \right\}}
  }
  {
    \sum_{j=1}^N
    \mathbf{1}
    _{
      A_N(x)
    }
    (x_j)
    }
  \,,
\end{align*}
where we keep to the convention $"0/0=0"$.
\begin{lemma}
  \quad
  \begin{enumerate}[label=(\roman*)]
\item
  $B(\cdot,(X_1,\ldots,X_N))(\omega)$ is 
  $\left(
    \mathcal{B}(\R^d),\mathcal{B}(\R^N)
  \right)$-measurable
  and
  constant on each cell $A_N\in\mathcal{P}_N$
  for all $\omega\in\Omega$. 
\item
  $B(X,(X_1,\ldots,X_N))$ is $\left(
    \mathcal{A},\mathcal{B}(\R^N)
  \right)$-measurable. 
  \end{enumerate}
\end{lemma}

\begin{proof}
Consider
for $k\in \left\{
  1,\ldots,N
\right\}$
and $\omega\in\Omega$
the indicator function
\begin{align*}
  \mathbf{1}
  _
  {A_N(X_k(\omega))}
  \colon \R^d\to \left\{
    0,1
  \right\}
  \,.
\end{align*}
Since 
$
  {A_N(X_k(\omega))}
  \in\mathcal{B}(\R^d)
$
this is a 
  $\left(
    \mathcal{B}(\R^d),\mathcal{B}(\R^N)
  \right)$-measurable
  function.
  Since it is 1 if $
  x\in
  {A_N(X_k(\omega))}
  $
  and 0 else, it is also constant on each cell. It follows statement (i).
Since
\begin{align*}
  \mathbf{1}
  _
  {A_N(X_k(\omega))}(X(\omega))
  \ 
  =
  \ 
  \mathbf{1}
  \bigcup_{i\in\mathbb{N}}
  \left\{
    X,X_k \in A_{N,i}
  \right\}
  (\omega)
  \qquad
  \text{for all}\ 
  \omega\in\Omega
  \,,
\end{align*}
and
$
  \bigcup_{i\in\mathbb{N}}
  \left\{
    X,X_k \in A_{N,i}
  \right\}
  \in\mathcal{A}
$
it follows statement (ii).
\end{proof}

\subsection{Weights Function}
Theorem~\ref{dual_solution_th} tells us that if an optimal solution
$
(\rho^\dagger,\lambda_0^\dagger,\lambda^\dagger)
$
to Problem~\ref{dual} exists,
then the unique optimal solutions to Problem~\ref{bw:1:primal} are 
\begin{gather*}
  w^\dagger_i
  \ 
  :=
  \ 
  (
  \varphi^{'}
  )^{-1}
  \left(
    \rho^\dagger_i
  \ 
    +
  \ 
\lambda_0^\dagger
  \ 
+
  \ 
\inner
{B(X_i)}
{
\lambda^{\dagger}
}
  \right)
  \qquad
  \text{for all}\ 
  i\in
  \left\{ 1,\ldots,n \right\}
  \,.
\end{gather*}

This point of view is sufficient from a practical point of view
(when we are interested in actually computing the optimal weights).

For the subsequent analysis we need to view the weights as random quantities. To this end, we shall introduce the weights process. 

\begin{definition}
  \label{def:weights_function}
  We define the weights process to be a stochastic process indexed 
  over
  $
  \mathcal{X}
  \,
  \times
  \,
  \R^N_{\ge 0}
  \,
  \times
  \,
  \R
  \times
  \,
  \R^{N}
  $
  with values in $\R^N$ such that
\begin{align*}
  w
  \left( 
  x
  ,
  \rho
  ,
  \lambda_0
  ,
  \lambda
  \right)
  \ 
  :=
  \ 
  \left[ 
    (
    \varphi^{'}
    )^{-1}
    \left( 
     \rho_i 
     +
      \lambda_0
      +
      \inner{B(x)}{\lambda}
    \right)
  \right]_{i\in \left\{ 1, \ldots,N \right\}}
    \,.
\end{align*}
\end{definition}
By the measurability of the basis functions (or the basis processes),
the weights processes is also measurable 
(the quantities 
$
  w
  \left( 
  x
  ,
  \rho
  ,
  \lambda_0
  ,
  \lambda
  \right)
$
are random variables).
The next question is if there are (plausible) assumptions such that the weights processes at a random parameter
$
(\rho^\dagger,\lambda_0^\dagger,\lambda^\dagger)
$
is measurable.
To answer this in a good way we need the argmax measurability theorem \cite[Theorem~18.19]{Aliprantis2007}.
\subsection{Argmax Measurability Theorem}
We follow \cite{Aliprantis2007}
A \textbf{correspondence}$ \psi$ from a set $X$ to a set $Y$ assigns to each $x\in X$ a subset $\psi(x)\subset Y$.
To clarify that we map $x$ to a set, we use the double arrow, that is,
$
  \psi
  \colon
  X
  \twoheadrightarrow
  Y
$.
  Let 
  $(S,\Sigma_S)$ be a measurable space and $X$ a topological space.
  We say, that a correspondence 
  $
  \psi
  \colon
  S
  \twoheadrightarrow
  X
  $
  is 
  \textbf{
  weakly measurable
  },
  if
  \begin{gather*}
    \left\{ 
      s\in S
      \ 
      |
      \ 
      \psi(s)
      \cap
      O
      \neq
      \emptyset
    \right\}
    \in
    \Sigma_S
    \qquad
    \text{for all open subsets}
    \ 
    O\subset X
    \,.
  \end{gather*}

A selector from a correspondence $\psi\colon X\twoheadrightarrow Y$ is a function $s\colon X\to Y$ that satisfies 
$
s(x)\in\psi(x)
$
for all $x\in X$.

\begin{definition}
  Let 
  $(S,\Sigma_S)$ be a measurable space, and let $X$ and $Y$  be topological space.
  A function 
  $f\colon S\times X \to Y$
  is a \textbf{Caratheodory function} if
  \begin{align*}
    f(\cdot,x)
    &
    \colon
    S\to Y
    \qquad
    \text{is}\ 
    (\Sigma_{S},\mathcal{B}(Y))-measurable
    \ 
    \text{for all}
    \ 
    x\in X
    \,,
    \intertext{and}
    f(s,\cdot)
    &
    \colon
    X\to Y
    \qquad
    \text{is continuous for all}\ 
    s\in S
    \,.
  \end{align*}
\end{definition}
\begin{theorem}
  Let $X$ be a separable metrizable space and
  $
  (S,\Sigma_S)
  $
  a measurable space.
  Let $\psi\colon S \twoheadrightarrow X$ be a weakly measurable correspondence with non-empty compact values, and suppose
  $f\colon S\times X \to \R$
  is a Caratheodory function. Define the value function 
  $m\colon S\to \R$ by
  \begin{gather*}
    m(s):=\max_{\psi(s)}f(s,x)
    \,,
  \end{gather*}
  and the correspondence 
  $mu\colon S\twoheadrightarrow X$ of maximizers by
  \begin{gather*}
    \mu(s):= \left\{ 
      x\in \psi(s)
      |
      f(s,x)=m(s)
    \right\}
    \,.
  \end{gather*}
  Then the value function $m$ is measurable, 
  the argmax correspondence $\mu$ has non-empty and compact values,
  is measurable and admits a measurable selector.
\end{theorem}
\begin{proof}
  \cite[Theorem~18.19]{Aliprantis2007}
\end{proof}

Our goal is to find (plausible) assumptions under which we can measurably select optimal solutions of Problem~\ref{dual}.

\begin{assumption}
  There exists $\underline{N}\in\mathbb{N}$ such that 
  for all $N\ge \underline{N}$ 
  there exists a compact and deterministic parameter space
  $
  \Theta_N
  \subset
  \R^N_{\ge 0}
  \times
  \R
  \times
  \R^{N}
  $
  such that for all data sets $D_N$
  a sequence of optimal solution
  $(\rho^\dagger,\lambda_0^\dagger,\lambda^\dagger)$
  satisfying the Lemma
exist and it holds
\begin{gather*}
  (
  \rho^\dagger,
  0_{N-n},
  \lambda_0^\dagger,\lambda^\dagger)
  \in
  \Theta_N
  \,.
\end{gather*}
\end{assumption}

With this assumption we can construct the (constant) correspondence
\begin{align*}
  \psi
  \colon
  (\Omega,\mathcal{A},\P)
  \ 
  \twoheadrightarrow
  \ 
  \R^N_{\ge 0}
  \times
  \R
  \times
  \R^{N}
  \,,
  \qquad
  \omega
  \ 
  \to
  \ 
  \Theta_N
  \,.
\end{align*}
This is weakly measurable, because $\Theta_N$ is deterministic and thus
$\Theta_N\cap O$ is either true or false, that is,
  \begin{gather*}
    \left\{ 
      \omega\in \Omega
      \ 
      |
      \ 
      \psi(\omega)
      \cap
      O
      \neq
      \emptyset
    \right\}
    \ 
    \in
    \ 
    \left\{ \Omega,\emptyset \right\}
    \subset
    \Sigma_S
    \qquad
    \text{for all open subsets}
    \ 
    O
    \subset
  \R^N_{\ge 0}
  \times
  \R
  \times
  \R^{N}
    \,.
  \end{gather*}
  Furthermore the objective function is a Caratheodory function.
  Thus, by Theorem there exists the argmax correspondence  of Problem~\ref{dual} and a measurable selector that selects
  $(\rho^\dagger,0_{N-n},\lambda_0^\dagger,\lambda^\dagger)$ 
  of Assumption.
  With this, we can define the optimal weights processes
  \begin{gather*}
    w^\dagger(x)
    \ 
    :=
    \ 
    w
    \left( 
    x,\rho^\dagger,0_{N-n},\lambda_0^\dagger,\lambda^\dagger
    \right)
    \qquad
    \text{indexed over}\ 
    x\in\mathcal{X}\,.
  \end{gather*}
  Due to the definition of the (general) weights processes and the measurability of the argmax selector, the random variables $w^\dagger(x)$ are measurable for all $x\in\mathcal{X}$.
  To end the measurability discussion, note that $w^\dagger(X)$ is a random variable.


\begin{lemma}
  \label{weights_l_inf}
  Let Assumption~\ref{asu:existence_sol}, 
  Assumption~\ref{asu:basis},
  and Assumption~\ref{asu:objective_f} hold true.
  Then it holds 
  \begin{gather*}
    w_i^\dagger(X)
  \ 
  \in
  \ 
  L^\infty(\P)
  \qquad
  \text{for all}
  \ 
  i\in \left\{ 1,\ldots,N \right\}
  \ 
  \text{
    and
  for all 
  }\ 
  N\ge \underline{N}
  \,.
  \end{gather*}
\end{lemma}
\begin{proof}
  Let $N\ge \underline{N}$.
  By Assumption~\ref{asu:basis} it holds 
  $
    \text{for all}\ 
  $ 
  \begin{gather}
    \label{3967}
    \left| 
     \rho_i 
     +
      \lambda_0
      +
      \inner{B(x)}{\lambda}
    \right|
    \ 
    \lesssim
    \ 
    \norm{(\rho,\lambda_0,\lambda)}_2
    \,.
  \end{gather}
  $
    \text{for all}
    \ 
    x\in\mathcal{X}
    \
    \text{and for all}\ 
    (\rho,\lambda,\lambda_0)
    \in
  \R^N_{\ge 0}
  \times
  \R
  \times
  \R^{N}
  $.
  By Assumption~\ref{asu:existence_sol}, 
  there exists a (compact) parameter space
  $\Theta_N$ around the origin, with $\mathrm{diam}\,  \Theta_N<\infty$, 
  such that for all data sets $D_N$ it holds  $(\rho^\dagger,\lambda^\dagger,\lambda_0^\dagger)\in\Theta_N$.
  By Assumption~\ref{asu:objective_f}, $(\varphi^{'})^{-1}$ is non-decreasing and continuous. Thus
    by \eqref{3967} it holds
  \begin{align*}
    \left| 
    w_i^\dagger(x)
    \right|
    \ 
    \lesssim
    \ 
    \left| 
    (\varphi^{'})^{-1}
    \left( 
      \,
      -
    \norm{(\rho^\dagger,\lambda_0^\dagger,\lambda^\dagger)}_2
    \right)
    \right|
    \ 
    +
    \ 
    \left| 
    (\varphi^{'})^{-1}
    \left( 
      \,
    \norm{(\rho^\dagger,\lambda_0^\dagger,\lambda^\dagger)}_2
    \right)
    \right|
    \ 
    \ 
    <
    \ 
    \infty
  \end{align*}
  for all $x\in\mathcal{X}$
  and all $i\in \left\{ 1,\ldots,n \right\}$
  .
\end{proof}
\begin{lemma}
  \label{w.Z=0}
  Let Assumption~\ref{aa:assumption:2} and Assumption~\ref{aa:assumption:3} hold true.
  Furthermore, 
  let
  $N\ge\underline{N}$, and
  let
  $Z\in L^1(\P)$
  be a random variable that is independent of $D_N$ 
  with
  $
\E
\left[
  Z
  \,
  |
  \, 
  X
\right]
= 0
  $
  almost surely.
  It holds
  \begin{gather*}
  \E
  \left[
  w_i(X,\rho^\dagger,\lambda^\dagger,\lambda_0^\dagger)
  \cdot Z
  \right]
  =0
  \qquad
  \text{for all}
  \ 
  i\in \left\{ 1,\ldots,n \right\}
  \,.
  \end{gather*}
\end{lemma}
\begin{proof}
  Let
  $N\ge\underline{N}$.
  We write
  \begin{gather*}
  w_i(X,\rho^\dagger,\lambda^\dagger,\lambda_0^\dagger)
  \ 
  =
  \ 
  w^\dagger(X)
  \end{gather*}
  and ignore the index $i$.
  By Lemma~\ref{weights_l_inf} and 
  $Z\in L^1(\P)$
  it holds
  \begin{gather}
    \label{9876}
    \norm{
  w^\dagger(X)\cdot Z
    }_{L^1(\P)}
    \ 
  \le
    \ 
  \norm{w^\dagger(X)}_{L^\infty(\P)}
  \norm{Z}_{L^1(\P)}
  \ 
  <
  \ 
  \infty
  \,.
  \end{gather}
  By 
  \eqref{9876},
  $Z\perp D_N$
  and
  $
\E
\left[
  Z
  \,
  |
  \, 
  X
\right]
= 0
  $
  almost surely
  it holds 
  \begin{align*}
    \E
  \left[
  w^\dagger(X)
  \cdot
  Z
  \,
  |
  \,
  D_N,X
  \right]
  &
  \ 
  =
  \ 
  w^\dagger(X)
  \cdot
  \E
  \left[
  Z
  \,
  |
  \,
  D_N,X
  \right]
  \\
  &
  \ 
  =
  \ 
  w^\dagger(X)
  \cdot
  \E
  \left[
  Z
  \,
  |
  \,
  X
  \right]
  \
  =
  \ 
  0
  \qquad
  \text{almost surely.}
  \end{align*}
  Thus
  \begin{gather*}
    \E
    \left[
  w^\dagger(X)
  \cdot
  Z
  \,
    \right]
    \ 
    =
    \ 
    \E
    \left[
 \E
  \left[
  w^\dagger(X)
  \cdot
  Z
  \,
  |
  \,
  D_N,X
  \right]
    \right]
    \ 
    =
    \ 
    0
    \,.
     \end{gather*}
\end{proof}




%  \begin{ftheorem}
  \label{dual_solution_th}
  The dual of 
  Problem~\ref{bw:1:primal}
  is the unconstrained optimization problem 
\begin{gather*}
  \underset
  {\lambda_0,\ldots,\lambda_{N}\in \R}
  {
    \mathrm{minimize}
  }
  \quad
  \frac{1}{N}
\sum_{i=1} 
  ^N
  \left[ 
    \,
  T_i
  \cdot
  f^*
  \!
  \left( 
\lambda_0
+
\inner
{B(X_i)}
{\lambda}
  \right)
  \ 
-
\ 
  \left( 
\lambda_0
+
\inner
{B(X_i)}
{\lambda}
  \right)
  \,
  \right]
  \ 
+
\ 
\inner
{\delta}
{|\lambda|}
  \,.
\end{gather*}
  where
  \begin{gather*}
  f^*
  \,
  \colon
  \, 
  \R
  \ 
  \to
  \ 
  \R
  \,
  ,
  \qquad 
  x^*
  \ 
  \mapsto
  \ 
    x^*
    \!
    \cdot
    (f^{'})^{-1}(x^*)
  \ 
    -
  \ 
    f
    \left( 
      (f^{'})^{-1}(x^*)
    \right)
  \end{gather*}
  is the Legendre transformation of $f$,
  the vector
  $
    B(X_i)
    =
    \left[ 
      B_1(X_i)
      ,
      \ldots
      ,
      B_n(X_i)
    \right]
    ^\top
  $
  denotes 
  the $N$ basis functions of the covariates 
  of unit $i\in \left\{ 1, \ldots, N \right\}$
  and
  $
    \left| \lambda \right|
    =
    \left[ 
      \left| \lambda_1 \right|
      ,
      \ldots
      ,
      \left| \lambda_N \right|
    \right]
    ^\top
    ,
  $
  where $\left| \,\cdot\, \right|$
  is the absolute value of a real-valued scalar.
  Moreover,
  if $\lambda^\dagger$
  is an optimal solution of the above problem 
  then the optimal solution to problem
  Problem~\ref{bw:1:primal}
  is given by
  \begin{gather*}
    w_i^\dagger
    \ 
    =
    \ 
    (f^{'})^{-1}
    \left( 
      \inner
      {B(X_i)}
      {\lambda^\dagger}
      +
      \lambda_0^\dagger
    \right)
    \qquad
    \text{for}\ 
    i\in \left\{ 1\ldots,n \right\}
    \,.
  \end{gather*}
\end{ftheorem}

\newpage
We introduce some more notation. Let $\mathbf{I}_n$ be the $n$-dimensional unit matrix, $\mathrm{0}_n$ and $\mathrm{1}_n$ the $n$-dimensional vectors containing only zeros or ones.
Also we define the vector of basis functions of the covariates of unit 
$i\in \left\{ 1,\ldots,N \right\}$ to be
\begin{gather*}
     B(X_i)
    \ 
    :=
    \ 
    \left[ 
      B_1(X_i)
      ,
      \ldots
      ,
      B_N(X_i)
    \right]
    ^\top
    \in
    \R^N
    \,.
\end{gather*}
Furthermore, we define the matrix of basis functions \textbf{for the treated} to be
\begin{gather*}
    \mathbf{B}(\mathbf{X})
    \ 
    :=
    \ 
    \begin{bmatrix}
      B(X_1), \ldots, B(X_n)
    \end{bmatrix}
    \in\R^{N\times n}
    \,.
\end{gather*}
Note, that these are random quantities and that the size of $\mathbf{B}(\mathbf{X})$ depends on the random size $n\in\mathbb{N}$ of the treatment group in the sample.
\begin{lemma}
  \label{matrix_notation}
  A matrix formulation of Problem~\ref{bw:1:primal} is 
\begin{align}
  \label{cv:ts:primal}
  %%%% objective %%%%
    &\underset{w \in \R^n}
    {\mathrm{minimize}}
    &&\qquad\qquad
    \varphi(w)
    &&&
    \\
    %%%% Ax >= b %%%%
    \nonumber
    &\mathrm{subject}\ \mathrm{to} 
    &&\qquad\qquad
    \mathbf{U}w
    \ 
    \ge
    \ 
    d
    \,,
    \\
    \nonumber
    &
    &&\qquad\qquad
    \mathbf{A}w
    \ 
    =
    \ 
    a
    \,,
\end{align}
with objective function
\begin{gather*}
  \varphi
  \ 
  :
  \ 
  \R^n
  \to
  \ 
  \overline{\R}
  \,
  ,
  \qquad
  [w_1,\ldots,w_n]^\top
  \ 
  \mapsto
  \ 
  \sum_{i=1}^n f(w_i)
  \,,
\end{gather*}
inequality matrix and vector
\begin{alignat*}{2}
    \mathbf{U}
    &
    \ 
    :=
    \ 
    \begin{bmatrix}
      \mathbf{I}_n
      \\
      \pm\,\mathbf{B}(\mathbf{X})
    \end{bmatrix}
    \in
    \R^{(n+2N)\times n}
        \qquad
    &&
d
    \ 
    :=
    \ 
    \begin{bmatrix}
      0_n
      \\
      -N\delta 
      \ 
      \pm\ 
      \sum_{i = 1}^{N} B(X_i)
    \end{bmatrix}
    \in
    \R^{(n+2N)}
    \,,
    \intertext{and equality matrix and vector}
    \mathbf{A}
    &
    \ 
    :=
    \ 
      \mathrm{1}_n
      ^\top
      \in\R^{1\times n}
      \qquad
    &&
    a
  \ 
    :=
    \ 
    N
    \in\mathbb{N}
    \,.
\end{alignat*}
\end{lemma}
\begin{proof}
  A simple calculation yields the result. We omit the details.
\end{proof}
\begin{remark}
  The inequality constraints of
  Lemma~\ref{matrix_notation} differ from its counterpart
  \cite[Proof of Lemma~1]{Wang2019}.
  We don't transform the variable $w$, but shift to $d$ what prevents us from keeping $w$.
  Note, that the choice of
  \cite[Proof of Lemma~1]{Wang2019} leads to a mistake on page 21.
  The mistake is most obvious in the second display, where the first implication follows from dividing by 0. I discussed this with the authors and proposed a version of Lemma\ref{matrix_notation} to solve the problem. I think it's best not to transform variables, because the mistake comes from (wrongly) calculating the convex conjugate of the (more complicated) transformed version of the objective function. The subsequent analysis even simplifies with my version.
\end{remark}


  \newpage
%  \section{Dual Formulation}
%  \begin{ftheorem*}
  The dual of 
  Problem~\ref{bw:1:primal}
  is the unconstrained optimization problem 
  \begin{gather*}
    \underset{\lambda \in \R^n}{\mathrm{minimize}}
    \qquad
    \frac{1}{n}
    \sum_{i = 1}^{n} 
    T_i 
    \cdot
    f^*
    (
      m_n(\lambda|X_i)
    )
    -
      m_n(\lambda|X_i)
    \ 
    +
    \ 
    \inner{\delta}{\left| \lambda \right|}
    \,,
  \end{gather*}
  where
  \begin{gather*}
  f^*
  \,
  \colon
  \, 
  \R
  \ 
  \to
  \ 
  \R
  \,
  ,
  \qquad 
  x^*
  \ 
  \mapsto
  \ 
    x^*
    \!
    \cdot
    (f^{'})^{-1}(x^*)
  \ 
    -
  \ 
    f
    \left( 
      (f^{'})^{-1}(x^*)
    \right)
  \end{gather*}
  is the Legendre transformation of $f$,
  the vector
  $
    B(X_i)
    =
    \left[ 
      B_1(X_i)
      ,
      \ldots
      ,
      B_n(X_i)
    \right]
    ^\top
  $
  denotes 
  the $n$ basis functions of the covariates 
  of unit $i\in \left\{ 1, \ldots, n \right\}$
  and
  $
    \left| \lambda \right|
    =
    \left[ 
      \left| \lambda_1 \right|
      ,
      \ldots
      ,
      \left| \lambda_n \right|
    \right]
    ^\top
    ,
  $
  where $\left| \,\cdot\, \right|$
  is the absolute value of a real-valued scalar.
  Moreover,
  if $\lambda^\dagger$
  is an optimal solution of the above problem 
  then
  for all $i$ with $T_i=1$
  \begin{gather*}
    w_i^\dagger
    \ 
    =
    \ 
    (f^{'})^{-1}
    \left( 
      m_n
      (
      \lambda^\dagger
      |
      X_i
      )
    \right)
  \end{gather*}
  is part of the optimal solution to 
  Problem~\ref{bw:1:primal}
  .
\end{ftheorem*}
\subsection*{Plan of proof}
  We bring 
  Problem~\ref{bw:1:primal}
  in the form of ts.
  Then we apply the results of the ts chapter.
  We wait to conclude on the weights.
  Then we eliminate the non-negativity constraints
  on the dual variables leveraging convexity and optimality.
%%%%%%%%%%%%%%%  
%%%% PROOF %%%%
%%%%%%%%%%%%%%%
\begin{proof}

  \subsubsection*{Form}
  We consider the vector of basis functions of the covariates 
  of unit $i\in \left\{ 1, \ldots, n \right\}$, that is,
\begin{align*}
    B(X_i)
    &
    \ 
    :=
    \ 
    \left[ 
      B_1(X_i)
      ,
      \ldots
      ,
      B_n(X_i)
    \right]
    ^\top
    \,,
    \intertext{the constraints vector}
d
    &
    \ 
    :=
    \ 
    \begin{bmatrix}
      0_n
      \\
      \pm n
      \\
      -n\cdot\delta 
      \pm\,
      \sum_{i = 1}^{n} B_k(X_i)
    \end{bmatrix}
    \,,
    \intertext{
  the matrix of the basis functions with indicators of treatment
  }
    T\mathbf{B}(\mathbf{X})
    &
    \ 
    :=
    \ 
    \begin{bmatrix}
      T_1B(X_1), \ldots, T_nB(X_n)
    \end{bmatrix}
    \intertext{
      and
  the constraint matrix with indicator of treatment
    }
    T\mathbf{A}
    &
    \ 
    :=
    \ 
    \begin{bmatrix}
      \mathrm{diag}
      [T_1,\ldots,T_n]
      \\
      \pm
      [T_1,\ldots,T_n]
      \\
      \pm\,T\mathbf{B}(\mathbf{X})
    \end{bmatrix}
        \,.
  \end{align*}
  Next we write
  Problem~\ref{bw:1:primal}
  in the form of ts.
 \begin{gather*}
    \underset{w_1, \ldots, w_n \in \R}{\mathrm{minimize}}
    \qquad
    \sum_{i = 1}^{n} T_i \cdot f(w_i)
    \\
    \mathbf{A}w 
    \ 
    \ge
    \ 
    d
    \,,
\end{gather*}
The convex conjugate of the objective function 
of 
  Problem~\ref{bw:1:primal}
is
\begin{gather*}
  [x^*_1,\ldots,x^*_n]^\top
  \ 
  \mapsto
  \ 
  \sum_{T_i=1} f^*(x^*_i)
  +
  \sum_{T_i=0} 
  \chi_{\left\{ 0 \right\}}(x^*_i)
  \,,
\end{gather*}
where we define the characteristic function $\chi$ by
\begin{gather*}
  \chi_{\left\{ 0 \right\}}
  (t)
  \ 
  =
  \ 
  \infty
  \cdot
  (
  1-\mathbf{1}_{\left\{ 0 \right\}}(t)
  )
  \ 
  =
  \ 
  \begin{cases}
    0,& \text{if}\ t=0,\\
    \infty,& \text{else}\,.
  \end{cases}
\end{gather*}
Note that the $i$-th column of $T\mathbf{A}$,
which we denote as 
$T\mathbf{A}_i$,
vanishes if 
$T_i=0$. 
In the subsecquent analysis
this prevents the characteristic function from blowing up.
We consider the form of the dual in ts to get
\begin{gather*}
  \sum_{T_i=1}  f^*(T\mathbf{A}_{ i}^\top\lambda)
  +
  \sum_{T_i=0} 
  \chi_{\left\{ 0 \right\}}
(T\mathbf{A}_{i}^\top\lambda)
  =
  \sum_{i=1}^n T_i f^*(\mathbf{A}_{ i}^\top\lambda)
  \,.
\end{gather*}
where
$\mathbf{A}$
 is as $T\mathbf{A}$ with the indicators of treatment replaced by $1$.
 It is important, that in the final form the indicators of treatment are outside the argument of $f^*$ and we have no singularity for 
 $T_i=0$.
The corresponding dual problem in \cite{Tseng1991} is then
\begin{gather*}
  \underset
  {\lambda_1,\ldots,\lambda_{K}\ge 0}
  {
  \mathrm{maximize}
  }
  \quad
  -
  \sum_{i=1} 
  ^n
  T_i
  \cdot
  f^*
(\mathbf{A}_{ i}^\top\lambda)
  \ 
  +
  \ 
  \inner{\lambda}{d}
  \,.
\end{gather*}

\subsubsection*{Non-negativity constraints}

  Next we want to remove the non-negativity constraints on $\lambda$.
  To this end we write
  \begin{gather}
    \lambda
    :=
    \begin{bmatrix}
      \rho_1,
      \ldots,
      \rho_n,
      \ 
      \lambda_0^+,
      \lambda_0^-,
      \ 
      \lambda_1^+,
      \ldots,
      \lambda_n^+,
      \ 
      \lambda_1^-,
      \ldots,
      \lambda_n^-
    \end{bmatrix}
    ^\top
    \,.
  \end{gather}
  We expand the objective function $G$ of the dual problem.
  \begin{align*}
    G
    (
    \rho,
    \lambda_0^\pm,
    \lambda^\pm
    )
    =
  &-
  \sum_{i=1} 
  ^n
  T_i
  \cdot
  f^*
  \left( 
\rho_i
\ 
+
\ 
\lambda_0^+
\!
-
\lambda_0^-
\ 
+
\ 
\inner
{B(X_i)}
{\,\lambda^+ \!\!- \lambda^-}
  \right)
  \\
  &+
  \ 
  n
  \cdot
  (
\lambda_0^+
\!
-
\lambda_0^-
  )
  \ 
+
  \ 
  \sum_{i=1} 
  ^n
\inner
{B(X_i)}
{\,\lambda^+ \!\!- \lambda^-}
  \ 
-
  \ 
  n
  \cdot
\inner
{\delta}
{\,\lambda^+ \!+ \lambda^-}
  \end{align*}
  To illustrate the procedure, we show 
  for all $i \in \left\{ 1,\ldots,n \right\}$


\begin{alignat*}{2}
  \text{either}
  &
  &&
  \qquad
      \lambda_i^+ > 0
  \\
  \text{or}
  &
  &&
  \qquad
      \lambda_i^- > 0
  \,.
\end{alignat*}
Assume towards a contradiction that 
there exists
$i \in \left\{ 1,\ldots,n \right\}$
such that
$
      \lambda_i^+ > 0
$
and
$
      \lambda_i^- > 0
$ 
and that 
$\lambda$ is optimal.
Consider
  \begin{gather}
    \tilde{\lambda}
    \ 
    :=
    \ 
    \begin{bmatrix}
      \rho
      ,
      \ 
      \lambda_0^\pm,
      \ 
      \lambda_1^+,
      \ldots,
      \ 
      \lambda_i^+
      \!
      -
      (
      \lambda_i^+
      \!
      \land
      \lambda_i^-
      )\,,
      \ 
      \ldots,
      \lambda_n^+,
      \ 
      \lambda_1^-,
      \ldots,
      \lambda_i^-
      \!
      -
      (
      \lambda_i^+
      \!
      \land
      \lambda_i^-
      )\,,
      \ 
      \ldots,
      \lambda_n^-
    \end{bmatrix}
    ^\top
    \,.
  \end{gather}
  Since 
  $
      \lambda_i^\pm
      -
      (
      \lambda_i^+
      \!
      \land
      \lambda_i^-
      )
      \ge 
      0
  $,
  the perturbed vector $\tilde{\lambda}$ is in the domain of the 
  optimization problem.
  But 
  \begin{align}
  G(\tilde{\lambda})
  -
  G(\lambda)
  \ 
  =
  \ 
  2
  n
  \cdot
  \delta_i
  \cdot
      (
      \lambda_i^+
      \!
      \land
      \lambda_i^-
      )
  \ 
  >
  \ 
  0
  \,,
  \end{align}
  which contradicts the optimality of $\lambda$.
  Likewise we can show
\begin{alignat*}{2}
  \text{either}
  &
  &&
  \qquad
      \lambda_i^+ > 0
  \\
  \text{or}
  &
  &&
  \qquad
      \lambda_i^- > 0
  \,.
\end{alignat*}
Let now
\begin{gather}
    \tilde{\lambda}
    \ 
    :=
    \ 
    \begin{bmatrix}
      \rho
      ,
      \ 
      \lambda_0^\pm
      \!
      -
      (
      \lambda_0^+
      \!
      \land
      \lambda_0^-
      )\,,
      \ 
      \lambda_1^+,
      \ldots,
      \lambda_n^+,
      \ 
      \lambda_1^-,
      \ldots,
      \lambda_n^-
    \end{bmatrix}
    ^\top
    \,.
  \end{gather}
We get
  \begin{align}
  G(\tilde{\lambda})
  -
  G(\lambda)
  \ 
  =
  \ 
  2
  n
  \cdot
  \delta_i
  \cdot
      (
      \lambda_i^+
      \!
      \land
      \lambda_i^-
      )
  \ 
  >
  \ 
  0
  \,,
  \end{align}
But then 
$
\lambda^\pm_i
\ge 0
$
collapses to
$
\lambda_i\in \R
$ 
for 
$i\in \left\{ 0,\ldots,n \right\}$, that is,
$ \lambda_i=\lambda_i^+\!-\lambda_i^- $.
Note that
$ |\lambda_i|=\lambda_i^+\!+\lambda_i^- $.
Likewise we can see, that 
$
\lambda_0
=
\lambda_0^+
-
\lambda_0^-
\in \R
$
removes the constraint on $\lambda_0^\pm$.
Let us take this into account for $G$. We get
  \begin{align*}
    G
    (
    \rho,
    \lambda_0,
    \lambda
    )
    =
  &-
  \sum_{i=1} 
  ^n
  T_i
  \cdot
  f^*
  \left( 
\rho_i
\ 
+
\ 
\lambda_0
\ 
+
\ 
\inner
{B(X_i)}
{\lambda}
  \right)
  \\
  &+
  \ 
  n
  \cdot
\lambda_0
  \ 
+
  \ 
  \sum_{i=1} 
  ^n
\inner
{B(X_i)}
{\lambda}
  \ 
-
  \ 
  n
  \cdot
\inner
{\delta}
{|\lambda|}
\,.
  \end{align*}
Next we show, that $\rho=0$.
Suppose there exists 
$
i\in \left\{ 1,\ldots, n \right\}
$
such that 
$
\rho_i>0
$
and
$
  T_i
  \cdot
  (f^{'})^{-1}
  \left( 
\rho_i
\ 
+
\ 
\lambda_0
\ 
+
\ 
\inner
{B(X_i)}
{\lambda}
  \right)
  <
  0
$.
It follows
\begin{gather}
  G(0,\lambda_0,\lambda)
  -
  G(\rho_i,\lambda_0,\lambda)
  \ge
  T_i
  \cdot
  (f^{'})^{-1}
  \left( 
\rho_i
\ 
+
\ 
\lambda_0
\ 
+
\ 
\inner
{B(X_i)}
{\lambda}
  \right)
(-\rho_i)
>0,
\end{gather}
which contradicts the optimality of $\lambda$.
Suppose
$
  T_i
  \cdot
  (f^{'})^{-1}
  \left( 
\rho_i
\ 
+
\ 
\lambda_0
\ 
+
\ 
\inner
{B(X_i)}
{\lambda}
  \right)
  >
  0
$.
Then the claim yields to a perturbation argument as in ts.
Thus
To eliminate the constraints for $\rho$ 
we use a similar argument as in the complementary slackness
section of the ts chapter.
Thus we have complementary slackness of 
$\rho_i$ and
$
  T_i
  \cdot
  (f^{'})^{-1}
  \left( 
\rho_i
\ 
+
\ 
\lambda_0
\ 
+
\ 
\inner
{B(X_i)}
{\lambda}
  \right)
$.
But then
every
optimal solution $\lambda$ remains optimal by taking $\rho=0$.

Dividing the optimization problem by $n$ and reversing it, we get

\begin{gather*}
  \underset
  {\lambda_0,\ldots,\lambda_{n}\in \R}
  {
    \mathrm{minimize}
  }
  \quad
  \frac{1}{n}
\sum_{i=1} 
  ^n
  \left[ 
    \,
  T_i
  \cdot
  f^*
  \!
  \left( 
\lambda_0
+
\inner
{B(X_i)}
{\lambda}
  \right)
  \ 
-
\ 
  \left( 
\lambda_0
+
\inner
{B(X_i)}
{\lambda}
  \right)
  \,
  \right]
  \ 
+
\ 
\inner
{\delta}
{|\lambda|}
  \,.
\end{gather*}

\end{proof}

%  \newpage
  %\section{Application of Convex Optimization}
  %\begin{assumption}
  \begin{enumerate}[label={(\roman*)}]
    Assume that the map 
    $
      f: \R \to \overline{\R}
    $
    has the following properties.
    \item
      $
        f 
        \ 
        \text{is strictly convex.}
      $
    \item
      $
        f
        \ 
        \text{is lower-semicontinuous and continuously differentiable on}
        \ 
        \mathrm{int}(\mathrm{dom}(f))
        .
      $
    \item
      $
        \text{The derivative of}\ 
        f\ 
        \text{on}\ 
        \mathrm{int}(\mathrm{dom}(f))
        \ 
        \text{is a diffeomorphism.}
      $
    \item
      $
        \text{The Legendre transformation}
        \ 
        f^*
        \ 
        \text{of}\ 
        f
        \ 
        \text{is finite}
        .
      $
    \item
      $
        \text{The function}\
        x\mapsto xt - f(x)
        \ 
        \text{takes its supremum on}
        \ 
        \mathrm{int}(\mathrm{dom}(f))
        \ 
        \text{for all}
        \ 
        t\in \R.
      $
  \end{enumerate}
\end{assumption}
Assumption (i),(ii) and (iv) are from \cite{Tseng1991}. See corresponding section in the thesis.
Assumption (iii) and (v) is needed to get a meaningful analytic expression for the weights after establishing the properties pointed out in \cite{Tseng1991}.

\todo[color=red!40, inline]{Elaborate on assumptions. Give as a counterexample 
  \[
  f(x):=
  (1+\delta_{x\ge 0})
  (x^2
  +x\cdot \lambda)
\]
}

We consider the following optimization problem.
\todo[color=orange!40, inline]{Consider Non-negativity constraints only for $f(x)=x\log x$. Incorporating Non-negativity constraint for general $f$ is complicated or impossible. Maybe find counterexample.}
  \begin{gather*}
    \underset{w_1, \ldots, w_n \in \R}{\mathrm{minimize}}
    \qquad
    \sum_{i = 1}^{n} T_i f(w_i)
  \end{gather*}
subject to the constraints
\begin{gather*}
%    w_i T_i \ge 0,
%   \qquad 
%    i = 1, \ldots, n,
%   \\
%  \sum_{ i = 1 }^{n}
%    w_i T_i
%  =
%  1
%  \\
    \left| 
      \frac{1}{n} 
      \sum_{i = 1}^{n} 
      (
      w_i T_i 
      - 
      1
      )
      \cdot
      B_k(X_i)
    \right|
    \ 
    \le 
    \ 
    \delta_k,
    \qquad
    k = 1, \ldots, K
\end{gather*}
\pagebreak
~\begin{ftheorem}
  Under Assumption,
  the dual of the above Problem is the unconstrained optimization problem 
  \begin{gather*}
    \underset{\lambda \in \R^K}{\mathrm{minimize}}
    \qquad
    \frac{1}{n}
    \sum_{i = 1}^{n} 
    T_i 
    \cdot
    f^*
    (
      \inner{B(X_i)}{\lambda}
    )
    -
    \inner{B(X_i)}{\lambda}
    \ 
    +
    \ 
    \inner{\delta}{\left| \lambda \right|}
    \,,
  \end{gather*}
  where
  \begin{gather*}
  f^*
  \,
  \colon
  \, 
  \R
  \ 
  \to
  \ 
  \R
  \,
  ,
  \qquad 
  t 
  \ 
  \mapsto
  \ 
    t\,(f^{'})^{-1}(t)
  \ 
    -
  \ 
    f
    \left( 
      (f^{'})^{-1}(t)
    \right)
  \end{gather*}
  is the Legendre transformation of $f$,
  $
    B(X_i)
    =
    \left[ 
      B_1(X_i)
      ,
      \ldots
      ,
      B_K(X_i)
    \right]
    ^\top
  $
  denotes the $K$ basis functions of the covariates 
  of unit $i\in \left\{ 1, \ldots, n \right\}$
  and
  $
    \left| \lambda \right|
    =
    \left[ 
      \left| \lambda_1 \right|
      ,
      \ldots
      ,
      \left| \lambda_K \right|
    \right]
    ^\top
    ,
  $
  where $\left| \,\cdot\, \right|$
  is the absolute value of a real-valued scalar.
  Moreover, if $\lambda^\dagger$
  is an optimal solution then
  \begin{gather*}
    w_i^*
    \ 
    =
    \ 
    (f^{'})^{-1}
    \left( 
      \inner{B(X_i)}{\lambda^\dagger}
    \right)
    \quad
    \text{for all}\ 
    i
    \ 
    \text{with}\ 
    T_i=1
  \end{gather*}
  are uniquely part of any optimal solution to (P)
  .
\end{ftheorem}
\begin{proof}
  We prove the following Lemma at the end of the section.
  \begin{lemma}
    The dual of the optimization problem is
  \begin{gather*}
    \underset{\lambda \in \R^{2K}}{\mathrm{minimize}}
    \qquad
    \frac{1}{n}
    \sum_{i = 1}^{n} 
    T_i 
    \cdot
    f^*
    (
    \inner{Q_{\bullet i}}{\lambda}
    )
    -
    \inner{Q_{\bullet i}}{\lambda}
    \ 
    +
    \ 
    \inner{d}{\lambda}
  \end{gather*}
  subject to
  \begin{gather*}
    \lambda_k \ge 0
    \quad
    \text{for all}\ 
    k \in \left\{ 1, \ldots, K \right\}
    ,
  \end{gather*}
  where
  \begin{gather*}
    \mathbf{Q}
    :=
    \begin{bmatrix}
     % \mathbf{I}_n\\
     % \mathbf{B}(\mathbf{X})\\
      \pm \, \mathbf{B}(\mathbf{X})
    \end{bmatrix}
    ,
    \qquad
    \mathbf{B}(\mathbf{X})
    :=
    \begin{bmatrix}
      B(X_1), \ldots, B(X_n)
    \end{bmatrix}
    ,
    \qquad
    \text{and}
    \qquad
    d
    :=
    \begin{bmatrix}
      %0_n\\
      \delta \\
      \delta
    \end{bmatrix}
    .
  \end{gather*}

  \end{lemma}
  \begin{proof}
    First we disentangle the box constraints.
    To this end, we get
    \begin{gather*}
    \underset{w_1, \ldots, w_n \in \R}{\mathrm{minimize}}
    \qquad
    \sum_{i = 1}^{n} T_i f(w_i)
    \\
      \pm\,\sum_{i = 1}^{n} w_i T_i B_k(X_i)
      \ 
      \ge
      \ 
      -n\cdot\delta_k
      \ 
      \pm 
      \ 
      \sum_{i = 1}^{n} B_k(X_i)
      \,,
      \qquad
      k=1,\ldots,K
      \,.
    \end{gather*}
    The corresponding matrix notation is
 \begin{gather*}
    \underset{w_1, \ldots, w_n \in \R}{\mathrm{minimize}}
    \qquad
    \sum_{i = 1}^{n} T_i \cdot f(w_i)
    \\
    \mathbf{Q}w 
    \ 
    \ge
    \ 
    d
    \,,
\end{gather*}
where
\begin{align*}
    \mathbf{Q}
    &
    \ 
    :=
    \ 
    \begin{bmatrix}
      \pm\,T\mathbf{B}(\mathbf{X})
    \end{bmatrix}
    ,
    \\
    T\mathbf{B}(\mathbf{X})
    &
    \ 
    :=
    \ 
    \begin{bmatrix}
      T_1B(X_1), \ldots, T_nB(X_n)
    \end{bmatrix}
    ,
    \\
    d
    &
    \ 
    :=
    \ 
    \begin{bmatrix}
      -n\cdot\delta 
      \pm\,
      \sum_{i = 1}^{n} B_k(X_i)
    \end{bmatrix}
    \,.
  \end{align*}
The convex conjugate is
\begin{gather*}
  \sum_{T_i=1} T_i f^*(\lambda_i)
  +
  \sum_{T_i=0} 
  \delta_{\left\{ 0 \right\}}(\lambda_i)
  \,,
\end{gather*}
where
\begin{gather*}
  \delta_{\left\{ 0 \right\}}
  (t)
  =
  \begin{cases}
    0,& \text{if}\, t=0,\\
    \infty,& \text{else}\,.
  \end{cases}
\end{gather*}
Note that
\begin{gather*}
  \sum_{T_i=0} 
  \delta_{\left\{ 0 \right\}}
  (
  T_i \cdot \inner{B(X_i)}{\lambda}
  )
  =
  0
  \,.
\end{gather*}

The corresponding dual problem in \cite{Tseng1991} is then
\begin{gather*}
  \underset
  {\lambda_1,\ldots,\lambda_{K}\ge 0}
  {
  \mathrm{minimize}
  }
  \quad
  \sum_{i=1} 
  ^n
  T_i
  \cdot
  f^*
  (
  T_i \cdot\inner{B(X_i)}{\lambda}
  )
  -
  \inner{B(X_i)}{\lambda}
  \ 
  +
  \ 
  n\inner{\lambda}{d}
  \,.
\end{gather*}
  If we keep only the outer $T_i$ and divide by $1/n$ the problem remains the same.
  This concludes the proof.
  \end{proof}

\todo[color=red!40, inline]{Continue with paraphrase of \cite[page 20]{Wang2019}}

\end{proof}


  \chapter{Asymptotic Analysis}
  \section{Consistency of Optimal Solutions}
  \subsection{Estimate of an Oracle Parameter by the Dual}
  \begin{assumption}
  \label{aa:assumption:overlap}
  The propensity score function satisfies
  $\pi(x)\in(0,1)$ for all $x\in\mathcal{X}$.
\end{assumption}
\begin{assumption}
  \label{aa:assumption:ps_continuous}
  The propensity score function is continuous in $\mathcal{X}$.
\end{assumption}
\subsection*{My Contribution}
I found out, that consistency for the dual variable is enough to prove later results. This simplifies the proof. 
In \cite{Wang2019} they use a quadratic Taylor expansion to obtain learning rates. I found out, that a simpler mean value result for differentiable convex functions is sufficient to proof consistency. 
Since I work with partitioning estimates, I found a suitable oracle parameter.
I prove an (extended) lemma which is central but the details were omited.

Throughout this section we assume for all $N\in\mathbb{N}$ the existence of an 
optimal solution 
$(\lambda_0^\dagger,\lambda^\dagger)$
to Problem?
We define the oracle parameter $\lambda^*\in\R^N$ to be the vector with coordinates
\begin{gather}
  \label{oracle_1}
  \lambda^*_k
  \ 
  :=
  \ 
  f^{'}
  \left( 
    \frac{1}{\pi(X_k)}
  \right)
  \qquad
  \text{for all}\ 
  k\in \left\{ 1,\ldots,N \right\}
  \,,
\end{gather}
where $\pi(x)=\P[T=1|X=x]$ is the \textbf{propensity score} at $x\in\mathcal{X}$. Why this choice? 
%First, the $\lambda_0^\dagger$ part is unimportant. We need it to eliminate the same factor in
%\begin{gather}
%  w(x):=
%  (
%  f^{'}
%  )^{-1}
%  \left( \inner{B(x)}{\lambda^\dagger}+\lambda_0^\dagger \right)
%\end{gather}
%that is 
%\begin{gather}
%  \inner{B(x)}{\lambda^*}+\lambda_0^\dagger 
%  =
%\sum_{k=1}^{N} 
%B_k(x)
%  f^{'}
%  \left( 
%    \frac{1}{\pi(X_k)}
%  \right)
%  \,.
%\end{gather}
With hindsight, we need
for all $i\in \left\{ 1,\ldots,N \right\}$
\begin{gather}
\left| 
\inner{B(X_i)}
{\lambda^*}
  -
  f^{'}
  \left( 
    \frac{1}{\pi(X_i)}
  \right)
\right|
\ 
\le
\ 
      \omega
      \left( f^{'}\circ (x\mapsto1/x)\circ \pi,h_N \right)
      \,.
\end{gather}
Consequently, if $\pi$ is continuous and positive (not 0) on $\mathcal{X}$ and the width of the partition $h_N$ converges to 0, we get 
\begin{gather}
 \left| 
  \inner{B(X_i)}{\lambda^*}
  -
  f^{'}
  \left( 
    \frac{1}{\pi(X_k)}
  \right)
 \right| 
 \
 \to
 \ 
 0
 \qquad
 \text{almost surely.}
\end{gather}
We need this to bound the term $M_N$ in \eqref{c:first:1}.
The first big goal is to prove the following proposition.
\begin{proposition}
  \label{bw:cd:th}
Define the oracle parameter $\lambda^*$ as in \eqref{oracle_1}.
Furthermore, assume that the propensity score function is continuous and positive on $\mathcal{X}$.
Then for all $\varepsilon>0$ with probability going to 1 there exists an optimal solution
     $(\lambda^\dagger,\lambda_0^\dagger)$
     to Problem ? with
     $
     \norm
     {
       (\lambda^\dagger,\lambda_0^\dagger)
      -
      (\lambda^*,0)
     }_2
     \le
     \varepsilon
     $.
\end{proposition}
\begin{remark}
  In the analysis of the next sections will assume existence of an optimal solution to Problem?
  The purpose of Proposition~\ref{bw:cd:th} is therefore twofold. 
  First, it ensures the existence of an optimal solution with probability going to 1 as $N\to\infty$.
  This shows that the assumption of existence is (at least for large $N$) likely met.
  Nevertheless, it is beyond the scope of my thesis to investigate feasibility of Problem ? in more detail. 
  Second, if a solution exists it converges to the oracle parameter $\lambda^*$ in probability. This ensures the consistency of the weights for the inverse propensity score.
\end{remark}

We use a hint from the last display of~\cite[p.22]{Wang2019}.
The high-level idea is that the existence of the optimal solution and its proximity to the oracle parameter can be analysed by the objective function.
\begin{lemma}
  \label{bw:cd:lem}
  Let $m\in\mathbb{N}$ and
  $g \,:\, \R^m \to \overline{\R}$ 
  be convex.
  Then 
  for all $y \in \R^m$ and $\varepsilon>0$ 
    \begin{gather}
      \label{7060_0}
      \inf_{\norm{\Delta}=\varepsilon} g(y+\Delta) - g(y) \ge 0 \quad
    \end{gather}
    implies
    the existence of  
    a global minimum
    $
    y^* \in \,\R^m
    $
    of $g$
    satisfying
    $
      \norm{y^* - y}_2 \le \varepsilon
    $.
\end{lemma}
\begin{proof}
  Let $B$ be the euclidian ball in $\R^m$.
  Since 
  $
  y
  \,
  +
  \,
  \varepsilon
  B
  $
  is convex, it contains a 
  local minimum  
  of $g$.
  Suppose towards a contradiction that
  $
    y^* 
    \ 
    \in 
    \ 
  y
  \,
  +
  \,
  \varepsilon
  B
  $
  is a local minimum, but not a global one, and
  \eqref{7060_0} is true.
  Then it holds
  \begin{gather}
    \label{7060_3}
    g(x) < g(y^*)
    \quad
    \text{for some}\ 
    x 
    \in 
    \R^m 
    \setminus 
    \left( 
  y
  \,
  +
  \,
  \varepsilon
  B
    \right)
  \,.
  \end{gather}
  Furthermore, since 
  $
  y
  \,
  +
  \,
  \varepsilon
  B
  $ is compact and contains $y^*$,
  the line segment connecting 
  $y^*$ and $x$
  intersects the boundary of 
  $y + \mathcal{C}$, that is,
  there exist
  $
    \theta \in (0,1)
  $
  and 
  $
    \Delta_x
  $
  with 
  $
    \norm{\Delta_x}_2=\varepsilon
  $
  such that
  \begin{gather}
    \label{7060_4}
    \theta x + (1 - \theta) y^* = y + \Delta_x
    \,.
  \end{gather}
    It follows
    \begin{align}
      \label{7060_5}
      \begin{split}
      g(y^*)
      \le
      g(y)
      \le
      g(y + \Delta_x)
      &=
      g(
        \theta x + (1 - \theta) y^*
      )
      \\
      &\le
      \theta g(x)
      + 
      (1 - \theta)
      g(y^*)
      <
      g(y^*)
      ,
      \end{split}
    \end{align}
    which is a contradiction.
    The first inequality is due to
    $y^*$ being a local minimum of $g$ in
    $
  y
  \,
  +
  \,
  \varepsilon
  B
    $,
    the second inequality is due to  
    \eqref{7060_0} being true,
    the equality is due to \eqref{7060_4},
    the third inequality is due to the convexity of $g$
    and the strict inequality is due to \eqref{7060_3}.
    Thus every local minimum of $g$ in
    $
  y
  \,
  +
  \,
  \varepsilon
  B
    $
    is also a global minimum.
    %It follows the right-hand side of \eqref{7060_0}.
\end{proof}
\begin{remark}
  The hint from \cite[page 22]{Wang2019}
  states
  \eqref{7060_0}
  with strict inequality.
  They base their subsequent analysis on a quadratic Taylor expansion and aim to prove \eqref{7060_0} with strict inequality.
  I show, that this approach is inefficient.  
  To do that, I need Lemma~\ref{bw:cd:lem} exactly as stated - with $\ge$ in \eqref{7060_0}.
\end{remark}
Lemma~\ref{bw:cd:lem} motivates the following lemma.
 \begin{lemma}
   \label{bw:cd:lem2}
   Under the conditions of Theorem~\ref{bw:cd:th} it holds
   for all $\varepsilon>0$
\begin{gather}
   \P
   \left[ 
     \inf _ { 
       \norm{
         (
     \Delta
     ,
     \Delta_0
         )
 } 
= \varepsilon }
     G
     (
     \lambda^*
      +
      \Delta
      ,
     \Delta_0
     )
     -
     G
     (
     \lambda^*,
     0
     )
     \ge 
     0
   \right]
   \ 
   \to
   \ 
   1
   \qquad
   \text{for}
   \ 
   N\to\infty
   \,.
\end{gather}
 \end{lemma}
 \begin{proof}
  Recall the objective function $G$ of Problem? 
\begin{align*}
  G(\lambda,\lambda_0)
  &
  \ 
  =
  \ 
    \frac{1}{N}
\sum_{i=1} 
  ^N
  \left[ 
    \,
  T_i
  \cdot
  f^*
  \!
  \left( 
\lambda_0
+
\inner
{B(X_i)}
{\lambda}
  \right)
  \ 
-
\ 
  \left( 
\lambda_0
+
\inner
{B(X_i)}
{\lambda}
  \right)
  \,
  \right]
  \ 
+
\ 
\inner
{\delta}
{|\lambda|}
\,
\\
  &
\ 
=
\ 
g(\lambda,\lambda_0)
\ 
+
\ 
\inner
{\delta}
{|\lambda|}
\,
\end{align*}
with
  \begin{gather}
    g
    \ 
    := 
    \ 
    (\lambda,\lambda_0)
    \ 
    \mapsto
    \ 
    \frac{1}{N}
\sum_{i=1} 
  ^N
  \left[ 
    \,
  T_i
  \cdot
  f^*
  \!
  \left( 
\lambda_0
+
\inner
{B(X_i)}
{\lambda}
  \right)
  \ 
-
\ 
  \left( 
\lambda_0
+
\inner
{B(X_i)}
{\lambda}
  \right)
  \,
  \right]
  \,.
  \end{gather}
   Since we assume the convex conjugate $f^*$ of $f$ to be differentiable
   (it always convex),
   $g$ is differentiable convex function with gradient 
  \begin{gather*}
    \nabla
    g
    \ 
    =
    \ 
    (\lambda,\lambda_0)
    \ 
    \mapsto
    \ 
\frac{1}{N}
\sum_{i=1} 
  ^N
  \left[ 
    \,
  T_i
  \cdot
  (f^{'})^{-1}
  \!
  \left( 
\lambda_0
+
\inner
{B(X_i)}
{\lambda}
  \right)
  \ 
-
\ 
1
  \right]
  [
  B(X_i)^\top
  ,
  1
  ]^\top
  \,.
  \end{gather*}
It is well know that a differentiable convex functions $g$ satisfies
  \begin{gather*}
    g(x)
    \ 
    -
    \ 
    g(y)
    \ 
    \ge
    \ 
    \nabla
    g(y)^\top
    (x-y)
    \qquad 
    \text{for all}\ 
    x,y\,.
  \end{gather*}
Thus 
\begin{align}
  \label{c:1}
  \begin{split}
     &
   G
     (
     \lambda^*
      +
      \Delta
      ,
     \Delta_0
     )
     \ 
     -
     \ 
     G
     (
     \lambda^*,
     0
     )
         %%%%%%%%%%%%% 1 %%%%%%%%%%%%%%
     \\
     &
     \quad
     \ge
     -
     \frac{1}{N}
     \sum_{i=1}^{N} 
     \left[ 
       B(X_i)^\top,
       1
     \right]
     \cdot
     \begin{bmatrix}
       \Delta\\
       \Delta_0
     \end{bmatrix}
     \left( 
       1
       \ 
       -
       \ 
     T_i
     \cdot
     (f^{'})^{-1}
     \left( 
       \inner
       {B(X_i)}
       {\lambda^*}
     \right)
     \right)
     \\
     &
     \qquad
     +
     \ 
     \inner
     {\delta}
     {
       |\lambda^*+\Delta|
       -
       |\lambda^*|
     }
     \,.
   \end{split}
\end{align}

Next, we fix $
\tilde{\varepsilon}
>0
$
and establish in \eqref{c:1} the lower bound
$
-
\tilde{\varepsilon}
$
with probability going to $1$ as $N\to\infty$.
Then we conclude that this holds for all $\tilde{\varepsilon}>0$.
The measurability of
$
G
     (
     \lambda^*
      +
      \Delta
      ,
     \Delta_0
     )
     \ 
     -
     \ 
     G
     (
     \lambda^*,
     0
     )
$
will give us the lower bound 0 in \eqref{c:1} with probability going to $1$.

In \eqref{c:1} we control the \textbf{first term} by the law of large numbers 
and the \textbf{second term} by $\norm{\delta}_1\to 0$.
\subsection*{First Term}
We note, that by $\norm{B(x)}_2\le 1$ for all $x\in\mathcal{X}$ and the Cauchy-Schwarz inequality
it holds
\begin{gather}
  \label{c:first:0}
      \left[ 
       B(X_i)^\top,
       1
     \right]
     \cdot
     \begin{bmatrix}
       \Delta\\
       \Delta_0
     \end{bmatrix}
     \ 
     \lesssim
     \ 
     \norm{(\Delta,\Delta_0)}
     =\varepsilon
     \,.
\end{gather}
Next, we see that
\begin{align}
  \label{c:first:1}
  \begin{split}
  &
     \frac{1}{N}
     \sum_{i=1}^{N} 
     \left( 
       1
       \ 
       -
       \ 
     T_i
     \cdot
     (f^{'})^{-1}
     \left( 
       \inner
       {B(X_i)}
       {\lambda^*}
     \right)
     \right)
     \\
     &
     \ 
     \lesssim
     \ 
     \frac{1}{N}
     \sum_{i=1}^{N} 
     \left|
     1
     -
     \frac{T_i}{\pi(X_i)}
     \right|
      \ 
     +
      \ 
     \frac{1}{N}
     \sum_{i=1}^{N} 
     \left| 
        \inner
       {B(X_i)}
       {\lambda^*}
      \ 
        -
        \ 
        f^{'}
        \left( 
          \frac{1}{\pi(X_i)}
        \right)
     \right|
     \\
     &
     \ 
     =:
     \ 
     S_N
     \
     +
     \ 
     M_N
     \,.
\end{split}
\end{align}
With $
\tilde{\varepsilon}
>0
$
fixed previously,
we want to
establish the upper bound
$
\tilde{\varepsilon}
/
(2\varepsilon)
$
with probability going to $1$ as $N\to\infty$.
First, we bound $S_N$.
By the properties of conditional expectation it holds
\begin{gather*}
  \E
  \left[ 
    \frac{T}{\pi(X)}
  \right]
  =
  \E
  \left[ 
    \frac{\E[T|X]}{\pi(X)}
  \right]
  =1
  \,.
\end{gather*}
Also
\begin{gather}
  \E
  \left[ 
    \left| 
    1
    -
    \frac{T}{\pi(X)}
    \right|
  \right]
  \le
  1
  +
  \E
  \left[ 
    \frac{T}{\pi(X)}
  \right]
  =
  2
  \,.
\end{gather}
Thus Etemadi's ($\mathcal{L}_1$ version) strong law of large numbers (cf.\cite[Theorem~5.17]{Klenke2020}) applies
to $S_N$, that is,
$S_N \le 
\tilde{\varepsilon}
/
(4\varepsilon)
$
with probability going to 1.

Next, we bound $M_N$.
Recall that $\sum_{k=1}^{N}B_k(x)=1$ for all $x\in\mathcal{X}$. Thus
\begin{align}
  \label{part_trick}
  \begin{split}
  &
\left| 
        \inner
       {B(X_i)}
       {\lambda^*}
       \ 
        -
        \ 
        f^{'}
        \left( 
          \frac{1}{\pi(X_i)}
     \right)
\right|
      \\
      &
      \ 
      =
      \ 
      \left| 
      \sum_{k=1}^{N} 
      B_k(X_i)
      \cdot
        f^{'}
        \left( 
          \frac{1}{\pi(X_k)}
          \right)
      \ 
      -
      \ 
        f^{'}
        \left( 
          \frac{1}{\pi(X_i)}
     \right)
      \right|
      \\
      &
      \ 
      =
      \ 
      \left| 
      \sum_{k=1}^{N} 
      B_k(X_i)
      \left( 
        f^{'}
        \left( 
          \frac{1}{\pi(X_k)}
     \right)
     -
        f^{'}
        \left( 
          \frac{1}{\pi(X_i)}
     \right)
      \right)
      \right|
      \\
      &
      \ 
      \le
      \ 
      \sum_{k=1}^{N} 
      \frac{1_{\left\{ X_k\in A_N(X_i) \right\}}}
      {
        \sum_{j=1}^{N} 
1_{\left\{ X_j\in A_N(X_i) \right\}}
      }
      \left| 
        f^{'}
        \left( 
          \frac{1}{\pi(X_k)}
     \right)
     -
        f^{'}
        \left( 
          \frac{1}{\pi(X_i)}
     \right)
      \right|
      \\
      &
      \ 
      \le
      \ 
      \omega
      \left( f^{'}\circ (x\mapsto1/x)\circ \pi,h_N \right)
      \ 
      \to
      \ 
      0
      \,,
\end{split}
\end{align}
where $\omega$ is the modulus of continuity.
The convergence to 0 is due to the $f^{'}$ being continuous, $\pi(x)\in(0,1)$ for all $x\in\mathcal{X}$, the (assumed) continuity of $\pi$ and $h_N\to 0$ for $N\to\infty$.
We conclude, that $
M_N\le
\tilde{\varepsilon}
/
(4\varepsilon)
$
with probability going to 1.

This establishes the desired bound of 
$
\tilde{\varepsilon}/(2\varepsilon)
$
in \eqref{c:first:1}.
Together with \eqref{c:first:0}
we conclude that the \textbf{first term} 
in
\eqref{c:1}
is bounded below by
$
-
\tilde{\varepsilon}/2
$
with probability going to $1$ as $N\to\infty$.
\subsection*{Second Term}
It holds
\begin{gather*}
  |x+y|-|x|\ge
  -|y|
  \qquad
  \text{for all}\ 
  x,y\in\R
  \,.
\end{gather*}
Since
$\delta\ge 0$
we get
\begin{align*}
  &
     \inner
     {\delta}
     {
       |\lambda^*+\Delta|
       -
       |\lambda^*|
     }
     \\
     &
     \ 
     \ge
     \ 
     -
     \inner{\delta}
     {|\Delta|}
     \ 
     \ge
     \ 
     -
     \norm{\delta}_1
     \norm{\Delta}_\infty
     \ 
     \ge
     \ 
     -
     \norm{\delta}_1
     \norm{(\Delta,\Delta_0)}_2
     \ 
     \ge
     \ 
     -
     \norm{\delta}_1
     \varepsilon
     \ 
     \ge
     \ 
     -
     \tilde{\varepsilon}/2
     \,,
\end{align*}
with probability going to $1$ as $N\to \infty$.
The convergence is due to $\norm{\delta}_1$ converging to $0$ in probability.
\subsection*{Conclusion}
With the analysis of the \textbf{first} and \textbf{second term} in
\eqref{c:1} we conclude
\begin{gather}
  \label{c:7}
  G
     (
     \lambda^*
      +
      \Delta
      ,
     \Delta_0
     )
     \ 
     -
     \ 
     G
     (
     \lambda^*,
     0
     )
     \ge
     -
     \tilde{\varepsilon}
\end{gather}
with probability going to $1$ as $N\to \infty$.
But this holds for all $\tilde{\varepsilon}>0$.
A closer look reveals that
$
  G
     (
     \lambda^*
      +
      \Delta
      ,
     \Delta_0
     )
     \ 
     -
     \ 
     G
     (
     \lambda^*,
     0
     )
$ 
is measurable.
Indeed, this holds because $X$,$T$,$B(X)$ and $\lambda^*$ are measurable and $f^*$ is continuous.
Since \eqref{c:7} holds true for all $\tilde{\varepsilon}>0$, it thus follows 
\begin{gather}
  G
     (
     \lambda^*
      +
      \Delta
      ,
     \Delta_0
     )
     \ 
     -
     \ 
     G
     (
     \lambda^*,
     0
     )
     \ge
     0
\end{gather}
with probability going to $1$ as $N\to \infty$.
But this holds for all 
$
(\Delta,\Delta_0)
$
with 
$
\norm{
(\Delta,\Delta_0)
}
=\varepsilon
$. Thus
\begin{gather}
   \inf _ { 
       \norm{
         (
     \Delta
     ,
     \Delta_0
         )
 } 
= \varepsilon }
     G
     (
     \lambda^*
      +
      \Delta
      ,
     \Delta_0
     )
     -
     G
     (
     \lambda^*,
     0
     )
     \ge 
     0
\end{gather}
with probability going to $1$ as $N\to \infty$.
Finally, we see, that this holds for all $\varepsilon>0$. This finish the proof.
 \end{proof}

\begin{proof}
  \emph{(Proposition~\ref{bw:cd:th})}
An immediate consequence of Lemma~\ref{bw:cd:lem} is
that for all $\varepsilon>0$ it holds
\begin{gather*}
   \P
   \left[ 
     \text{There exists an optimal solution $(\lambda^\dagger,\lambda_0^\dagger)$ to Problem? with}
     \ 
     \norm
     {
       (\lambda^\dagger,\lambda_0^\dagger)
      -
      (\lambda^*,0)
     }_2
     \le
     \varepsilon
   \right]
   \\
   \
   \ge
   \ 
   \P
   \left[ 
     \inf _ { 
       \norm{
         (
     \Delta
     ,
     \Delta_0
         )
 } 
= \varepsilon }
     G
     (
     \lambda^*
      +
      \Delta
      ,
     \Delta_0
     )
     -
     G
     (
     \lambda^*,
     0
     )
     \ge 
     0
   \right]
   \,.
 \end{gather*}
 Applying Lemma~\ref{bw:cd:lem2} finishes the proof.
\end{proof}


  \subsection{Estimate of the Inverse Propensitiy Score by the Weights}
  The following theorem is an easy consequence of Theorem~\ref{bw:cd:th}.
\begin{ftheorem}
  Consider the weights function defined by
  \begin{gather}
    w(x)
    :=
    (
    f^{'}
    )^{-1}
    \left( 
      \inner{B(x)}{\lambda^\dagger}
      +
      \lambda_0^\dagger
    \right)
    \qquad
    \text{for all}\ 
    x\in\mathcal{X}
    \,.
  \end{gather}
  Under the conditions of Theorem~\ref{bw:cd:th} 
  it holds
  $w(X)\ 
  \overset{\P}{\to}
  \ 
  1/\pi(X)
  $
\end{ftheorem}
\begin{proof}
  For all $\varepsilon>0$ it holds
\begin{align}
  \begin{split}
  \left| 
  w(X)
  -
  \frac{1}{\pi(X)}
  \right|
  &
  \ 
  =
  \ 
  \left| 
  (f^{'})^{-1}
  \left( 
    \inner{B(X)}
    {\lambda^\dagger}
    +
    \lambda_0^\dagger
  \right)
  -
  \frac{1}{\pi(X)}
  \right|
  \\
  &
  \ 
  \lesssim
  \ 
  \left| 
  \inner
  {B(X)}
  {\lambda^\dagger-\lambda^*}
  \right|
  +
  \left| 
    \inner{B(X)}
    {\lambda^\dagger}
    +
    \lambda_0^\dagger
    -
    f^{'}
    \left( 
  \frac{1}{\pi(X)}
    \right)
  \right|
  \\
  &
  \ 
  \lesssim
  \ 
  \norm{\lambda^\dagger-\lambda^*}_2
  +
  \left| 
  \sum_{i=1}^{N} 
  B_k(X)
  \cdot
    f^{'}
    \left( 
  \frac{1}{\pi_k}
    \right)
    -
    f^{'}
    \left( 
  \frac{1}{\pi(X)}
    \right)
  \right|
  \\
  &
  \ 
  \le
  \ 
  \frac{\varepsilon}{2}
+
  \frac{\varepsilon}{2}
  \le
  \varepsilon
  \,,
\end{split}
\end{align}
with probability going to $1$ as $N\to\infty$.
\end{proof}

  %Throughout this section assume the existence of an 
optimal solution 
$(\lambda^\dagger,\lambda_0^\dagger)$.
We use a hint from the last display of~\cite[p.22]{Wang2019}.
The high-level idea is, 
to connect the optimality of a dual solution 
to 
being in the 
neighborhood of an oracle parameter
by looking at the objective function of the dual.
We deliver the omitted technical details.
\subsection*{Neighbourhood of Oracle Parameter}
Let $\lambda^*$ denote the vector with coordinates
\begin{gather}
  \lambda^*_i
  :=
  f^{'}(1/\pi_i)
  -
  \lambda^\dagger_0
  \,,
\end{gather}
where $\pi_i=\P[T_i=1|X_i]$ is the \textbf{propensity score} of the 
$i$-th unit.

\begin{theorem}
  \label{bw:cd:th}
  For all
  $\varepsilon>0$ it holds
  \begin{gather}
    \P
    \left[ 
    \norm{
      \lambda^\dagger
      -
      \lambda^*
    }
    \ge
    \varepsilon
    \right]
    \ 
    \to
    \ 
    0
    \qquad
    \text{for}\ 
    N
    \to 
    \infty
    \,.
  \end{gather}
\end{theorem}
We want to leverage the convexity of
the objective function of the dual to get
 \begin{gather*}
   \P
   \left[ 
     \norm
     {
      \lambda ^ \dagger
      -
      \lambda^*
     }
     \le
     \varepsilon
   \right]
   =
   \P
   \left[ 
     \inf _ { 
       \norm{
         (
     \Delta
     ,
     \Delta_0
         )
 } 
= \varepsilon }
     G
     (
     \lambda^*
      +
      \Delta
      ,
      \lambda^\dagger_0
      +
     \Delta_0
     )
     -
     G
     (
     \lambda^*,
      \lambda^\dagger_0
     )
     \ge 
     0
   \right]
   \,.
 \end{gather*}
We learned about a similar idea from \cite[p.22]{Wang2019}. The next Lemma makes this rigorous. 
\begin{lemma}
  \label{bw:cd:lem}
  Let $m\in\mathbb{N}$ and
  $g \,:\, \R^m \to \overline{\R}$ 
  be convex.
  Then 
  for all $y \in \R^m$ and $\varepsilon>0$ 
    \begin{gather}
      \label{7060_0}
      \inf_{\norm{\Delta}=\varepsilon} g(y+\Delta) - g(y) \ge 0 \quad
    \end{gather}
    implies
    the existence of  
    a global minimum
    $
    y^* \in \,\R^m
    $
    of $g$
    satisfying
    $
      \norm{y^* - y} \le \varepsilon
    $.
\end{lemma}
\begin{proof}
  Since 
  $
  y
  \,
  +
  \,
  \varepsilon
  B
  $
  is convex, it contains a 
  local minimum  
  of $g$.
  Suppose towards a contradiction that
  $
    y^* 
    \ 
    \in 
    \ 
  y
  \,
  +
  \,
  \varepsilon
  B
  $
  is a local minimum, but not a global one, and
  \eqref{7060_0} is true.
  Then it holds
  \begin{gather}
    \label{7060_3}
    g(x) < g(y^*)
    \quad
    \text{for some}\ 
    x 
    \in 
    \R^m 
    \setminus 
    \left( 
  y
  \,
  +
  \,
  \varepsilon
  B
    \right)
  \,.
  \end{gather}
  Furthermore, since 
  $
  y
  \,
  +
  \,
  \varepsilon
  B
  $ is compact and contains $y^*$,
  the line segment connecting 
  $y^*$ and $x$
  intersects the boundary of 
  $y + \mathcal{C}$, that is,
  there exist
  $
    \theta \in (0,1)
  $
  and 
  $
    \Delta_x
  $
  with 
  $
    \norm{\Delta_x}=\varepsilon
  $
  such that
  \begin{gather}
    \label{7060_4}
    \theta x + (1 - \theta) y^* = y + \Delta_x
    \,.
  \end{gather}
    It follows
    \begin{align}
      \label{7060_5}
      \begin{split}
      g(y^*)
      \le
      g(y)
      \le
      g(y + \Delta_x)
      &=
      g(
        \theta x + (1 - \theta) y^*
      )
      \\
      &\le
      \theta g(x)
      + 
      (1 - \theta)
      g(y^*)
      <
      g(y^*)
      ,
      \end{split}
    \end{align}
    which is a contradiction.
    The first inequality is due to
    $y^*$ being a local minimum of $g$ in
    $
  y
  \,
  +
  \,
  \varepsilon
  B
    $,
    the second inequality is due to  
    \eqref{7060_0} being true,
    the equality is due to \eqref{7060_4},
    the third inequality is due to the convexity of $g$
    and the strict inequality is due to \eqref{7060_3}.
    Thus every local minimum of $g$ in
    $
  y
  \,
  +
  \,
  \varepsilon
  B
    $
    is also a global minimum.
    %It follows the right-hand side of \eqref{7060_0}.
\end{proof}
\begin{proof}
  The objective function $G$ of the dual satisfies
\begin{gather*}
  G(\lambda,\lambda_0)
  \ 
  :=
  \ 
    \frac{1}{N}
\sum_{i=1} 
  ^N
  \left[ 
    \,
  T_i
  \cdot
  f^*
  \!
  \left( 
\lambda_0
+
\inner
{B(X_i)}
{\lambda}
  \right)
  \ 
-
\ 
  \left( 
\lambda_0
+
\inner
{B(X_i)}
{\lambda}
  \right)
  \,
  \right]
  \ 
+
\ 
\inner
{\delta}
{|\lambda|}
\,.
\end{gather*}
Without the last term, this is a differentiable convex function.

It is well know that a differentiable convex functions $g$ satisfies
  \begin{gather}
    \label{cv:ts:concD}
    g(x)-g(y)
    \ge
    \nabla
    g(y)^\top
    (x-y)
    \qquad 
    \text{for all}\ 
    x,y\,.
  \end{gather}
  The gradient of
  \begin{gather}
    g := 
    (\lambda,\lambda_0)
    \mapsto
    \frac{1}{N}
\sum_{i=1} 
  ^N
  \left[ 
    \,
  T_i
  \cdot
  f^*
  \!
  \left( 
\lambda_0
+
\inner
{B(X_i)}
{\lambda}
  \right)
  \ 
-
\ 
  \left( 
\lambda_0
+
\inner
{B(X_i)}
{\lambda}
  \right)
  \,
  \right]
  \end{gather}
  is 

  \begin{gather}
    \nabla
    g
    =
    (\lambda,\lambda_0)
    \mapsto
\frac{1}{N}
\sum_{i=1} 
  ^N
  \left[ 
    \,
  T_i
  \cdot
  (f^{'})^{-1}
  \!
  \left( 
\lambda_0
+
\inner
{B(X_i)}
{\lambda}
  \right)
  \ 
-
\ 
1
  \right]
  [
  B(X_i)^\top
  ,
  1
  ]^\top
  \end{gather}
  Thus
\begin{align}
  \label{c:1}
  \begin{split}
     &
   G
     (
     \lambda^*
      +
      \Delta
      ,
      \lambda^\dagger_0
      +
     \Delta_0
     )
     \ 
     -
     \ 
     G
     (
     \lambda^*,
      \lambda^\dagger_0
     )
         %%%%%%%%%%%%% 1 %%%%%%%%%%%%%%
     \\
     &
     \quad
     \ge
     -
     \frac{1}{N}
     \sum_{i=1}^{N} 
     \left[ 
       B(X_i)^\top,
       1
     \right]
     \cdot
     \begin{bmatrix}
       \Delta\\
       \Delta_0
     \end{bmatrix}
     \left( 
       1
       \ 
       -
       \ 
     T_i
     \cdot
     (f^{'})^{-1}
     \left( 
       \inner
       {B(X_i)}
       {\lambda^*}
       +
      \lambda^\dagger_0
     \right)
     \right)
     \\
     &
     \qquad
     +
     \ 
     \inner
     {\delta}
     {
       |\lambda^*+\Delta|
       -
       |\lambda^*|
     }
     \,.
   \end{split}
\end{align}
We fix $
\tilde{\varepsilon}
>0
$
and establish the lower bound
$
-
\tilde{\varepsilon}
$
with probability going to $1$ as $N\to\infty$.
We control the \textbf{first term} by (what?) 
and the \textbf{second term} by $\norm{\delta}$.
\subsection*{First Term}
We note, that by $\norm{B(x)}\le 1$ and the Cauchy-Schwarz inequality
it holds
\begin{gather}
  \label{c:first:0}
      \left[ 
       B(X_i)^\top,
       1
     \right]
     \cdot
     \begin{bmatrix}
       \Delta\\
       \Delta_0
     \end{bmatrix}
     \ 
     \lesssim
     \ 
     \norm{(\Delta,\Delta_0)}
     =\varepsilon
     \,.
\end{gather}
Next, we see that
\begin{align}
  \label{c:first:1}
  \begin{split}
  &
     \frac{1}{N}
     \sum_{i=1}^{N} 
     \left( 
       1
       \ 
       -
       \ 
     T_i
     \cdot
     (f^{'})^{-1}
     \left( 
       \inner
       {B(X_i)}
       {\lambda^*}
       +
      \lambda^\dagger_0
     \right)
     \right)
     \\
     &
     \ 
     \lesssim
     \ 
     \frac{1}{N}
     \sum_{i=1}^{N} 
     \left|
     1
     -
     \frac{T_i}{\pi_i}
     \right|
      \ 
     +
      \ 
     \frac{1}{N}
     \sum_{i=1}^{N} 
     \left| 
        \inner
       {B(X_i)}
       {\lambda^*}
       +
      \lambda^\dagger_0
      \ 
        -
        \ 
        f^{'}
        \left( 
          \frac{1}{\pi_i}
        \right)
     \right|
     \\
     &
     \ 
     =:
     \ 
     S_N
     \
     +
     \ 
     M_N
     \,.
\end{split}
\end{align}
With $
\tilde{\varepsilon}
>0
$
fixed previously,
we want to
establish the upper bound
$
\tilde{\varepsilon}
/
(2\varepsilon)
$
with probability going to $1$ as $N\to\infty$.

First, we bound $S_N$.
By the properties of conditional expectation it holds
\begin{gather*}
  \E
  \left[ 
    \frac{T}{\pi(X)}
  \right]
  =
  \E
  \left[ 
    \frac{\E[T|X]}{\pi(X)}
  \right]
  =1
  \,.
\end{gather*}
By the weak law of large numbers (L1 version ? some assumption on 1/pi?)
\begin{gather}
  \P
  \left[ 
    S_N
    \ge
\tilde{\varepsilon}
/(4\varepsilon)
  \right]
  \ 
  \to 
  \ 
  0
  \qquad 
  \text{for}\ 
  N\to \infty
  \,.
\end{gather}
Next, we bound $M_N$.
Recall that $\sum_{k=1}^{N}B_k(x)=1$. Thus
\begin{gather*}
        \inner
       {B(X)}
       {\lambda^*}
      \ 
       +
      \ 
      \lambda^\dagger_0
      \ 
      =
      \ 
      \sum_{k=1}^{N} 
      B_k(X)
      \left( 
        f^{'}
        \left( 
          \frac{1}{\pi_k}
      \right)
      -
      \lambda_0^\dagger
      \right)
      \ 
      +
      \ 
      \lambda_0^\dagger
      \ 
      =
      \ 
      \sum_{k=1}^{N} 
      B_k(X)
      \cdot
        f^{'}
        \left( 
          \frac{1}{\pi_k}
      \right)
      \,.
\end{gather*}
By Markov's inequality it holds
\begin{align*}
  &
  \P \left[ 
    M_N \ge 
\tilde{\varepsilon}
/(4\varepsilon)
  \right]
  \\
  &
  \ 
  \le
  \ 
  \frac{4\varepsilon}{
\tilde{\varepsilon}
  }
  \, 
  \frac{1}{N}
  \sum_{i=1}^{N} 
  \E
  \left[ 
    \left| 
      \sum_{k=1}^{N} 
      B_k(X_i)
      \cdot
        f^{'}
        \left( 
          \frac{1}{\pi_k}
      \right)
      -
        f^{'}
        \left( 
          \frac{1}{\pi_i}
      \right)
    \right|
  \right]
  \\
  &
  \ 
  \le
  \ 
  \frac{4\varepsilon}{
\tilde{\varepsilon}
  }
  \, 
  \E
  \left[ 
    \left| 
      \sum_{k=1}^{N} 
      B_k(X)
      \cdot
        f^{'}
        \left( 
          \frac{1}{\pi_k}
      \right)
      -
        f^{'}
        \left( 
          \frac{1}{\pi(X)}
      \right)
    \right|
  \right]
  \\
  &
  \ 
  \le
  \ 
  \frac{4\varepsilon}{
\tilde{\varepsilon}
}
  \, 
  \E
  \left[ 
    \left| 
      \sum_{k=1}^{N} 
      B_k(X)
      \cdot
        f^{'}
        \left( 
          \frac{1}{\pi_k}
      \right)
      -
        f^{'}
        \left( 
          \frac{1}{\pi(X)}
      \right)
    \right|
    ^2
  \right]
  ^{1/2}
  \ 
  \to 
  \ 
  0
  \qquad 
  \text{for}\ 
  N\to \infty
  \,.
\end{align*}
The convergence is due to the universal consistency of $B$.
This establishes the desired bound of 
$
\tilde{\varepsilon}/(2\varepsilon)
$
in \eqref{c:first:1}.
Together with \eqref{c:first:0}
we conclude that the \textbf{first term} 
in
\eqref{c:1}
is bounded below by
$
-
\tilde{\varepsilon}/2
$
with probability going to $1$ as $N\to\infty$.
\subsection*{Second Term}
It holds
\begin{gather*}
  |x+y|-|x|\ge
  -|y|
  \qquad
  \text{for all}\ 
  x,y
  \,.
\end{gather*}
Since
$\delta\ge 0$
we get
\begin{align*}
  &
     \inner
     {\delta}
     {
       |\lambda^*+\Delta|
       -
       |\lambda^*|
     }
     \\
     &
     \ 
     \ge
     \ 
     -
     \inner{\delta}
     {|\Delta|}
     \ 
     \ge
     \ 
     -
     \norm{\delta}
     \norm{\Delta}
     \ 
     \ge
     \ 
     -
     \norm{\delta}
     \norm{(\Delta,\Delta_0)}
     \ 
     \ge
     \ 
     -
     \norm{\delta}
     \varepsilon
     \ 
     \ge
     \ 
     -
     \tilde{\varepsilon}/2
     \,,
\end{align*}
with probability going to $1$ as $N\to \infty$.
The convergence is due to $\norm{\delta}$ converging to $0$ in probability.
\subsection*{Conclusion}
With the analysis of the \textbf{first} and \textbf{second term} in
\eqref{c:1} we conclude
\begin{gather}
  G
     (
     \lambda^*
      +
      \Delta
      ,
      \lambda^\dagger_0
      +
     \Delta_0
     )
     \ 
     -
     \ 
     G
     (
     \lambda^*,
      \lambda^\dagger_0
     )
     \ge
     -
     \tilde{\varepsilon}
\end{gather}
with probability going to $1$ as $N\to \infty$.
Since this holds true for all $\tilde{\varepsilon}>0$ we get
\begin{gather}
  G
     (
     \lambda^*
      +
      \Delta
      ,
      \lambda^\dagger_0
      +
     \Delta_0
     )
     \ 
     -
     \ 
     G
     (
     \lambda^*,
      \lambda^\dagger_0
     )
     \ge
     0
\end{gather}
with probability going to $1$ as $N\to \infty$.
But this holds for all 
$
(\Delta,\Delta_0)
$
with 
$
\norm{
(\Delta,\Delta_0)
}
=\varepsilon
$. Thus
\begin{gather}
   \inf _ { 
       \norm{
         (
     \Delta
     ,
     \Delta_0
         )
 } 
= \varepsilon }
     G
     (
     \lambda^*
      +
      \Delta
      ,
      \lambda^\dagger_0
      +
     \Delta_0
     )
     -
     G
     (
     \lambda^*,
      \lambda^\dagger_0
     )
     \ge 
     0
\end{gather}
with probability going to $1$ as $N\to \infty$.
Thus by Lemma~\ref{bw:cd:lem}
\begin{gather}
  \P
    \left[ 
    \norm{
      \lambda^\dagger
      -
      \lambda^*
    }
    \ge
    \varepsilon
    \right]
    \ 
    \to
    \ 
    0
    \qquad
    \text{for}\ 
    N
    \to 
    \infty
    \,.
\end{gather}
Finally, note that this holds for all $\varepsilon>0$. This finishes the proof.
\end{proof}
\section*{Consistency for Inverse Propensitiy Score}
\begin{theorem}
  For all $\varepsilon>0$ it holds
  \begin{align*}
    \P
    \left[ 
  \left| 
  w(X)
  -
  1/\pi(X)
  \right|
  \ge
  \varepsilon
    \right]
    \ 
    \to
    \ 
    0
    \qquad
    \text{for}
    \ 
    N
    \to
    \infty
    \,.
    \intertext{
  Furthermore, it holds
    }
  \E
    \left[ 
  \left| 
  w(X)
  -
  1/\pi(X)
  \right|
  ^2
    \right]
    ^{1/2}
    \ 
    \to
    \ 
    0
    \qquad
    \text{for}
    \ 
    N
    \to
    \infty
    \,.
  \end{align*}
\end{theorem}
\begin{proof}
  We employ the consistency of the dual variable, the universal consistency and boundedness of the regression basis and the constraint on the arithmetic mean of the weights. 
\begin{align}
  \begin{split}
  \left| 
  w(X)
  -
  \frac{1}{\pi(X)}
  \right|
  &
  \ 
  =
  \ 
  \left| 
  (f^{'})^{-1}
  \left( 
    \inner{B(X)}
    {\lambda^\dagger}
    +
    \lambda_0^\dagger
  \right)
  -
  \frac{1}{\pi(X)}
  \right|
  \\
  &
  \ 
  \lesssim
  \ 
  \left| 
  \inner
  {B(X)}
  {\lambda^\dagger-\lambda^*}
  \right|
  +
  \left| 
    \inner{B(X)}
    {\lambda^\dagger}
    +
    \lambda_0^\dagger
    -
    f^{'}
    \left( 
  \frac{1}{\pi(X)}
    \right)
  \right|
  \\
  &
  \ 
  \lesssim
  \ 
  \norm{\lambda^\dagger-\lambda^*}
  +
  \left| 
  \sum_{i=1}^{N} 
  B_k(X)
  \cdot
    f^{'}
    \left( 
  \frac{1}{\pi_k}
    \right)
    -
    f^{'}
    \left( 
  \frac{1}{\pi(X)}
    \right)
  \right|
  \\
  &
  \ 
  \le
  \ 
  \frac{\varepsilon}{2}
+
  \frac{\varepsilon}{2}
  \le
  \varepsilon
  \,,
\end{split}
\end{align}
with probability going to $1$ as $N\to\infty$.

If we prove boundedness, $L_2$-convergence follows readily.
\begin{align*}
    \left| 
  w(X)
  -
  \frac{1}{\pi(X)}
  \right|
  ^2
  &
  \ 
  \le
  \ 
    \left| 
  w(X)
  -
  \frac{1}{C_\pi}
  \right|
  ^2
  \\
  &
  \ 
  \lesssim
  \ 
    \left| 
    \inner
    {B(X)}
    {\lambda^\dagger}
    +
    \lambda_0^\dagger
  -
  f^{'}
  \left( 
  \frac{1}{C_\pi}
  \right)
  \right|
  ^2
  \\
  &
  \ 
  \lesssim
  \left( 
    \mathrm{diam}\,\Theta
  +
  f^{'}
  \left( 
  \frac{1}{C_\pi}
  \right)
  \right)
  ^2
\end{align*}
\end{proof}

\subsection*{Consistency of the Weighted Mean}
\begin{theorem}
  \label{bw:ce:th}
  For all
  $\varepsilon>0$ it holds
  \begin{gather*}
    \P
    \left[ 
      \,
  \left| 
    \frac{1}{N}
    \sum_{i=1}^{n} 
    w_i
    Y_i
    \ 
    -
    \ 
    \E[Y(1)]
  \right|
    \ 
    \ge
    \ 
    \varepsilon
    \right]
    \ 
    \to
    \ 
    0
    \qquad
    \text{for}\ 
    N
    \to 
    \infty
    \,,
  \end{gather*}
  that is, the weighted mean is a consistent estimator.
  Furthermore, it holds for all $p\in [1,\infty)$
  \begin{gather*}
    \E
  \left| 
    \frac{1}{N}
    \sum_{i=1}^{n} 
    w_i
    Y_i
    \ 
    -
    \ 
    \E[Y(1)]
  \right|
  ^p
    \ 
    \to
    \ 
    0
    \qquad
    \text{for}\ 
    N
    \to 
    \infty
    \,.
  \end{gather*}
\end{theorem}
\begin{proof}
  Let
  $
    \mathbf{Y}(1)
  $
  be the vector with $i$-th coordinate
  $
  Y_i(1)
  $,
  that is, the vector of marginal potential outcomes under treatment.
  Note, that $\mathbf{Y}(1)$ is usually unknown. Nevertheless, we can leverage its existence in the following error decomposition.
  Also note, that for $i>n$ we may choose $w_i=1/\pi_i$.
\begin{align*}
  \left| 
    \frac{1}{N}
    \sum_{i=1}^{n} 
    w_i
    Y_i
    -
    \E[Y(1)]
  \right|
  &
  \ 
  \le
  \ 
  \left|  
  \frac{1}{N}
  \left( 
    \sum_{i=1}^{n} 
    w_i
    B(X_i)
    -
    \sum_{i=1}^{N} 
    B(X_i)
  \right)
  ^\top
     \mathbf{Y}(1) 
  \right|
  %%%% 1 %%%%
  \\
  &
  \qquad
  +
  \ 
  \left|  
  \frac{1}{N}
    \sum_{i=1}^{N} 
    \left( T_i\cdot w_i -1 \right)
    \left( 
    \E[Y(1)| X_i]
    -
    \inner
    {B(X_i)}
    { \mathbf{Y}(1) }
    \right)
  \right|
  %%%% 2 %%%%
  \\
  &
  \qquad
  +
  \ 
  \left|  
  \frac{1}{N}
    \sum_{i=1}^{N} 
    T_i
    (
    w_i 
    -
    1/\pi_i
    )
    \left( 
      Y_i
    -
    \E[Y(1)| X_i]
    \right)
  \right|
  %%%% 3 %%%%
  \\
  &
  \qquad
  +
  \ 
  \left|  
  \frac{1}{N}
    \sum_{i=1}^{N} 
    T_i
    /\pi_i
    \left( 
      Y_i
    -
    \E[Y(1)| X_i]
    \right)
    +
    \left( 
    \E[Y(1)| X_i]
    -
    \E[Y(1)]
    \right)
  \right|
  \\
  %%%%%%%%%%%%%%%
  &
  \ 
  =:
  \ 
  R_1
  +
  R_2
  +
  R_3
  +
  R_4
  \end{align*}
  \subsubsection*{Analysis of $R_1$}
  By the Cauchy-Schwarz inequality it holds
  \begin{align*}
   \left|  
  \frac{1}{N}
  \left( 
    \sum_{i=1}^{n} 
    w_i
    B(X_i)
    -
    \sum_{i=1}^{N} 
    B(X_i)
  \right)
  ^\top
     \mathbf{Y}(1) 
  \right|
    \ 
  \le
    \ 
  \norm{\delta}
  \norm{
     \mathbf{Y}(1) 
  }
  \ 
  \lesssim 
  \ 
  \norm{\delta}
  N
  \to
  0
  \qquad
  \text{for}\ 
  N\to \infty
  \,.
  \end{align*}
  \subsubsection*{Analysis of $R_2$}
This calculation will be central to the asymptotic normality.
\begin{align*}
  &
  \P[R_2\ge \varepsilon]
  \\
  &
  \ 
  \le
  \ 
  \varepsilon^{-1}
  \frac{1}{N}\sum_{i=1}^{N} 
  \E
  \left[ 
\left| 
    \left( T_i\cdot w_i -1 \right)
    \left( 
    \E[Y(1)| X_i]
    -
    \inner
    {B(X_i)}
    { \mathbf{Y}(1) }
    \right)
\right|
  \right]
  \\
  &
  \ 
  \le
  \ 
  \varepsilon^{-1}
  \E
  \left[ 
\left| 
    w(X) - 1/\pi(X)
\right|
^2
  \right]
  ^{1/2}
  \E
  \left[ 
\left| 
    \E[Y(1)| X_i]
    -
    \inner
    {B(X_i)}
    { \mathbf{Y}(1) }
\right|
^2
  \right]
  ^{1/2}
  \ 
  \to
  \ 
  0
  \\
  &
\end{align*}
for $N\to\infty$. 
Notice that the rates multiply. This is important for later.
\subsection*{Analysis of $R_3$}
By Theorem?
it holds
\begin{gather*}
  \P[R_3\ge \varepsilon]\le
  \frac{\E[\left| w(X) - 1/\pi(X)\right|^2]^{1/2}}{\varepsilon}
  \ 
  \to
  0
  \qquad
  \text{for}\ 
  N
  \to
  \infty
  \,.
\end{gather*}
\subsection*{Analysis of $R_4$}
\begin{align}
  \begin{split}
  \E
  \left[ 
    Y(T)
    \cdot
    (T\,/\,\pi(X))
  \right]
  &\ =\ 
  \E
  \left[ 
    Y(1)
    \,
    /
    \,
    \pi(X)
    \,
    \vert
    \,
    T=1
  \right]
  \cdot
  \P[T=1]
  \\
  &\ =\ 
  \int_\mathcal{X}
  \E
  \left[ 
    Y(1)
    \,
    \vert
    \,
    X=x
    ,
    T=1
  \right]
  \cdot
  (
  \P[T=1]
  \,
  /
  \,
  \pi(x)
  )
  \,
  \P_{X|T}(dx\,|\,1)
  \\
  &\ =\ 
  \int_\mathcal{X}
  \left[ 
    Y(1)
    \vert
    X=x
  \right]
  \P_X(dx)
  \\
  &
 \  
  =
  \ 
  \E[Y(1)]
  \,.
\end{split}
\end{align}
The first equality holds because of the definition of the weights.
The second, third and last equality stem from 
$
  T\in \left\{ 0,1 \right\}
$,
and the law of total expectation, applied with $T$ and $X$.
The fourth equality is justified by the assumption of conditional unconfoundedness.
The density transformation is due to Bayes's Theorem.
Thus the expectation of the summands in $R_4$ is 0.
Convergence follows from the WLLN.
\subsection*{Conclusion}
We conclude that 
\begin{gather}
  \left| 
    \frac{1}{N}
    \sum_{i=1}^{n} 
    w_i
    Y_i
    -
    \E[Y(1)]
  \right|
  \le \varepsilon
\end{gather}
with probability going to 1 as $N\to \infty$. To show the second statement we calculate
\begin{align}
  \left| 
    \frac{1}{N}
    \sum_{i=1}^{n} 
    w_i
    Y_i
    -
    \E[Y(1)]
  \right|
  ^p
  &
  \ 
  =
  \ 
  \left| 
    \frac{1}{N}
    \sum_{i=1}^{n} 
    w_i
    \left( 
    Y_i
    -
    \E[Y(1)]
    \right)
  \right|
  ^p
  \\
  &
  \ 
  \le
  \ 
    \frac{1}{N}
    \sum_{i=1}^{n} 
    w_i
  \left| 
    \left( 
    Y_i
    -
    \E[Y(1)]
    \right)
  \right|
  ^p
  \ 
  \le
  \ 
  (2M)^p
\end{align}
\end{proof}



%  \subsection*{Consistency of the Dual}
%  \begin{lemma}
  \label{bw:cd:lem}
  Let $m\in\mathbb{N}$ and
  $g \,:\, \R^m \to \overline{\R}$ 
  be convex.
  Then 
  for all $y \in \R^m$ and $\varepsilon>0$ 
    \begin{gather}
      \label{7060_0}
      \inf_{\norm{\Delta}=\varepsilon} g(y+\Delta) - g(y) \ge 0 \quad
    \end{gather}
    implies
    the existence of  
    a global minimum
    $
    y^* \in \,\R^m
    $
    of $g$
    satisfying
    $
      \norm{y^* - y} \le \varepsilon
    $.
\end{lemma}
\begin{proof}
  Since 
  $
  y
  \,
  +
  \,
  \varepsilon
  B
  $
  is convex, it contains a 
  local minimum  
  of $g$.
  Suppose towards a contradiction that
  $
    y^* 
    \ 
    \in 
    \ 
  y
  \,
  +
  \,
  \varepsilon
  B
  $
  is a local minimum, but not a global one, and
  \eqref{7060_0} is true.
  Then it holds
  \begin{gather}
    \label{7060_3}
    g(x) < g(y^*)
    \quad
    \text{for some}\ 
    x 
    \in 
    \R^m 
    \setminus 
    \left( 
  y
  \,
  +
  \,
  \varepsilon
  B
    \right)
  \,.
  \end{gather}
  Furthermore, since 
  $
  y
  \,
  +
  \,
  \varepsilon
  B
  $ is compact and contains $y^*$,
  the line segment connecting 
  $y^*$ and $x$
  intersects the boundary of 
  $y + \mathcal{C}$, that is,
  there exist
  $
    \theta \in (0,1)
  $
  and 
  $
    \Delta_x
  $
  with 
  $
    \norm{\Delta_x}=\varepsilon
  $
  such that
  \begin{gather}
    \label{7060_4}
    \theta x + (1 - \theta) y^* = y + \Delta_x
    \,.
  \end{gather}
    It follows
    \begin{align}
      \label{7060_5}
      \begin{split}
      g(y^*)
      \le
      g(y)
      \le
      g(y + \Delta_x)
      &=
      g(
        \theta x + (1 - \theta) y^*
      )
      \\
      &\le
      \theta g(x)
      + 
      (1 - \theta)
      g(y^*)
      <
      g(y^*)
      ,
      \end{split}
    \end{align}
    which is a contradiction.
    The first inequality is due to
    $y^*$ being a local minimum of $g$ in
    $
  y
  \,
  +
  \,
  \varepsilon
  B
    $,
    the second inequality is due to  
    \eqref{7060_0} being true,
    the equality is due to \eqref{7060_4},
    the third inequality is due to the convexity of $g$
    and the strict inequality is due to \eqref{7060_3}.
    Thus every local minimum of $g$ in
    $
  y
  \,
  +
  \,
  \varepsilon
  B
    $
    is also a global minimum.
    %It follows the right-hand side of \eqref{7060_0}.
\end{proof}

Let $\lambda^*$ denote the vector with coordinates
\begin{gather}
  \lambda^*_i
  :=
  f^{'}(1/\pi_i)
  -
  \lambda^\dagger_0
\end{gather}
It follows from Lemma~\ref{bw:cd:lem}
 \begin{gather*}
   \P
   \left[ 
     \norm
     {
      \lambda ^ \dagger
      -
      \lambda^*
     }
     \le
     \varepsilon
   \right]
   =
   \P
   \left[ 
     \inf _ { 
       \norm{
         (
     \Delta
     ,
     \Delta_0
         )
 } 
= \varepsilon }
     G
     (
     \lambda^*
      +
      \Delta
      ,
      \lambda^\dagger_0
      +
     \Delta_0
     )
     -
     G
     (
     \lambda^*,
      \lambda^\dagger_0
     )
     \ge 
     0
   \right]
   \,.
 \end{gather*}
 
 We separate the differentiable part in $G$ to get

 \begin{align*}
   &
   G
     (
     \lambda^*
      +
      \Delta
      ,
      \lambda^\dagger_0
      +
     \Delta_0
     )
     -
     G
     (
     \lambda^*,
      \lambda^\dagger_0
     )
         %%%%%%%%%%%%% 1 %%%%%%%%%%%%%%
     \\
     &
     \quad
     \ge
     -
     \frac{1}{N}
     \sum_{i=1}^{N} 
     \left[ 
       B(X_i)^\top,
       1
     \right]
     \cdot
     \begin{bmatrix}
       \Delta\\
       \Delta_0
     \end{bmatrix}
     \left( 
       1
       \ 
       -
       \ 
     T_i
     \cdot
     (f^{'})^{-1}
     \left( 
       \inner
       {B(X_i)}
       {\lambda^*}
       +
      \lambda^\dagger_0
     \right)
     \right)
     \\
     &
     \qquad
     +
     \ 
     \inner
     {\delta}
     {
       |\lambda^*+\Delta|
       -
       |\lambda^*|
     }
     %%%% 2 %%%%
     \\
     &
     \quad
     \ge
     \ 
     -
       \norm{
         (
     \Delta
     ,
     \Delta_0
         )
 } 
     \left( 
     \norm{B(X_i)}
     \cdot
     \frac{1}{N}
     \sum_{i=1}^{N} 
     \left| 
       1
       \ 
       -
       \ 
     T_i
     \cdot
     (f^{'})^{-1}
     \left( 
       \inner
       {B(X_i)}
       {\lambda^*}
       +
      \lambda^\dagger_0
     \right)
     \right|
     \ 
     +
     \ 
     \norm{\delta}
     \right)
     %%%%% 2 %%%%
     \\
     &
     \quad
     \ge
     \ 
     -
     \varepsilon
     \left( 
     \frac{1}{N}
     \sum_{i=1}^{N} 
     \left| 
       1
       \ 
       -
       \ 
     T_i
     \cdot
     (f^{'})^{-1}
     \left( 
       \inner
       {B(X_i)}
       {\lambda^*}
       +
      \lambda^\dagger_0
     \right)
     \right|
     \ 
     +
     \ 
     \norm{\delta}
     \right)
     \ 
     =:
     -\varepsilon
     (S+\norm{\delta})
 \end{align*}
 \subsubsection*{Analysis of $S$}
 By the triangle inequality we get
 \begin{gather}
  S
  \le
  \frac{1}{N}
  \sum_{i=1}^{N} 
  |
  1
  -
  T_i
  /
  \pi_i
  |
  +
  \max_{i=1,\ldots,n}
     \left| 
     (f^{'})^{-1}
     \left( 
       \inner
       {B(X_i)}
       {\lambda^*}
       +
      \lambda^\dagger_0
     \right)
     \ 
      -
      \ 
      1/\pi_i
     \right|
     \\
     =:
     S_1
     +
     M
 \end{gather}

 \subsubsection*{Analysis of $S_1$}
 Since
 $X_i$ and $T_i$ are i.i.d. 
 and
 \begin{gather}
   \E[
  T_i
  /
  \pi_i
   ]
   =
   \E[
   \E[
  T_i
  |X_i
   ]
  /
  \pi_i
   ]
   =1
 \end{gather}
 it holds by the weak law of large numbers
 \begin{gather}
   \P[S_1\le 
   \tilde{\varepsilon}/(4\varepsilon)
   ]\to 1
   \quad
   \text{for}\ 
   n\to \infty
 \end{gather}
 \subsubsection*{Analysis of $M$}
 Since  
 \begin{gather}
   \inner{B(X_i)}{\mathrm{1_N}}=1
 \end{gather}
 and
 \begin{gather}
   \lambda^*_i
   =
  f^{'}(1/\pi_i)
  -
  \lambda^\dagger_0
 \end{gather}
 it holds
 \begin{gather}
       \inner
       {B(X_i)}
       {\lambda^*}
       +
      \lambda^\dagger_0
     =
    \sum_{k=1}^{N}  
  f^{'}(1/\pi_i)
  \cdot
       B_k(X_i)
 \end{gather}

By the universal consistency of $B$
this converges to $
f^{'}
(
1/\pi_i
)
$ in probability.
By the continuity of 
$
(
f^{'}
)^{-1}
$
it follows

 \begin{gather}
   \P[M\le 
   \tilde{\varepsilon}/(4\varepsilon)
   ]\to 1
   \quad
   \text{for}\ 
   n\to \infty
 \end{gather}
 \subsubsection*{Conclusion}
 It follows
 \begin{gather}
   \P[S\le
   \tilde{\varepsilon}/(2\varepsilon)
   ]\to 1
   \quad
   \text{for}\ 
   n\to \infty
 \end{gather}
We get
for all $\tilde{\varepsilon}>0$
\begin{gather}
   G
     (
     \lambda^*
      +
      \Delta
      ,
      \lambda^\dagger_0
      +
     \Delta_0
     )
     -
     G
     (
     \lambda^*,
      \lambda^\dagger_0
     )
     \ge
     -
     \tilde{\varepsilon}
\end{gather}
with probability going to 1 for $n\to \infty$.
Thus it also holds
\begin{gather}
  \P
  \left[ 
   G
     (
     \lambda^*
      +
      \Delta
      ,
      \lambda^\dagger_0
      +
     \Delta_0
     )
     -
     G
     (
     \lambda^*,
      \lambda^\dagger_0
     )
     \ge
     0
  \right]
  \to
  0
  \qquad
  \text{for}
  \ 
  n\to \infty
  \,.
\end{gather}
This finishes the proof.




%  \subsection*{Consistency of the Primal}
%   \begin{align*}
   \left| 
   1 / \pi_i
   -
   w_i
   \right|
   & 
   \ 
   =
   \ 
   \left| 
   1 / \pi_i
   -
   (
      f ^ { ' }
   )
   ^ {-1}
   \left( 
   \inner
   { B(X_i)}
   { \lambda ^ \dagger }
   +
   \lambda^\dagger_0
   \right)
   \right|
   %%%% 1 %%%%
   \\
   & 
   \ 
   \le
   \ 
   \left| 
   1 / \pi_i
   -
   (
      f ^ { ' }
   )
   ^ {-1}
   \left( 
   \inner
   { B(X_i)}
   { \lambda ^ * }
   +
   \lambda^\dagger_0
   \right)
   \right|
      %%%% 1 a %%%%
   \\
   &
   \qquad
   +
   \ 
\left| 
   (
      f ^ { ' }
   )
   ^ {-1}
   \left( 
   \inner
   { B(X_i)}
   { \lambda ^ \dagger }
   +
   \lambda^\dagger_0
   \right)
   -
   (
      f ^ { ' }
   )
   ^ {-1}
   \left( 
   \inner
   { B(X_i)}
   { \lambda ^ *}
   +
   \lambda^\dagger_0
   \right)
   \right|
      %%%% 2 %%%%
   \\
   &
   \ 
   =:
   I_1 +
   I_2
 \end{align*}
 By the universal consistency of $B$ and the continuity of $(f^{'})^{-1}$ the first term converges to $0$ in probability.
 By the continuity of $(f^{'})^{-1}$,
 the uniform boundedness of $\norm{B}$ and the dual consistency
 the second term converges to $0$ in probability.


%  \subsection*{Consistency of the Weighted Mean}
%  \begin{align*}
  \left| 
    \frac{1}{N}
    \sum_{i=1}^{n} 
    w_i
    Y_i
    -
    \E[Y(1)]
  \right|
  &
  \ 
  \le
  \ 
  \left|  
  \frac{1}{N}
  \left( 
    \sum_{i=1}^{n} 
    w_i
    B(X_i)
    -
    \sum_{i=1}^{N} 
    B(X_i)
  \right)
  ^\top
     \mathbf{Y}(1) 
  \right|
  %%%% 1 %%%%
  \\
  &
  \qquad
  +
  \ 
  \left|  
  \frac{1}{N}
    \sum_{i=1}^{n} 
    w_i
    \left( 
    \E[Y(1)| X_i]
    -
    \inner
    {B(X_i))}
    { \mathbf{Y}(1) }
    \right)
  \right|
  %%%% 1 %%%%
  \\
  &
  \qquad
  +
  \ 
  \left|  
  \frac{1}{N}
    \sum_{i=1}^{N} 
    \left( 
    \E[Y(1)| X_i]
    -
    \inner
    {B(X_i))}
    { \mathbf{Y}(1) }
    \right)
  \right|
  %%%% 2 %%%%
  \\
  &
  \qquad
  +
  \ 
  \left|  
  \frac{1}{N}
    \sum_{i=1}^{n} 
    (
    w_i 
    -
    1/\pi_i
    )
    \left( 
      Y_i
    -
    \E[Y(1)| X_i]
    \right)
  \right|
  %%%% 3 %%%%
  \\
  &
  \qquad
  +
  \ 
  \left|  
  \frac{1}{N}
    \sum_{i=1}^{N} 
    T_i
    /\pi_i
    \left( 
      Y_i
    -
    \E[Y(1)| X_i]
    \right)
    +
    \left( 
    \E[Y(1)| X_i]
    -
    \E[Y(1)]
    \right)
  \right|
  \\
  %%%%%%%%%%%%%%%
  &
  \ 
  =:
  \ 
  R_1
  +
  R_2
  +
  R_3
  +
  R_4
  +
  R_5
  \end{align*}
  \subsubsection*{Analysis of $R_1$}
  By the Cauchy-Schwarz inequality it holds
  \begin{gather}
    R_1
    \ 
    =
    \ 
   \left|  
  \frac{1}{N}
  \left( 
    \sum_{i=1}^{n} 
    w_i
    B(X_i)
    -
    \sum_{i=1}^{N} 
    B(X_i)
  \right)
  ^\top
     \mathbf{Y}(1) 
  \right|
  \le
  \norm{\delta}
  \norm{
     \mathbf{Y}(1) 
  }
  \end{gather}

  \subsubsection*{Analysis of $R_2$ and $R_3$}
By the universal consistency of $B$ and 
$
  \frac{1}{N}
    \sum_{i=1}^{n} 
    w_i
    =1
$
the terms
$R_2$ and $R_3$ converge to $0$ in probability.
\subsubsection*{Analysis of $R_4$}
By the boundedness of $Y$ and the primal consistency 
$R_4$ 
converges to $0$ in probability.
\subsubsection*{Analysis of $R_5$}
Since the expectation is $0$, the term $R_5$ 
converges to $0$ in probability by the weak law of large numbers.


%
%  \chapter{Asymptotic Normality and Convergence to Gaussian Bridge}
%  \begin{theorem}
\cite[Theorem~4.3]{Gyorfi2002}
  Under conditions
  the partition estimate has
  \begin{gather}
    \E\norm{m_N-m}^2
    \le
    C_\P
    N^{- \frac{2}{d+2}}
  \end{gather}
\end{theorem}
\begin{theorem}
\cite[Theorem~5.2]{Gyorfi2002}
  Under conditions
  the kernel estimate has
  \begin{gather}
    \E\norm{m_N-m}^2
    \le
    C_\P
    N^{- \frac{2}{d+2}}
  \end{gather}
\end{theorem}

%  \subsection*{Learning Rates for the Dual}
%  \begin{theorem}
  Under conditions
  \begin{gather}
    \P
    \left[ 
      \norm{\lambda^\dagger - \lambda^*}
      \le
      C_\P
      C_\tau
      \varepsilon_n
    \right]
    \ge
    1-\tau
    \,,
  \end{gather}
  where 
  $\varepsilon_n$ is the square root of the basis function Learning rate and $C_\tau$ depends on the Concentration Inequality.
  We need bernstein confidence $\sqrt{\log (1/\tau)}$ to preserve minimal Learning rate for $d=1$.
\end{theorem}

%  \subsection*{Learning Rates for the Primal}
%  \begin{theorem}
  Under conditions
  the weights satisfy
  \begin{gather}
    \E
    [
    |
    w^\dagger(X)
    -
    1/\pi(X)
    |
    ^2
    ]
    ^{1/2}
    \le
    C_\P
  \sqrt{\log(n)}
  n^{-1/(2+d)}
  \end{gather}
  where 
  $varepsilon_n$ depends on the Learning rate of the basis functions
  and the confidence of the dual.
  $C_\P$ depends on the size of the parameter space.
\end{theorem}
\begin{proof}
  \begin{align}
    \E
    [
    |
    w^\dagger(X)
    -
    1/\pi(X)
    |
    ^2
    ]
    ^{1/2}
    &
    \ 
    =
    \ 
    \E
    \left[ 
    \left| 
    (f^{'})^{-1}
    \left( 
      \inner
      {B(X)}
      {\lambda^\dagger}
      +
      \lambda_0^\dagger
    \right)
    -
    1/\pi(X)
    \right|
    ^2
    \right]
    ^{1/2}
    \\
    &
    \ 
    \le
    \ 
    \left| 
    (f^{'})^{-1}
    \right|_L
    \left( 
      I_1 + I_2
    \right)
  \end{align}
  where
  \begin{align}
    I_1
    & 
    \ 
    :=
    \ 
    \left( 
    \E\norm{\lambda^\dagger-\lambda^*}^2
    \right)
    ^{1/2}
    \\
    I_2
    & 
    \ 
    :=
    \ 
      \E
      \left[ 
        \left| 
        \sum_{k=1}^{N} 
        B_k(X)\cdot 
        f^{'}(X_k)
        -
        f^{'}(X)
        \right|
        ^2
      \right]
    ^{1/2}
  \end{align}
  It holds 
  $I_2\le n^{-1/(d+2)}$
  by the lr of the basis.
  To analyse $I_1$ we use the lr of the dual.
  \begin{align}
   I_1
   \le
   C_\tau
   n^{-1/(d+2)}
   +
   \sqrt{\tau}
   \cdot
   \mathrm{diam}\,\Theta
  \end{align}
Note that the Markov confidence $1/\sqrt{\tau}$ is insufficient. 
We need the Bernstein confidence $\sqrt{\log(1/\tau)}$. 
With Bernstein confidence, bounded diameter and $\tau=n^{-2/(d+2)}$
we get
\begin{gather}
  I_1
  \le
  C
  \cdot
  \sqrt{\log(n)}
  n^{-1/(2+d)}
\end{gather}
Thus
\begin{gather}
    \E
    [
    |
    w^\dagger(X)
    -
    1/\pi(X)
    |
    ^2
    ]
    ^{1/2}
    \le
    C
    \cdot
  \sqrt{\log(n)}
  n^{-1/(2+d)}
\end{gather}
\end{proof}

%  \subsection*{Asymptotic Normality of the Weighted Mean}
  \subsection*{Gaussian Bridge}
  \subsection{Tools and Assumptions}

To control the remaining terms $R_3$ and $R_4$ we use empirical processes. We introduce the concept and the results we need in the next paragraphs. 
For an introduction to empirical processes see \cite{Vaart2000}. More advanced techniques are in \cite{vaart2013}.

Let 
$
  \left( 
    \Omega,
    \mathcal{A},
    \P
  \right)
$
be a probability space,
$
  \left( 
    \mathcal{Z},
    \Sigma
  \right)
$
a measurable space, and 
$
  \xi_1,\ldots,\xi_N
  :
  \left( 
    \Omega,
    \mathcal{A},
    \P
  \right)
  \to
  \left( 
    \mathcal{Z},
    \Sigma
  \right)
$
a sample 
of independent and identically-distributed
random variables
with probability distribution $\P_{\!\xi}$.
A family $\mathcal{F}$ of measurable functions 
$
  f:
  \left( 
    \mathcal{Z},
    \Sigma
  \right)
    \to
  \left( 
    \R,
    \mathcal{B}(\R)
  \right)
$
induces a stochastic process by
\begin{gather}
  f
  \ 
  \mapsto
  \ 
  \G_N f 
  \ 
  :=
  \ 
  \frac{1}{\sqrt{n}}
  \sum_{i=1}^{N} 
  (
    f(\xi_i)
    -
    \E_\xi[f]
  \,.
\end{gather}
We call this the  \textbf{empirical process} $\G_N$ indexed by $\mathcal{F}$.
We define the (random) norm
\begin{gather}
  \norm{\G_n}_\mathcal{F}
  :=
  \sup_
        { f \in \mathcal{F}}
        \left|
          \G_N f
        \right|
        .
\end{gather}
\begin{remark}
We stress that 
$
  \norm{\G_n}_\mathcal{F}
$
often ceases to be measurable, even in simple situations~\cite[page 3]{vaart2013}.
To deal with this, we introduce the notion of \textbf{outer expectation} $\E^*$, that is,
\begin{gather}
  \E^*[Z]
  \ 
  :=
  \ 
    \inf
  \left\{ 
    \E[U]
  \ 
  \lvert
  \ 
    U\ge Z,
    \ 
    U:
  \left( 
    \Omega,
    \mathcal{A},
    \P
  \right)
  \to 
  \left( 
    \overline{\R},
    \mathcal{B}(\overline{\R})
  \right)
  \text{measurable and}
  \ 
  \E[U]<\infty
  \right\}
  \,.
\end{gather}
In our application the technical difficulties halt at this point, because we only consider $Z$ with $\E^*[Z]<\infty$. Then there exists a smallest measurable function $Z^*$ dominating $Z$ with
$\E^*[Z]=\E[Z^*]$. Thus, we may assume $Z$ to be measurable in this regard.
\begin{definition}
  We call a class 
  $\mathcal{F}$ of measurable functions 
$\P$-Donsker
if the sequence of processes 
$\left\{ G_N f \colon f\in\mathcal{F}\right\}$
converges in
$l^\infty(\mathcal{F})$
to a tight limit process.
\end{definition}

\begin{theorem}
  Every class $\mathcal{F}$ of measurable functions 
  with
  $
    J
    _{[]}
    (
    1
    ,
    \mathcal{F}
    ,
    L_2(\P)
    )
    <\infty
  $
  is

  $\P$-Donsker, that is,
  the sequence of processes 
$\left\{ G_N f \colon f\in\mathcal{F}\right\}$
  converges 
  in
$l^\infty(\mathcal{F})$
to a Gaussian process with mean 0 and covariance function given by
\begin{gather}
  \mathrm{Cov}(f,g)
  :=
  \E[fg]-\E[f]\E[g]
  \,.
\end{gather}
\end{theorem}
\begin{proof}
  \cite[Theorem~19.5]{Vaart2000}
\end{proof}



\end{remark}
In our application we need concentration inequalities for 
$
  \norm{\G_n}_\mathcal{F}
$.
One easy way is to use maximal inequalities for the expectation together with Markov's inequality. There are also Bernstein-like inequalities for empirical processes. We need to introduce more concepts.

Given two functions $\underline{f}$ and $\overline{f}$, the bracket
$[\underline{f},\overline{f}]$ 
is the set of all functions $f$ with 
$\underline{f}\le f \le \overline{f}$.
We define a
$(\varepsilon, L^{r}(\P))$-bracket
to be a bracket
$[\underline{f},\overline{f}]$ with
$\norm{\overline{f}-\underline{f}}_{ L^r(\P)}< \varepsilon$.
The bracketing number 
$
N_{[\,]}(\varepsilon, \mathcal{F}, L^r(\P))
$
is 
the minimum number of 
$(\varepsilon, L^{r}(\P))$-brackets needed to cover $\mathcal{F}$.
For most classes $\mathcal{F}$ the bracketing number grows to infinity for $\varepsilon\to 0$.
To measure the speed of convergence we introduce the bracketing integral
\begin{gather}
     J
    _{[]}
    (
    \delta
    ,
    \mathcal{F}
    ,
    L_r(\P)
    )
    =
  \int_0^{\delta}
      \sqrt{
        \log 
      N_{[\,]}
\left( \varepsilon, \mathcal{F}_N, L^r(\P) \right)
    }
    \,
    d\varepsilon
    \,.
\end{gather}
An envelope function $F$ of a class $\mathcal{F}$ satisfies 
$|f(x)|\le F(x)< \infty$ for all $f\in\mathcal{F}$ and all $x$.
\begin{theorem}
  \label{th:max_ineq}
  \emph{(Maximal inequality)}
  For any class $\mathcal{F}$ of measurable functions with envelope function $F,$
  \begin{gather}
    \E^*
    _\P
    [
    \norm{
      \G
      _n
      }
      _\mathcal{F}
    ]
    \lesssim
    J
    _{[]}
    (
    \norm{
      F
    }
    _{ L^2(\P)}
    ,
    \mathcal{F}
    ,
    L_2(\P)
    )
    .
  \end{gather}
\end{theorem}
\begin{proof}
  \cite[Corollary~19.35]{Vaart2000}
\end{proof}

\begin{lemma}
  \label{markov_max_lemma}
  Let $(\mathcal{H}_N)$ be a sequence of measurable function classes with envelope functions $(H_N)$.
  If
  \begin{gather}
    J_{[\, ]}
    \left( 
    \norm{F_N}_{ L^2(\P)}
    ,
    \mathcal{H}_N
    ,
     L^2(\P)
    \right)
    \ 
    \to
    \ 
    0
    \qquad
    \text{for}
    \ 
    N
    \to
    \infty
  \end{gather}
  it holds 
  $
  \norm{\G_N}^*_{\mathcal{H}_N}\overset{\P}{\to}0
  $.
\end{lemma}
\begin{proof}
  By Markov's inequality and Theorem~\ref{th:max_ineq} it holds for all $\varepsilon>0$
  \begin{align*}
    &
    \P
    [
  \norm{\G_N}^*_{\mathcal{H}_N}
  \ge
  \varepsilon
    ]
    \\
    &
    \ 
    \le
    \ 
    \varepsilon^{-1}
    \E
    [
  \norm{\G_N}^*_{\mathcal{H}_N}
    ]
    \ 
    =
    \ 
    \varepsilon^{-1}
    \E^*
    [
  \norm{\G_N}_{\mathcal{H}_N}
    ]
    \ 
    \lesssim
    \ 
    \varepsilon^{-1}
    J_{[\, ]}
    \left( 
    \norm{H_N}_{ L^2(\P)}
    ,
    \mathcal{H}_N
    ,
     L^2(\P)
    \right)
    \ 
    \to
    \ 
    0
  \end{align*}
  for $N\to\infty$. 
\end{proof}
A technical lemma for products of function classes.

Define the product of two function classes
\begin{gather*}
  \mathcal{F}\cdot \mathcal{G}
  :=
  \left\{ 
    f\cdot g
    \colon
    f\in\mathcal{F},
    g\in\mathcal{G}
  \right\}\,.
\end{gather*}
\begin{lemma}
  \label{lem_prod_br}
  Let
  $\mathcal{F}$ and $\mathcal{G}$ be two function classes 
  with envelope functions $F$ and $G$ satisfying
  $\norm{F}_\infty,\norm{G}_\infty\le 1$.
  For all $\varepsilon>0$ and all $r\in [1,\infty)$ it holds
  \begin{gather*}
    N_{[\,]}(2\varepsilon,\mathcal{F}\cdot\mathcal{G},\mathrm{L}_r(\P))
    \
    \le
    \ 
    N_{[\,]}(\varepsilon,\mathcal{F},\mathrm{L}_r(\P))
    \cdot
    N_{[\,]}(\varepsilon,\mathcal{G},\mathrm{L}_r(\P))
    \,.
  \end{gather*}
\end{lemma}
\begin{proof}
  Let $f\in\mathcal{F}$ and $g\in\mathcal{G}$.
  We can choose two 
  $(\varepsilon,L^r(\P))$
  brackets
  $[\underline{f},\overline{f}]$
  and
  $[\underline{g},\overline{g}]$
  containing $f$ and $g$ with 
  $\norm{\underline{f}}_\infty,\norm{\overline{f}}_\infty\le\norm{F}_\infty\le 1$
  and
  $\norm{\underline{g}}_\infty,\norm{\overline{g}}_\infty\le\norm{G}_\infty\le 1$.
  We the get an 
  $(2\varepsilon,L^r(\P))$
  $[\underline{h},\overline{h}]$
  bracket, containing $f\cdot g$, by
\end{proof}

\begin{lemma}
  \label{aa:r3:lemma:1}
  Let
$(\varepsilon_N)\subset(0,1]$
be 
a decreasing sequence
with $\varepsilon_N\to 0$ for $N\to\infty$ 
and let Assumption~\ref{aa:assumption:1} hold true
for a sequence of function classes $\mathcal{F}_N$.
Then
\begin{gather*}
  J_{[\,]}(
\norm{F_N}_{L^2(\P)}
,\mathcal{F}_N\cdot\mathcal{F},\mathrm{L}_2(\P))
  \to 0
  \quad
  \text{and}
  \quad
  \norm{\G_N}^*_{\mathcal{F}_N\cdot\mathcal{F}}\overset{\P}{\to}0
  \qquad
  \text{for}\ 
  N\to\infty
  \,.
\end{gather*}
\end{lemma}
\begin{proof}
  By Assumption~\ref{aa:assumption:1}
  and Lemma~\ref{aa:mean:l:br} it holds
for some $k<2$
\begin{gather*}
\norm{F_N}_{L^2(\P)}
\ 
\le
\ 
\varepsilon_N
\quad
\text{and}
\quad
  \log
  N_{[\,]}(\varepsilon,\mathcal{F}_N,\mathrm{L}_2(\P))
  \ 
  \lesssim
  \ 
  \left( 
  \frac{1}{\varepsilon}
  \right)^k
  \quad
  \text{for all}
  \ 
  N\in\mathbb{N}
  \,,
\end{gather*}
and
  \begin{gather*}
    N_{[\,]}
    (
    \varepsilon
    ,
    \mathcal{F}, L^2(\P))
    \ 
    \lesssim
    \ 
    \left( 
      \frac{1}{\varepsilon}
    \right)^2
    \qquad
    \text{for all}
    \ 
    \varepsilon>0
    \,.
  \end{gather*}
  Since $\mathcal{F}_N$ and $\mathcal{F}$ have envelope function smaller 1, we can apply Lemma~\ref{lem_prod_br} to get
  \begin{gather*}
  \log
  N_{[\,]}(\varepsilon,\mathcal{F}_N\cdot\mathcal{F},\mathrm{L}_2(\P))
  \ 
  \lesssim
  \ 
  \left( 
  \frac{1}{\varepsilon}
  \right)^k
  +
  \log
  (1/\varepsilon)
  \ 
  \lesssim
  \ 
  \left( 
  \frac{1}{\varepsilon}
  \right)^k
  \quad
  \text{for all}\ 
  \varepsilon>0
  \,.
  \end{gather*}
  Since 
  $k/2\in(0,1)$
  it holds
\begin{align*}
  J_{[\,]}(
\norm{F_N}_{L^2(\P)}
,\mathcal{F}_N\cdot\mathcal{F},\mathrm{L}_2(\P))
  &
  \ 
=
  \ 
\int_0^{
\norm{F_N}_{L^2(\P)}
}
\sqrt{
  \log
  N_{[\,]}(\varepsilon,\mathcal{F}_N\cdot\mathcal{F},\mathrm{L}_2(\P))
}
\,d\varepsilon
\\
&
\ 
\lesssim
\ 
\int_0^{
  \varepsilon_N
}
  \left( 
  \frac{1}{\varepsilon}
\right)^{k/2}
\,d\varepsilon
\\
&
\ 
=
\ 
\frac{
\varepsilon_N^{1-k/2}
}{1-k/2}
\ 
\to 0
\ 
\qquad
\text{for}
\ 
N\to\infty
\,.
\end{align*}
The second statement follows from Lemma~\ref{markov_max_lemma}
for 
$\mathcal{H}_N:=\mathcal{F}_N\cdot\mathcal{F}$ and
$H_N:=F_N$.
\end{proof}

\begin{lemma}
  \label{ps_weights_lemma}
  Let
  $
  g_1\colon
  \mathcal{X}\to\R
  $
  and
  $
  g_2\colon
  \mathcal{Y}\to\R
  $
  be a measurable functions.
  It holds
  \begin{gather*}
    \E
    \left[
    \frac{T}{\pi(X)}
    g_1(X)
    \right]
    \ 
    =
    \ 
    \E
    \left[
    g_1(X)
    \right]
    \,.
  \end{gather*}
  If Assumption~\ref{aa:assumption:treatment_str_ign} holds true, then
  \begin{gather*}
    \E
    \left[
    \frac{T}{\pi(X)}
    g_2(Y(T))
    \right]
    \ 
    =
    \ 
    \E
    \left[
    f(Y(1))
    \right]
    \,.
  \end{gather*}
\end{lemma}

\begin{assumption}
  \label{aa:assumption:1}
  For any decreasing sequence
  $(\varepsilon_N)$ with $\varepsilon_N\to 0$ for $N\to\infty$,
there exists a sequence of (measurable) function classes
$(\mathcal{F}_N)$
with envelope functions
$(F_N)$,
satisfying 
for some $k<2$
\begin{gather*}
\norm{F_N}_{L^2(\P)}
\ 
\le
\ 
\varepsilon_N
\quad
\text{and}
\quad
  \log
  N_{[\,]}(\varepsilon,\mathcal{F}_N,\mathrm{L}_2(\P))
  \ 
  \lesssim
  \ 
  \left( 
  \frac{1}{\varepsilon}
  \right)^k
  \quad
  \text{for all}
  \ 
  N\in\mathbb{N}
  \,,
\end{gather*}
such that for all $(\lambda,\lambda_0)\in\R^{N+1}$ and for all $N\in\mathbb{N}$ the function
\begin{gather}
  \label{ghost_function}
  \mathcal{X}
  \ 
  \to
  \ 
  \R
  \,,
  \qquad
  x
  \ 
  \mapsto
  \ 
  \mathbf{1}{
    \left\{ 
      \sup_{y\in A_N(x)}
      \left| 
      w(y,\lambda,\lambda_0)
      -
      \frac{1}{\pi(y)}
      \right|
      \,
      \le
      \,
      \varepsilon_N
    \right\}
  }
  \left( 
      w(x,\lambda,\lambda_0)
      -
      \frac{1}{\pi(x)}
  \right)
\end{gather}
is contained in $\mathcal{F}_N$.
\end{assumption}
The next Lemma shows that mild assumptions on the regularity of the inverse propensity score implies Assumption~\ref{aa:assumption:1}.
To this end, we need notation from \cite[§2.7.1]{vaart2013}.
To this end, let for any vector $k\in\mathbb{N}_0^d \ (d\in\R)$
\begin{gather*}
  D^k
  \ :=\ 
  \frac
  {\partial^{\norm{k}_1}}
  {
    \partial^{k_1}x_1
    \cdots
    \partial^{k_d}x_d
  }
  \,,
\end{gather*}
and let $\lfloor a \rfloor$ be the greatest integer smaller than $a>0$.
For $\alpha>0$, a bounded set 
$\mathcal{X}\subset\R^d\ (d\in\mathbb{N})$
and
$M>0$, we define $C^\alpha_M(\mathcal{X})$ to be the space of all continuous functions $f\colon \mathcal{X}\to\R$ with
\begin{gather*}
  \max_{\norm{k}_1\le \alpha}\sup_{x\in\mathcal{X}}
  \left| D^k f(x) \right|
  \ 
  +
  \ 
  \max_{\norm{k}_1=\lfloor \alpha \rfloor}\sup_{x,y}
  \frac
  {
  \left|
  D^k f(x) 
  -
  D^k f(y) 
  \right|
  }
  {
    \norm{x-y}_2^{\alpha-\lfloor \alpha \rfloor}
  }
  \
  \le
  \ 
  M
  \,.
\end{gather*}
where the suprema in the second term are taken over all $x,y$ in the interior of $\mathcal{X}$ with $x\neq y$.
Furthermore, let
\begin{gather*}
  \mathcal{X}^1
  :=
  \left\{ 
    y\in\R^d
    \colon
    \norm{x-y}_2 <1
    \ 
    \text{for some}\ x\in\mathcal{X}
  \right\}
  \,.
\end{gather*}
We will employ the following result.
\begin{lemma}
  \label{vdv_coro}
  Let $\mathcal{P}=\left\{ A_1,A_2,\ldots \right\}$ be a partition of $\R^d$ into bounded, convex sets with non-empty interior, and let $\mathcal{F}$ be a class of functions $f\colon\R^d\to\R$ such that the restrictions $\mathcal{F}_{|A_j}$ belong to $C^\alpha_{M_j}(A_j)$
  for all $j\in\mathbb{N}$.
  Then there exists a constant $K$, depending only on $\alpha$, $V$, $r$ and $d$
  such that
  \begin{gather}
    \label{667}
    \log
    N_{[\,]}
    (
    \varepsilon
    ,
    \mathcal{F}
    ,
    L^r(\mathbf{Q})
    )
    \le
    K
    \left( \frac{1}{\varepsilon} \right)^V
    \left( 
      \sum_{j=1}^{\infty}
      \lambda(A_j^1)^{r/(V+r)}
      M_j^{Vr/(V+r)}
      \mathbf{Q}(A_j)^{V/(V+r)}
    \right)
    ^{(V+r)/r}
  \end{gather}
  for every $\varepsilon>0$, $V\ge d/\alpha$, and probability measure $\mathbf{Q}$.
\end{lemma}
\begin{proof}
  \emph{\cite[Corollary~2.7.4]{vaart2013}}
\end{proof}
\pagebreak
\begin{lemma}
  Let $(\mathcal{P}_N)$ denote a sequence of qubic partitions
  $\mathcal{P}_N=\left\{ A_{N,1},A_{N,2},\ldots \right\}$ 
  of $\R^d$ 
  with decreasing width $(h_N)\subset(0,1]$ such that $h_N\to 0$ for $N\to\infty$.
  Furthermore, assume that there exists
  $\alpha>d/2$, where $\mathcal{X}\subseteq \R^d$, such
  that for 
  $V:=d/\alpha$
 and for all 
$
(j,N)\in\mathbb{N}^2
$
there exists 
$M_{N,j}\ge 1$ such that 
\begin{gather}
  \label{897}
  \frac{1}{\pi(\cdot)}
  \in C^\alpha_{M_{N,j}}(A_{N,j})
  \quad
  \text{and}
  \quad
  \sum_{j=1}^{\infty} 
  M_{N,j}^{2V/(V+2)}
  \P
  [
  X\in A_{N,j}
  ]^{V/(V+2)}
  \ 
  \lesssim
  \ 
  1
  \,.
\end{gather}
It then holds the statement of Assumption~\ref{aa:assumption:1}.
\end{lemma}
\begin{proof}
  We want to employ Lemma~\ref{vdv_coro}. 
  To do this, the crucial observation is, that
  \begin{gather*}
    w(\cdot,\lambda,\lambda_0)
    \quad
    \text{is constant on each cell}\ 
    A_N\in\mathcal{P}_N
    \,.
  \end{gather*}
  Thus, the regularity of the function \eqref{ghost_function}
  on each cell $A_N\in\mathcal{P}_N$ is decided by 
  $1/\pi(\cdot)$.
  Indeed, \eqref{ghost_function} is either 0 if the threshold of $\varepsilon_N$ is exceeded somewhere in the cell, or has the form constant-minus-smooth-function.
In any case, it is continuous and bounded by $\varepsilon_N$.
All its derivatives are 0 (if the threshold is exceeded) or are governed by $1/\pi(\cdot)$.
Thus, it follows from \eqref{897}
\begin{gather}
  \eqref{ghost_function}
  \in C^\alpha_{M_{N,j}}(A_{N,j})
  \quad
  \text{and}
  \quad
  \sum_{j=1}^{\infty} 
  M_{N,j}^{2V/(V+2)}
  \P
  [
  X\in A_{N,j}
  ]^{V/(V+2)}
  \ 
  \lesssim
  \ 
  1
  \,.
\end{gather}
To bound the right-hand-side in \eqref{667} we note
that 
$
\lambda(A_{N,j})=h^d_N
$ 
and thus
$
\lambda(A_{N,j}^1)\lesssim 1
$
for all $(j,N)\in\mathbb{N}^2$.
Thus
\begin{gather*}
  \sum_{j=1}^{\infty} 
  \lambda(A_{N,j}^1)^{2/(V+2)}
  M_{N,j}^{2V/(V+2)}
  \P
  [
  X\in A_{N,j}
  ]^{V/(V+2)}
  \ 
  \lesssim
  \ 
  1
  \,.
\end{gather*}
We get $\eqref{ghost_function}\in\mathcal{F}_N$, where 
$
\mathcal{F}_N
$
restricted to $A_{N,j}$ is 
$
  C^\alpha_{M_{N,j}}(A_{N,j})
$
and satisfies the requirements of Lemma~\ref{vdv_coro}.
Since $V=d/\alpha \in (0,2)$ by $\alpha>d/2$, 
applying Lemma~\ref{vdv_coro} finishes the proof.
\end{proof}
In the next examples we show the concrete applications of this Lemma.
\begin{example}
Let $d\in\R$ and assume $\pi(\cdot)$ follows a logistic regression model. Then there exist $(\beta_0,\beta)\in\R^{d+1}$ such that
\begin{gather*}
  \pi(x)
  \ 
  =
  \ 
  \frac{1}
  {
    1
    +
    \exp
    \left( 
      -
        \beta_0
        -
        \inner{\beta}{x}
    \right)
  }
  \quad
  \text{and}
  \quad
  \frac
  {1}
  {
  \pi(x)
  }
  \ 
  =
  \ 
    1
    +
    \exp
    \left( 
      -
        \beta_0
        -
        \inner{\beta}{x}
    \right)
    \quad
    \text{for all}\ 
    x\in\R^d
    \,.
\end{gather*}
Clearly, by the smoothness of the exponential function, for all 
$\alpha>0$ there exist $(M_{N,j})\ge 1$ such that
$
  1/\pi(\cdot)\in C^\alpha_{M_{N,j}}(A_{N,j})
$.
Assume $\#\mathcal{X}<\infty$, that is, $X$ can take only finitely many values with positive probability.
We write
\begin{gather*}
  J_N
  :=
  \left\{ 
    j\in\mathbb{N}
    \colon
    \P[X\in A_{N,j}]>0
  \right\}
  \,.
\end{gather*}
It holds
$\#J_N\le \# \mathcal{X}<\infty$.
Thus, the following maximum is attained
\begin{gather*}
  \max_{j\in J_N} M_{N,j}
  =:M^*_N
  \,.
\end{gather*}
But the partitions increasingly better fit the support of $X$. Thus
$M^*_N$ is decreasing in $N$, that is,  $\infty>M^*_1\ge M^*_N$.
It follows
\begin{align*}
  \sum_{j=1}^{\infty} 
  M_{N,j}
  ^{2V/(V+2)}
  \P
  [
  X\in A_{N,j}
  ]^{V/(V+2)}
  \ 
  \le
  \ 
  \left( 
  M^*_1
  \right)
  ^{2V/(V+2)}
  \cdot
  \# J_N
  \ 
  \lesssim
  \ 
  1
  \,.
\end{align*}
\end{example}
The next question is if we can extend this to countable sets
$\mathcal{X}$. To focus our attention, we assume (the best case)
$M_{N,j}=1$.
\begin{example}
  We give three examples. Let $\propto$ denote the equal-up-to-a-constant order.
  Without loss of generality we assume $\mathcal{X}=\mathbb{N}$.
  First consider
  for some $\varepsilon\in(0,1)$
  \begin{gather*}
    \P[X=k]
    \propto
    k^{-(1+\varepsilon)}
    \,.
  \end{gather*}
  This (barely) is a distribution. But scaling with $V/(V+2)\in(0,1)$ we cause trouble. Indeed, by $V<2$ it holds
  \begin{gather*}
    (
    1+\varepsilon
    )
    \cdot
    \frac{V}{V+2}
    <1
    \,,
  \end{gather*}
  and thus
  \begin{align*}
  \sum_{j=1}^\infty
  \P
  [
  X\in A_{N,j}
  ]^{V/(V+2)}
  &
  \ 
  =
  \ 
  \lim_{N\to\infty}
  \sum_{j\in J_N}
  \P
  [
  X\in A_{N,j}
  ]^{V/(V+2)}
  \\
  &
  \ 
  =
  \ 
  \sum_{k=1}^{\infty} 
  \P
  [
  X=k
  ]^{V/(V+2)}
  \\
  &
  \ 
  \propto
  \ 
  \sum_{k=1}^{\infty} 
  \left( 
  \frac{1}{k}
  \right)
  ^{(1+\varepsilon)\cdot V/(V+2)}
  =\infty
  \,.
  \end{align*}
  On the other hand, the same arguments applied for
  \begin{gather*}
    \P[X=k]
    \propto
    k^{-((V+2)/V+\varepsilon)}
    \,,
  \end{gather*}
  with $\varepsilon>0$ yield
\begin{gather*}
  \sum_{j=1}^\infty
  \P
  [
  X\in A_{N,j}
  ]^{V/(V+2)}
  \lesssim 1
  \,.
\end{gather*}
Finally, let $X$ be Poisson distributed with parameter $\lambda\in(0,1)$.
Then
\begin{align*}
  \sum_{j=1}^\infty
  \P
  [
  X\in A_{N,j}
  ]^{V/(V+2)}
  &
  \ 
  =
  \ 
  \sum_{k=0}^{\infty} 
  \P
  [
  X=k
  ]^{V/(V+2)}
  \\
  &
  \ 
  =
  \ 
  e^{-\lambda\cdot V/(V+2)}
  \sum_{k=0}^{\infty} 
  \left( 
    \frac{\lambda^k}{k\text{!}}
  \right)
  ^{V/(V+2)}
  \\
  &
  \ 
  \le
  \ 
  \sum_{k=0}^{\infty} 
\left( 
    \lambda
  ^{V/(V+2)}
\right)^k
  =
  \frac{1}
  {
    1-
    \lambda
  ^{V/(V+2)}
  }
  \lesssim 1
  \,.
  \end{align*}

\end{example}
In the next example we show, that 
\begin{gather*}
  \sum_{j=1}^\infty
  \P
  [
  X\in A_{N,j}
  ]^{V/(V+2)}
  \to
  \infty
\end{gather*}
for all continuous distributions of $X$.
\begin{example}
  Let $f_X$ be the probability density of $X$. 
  Then there exists a compact set $K\subset\mathcal{X}\subset \R^d$, such that
  $
  \inf_{x\in K}f_X(x)
  >0
  $. Since $\mathcal{P}_N$ are cubic partitions, it holds for 
  \begin{gather*}
    I_N
    :=
    \left\{ 
      i\in\mathbb{N}\colon
      A_{N,i}\subset K
    \right\}
    \qquad
    \text{that}
    \qquad
    \bigcup_{i\in I_N}A_{N,i}\nearrow K
    \,.
  \end{gather*}
Thus
\begin{align*}
  \sum_{i=1}^\infty
  \P
  [
  X\in A_{N,i}
  ]^{V/(V+2)}
  &
  \ 
  \ge
  \ 
  \sum_{i\in I_N}
  \P
  [
  X\in A_{N,i}
  ]^{V/(V+2)}
  \\
  &
  \ 
  \ge
  \ 
  \inf_{x\in K}f_X(x)
  ^{V/(V+2)}
  \cdot 
  h_N^{d\cdot(V/(V+2)-1)}
  \sum_{i\in I_N}
  \lambda
  \left( 
  A_{N,i}
  \right)
  \\
  &
  \ 
  \to
  \ 
  \infty
  \,.
\end{align*}
This follows from
$
  \sum_{i\in I_N}
  \lambda
  \left( 
  A_{N,i}
  \right)
  \to 
  \lambda(K)>0
$,
  $
  \inf_{x\in K}f_X(x)
  >0
  $,
  $V/(V+2)-1<0$ and $h_N\to 0$.

\end{example}


  $
(\lambda^\dagger,\lambda_0^\dagger)
$
We adapt the error decomposition in \cite[page 27]{Wang2019} 
to estimates of the distribution function $F_{Y(1)}$ of $Y(1)$, 
that is,
\begin{gather}
  F_{Y(1)}
  \ 
  \colon
  \ 
  \R
  \ 
  \to
  \ 
  [0,1]
  \, 
  , 
  \qquad
  z
  \ 
  \mapsto
  \ 
  \P
  [
  Y(1)
  \le
  z
  ]
  \,.
\end{gather}
\begin{assumption}
  \label{aa:assumption:1}
  For any decreasing sequence
  $(\varepsilon_N)$ with $\varepsilon_N\to 0$ for $N\to\infty$,
there exists a sequence of (measurable) function classes
$(\mathcal{F}_N)$
with envelope functions
$(F_N)$,
satisfying 
for some $k<2$
\begin{gather*}
\norm{F_N}_{L^2(\P)}
\ 
\le
\ 
\varepsilon_N
\quad
\text{and}
\quad
  \log
  N_{[\,]}(\varepsilon,\mathcal{F}_N,\mathrm{L}_2(\P))
  \ 
  \lesssim
  \ 
  \left( 
  \frac{1}{\varepsilon}
  \right)^k
  \quad
  \text{for all}
  \ 
  N\in\mathbb{N}
  \,,
\end{gather*}
such that for all $(\lambda,\lambda_0)\in\R^{N+1}$ and for all $N\in\mathbb{N}$ the function
\begin{gather}
  \label{ghost_function}
  \mathcal{X}
  \ 
  \to
  \ 
  \R
  \,,
  \qquad
  x
  \ 
  \mapsto
  \ 
  \mathbf{1}_{
    \left\{ 
      \sup_{y\in A_N(x)}
      \left| 
      w(y,\lambda,\lambda_0)
      -
      \frac{1}{\pi(y)}
      \right|
      \,
      \le
      \,
      \varepsilon_N
    \right\}
  }
  \left( 
      w(x,\lambda,\lambda_0)
      -
      \frac{1}{\pi(x)}
  \right)
\end{gather}
is contained in $\mathcal{F}_N$.
\end{assumption}
\begin{remark}
  We can derive Assumption~\ref{aa:assumption:1} from regularity assumptions on the inverse propensity score 
  $1/\pi(\cdot)$ and the distribution of $X$.
  To this end, we want to employ \cite[Corollary~2.7.4]{vaart2013}.
  To do this, the crucial observation is, that
  \begin{gather*}
    w(\cdot,\lambda,\lambda_0)
    \quad
    \text{is constant on each cell}\ 
    A_N\in\mathcal{P}_N
    \,.
  \end{gather*}
  Thus, the regularity of the function \eqref{ghost_function}
  on each cell $A_N\in\mathcal{P}_N$ is decided by 
  $1/\pi(\cdot)$.
  Thus, we assume that there exists
  $\alpha>d/2$, where $\mathcal{X}\subseteq \R^d$, such
  that for 
  $V:=d/\alpha$.
 and for all 
$
(j,N)\in\mathbb{N}^2
$
there exists 
$M_{N,j}\ge \varepsilon_N$ such that 
\begin{gather*}
  \frac{1}{\pi(\cdot)}
  \in C^\alpha_{M_{N,j}}(A_{N,j})
  \quad
  \text{and}
  \quad
  \sum_{j=1}^{\infty} 
  M_{N,j}^{2V/(V+2)}
  \P
  [
  X\in A_{N,j}
  ]^{V/(V+2)}
  \ 
  \lesssim
  \ 
  1
  \,.
\end{gather*}
Then the function \ref{ghost_function}, restricted to $A_{N,j}$, is in 
$
  C^\alpha_{M_{N,j}}(A_{N,j})
$
for all 
$
(j,N)\in\mathbb{N}^2
$.
Then \cite[Corollary~2.7.4]{vaart2013} gives us the statement 
in
Assumption~\ref{aa:assumption:1}.
\end{remark}
\begin{ftheorem}
  \label{aa:mean:th}
  Under conditions 
the stochastic process
\begin{gather}
    \sqrt{N}
    \left( 
  \frac{1}{N}
    \sum_{i=1}^{n} 
    w_i^\dagger
    \mathbf{1}
    _{\left\{ Y_i\,\le\, z \right\}}
    \ 
    -
    \ 
    F_{Y(1)}(z)
    \right)
    _{z\in\R}
    \,
  \end{gather}
  converges in
  $l^\infty(\R)$
  to a Gaussian process with mean 0 and covariance
\begin{align}
  \label{cov:lp}
 \begin{split}
  &
  \mathbf{Cov}
  (z_1,z_2)
  \\
  &
  =\ 
  \E
  \left[ 
 \frac{
 F_{Y(1)}(z_1 \land z_2\,|\,X)
}{\pi(X)}
\ 
-
\ 
 \frac{1-\pi(X)}{\pi(X)}
 F_{Y(1)}(z_1|X)
 \cdot
 F_{Y(1)}(z_2|X)
  \right]
  \\
  &
  \qquad 
 -
 \ 
 F_{Y(1)}(z_1)
 \cdot
 F_{Y(1)}(z_2)
 \end{split}
\end{align}
\end{ftheorem}

\pagebreak
\begin{lemma}
  \label{aa:mean:lemma_decomp}
  It holds
  \begin{gather}
  \sqrt{N}
\left( 
    \frac{1}{N}
    \sum_{i=1}^{n} 
    w(X_i)
    \mathbf{1}_{\left\{ Y_i \le z \right\}}
    \ 
    -
    \ 
    F_{Y(1)}(z)
    \right)
    _{z\in\R}
    \ 
    =
    \ 
    R_1
    \ 
    +
    \ 
    R_2
    \ 
    +
    \ 
    R_3
    \ 
    +
    \ 
    R_4
  \end{gather}
  with
\begin{align*}
  R_1
  &
  \ 
  :=
  \ 
  \sqrt{N}
  \sum_{k=1}^{N} 
  \left[ 
  \frac{1}{N}
  \left( 
    \sum_{i=1}^{n} 
    w(X_i)
    B_k(X_i)
    -
    \sum_{i=1}^{N} 
    B_k(X_i)
  \right)
  F_{Y(1)}(z|X_k)
  \right]
  _{z\in\R}
  \,,
  %%%% 1 %%%%
  \\
  R_2
  &
  \
  :=
  \ 
  \sqrt{N}
    \sum_{i=1}^{N} 
    \left[ 
  \frac{
    T_i\cdot w(X_i) -1 }{N}
    \left( 
  F_{Y(1)}(z|X_i)
    \ 
    -
    \ 
    \sum_{k=1}^{N} 
    B_k(X_i)
    \cdot
  F_{Y(1)}(z|X_k)
    \right)
    \right]
  _{z\in\R}
  \,,
  %%%% 2 %%%%
  \\
  R_3
  &
  \
  :=
  \ 
  \sqrt{N}
  \left( 
  \frac{1}{N}
    \sum_{i=1}^{N} 
    \left[ 
    T_i
    \left( 
    w(X_i) 
    -
    \frac{1}{\pi(X_i)}
    \right)
    \left( 
    \mathbf{1}_{\left\{ Y_i \le z \right\}}
    -
  F_{Y(1)}(z|X_i)
    \right)
    \right]
  \right)
  _{z\in\R}
  \,,
  %%%% 3 %%%%
  \\
  R_4
  &
  \
  :=
  \ 
  \sqrt{N}
  \left( 
  \frac{1}{N}
    \sum_{i=1}^{N} 
    \frac{T_i}{\pi(X_i)}
    \left( 
    \mathbf{1}_{\left\{ Y_i \le z \right\}}
    -
  F_{Y(1)}(z|X_i)
    \right)
    \ 
    +
    \ 
    \left( 
  F_{Y(1)}(z|X_i)
    -
  F_{Y(1)}(z)
    \right)
  \right)
  _{z\in\R}
  \,.
  \end{align*}
\end{lemma}
\nopagebreak
\begin{proof}
  We fix $z\in\R$.
  It holds
  \begin{align*}
    &
    \frac{1}{N}
    \sum_{i=1}^{N} 
    w^\dagger(X_i)
    \cdot
    T_i
    \cdot
    \mathbf{1}_{\left\{ Y_i(T_i)\, \le\, z \right\}}
    \\
    &
    \ 
    =
    \ 
    \frac{1}{N}
    \sum_{i=1}^{N} 
    \left( 
    w^\dagger(X_i)
    -
    \frac{1}{\pi(X_i)}
    \right)
    T_i
    \cdot
    \mathbf{1}_{\left\{ Y_i(T_i)\, \le\, z \right\}}
    \\
    &
    \quad 
    +
    \ 
    \frac{1}{N}
    \sum_{i=1}^{N} 
    \frac{T_i}{\pi(X_i)}
    \mathbf{1}_{\left\{ Y_i(T_i)\, \le\, z \right\}}
    \\
    &
    \ 
    =
    \ 
    \frac{1}{N}
    \sum_{i=1}^{N} 
    \left( 
    w^\dagger(X_i)
    -
    \frac{1}{\pi(X_i)}
    \right)
    T_i
    \left( 
    \mathbf{1}_{\left\{ Y_i(T_i)\, \le\, z \right\}}
    -
    F_{Y(1)}(z|X_i)
    \right)
    \\
    &
    \quad 
    +
    \ 
    \frac{1}{N}
    \sum_{i=1}^{N} 
    \frac{T_i}{\pi(X_i)}
    \left( 
    \mathbf{1}_{\left\{ Y_i(T_i)\, \le\, z \right\}}
    -
    F_{Y(1)}(z|X_i)
    \right)
    \\
    &
    \qquad 
    +
    \ 
    \frac{1}{N}
    \sum_{i=1}^{N} 
    w^\dagger(X_i)\cdot T_i\cdot
    F_{Y(1)}(z|X_i)
    \\
    &
    \ 
    =
    \ 
    R_3(z)/\sqrt{N}
    \\
    &
    \quad 
    +
    \ 
    \frac{1}{N}
    \sum_{i=1}^{N} 
    \frac{T_i}{\pi(X_i)}
    \left( 
    \mathbf{1}_{\left\{ Y_i(T_i)\, \le\, z \right\}}
    -
    F_{Y(1)}(z|X_i)
    \right)
    +
    \left( 
    F_{Y(1)}(z|X_i)
    -
    F_{Y(1)}(z)
    \right)
    \\
    &
    \qquad
    +
    \ 
    \frac{1}{N}
    \sum_{i=1}^{N} 
    \left( 
    w^\dagger(X_i)\cdot T_i
    \ 
    -
    \ 
    1
    \right)
    F_{Y(1)}(z|X_i)
    \\
    &
    \quad\qquad
    +
    \ 
    F_{Y(1)}(z)
    \\
    &
    \ 
    =
    \ 
    R_3(z)/\sqrt{N}
    \\
    &
    \quad
    +
    \ 
    R_4(z)/\sqrt{N}
    \\
    &
    \qquad
    +
    \ 
    \frac{1}{N}
    \sum_{i=1}^{N} 
    \left( 
    w^\dagger(X_i)\cdot T_i
    \ 
    -
    \ 
    1
    \right)
    \left( 
    F_{Y(1)}(z|X_i)
    -
    \sum_{k=1}^{N} 
    B_k(X_i)
    \cdot
  F_{Y(1)}(z|X_k)
    \right)
    \\
    &
    \quad\qquad
    +
    \ 
    \frac{1}{N}
    \sum_{i=1}^{N} 
    \left( 
    w^\dagger(X_i)\cdot T_i
    \ 
    -
    \ 
    1
    \right)
    \sum_{k=1}^{N} 
    B_k(X_i)
    \cdot
  F_{Y(1)}(z|X_k)
    \\
    &
    \qquad\qquad
    +
    \ 
    F_{Y(1)}(z)
\\
    &
    \ 
    =
    \ 
    R_3(z)/\sqrt{N}
    \\
    &
    \quad
    +
    \ 
    R_4(z)/\sqrt{N}
    \\
    &
    \qquad
    +
    \ 
    R_2(z)/\sqrt{N}
    \\
    &
    \quad\qquad
    +
    \ 
    \sum_{k=1}^{N} 
    \frac{1}{N}
    \sum_{i=1}^{N} 
    \left( 
    w^\dagger(X_i)\cdot T_i
    B_k(X_i)
    \ 
    -
    \ 
    B_k(X_i)
    \right)
    \cdot
  F_{Y(1)}(z|X_k)
    \\
    &
    \qquad\qquad
    +
    \ 
    F_{Y(1)}(z)
    \\
    &
    \ 
    =
    \ 
    \left( 
R_3(z)
    \ 
    +
    \ 
    R_4(z)
    \ 
    +
    \ 
    R_2(z)
    \ 
    +
    \ 
    R_1(z)
    \right)
    /\sqrt{N}
    \ 
    +
    \ 
    F_{Y(1)}(z)
    \,.
  \end{align*}
  This holds for all $z\in\R$.
  Multiplying with $\sqrt{N}$ yields the result.
\end{proof}
Let $F_{Y(1)}(z|x):=\P[Y(1)\le z|X=x]$ denote a conditional version of the distribution function of $Y(1)$ at $x\in\mathcal{X}$.
We also need the propensity score
$
  \pi(x):=\P[T=1|X=x]
$
and the weights function
$
  w(x)
  :=
  (
  f^{'}
  )^{-1}
  \left( 
    \inner{B(x)}{\lambda^\dagger}
    +
    \lambda_0^\dagger
  \right)
$.
\begin{lemma}
  \label{aa:mean:l:r1}
Let the weights function $w$ satisfy the box constraints in 
Problem~\ref{bw:1:primal} and 
$\sqrt{N}\norm{\delta}_1\overset{\P}{\to}0$.
Then it holds
$\sup_{z\in\R}|R_1(z)|\overset{\P}{\to}0$.
  \end{lemma}
\begin{proof}
%  By Theorem~\ref{dual_solution_th}
%  it holds $w_i^\dagger=w(X_i)$ for $i\in \left\{ 1,\ldots,n \right\}$, that is, for $i\le n$ we can identify $w(X_i)$ with the optimal
%  solution to 
%  problem~\ref{bw:1:primal}. 
%  Thus the constraints of the problem apply.
%   Let's bound $R_1$.
It holds
  \begin{align}
    \label{R_1:1}
    \begin{split}
    \sup_{z\in\R}
    \left| 
    R_1(z)
    \right|
    &
    \ 
    =
    \ 
    \sup_{z\in\R}
    \left| 
  \sqrt{N}
  \sum_{k=1}^{N} 
  \left[ 
  \frac{1}{N}
  \left( 
    \sum_{i=1}^{n} 
    w(X_i)
    B_k(X_i)
    -
    \sum_{i=1}^{N} 
    B_k(X_i)
  \right)
  F_{Y(1)}(z|X_k)
  \right]
    \right|
    \\
    &
    \ 
    \le
    \ 
  \sqrt{N}
  \sum_{k=1}^{N} 
  \left| 
  \frac{1}{N}
  \left( 
    \sum_{i=1}^{n} 
    w(X_i)
    B_k(X_i)
    -
    \sum_{i=1}^{N} 
    B_k(X_i)
  \right)
  \right|
    \sup_{z\in\R}
  F_{Y(1)}(z|X_k)
  \\
    &
    \ 
    \le
    \ 
  \sqrt{N}
  \norm{\delta}_1
    \end{split}
  \end{align}
  The last inequality is due to $F_{Y(1)}\in[0,1]$ and the assumption that $(w(X_i))$ satisfies the box constraints of Problem~\ref{bw:1:primal}.
  Since we assume 
$\sqrt{N}\norm{\delta}_1\overset{\P}{\to}0$
it holds
$\sup_{z\in\R}|R_1(z)|\overset{\P}{\to}0$.
\end{proof}
\begin{remark}
  We want to comment on the box constraints of Problem~\ref{bw:1:primal}, that is,
 \begin{gather*}
      \left| 
      \frac{1}{N} 
      \left( 
      \sum_{i = 1}^{n} 
      w(X_i)
      B_k(X_i)
      -
      \sum_{i=1}^{N} 
      B_k(X_i)
      \right)
    \right|
    \ 
    \le 
    \ 
    \delta_k
    \qquad
    \text{for all}\ 
    k \in \left\{ 1, \ldots, N \right\}
    \,.
  \end{gather*}
  Note, that the first sum goes over $\left\{ 1,\ldots,n \right\}$ while the second sum goes over $\left\{ 1,\ldots,N \right\}$.
  A second, equivalent version of the constraints is
  \begin{gather*}
      \left| 
      \frac{1}{N} 
      \left( 
      \sum_{i = 1}^{N} 
      T_i
      w(X_i)
      B_k(X_i)
      -
      \sum_{i=1}^{N} 
      B_k(X_i)
      \right)
    \right|
    \ 
    \le 
    \ 
    \delta_k
    \qquad
    \text{for all}\ 
    k \in \left\{ 1, \ldots, N \right\}
    \,.
  \end{gather*}
  Now both sums go over $\left\{ 1,\ldots,N \right\}$ and the
  indicator of treatment $T_i$ takes care that in the first sum only the terms with $i\le n$ are effective. 
  Having this flexibility with the versions helps. I regard the first version as suitable for non-probabilistic computations, although $n$ is of course a random variable. On the other hand, the second version is more honest, exactly telling the dependence on the indicator of treatment. This version is useful in probabilistic computations. 

  Also we want to comment on the assumption on $\norm{\delta}$.
  Playing around with norm equivalences we discover that 
  $\sqrt{N}\norm{\delta}_1\overset{\P}{\to}0$ for $N\to \infty$ is the weakest
  (natural) assumption to
  control $R_1$.
  Indeed, other ways to continue the second row in \eqref{R_1:1} are
  \begin{gather*}
    (\,\cdots)
    \ 
  \le
    \ 
  \sqrt{N}
  \norm{\delta}_2
  \left( 
  \sum_{k=1}^{N} 
  \left( 
    \sup_{z\in\R}
  F_{Y(1)}(z|X_k)
  \right)^2
\right)^{1/2}
\ 
\le
\ 
N
  \norm{\delta}_2\,,
  \end{gather*}
  by the Cauchy-Schwarz inequality and
  $
  F_{Y(1)}\in [0,1]
  $,
or
\begin{gather*}
  (\,\cdots)
  \ 
  \le
  \ 
  \sqrt{N}
  \norm{\delta}_\infty
  \sum_{k=1}^{N} 
    \sup_{z\in\R}
  F_{Y(1)}(z|X_k)
  \ 
  \le
  \ 
  N^{3/2}
  \norm{\delta}_\infty
  \,.
\end{gather*}
Since $\delta\in \R^N$, however, it holds
\begin{gather*}
  \sqrt{N}\norm{\delta}_1
  \ 
  \le
  \ 
  N\norm{\delta}_2
  \ 
  \le
  \ 
  N^{3/2}\norm{\delta}_\infty
  \,.
\end{gather*}
With hindsight, the assumption 
$\sqrt{N}\norm{\delta}_1\overset{\P}{\to}0$ for $N\to \infty$ 
  also 
  suffices 
  to control the second (or first) occurrence of a term, that we control by assumptions on $\norm{\delta}$.
This is the \textbf{second term} of \eqref{c:1}, where we estimate
\begin{gather*}
  \inner{\delta}{\left| \Delta \right|}
  \ 
  =
  \ 
  \sum_{k=1}^{N} 
  \delta_k
  \left| \Delta_k \right|
  \ 
  \le
  \ 
  \norm{\delta}_1
  \norm{\Delta}_\infty
  \ 
  \le
  \ 
  \norm{\delta}_1
  \norm{\Delta}_2
  \ 
  \le
  \ 
  \norm{\delta}_1
  \varepsilon
  \ 
  \overset{\P}{\to}
  \ 
  0
  \quad
  \text{for}\ 
  N\to \infty
  \,.
\end{gather*}

\end{remark}
  \begin{lemma}
    \label{aa:mean:l:r2}
    Let the conditions of Theorem~\ref{aa:weights:th} hold true.
    Furthermore assume, that the width of the partitioning estimate $h_N$ and a conditional version of the distribution function of $Y(1)$ satisfy
\begin{gather*}
  \sqrt{N}
  \sup_{z\in\R}
  \omega
  \left( 
    F_{Y(1)}(z|\cdot)
    ,h_N
  \right)
  \to
  0
  \qquad
  \text{for}\ 
  N\to \infty
  \,,
\end{gather*}
where $\omega$ is the modulus of continuity.
Then it holds
$\sup_{z\in\R}|R_2(z)|\overset{\P}{\to}0$.
  \end{lemma}
\begin{proof}
\begin{align*}
&
  \sup_{z\in\R}
  \left| 
  R_2(z)
  \right|
  \\
  &
  \ 
  \le
  \ 
  \sqrt{N}
    \sum_{i=1}^{N} 
    \left[ 
  \frac{
    T_i\cdot w(X_i) -1 }{N}
  \sup_{z\in\R}
    \left| 
  F_{Y(1)}(z|X_i)
    \ 
    -
    \ 
    \sum_{k=1}^{N} 
    B_k(X_i)
    \cdot
  F_{Y(1)}(z|X_k)
    \right|
    \right]
    \\
    &
  \ 
    \le
  \ 
    \sqrt{N}
  \sup_{z\in\R}
  \omega
  \left( 
    F_{Y(1)}(z|\cdot)
    ,h_N
  \right)
  \sum_{i=1}^{N} 
  \frac{
    T_i\cdot w(X_i) +1 }{N}
    \\
    &
  \ 
    =
  \ 
    2
    \sqrt{N}
  \sup_{z\in\R}
  \omega
  \left( 
    F_{Y(1)}(z|\cdot)
    ,h_N
  \right)
  \,.
\end{align*}
The equality is due to 
\begin{gather}
  1
  \ 
  =
  \ 
\frac{1}{N}\sum_{i=1}^{n}w_i^\dagger
  \ 
  =
  \ 
\frac{1}{N}\sum_{i=1}^{n}w(X_i)
  \ 
=
  \ 
\frac{1}{N}\sum_{i=1}^{N}T_iw(X_i)
\,,
\end{gather}
that is, $w(X_i)$ satisfy the second constraint of Problem~\ref{bw:1:primal}.
The second inequality follows from 
$\sum_{k=1}^{N}B_k(X)=1$ and the convexity of the absolute value. 
Indeed,
\begin{align*}
&
  \sup_{z\in\R}
  \left| 
  F_{Y(1)}(z|X_i)
  -
  \sum_{k=1}^{N} 
  B_k(X_i)
  \cdot
  F_{Y(1)}(z|X_k)
  \right|
  \\
  &
  \ 
  \le
  \ 
  \sum_{k=1}^{N} 
  \frac{\mathbf{1}_{\left\{ X_k\in A_N(X_i) \right\}}}
  {
    \sum_{j=1}^{N} 
\mathbf{1}_{\left\{ X_j\in A_N(X_i) \right\}}
  }
  \sup_{z\in\R}
  \left| 
  F_{Y(1)}(z|X_i)
  -
  F_{Y(1)}(z|X_k)
  \right|
  \\
  &
  \ 
  \le
  \ 
  \sup_{z\in\R}
  \omega
  \left( 
    F_{Y(1)}(z|\cdot)
    ,h_N
  \right)
  \,.
\end{align*}

Since we assume
\begin{gather*}
  \sqrt{N}
  \sup_{z\in\R}
  \omega
  \left( 
    F_{Y(1)}(z|\cdot)
    ,h_N
  \right)
  \to
  0
  \qquad
  \text{for}\ 
  N\to \infty
  \,,
\end{gather*}
it follows
$\sup_{z\in\R}|R_2(z)|\overset{\P}{\to}0$.
\end{proof}
\begin{remark}
In the original paper \cite{Wang2019} the authors derive concrete learning rates for the weights and employ them in bounding this term. They obtain a multiplied learning rate, which is sufficiently fast. Their approach, however, calls for concrete learning rates of the weights. Arguably, the process of deriving such rates is the most complicated part of the paper. 
I found out, that we don't need concrete rates for the weights. 
Consistency of the weights is enough and gives us an (arbitrarily slow but sufficient) learning rate to establish the results.
We don't even need rates for the weights to control $R_2$.
They only play a role in bounding $R_3$.

We also want to comment on the assumption
\begin{gather*}
  \sqrt{N}
  \sup_{z\in\R}
  \omega
  \left( 
    F_{Y(1)}(z|\cdot)
    ,h_N
  \right)
  \to
  0
  \qquad
  \text{for}\ 
  N\to \infty
  \,,
\end{gather*}
I decided to keep this more general (and abstract) assumption, althogh
there are many (more concrete, yet stronger) assumptions on the regularity of
$
    F_{Y(1)}(z|\cdot)
$
and the convergence speed of $h_N$.
If for example 
$
    F_{Y(1)}(z|\cdot)
$
is $\alpha$-Hölder continuous with $\alpha\in(0,1]$ for all $z\in\R$, it suffices $\sqrt{N}h_N^\alpha\to0$ to control $R_2$.
\end{remark}

To control the remaining terms $R_3$ and $R_4$ we use empirical processes. We introduce the concept and the results we need in the next paragraphs. 
For an introduction to empirical processes see \cite{Vaart2000}. More advanced techniques are in \cite{vaart2013}.

Let 
$
  \left( 
    \Omega,
    \mathcal{A},
    \P
  \right)
$
be a probability space,
$
  \left( 
    \mathcal{Z},
    \Sigma
  \right)
$
a measurable space, and 
$
  \xi_1,\ldots,\xi_N
  :
  \left( 
    \Omega,
    \mathcal{A},
    \P
  \right)
  \to
  \left( 
    \mathcal{Z},
    \Sigma
  \right)
$
a sample 
of independent and identically-distributed
random variables
with probability distribution $\P_{\!\xi}$.
A family $\mathcal{F}$ of measurable functions 
$
  f:
  \left( 
    \mathcal{Z},
    \Sigma
  \right)
    \to
  \left( 
    \R,
    \mathcal{B}(\R)
  \right)
$
induces a stochastic process by
\begin{gather}
  f
  \ 
  \mapsto
  \ 
  \G_N f 
  \ 
  :=
  \ 
  \frac{1}{\sqrt{n}}
  \sum_{i=1}^{N} 
  (
    f(\xi_i)
    -
    \E_\xi[f]
  \,.
\end{gather}
We call this the  \textbf{empirical process} $\G_N$ indexed by $\mathcal{F}$.
We define the (random) norm
\begin{gather}
  \norm{\G_n}_\mathcal{F}
  :=
  \sup_
        { f \in \mathcal{F}}
        \left|
          \G_N f
        \right|
        .
\end{gather}
\begin{remark}
We stress that 
$
  \norm{\G_n}_\mathcal{F}
$
often ceases to be measurable, even in simple situations~\cite[page 3]{vaart2013}.
To deal with this, we introduce the notion of \textbf{outer expectation} $\E^*$, that is,
\begin{gather}
  \E^*[Z]
  \ 
  :=
  \ 
    \inf
  \left\{ 
    \E[U]
  \ 
  \lvert
  \ 
    U\ge Z,
    \ 
    U:
  \left( 
    \Omega,
    \mathcal{A},
    \P
  \right)
  \to 
  \left( 
    \overline{\R},
    \mathcal{B}(\overline{\R})
  \right)
  \text{measurable and}
  \ 
  \E[U]<\infty
  \right\}
  \,.
\end{gather}
In our application the technical difficulties halt at this point, because we only consider $Z$ with $\E^*[Z]<\infty$. Then there exists a smallest measurable function $Z^*$ dominating $Z$ with
$\E^*[Z]=\E[Z^*]$. Thus, we may assume $Z$ to be measurable in this regard.


\end{remark}
In our application we need concentration inequalities for 
$
  \norm{\G_n}_\mathcal{F}
$.
One easy way is to use maximal inequalities for the expectation together with Markov's inequality. There are also Bernstein-like inequalities for empirical processes. We need to introduce more concepts.

Given two functions $\underline{f}$ and $\overline{f}$, the bracket
$[\underline{f},\overline{f}]$ 
is the set of all functions $f$ with 
$\underline{f}\le f \le \overline{f}$.
We define a
$(\varepsilon, L^{r}(\P))$-bracket
to be a bracket
$[\underline{f},\overline{f}]$ with
$\norm{\overline{f}-\underline{f}}_{ L^r(\P)}< \varepsilon$.
The bracketing number 
$
N_{[\,]}(\varepsilon, \mathcal{F}, L^r(\P))
$
is 
the minimum number of 
$(\varepsilon, L^{r}(\P))$-brackets needed to cover $\mathcal{F}$.
For most classes $\mathcal{F}$ the bracketing number grows to infinity for $\varepsilon\to 0$.
To measure the speed of convergence we introduce the bracketing integral
\begin{gather}
     J
    _{[]}
    (
    \delta
    ,
    \mathcal{F}
    ,
    L_r(\P)
    )
    =
  \int_0^{\delta}
      \sqrt{
        \log 
      N_{[\,]}
\left( \varepsilon, \mathcal{F}_N, L^r(\P) \right)
    }
    \,
    d\varepsilon
    \,.
\end{gather}
An envelope function $F$ of a class $\mathcal{F}$ satisfies 
$|f(x)|\le F(x)< \infty$ for all $f\in\mathcal{F}$ and all $x$.
\begin{theorem}
  \label{th:max_ineq}
  \emph{(Maximal inequality)}
  For any class $\mathcal{F}$ of measurable functions with envelope function $F,$
  \begin{gather}
    \E^*
    _\P
    [
    \norm{
      \G
      _n
      }
      _\mathcal{F}
    ]
    \lesssim
    J
    _{[]}
    (
    \norm{
      F
    }
    _{ L^2(\P)}
    ,
    \mathcal{F}
    ,
    L_2(\P)
    )
    .
  \end{gather}
\end{theorem}
\begin{proof}
  \cite[Corollary~19.35]{Vaart2000}
\end{proof}
%
%et 
%  $(\varepsilon_N)\subset(0,1]$ be a decreasing sequence with 
%  $\varepsilon_N\to 0$ as $N\to\infty$ and define for
%  $N\in\mathbb{N}$ and $z\in\R$ the function
% \begin{align*}
%    f_{\varepsilon_N,\pm}^z
%    \ 
%    :
%    \ 
%    &
%      \left\{ 0,1 \right\}
%      \times
%      \mathcal{X}
%      \times
%      \mathcal{Y}
%    \ 
%    \to
%    \ 
%    \R
%    \\
%    &
%      (T,X,Y(T))
%      \ 
%      \mapsto
%      \ 
%      \mathbf{1}
%      _
%      {\left\{ 
%          \left|
%      w(X)-\frac{1}{\pi(X)}
%          \right|
%          \,
%          \le
%          \,
%          \varepsilon_N/2
%      \right\}}
%      \left( 
%      w(X)-\frac{1}{\pi(X)}
%      \right)
%      ^{\pm}
%      T
%      \left( 
%        \mathbf{1}
%        _{\left\{  Y(T)\,\le\,z \right\}}
%        -
%        F_{Y(1)}(z|X)
%      \right)
%      \,,
%  \end{align*}
%  and define 
%  \begin{gather*}
%    \mathcal{F}_N^\pm
%    :=
%    \left\{ 
%    f_{\varepsilon_N,\pm}^z
%    |
%    z\in\R
%    \right\}
%    \,.
%  \end{gather*}
%\begin{lemma}
%  Consider the weights function
%  \begin{align*}
%    w
%    \colon
%    \mathcal{X}
%    \times
%    \R^{N+1}
%    \ 
%    \to
%    \ 
%    \R
%    \,,
%    \quad
%    (X,(\lambda,\lambda_0))
%    \ 
%    \mapsto
%    \ 
%    (
%    f^{'}
%    )^{-1}
%    \left( 
%      \inner{B(X)}{\lambda}
%      +
%      \lambda_0
%    \right)
%    \,.
%  \end{align*}
%    If the optimal solution 
%    $
%    (\lambda^\dagger,\lambda_0^\dagger)
%    $ to Problem? exists and is measurable for all $N\in\mathbb{N}$ 
%    then $
%    w
%    \circ
%    (
%    X
%    ,
%    (\lambda^\dagger,\lambda_0^\dagger)
%    )
%    $
%    is measurable.
%\end{lemma}
%\begin{lemma}
%  \label{bounded_f_lemma}
%    If the optimal solution $(\lambda^\dagger,\lambda_0^\dagger)$ to Problem? exists and is measurable for all $N\in\mathbb{N}$ 
%  the function class
%  $
%    \mathcal{F}_N^\pm
%  $
%    is measurable for all $N\in\mathbb{N}$.
%    Furthermore, $\mathcal{F}_N^\pm$ is bounded above by $\varepsilon_N$ and
%    \begin{gather*}
%     J_{[\, ]}
%    \left( 
%      \varepsilon_N
%    ,
%    \mathcal{F}_N^\pm
%    ,
%     L^2(\P)
%    \right)
%    \ 
%    \to
%    \ 
%    0
%    \qquad
%    \text{for}
%    \ 
%    N
%    \to
%    \infty
%\,. 
%    \end{gather*}
%\end{lemma}
%\begin{proof}
%  If the optimal solution $(\lambda^\dagger,\lambda_0^\dagger)$ to Problem? exists and is measurable for all $N\in\mathbb{N}$ 
%  then the weights function $w$ is measurable. Indeed,
%  this follows from
%  \begin{gather*}
%    w(X)=
%    (
%    f^{'}
%    )
%    ^{-1}
%    \left( \inner{B(X)}{\lambda^\dagger}+\lambda^\dagger_0 \right)
%    \,,
%  \end{gather*}
%  the measurability of the basis functions $B$ and the continuity of $
%    (
%    f^{'}
%    )
%    ^{-1}
%    $. Thus $f_{\varepsilon_N,\pm}^z$ are measurable for all $N\in\mathbb{N}$ and all $z\in\R$. Clearly, 
%    $\mathcal{F}_N^\pm$ is bounded above by $\varepsilon_N$.
%    Since 
%    \begin{gather*}
%       \mathbf{1}
%      _
%      {\left\{ 
%          \left|
%      w(X)-\frac{1}{\pi(X)}
%          \right|
%          \,
%          \le
%          \,
%          \varepsilon_N/2
%      \right\}}
%      \left( 
%      w(X)-\frac{1}{\pi(X)}
%      \right)
%      ^{\pm}
%    \end{gather*}
%    is non-negative and bounded above by $\varepsilon_N$
%    it follows from Lemma~\ref
%    {lem:brack_n}
%  \begin{gather*}
%    N_{[\,]}
%    (
%    \varepsilon
%    ,
%    \mathcal{F}_N^\pm, L^2(\P))
%    \ 
%    \lesssim
%    \ 
%    \left( 
%    \frac
%    {\varepsilon_N}
%    {\varepsilon}
%    \right)^2
%    \ 
%    \lesssim
%    \ 
%    \left( 
%    \frac
%    {1}
%    {\varepsilon}
%    \right)^2
%    \qquad
%    \text{for all}
%    \ 
%    \varepsilon>0
%    \,.
%  \end{gather*}
%Thus
%    \begin{align*}
%     J_{[\, ]}
%    \left( 
%      \varepsilon_N
%    ,
%    \mathcal{F}_N^\pm
%    ,
%     L^2(\P)
%    \right)
%    &
%    \ 
%=
%\  
%\int^{\varepsilon_N}_0
%\sqrt{
%\log
%    N_{[\,]}
%    (
%    \varepsilon
%    ,
%    \mathcal{F}_N^\pm, L^2(\P))
%}
%\, 
%d\varepsilon
%\\
%&
%\ 
%\lesssim
%\ 
%\int^{\varepsilon_N}_0
%\log
%\left( 
%  \frac{1}{\varepsilon}
%\right)
%\, 
%d\varepsilon
%\ 
%=
%\ 
%\varepsilon_N
%-
%\varepsilon_N \log(\varepsilon_N)
%\ 
%\to
%\ 
%0 
%\qquad
%\text{for}\ 
%N\to\infty
%\,.
%    \end{align*}
%\end{proof}
%In the next lemma we get rid of the $\varepsilon_N$ bound
%\begin{lemma}
%
%  \label{lemma_fpm}
% Consider for $z\in\R$ the function
% \begin{align*}
%    f_{\pm}^z
%    \ 
%    :
%    \ 
%      \left\{ 0,1 \right\}
%      \times
%      \mathcal{X}
%      \times
%      \mathcal{Y}
%    &
%    \ 
%    \to
%    \ 
%    \R
%    \\
%      (T,X,Y(T))
%    &
%      \ 
%      \mapsto
%      \ 
%      \left( 
%      w(X)-\frac{1}{\pi(X)}
%      \right)
%      ^{\pm}
%      T
%      \left( 
%        \mathbf{1}
%        _{\left\{  Y(T)\,\le\,z \right\}}
%        -
%        F_{Y(1)}(z|X)
%      \right)
%      \,.
%  \end{align*}
%  It holds
%  $\sup_{z\in\R}\left| \G_N f^z_\pm \right|\overset{\P}{\to}0$.
%  Furthermore, for
% \begin{align*}
%    f_z
%    \ 
%    :
%    \ 
%      \left\{ 0,1 \right\}
%      \times
%      \mathcal{X}
%      \times
%      \mathcal{Y}
%    &
%    \ 
%    \to
%    \ 
%    \R
%    \\
%      (T,X,Y(T))
%    &
%      \ 
%      \mapsto
%      \ 
%      \left( 
%      w(X)-\frac{1}{\pi(X)}
%      \right)
%      T
%      \left( 
%        \mathbf{1}
%        _{\left\{  Y(T)\,\le\,z \right\}}
%        -
%        F_{Y(1)}(z|X)
%      \right)
%      \,.
%  \end{align*}
%  it holds
%  $\sup_{z\in\R}\left| \G_N f_z \right|\overset{\P}{\to}0$.
%\end{lemma}
%\begin{proof}
%  It holds for all $z\in\R$
%  $
%  $
%  \begin{gather*}
%  f^z_\pm
%      (T,X,Y(T))
%      \ 
%      \lesssim
%      \ 
%     \left(
%      w(X)- \frac{1}{\pi(X)}
%      \right)
%      ^\pm
%      \,.
%  \end{gather*}
%  By Theorem~\ref{aa:weights:th}
%  there exists a decreasing sequence 
%  $(\varepsilon_N)\subset(0,1]$
%  with $\varepsilon_N\to 0$
%  and
%  $
%  \P[
%\left| 
%      w(X)- 1/\pi(X)
%\right|
%\le
%\varepsilon_N/2
%  ]
%  \to 1
%  $
%  for $N\to\infty$.
%  Therefore
%  \begin{gather*}
%  \P
%  \left[
%  f^z_\pm
%  \in
%  \mathcal{F}_N^\pm
%  \
%  \forall
%  z\in\R
%  \right]
%    =
%  \P[
%\left| 
%      w(X)- 1/\pi(X)
%\right|
%\le
%\varepsilon_N/2
%  ]
%  \to 1
%  \,.
%  \end{gather*}
%Then, for all $\varepsilon>0$ it holds
%\begin{align*}
%  \P
%  \left[
%  \sup_{z\in\R}
%  \left| 
%  \G_N
%  f^z_\pm
%  \right|
%  \le \varepsilon
%  \right]
%  &
%  \ 
%  \ge
%  \ 
%  \P
%  \left[
%  f^z_\pm
%  \in
%  \mathcal{F}_N^\pm
%  \
%  \forall
%  z\in\R
%  \,,
%  \ 
%  \norm{\G_N}^*_{\mathcal{F}_N^\pm}
%  \le \varepsilon
%  \right]
%  \\
%  &
%  \ 
%  \ge
%  \ 
%  \P
%  \left[
%  f^z_\pm
%  \in
%  \mathcal{F}_N^\pm
%  \
%  \forall
%  z\in\R
%  \right]
%  \ 
%  -
%  \ 
%  \P
%  \left[
%  \norm{\G_N}^*_{\mathcal{F}_N^\pm}
%  \le \varepsilon
%  \right]
%  \\
%  &
%  \ 
%  \to
%  \ 
%  1
%  \,.
%\end{align*}
%The convergence of the last term to 0 is due to 
%Lemma~\ref{bounded_f_lemma} and Lemma~\ref{markov_max_lemma}.
%Thus
%  $\sup_{z\in\R}\left| \G_N f^z_\pm \right|\overset{\P}{\to}0$.
%  Note, that 
%  \begin{gather*}
%\G_N f_z 
%=
%\G_N f^z_+
%-
%\G_N f^z_-
%\qquad
%\text{for all}\ 
%z\in\R
%\,.
%  \end{gather*}
%  Thus by Slutzky's Theorem\cite[Theorem~13.18]{Klenke2020}
%  it holds
%  $\sup_{z\in\R}\left| \G_N f_z \right|\overset{\P}{\to}0$.
%\end{proof}
%%\begin{lemma}
%%  Define for
%%  $z\in\R$
%% \begin{align*}
%%    f_z
%%    \ 
%%    :
%%    \ 
%%      \left\{ 0,1 \right\}
%%      \times
%%      \mathcal{X}
%%      \times
%%      \mathcal{Y}
%%&
%%    \ 
%%    \to
%%    \ 
%%      \R
%%    \\
%%    \left(
%%      T,X,Y(T)
%%    \right)
%%    &
%%      \ 
%%      \mapsto
%%      \ 
%%      \frac{T}{\pi(X)}
%%      \left( 
%%        \mathbf{1}
%%        _{\left\{  Y(T)\,\le\,z \right\}}
%%        -
%%        F_{Y(1)}(z|X)
%%      \right)
%%      \ 
%%      +
%%      \ 
%%      \left( 
%%        F_{Y(1)}(z|X)
%%        -
%%        F_{Y(1)}(z)
%%      \right)
%%      \,,
%%  \end{align*}
%%and consider the function class
%%$
%%  \mathcal{G}
%%  \ 
%%  :=
%%  \ 
%%  \left\{ 
%%    f_z
%%  \ 
%%    |
%%  \ 
%%    z\in\R
%%  \right\}
%%$
%%\begin{gather}
%%    \,.
%%\end{gather}
%%  Then $\mathcal{G}$ is a class of measurable functions 
%%  and it holds
%%  \begin{gather}
%%    N_{[\,]}(\varepsilon,\mathcal{G}, L^2(\P))
%%    \le
%%    ??
%%    \qquad
%%    \text{for all}
%%    \ 
%%    \varepsilon>0
%%    \,.
%%  \end{gather}
%%\end{lemma}
%%\begin{lemma}
%%  \label{lemma_max_ineq}
%%  Consider a function class $\mathcal{F}$ with unit ball
%%  $
%%  B_{\mathcal{F}}
%%  :=
%%  \left\{ 
%%    f\in \mathcal{F}
%%    \colon
%%    \norm{f}_\infty
%%    \le
%%    1
%%  \right\}
%%  $.
%%  Let $(\varepsilon_N)$ be a sequence converging to 0
%%  and let 
%%  $
%%  \left( 
%%    \mathcal{F}_N
%%  \right)
%%    :=
%%    \left( 
%%      C\cdot
%%    \varepsilon_N\cdot B_\mathcal{F}
%%    \right)
%%  $
%%  denote the sequence of scaled unit balls in $\mathcal{F}$.
%%  Assume that 
%%  there exists
%%  $k<2$ such that 
%%  the covering number of the unit ball in $\mathcal{F}$
%%  satisfies
%%  \begin{gather}
%%        \log 
%%      N_{[\,]}
%%\left( \varepsilon, B_\mathcal{F}, L^2(\P) \right)
%%\ 
%%\lesssim
%%\ 
%%\left( \frac{1}{\varepsilon} \right)^k
%%\qquad
%%\text{for all}\ 
%%\varepsilon>0
%%\,.
%%  \end{gather}
%%  Then it holds
%%  $
%%      \norm{G_N}^*_{\mathcal{F}_N}
%%    \ 
%%    \overset{\P}
%%    {\to}
%%    \ 
%%    0
%%  $
%%  for $N\to \infty$. 
%%\end{lemma}
%%\begin{proof}
%%  By maximal inequalities it holds
%%  \begin{align*}
%%    \E^*
%%    \left[ 
%%      \norm{G_N}_{\mathcal{F}_N}
%%    \right]
%%    &
%%      \ 
%%      \lesssim
%%      \ 
%%      J_{[\,]}\left( \varepsilon_N, \mathcal{F}_N, L^2(\P) \right)
%%      \\
%%    &
%%      \ 
%%      =
%%      \ 
%%      \int_0^{\varepsilon_N}
%%      \sqrt{
%%        \log 
%%      N_{[\,]}
%%\left( \varepsilon, \mathcal{F}_N, L^2(\P) \right)
%%    }
%%    \,
%%    d\varepsilon
%%    \\
%%    &
%%    \ 
%%    =
%%    \ 
%%      \int_0^{\varepsilon_N}
%%      \sqrt{\log 
%%      N_{[\,]}
%%\left( \varepsilon/(C\cdot \varepsilon_N), B_\mathcal{F}, L^2(\P) \right)
%%    }
%%    \,
%%    d\varepsilon
%%    \\
%%    &
%%    \ 
%%    \lesssim
%%    \ 
%%      \int_0^{\varepsilon_N}
%%      \left( 
%%      \frac{\varepsilon_N}{\varepsilon}
%%    \right)^{k/2}
%%    \,
%%    d\varepsilon
%%    \\
%%    &
%%    \ 
%%    =
%%    \ 
%%  \varepsilon_N^{k/2}
%%  \frac{1}{1-k/2}
%%  \varepsilon_N^{1-k/2}
%%  \\
%%    &
%%    \ 
%%  \lesssim
%%    \ 
%%  \varepsilon_N
%%  \\
%%    &
%%    \ 
%%  \to
%%  \ 
%%  0
%%  \qquad
%%  \text{for}\ 
%%  N\to
%%  \infty
%%  \,.
%%  \end{align*}
%%  Note, that $k<2$.
%%  By the boundedness of $\E^*$ there is no measurability problem.
%%  By Markov's Inequality it holds
%%  \begin{align*}
%%    \P
%%    \left[ 
%%      \norm{G_N}^*_{\mathcal{F}_N}\ge \varepsilon
%%    \right]
%%    \le
%%    \varepsilon^{-1}
%%    \,
%%    \E^*
%%    \left[ 
%%      \norm{G_N}_{\mathcal{F}_N}
%%    \right]
%%    \ 
%%    \to
%%    \ 
%%    0
%%    \qquad
%%    \text{for}\ 
%%    N\to\infty
%%    \,.
%%  \end{align*}
%%\end{proof}
%%The next two lemmas connect $R_3$ to the theory of empirical processes.
%%\begin{lemma}
%%  \label{aa:mean:l:fz}
%%  Consider 
%%  the (random) function
%%  $
%%  f_D^z
%%  $ given by
%%  \begin{gather}
%%    f_{D}^z(T,X,Y(T))
%%    :=
%%    T
%%    \left( 
%%    w(D,X)- \frac{1}{\pi(X)}
%%    \right)
%%    \left( 
%%    \mathbf{1}_{\left\{ Y(T) \le z \right\}}
%%    -
%%  F_{Y(1)}(z|X)
%%    \right)
%%    \,.
%%  \end{gather}
%%  Assume that 
%%  there exists a function class $\mathcal{F}$ satisfying the requirements of Lemma~\ref{lemma_max_ineq} and that
%%  $
%%  f_D^z
%%  \in \mathcal{F}
%%  $
%%  for all $z\in\R$ almost surely.
%%  It then holds
%%  $
%%  \sup_{z\in\R} \left| G_Nf_D^z \right|
%%  \overset{\P}{\to}0$ for 
%%  $N\to\infty$.
%%\end{lemma}
%%\begin{proof}
%%  By the consistency of the weights there exists a learning rate $(\varepsilon_N)$ such that
%%  \begin{gather}
%%    \P
%%    \left[ 
%%      \left| 
%%      w(X,D)
%%      -
%%      \frac{1}{\pi(X)}
%%      \right|
%%      \le
%%      \varepsilon_N
%%    \right]
%%    \to 1 
%%    \qquad
%%    \text{for}\ 
%%    N\to \infty
%%    \,.
%%  \end{gather}
%%  Let
%%  $\mathcal{F}_N:=\varepsilon_NB_\mathcal{F}$ as in 
%%  Lemma~\ref{lemma_max_ineq}.
%%  It holds
%%\begin{gather}
%%      \sup_{z\in\R}
%%      \left| 
%%    f_D^z
%%      \right|
%%      \lesssim
%%      \left| 
%%      w(X,D)
%%      -
%%      \frac{1}{\pi(X)}
%%      \right|
%%      \le
%%      \varepsilon_N
%%\end{gather}
%%with probability going to 1 as $N\to\infty$.
%%Thus
%% \begin{gather}
%%    \P
%%    \left[ 
%%    f_D^z
%%    \in
%%  \mathcal{F}_N
%%  \ \forall\,z\in\R
%%    \right]
%%    =
%%    \P
%%    \left[ 
%%      \sup_{z\in\R}
%%      \left| 
%%    f_D^z
%%      \right|
%%      \lesssim
%%      \varepsilon_N
%%    \right]
%%    \to 1
%%    \qquad
%%    \text{as}
%%    \ 
%%    N\to\infty
%%    \,.
%%  \end{gather}
%%  Then it holds
%%  for all $\varepsilon>0$ 
%%  \begin{align*}
%%    \P
%%    \left[ 
%%      \sup_{z\in\R}
%%  \left| G_Nf_D^z \right|
%%  \le
%%  \varepsilon
%%    \right]
%%    &
%%    \ 
%%    \ge
%%    \ 
%%    \P
%%    \left[ 
%%      \sup_{z\in\R}
%%  \left| G_Nf_D^z \right|
%%  \le
%%  \norm{G_N}^*_{\mathcal{F}_N}
%%  \le
%%  \varepsilon
%%    \right]
%%    \\
%%    &
%%    \ 
%%    \ge
%%    \ 
%%    \P
%%    \left[ 
%%    f_D^z
%%    \in
%%  \mathcal{F}_N
%%  \ \forall\,z\in\R
%%      \ 
%%      \text{and}\ 
%%  \norm{G_N}^*_{\mathcal{F}_N}
%%  \le
%%  \varepsilon
%%    \right]
%%    \\
%%    &
%%    \ 
%%    \ge
%%    \ 
%%    \P
%%    \left[ 
%%    f_D^z
%%    \in
%%  \mathcal{F}_N
%%  \ \forall\,z\in\R
%%    \right]
%%    -
%%    \P
%%    \left[ 
%%  \norm{G_N}^*_{\mathcal{F}_N}
%%    \ 
%%  \ge
%%    \ 
%%  \varepsilon
%%    \right]
%%    \\
%%    &
%%    \ 
%%    \to
%%    \ 
%%    1
%%    \,.
%%  \end{align*}
%%  The convergence of the second term is due to Lemma~\ref{lemma_max_ineq}.
%%
%%\end{proof}
%%
%\begin{lemma}
%  \label{aa:mean:l:r3}
%  Assume conditional unconfoundedness, that is,
%  \begin{gather}
%  (Y(0),Y(1))\perp T \ |\ X
%  \,.
%  \end{gather}
%  Then for all
%  $z\in\R$
%  it holds
%  $
%  G_Nf_z=R_3(z)
%  $.
%  Furthermore, under conditions it holds
%  $\sup_{z\in\R}\left| R_3(z) \right|\overset{\P}{\to}0$.
%\end{lemma}
%\begin{proof}
%  A standard computation shows
%  \begin{gather}
%    \E
%    \left[ 
%    \frac{T}{\pi(X)}
%    \left( 
%    \mathbf{1}_{\left\{ Y(T) \le z \right\}}
%    -
%  F_{Y(1)}(z|X)
%    \right)
%    \right]
%    =0
%    \,.
%  \end{gather}
%  Furthermore
%  \begin{align*}
%    &
%    \E
%    \left[ 
%      Tw(X,D)
%    \left( 
%    \mathbf{1}_{\left\{ Y(T) \le z \right\}}
%    -
%  F_{Y(1)}(z|X)
%    \right)
%    \right]
%    \\
%    &
%    \ 
%    =
%    \ 
%    \E
%    \left[ 
%      \E
%      \left[ 
%      w(X,D)
%      \left( 
%    \mathbf{1}_{\left\{ Y(1) \le z \right\}}
%    -
%  F_{Y(1)}(z|X)
%      \right)
%  |T=1,X,D
%      \right]
%    \right]
%    \\
%    &
%    \ 
%    =
%    \ 
%    \E
%    \left[ 
%      w(X,D)
%      \E
%      \left[ 
%    \mathbf{1}_{\left\{ Y(1) \le z \right\}}
%    -
%  F_{Y(1)}(z|X)
%  |X,D
%      \right]
%    \right]
%    \\
%    &
%    \ 
%    =
%    \ 
%    \E
%    \left[ 
%      w(X,D)
%      \E
%      \left[ 
%    \mathbf{1}_{\left\{ Y(1) \le z \right\}}
%    -
%  F_{Y(1)}(z|X)
%  |X
%      \right]
%    \right]
%    \\
%    &
%    \ 
%    =
%    \ 
%    0
%  \end{align*}
%  The second equality is due to the assumption of 
%  $(Y(0),Y(1))\perp T |X$.
%  The third equality is due to $X\perp D$.
%  Thus
%  $
%    \E
%    f_D^z
%    =
%    0
%  $
%\end{proof}
%

Next, we define some auxiliary functions.
  For $z\in\R$ we define the function
  \begin{align*}
    f_z
    \ 
    :
    \ 
      \left\{ 0,1 \right\}
      \times
      \mathcal{X}
      \times
      \mathcal{Y}
    &
    \ 
    \to
    \ 
    \R
    \\
      (t,x,y)
    &
      \ 
      \mapsto
      \ 
      t
      \left( 
        \mathbf{1}
        _{\left\{  y\,\le\,z \right\}}
        -
        F_{Y(1)}(z|x)
      \right)
      \,,
  \end{align*}
  and the function classes
  \begin{gather}
    \label{F_g}
    \begin{split}
    \mathcal{F}
    &
    \ 
    :=
    \ 
    \left\{ 
      f_z
      \ 
      |
      \ 
      z\in\R\ 
    \right\}
    \\
    \mathcal{G}
    &
    \ 
    :=
    \ 
    \left\{ 
      \frac{f_z}{\pi(\cdot)}
      +
      F_{Y(1)}(z|\cdot)
      -
      F_{Y(1)}(z)
      \ 
      \colon
      \ 
      z\in\R\ 
    \right\}
    \,.
    \end{split}
  \end{gather}
  \begin{assumption}
    \label{aa:assumption:treatment_str_ign}
    It holds 
    \begin{gather*}
    (Y(0),Y(1))\, \perp\,  T \ |\  X
    \qquad
    \text{and}
    \qquad
      0
      \ 
      <
      \ 
      \pi(x)
      \ 
      <
      \ 
      1
      \quad
      \text{for all}\ 
      x\in\mathcal{X}
      \,.
    \end{gather*}
  \end{assumption}
  \begin{lemma}
    \label{aa:mean:r3:lem:fz_E}
    It holds
    $f_z(T,X,Y(T))\in L^1(\P)$
    and 
    $f_z(T,X,Y(T))\perp D_N$
    for all $z\in\R$.
    If also Assumption~\ref{aa:assumption:treatment_str_ign} holds,
    then 
    for all $z\in\R$
    \begin{gather*}
      \E
      \left[
        f_z
        \left( 
          T,
          X,
          Y(T)
        \right)
        \,
        |
        \,
        X
      \right]
      \ 
      =
      \ 
      0
      \qquad
      \text{almost surely.}
    \end{gather*}
  \end{lemma}
  \begin{proof}
    Since $f_z$ is bounded by 1,
    it holds
    $f_z(T,X,Y(T))\in L^1(\P)$.
    Since
    \begin{gather*}
    (T,X,Y(T))
    \ 
    \perp 
    \ 
    D_N
    \ 
    =
    \ 
    (T_i,X_i)_{i\in \left\{ 1,\ldots,N \right\}}
    \end{gather*}
    it holds
    $f_z(T,X,Y(T))\perp D_N$
    for all $z\in\R$.
    For the third statement, note that
    \begin{align*}
      \E
      \left[
        f_z
        \left( 
          T,
          X,
          Y(T)
        \right)
        \,
        |
        \,
        X
      \right]
      &
      \ 
      =
      \ 
      \E
      \left[
      T
      \left( 
        \mathbf{1}
        _{\left\{  Y(T)\,\le\,z \right\}}
        -
        F_{Y(1)}(z|X)
      \right)
        \,
        |
        \,
        X
      \right]
      \\
      &
      \ 
      =
      \ 
      \E
      \left[
        \mathbf{1}
        _{\left\{  Y(1)\,\le\,z \right\}}
        -
        F_{Y(1)}(z|X)
        \,
        |
        \,
        X
        ,
        T=1
      \right]
      \pi(X)
      \\
      &
      \ 
      =
      \ 
      \left( 
      \E
      \left[
        \mathbf{1}
        _{\left\{  Y(1)\,\le\,z \right\}}
        \,
        |
        \,
        X
      \right]
      \ 
        -
      \ 
        F_{Y(1)}(z|X)
      \right)
      \pi(X)
      \\
      &
      \ 
      =
      \ 
      0
      \qquad
      \text{almost surely.}
    \end{align*}
    The third equality is due to Assumption~\ref{aa:assumption:treatment_str_ign}.
  \end{proof}
The next lemma provides bracketing numbers for specific function classes needed to control $R_3$ and $R_4$.
\newpage
\begin{lemma}
  \label{aa:mean:l:br}
  The function class $\mathcal{F}$ and $\mathcal{G}$ defined in \eqref{F_g} are measurable.
  Furthermore, 
  \begin{gather*}
    N_{[\,]}
    (
    \varepsilon
    ,
    \mathcal{F}, L^2(\P))
    \ 
    \lesssim
    \ 
    \left( 
      \frac{1}{\varepsilon}
    \right)^2
    \qquad
    \text{for all}
    \ 
    \varepsilon>0
    \,.
  \end{gather*}
  If $1/\pi(X)\in L^2(\P)$, it also holds 
  \begin{gather*}
    N_{[\,]}
    (
    \varepsilon
    ,
    \mathcal{G}, L^2(\P))
    \ 
    \lesssim
    \ 
    \left( 
    \frac{
      1+
    \norm{1/\pi(X)}_{ L^2(\P)}
    }
    {\varepsilon}
    \right)^4
    \qquad
    \text{for all}
    \ 
    \varepsilon>0
    \,.
  \end{gather*}
%  Furthermore, 
%  consider the function class
%  \begin{gather*}
%    \mathcal{G}
%    :=
%    \left\{ 
%      f_{1/\pi}^z
%      +
%        F_{Y(1)}(z|\cdot)
%        -
%        F_{Y(1)}(z)
%      \ 
%      |
%      \ 
%      z\in\R
%    \right\}
%    \,.
%  \end{gather*}
%  If $1/\pi(X)\in  L^2(\P)$
%  it holds
%  \begin{gather}
%    N_{[\,]}(\varepsilon,\mathcal{G}, L^2(\P))
%    \le
%    ??
%    \qquad
%    \text{for all}
%    \ 
%    \varepsilon>0
%    \,.
%  \end{gather}
\end{lemma}
\begin{proof}
  As in \cite[Example~19.6]{Vaart2000}
  we choose for
  $\varepsilon>0$ and $m\in\mathbb{N}$
  \begin{gather*}
  -\infty=z_0\ <\ z_1\ <\ \cdots\ <\ z_{m-1}\ <\ z_m=\infty
  \,
  \end{gather*}
  such that
  \begin{gather}
    \label{size_z}
    \P
    \left[ 
      Y(1)\in \left[ z_{l-1},z_l \right]\,
    \right]
    \ 
    \le
    \ 
    \varepsilon
    \qquad
    \text{for all}\ 
    l\in \left\{ 1,\ldots,m \right\}
  \end{gather}
  and $m \le 2/\varepsilon$.
  Next, we define $m$ brackets by
\begin{align*}
  \overline{f_l}
  (t,x,y)
  &
  \ 
  :=
  \ 
      t
      \left( 
        \mathbf{1}
        _{\left\{  y\,\le\,z_{l} \right\}}
        -
        F_{Y(1)}(z_{l-1}|x)
      \right)
      \,,
      \\
  \underline{f_l}
  (t,x,y)
  &
  \ 
  :=
  \ 
      t
      \left( 
        \mathbf{1}
        _{\left\{  y\,\le\,z_{l-1} \right\}}
        -
        F_{Y(1)}(z_l|x)
      \right)
      \,,
\end{align*}
for $l\in \left\{ 1,\ldots,m \right\}$.
These brackets cover $\mathcal{F}$.
Indeed,
\begin{gather*}
  \text{for all}\ 
  z\in\R
  \ 
  \text{there exists} \ 
l\in \left\{ 1,\ldots,m \right\}
\qquad 
\text{such that}\qquad
z_{l-1}
\ 
\le
\ 
z
\ 
\le
\ 
z_l
\,.
\end{gather*}
By the monotonicity of 
$
        \mathbf{1}
      _{\left\{  y\,\le\,(\cdot) \right\}}
$
and
$
        F_{Y(1)}(\cdot|x)
$
and the non-negativity of $T$ it follows
\begin{gather*}
  \text{for all}\ 
  z\in\R
  \ 
  \text{there exists} \ 
l\in \left\{ 1,\ldots,m \right\}
\qquad 
\text{such that}\qquad
  \underline{f_l}
  \ 
  \le
  \ 
  f_z
  \ 
  \le
  \ 
  \overline{f_l}
  \,.
\end{gather*}
Thus, the $m$ brackets 
$
[
  \underline{f_l}
  ,
  \overline{f_l}
]
$
cover $\mathcal{F}$.

Let's calculate the size of the brackets.
It holds
\begin{align*}
  &
\E
\left[ 
      T
      \cdot
      \left( 
        \mathbf{1}
        _{\left\{  Y(T)\,\le\,z_{l} \right\}}
        -
        F_{Y(1)}(z_{l-1}|X)
        \ 
        -
        \ 
        \mathbf{1}
        _{\left\{  Y(T)\,\le\,z_{l-1} \right\}}
        +
        F_{Y(1)}(z_{l}|X)
      \right)
      \,
\right]
\\
  &
  \ 
=
  \ 
\E
\left[ 
      T
      \cdot
      \left( 
        \mathbf{1}
        _{\left\{
        Y(T)
        \,
        \in 
        \,
    [z_{l-1},z_l]
\right\}}
\ 
        +
\ 
        \P
        \left[ 
          Y(1)
          \in
    [z_{l-1},z_l]
        \,
    |
        \,
    X
        \right]
      \right)
      \,
\right]
\\
  &
  \ 
\le
  \ 
\E
\left[ 
  \,
  \pi(X)
  \cdot
        \P
        \left[ 
          Y(1)
          \in
    [z_{l-1},z_l]
          \,
    |
          \,
    X
        \right]
          \ 
\right]
\ 
+
\ 
\varepsilon
\\
  &
  \ 
\le
  \ 
2
\,
\varepsilon
\,.
\end{align*}
We used \eqref{size_z}, $0\le T,\pi(X)\le 1$ and Lemma~\ref{ps_weights_lemma}.
It follows
\begin{align*}
  &
  \norm{
    \left( 
  \overline{f_l}
-
  \underline{f_l}
    \right)
  (T,X,Y(T))
}_
{ L^2(\P)}
\\
&
\ 
\lesssim
\ 
\E
\left[ 
  \,
      T
      \cdot
      \left( 
        \mathbf{1}
        _{\left\{
        Y(T)
        \,
        \in 
        \,
    [z_{l-1},z_l]
\right\}}
\ 
        +
\ 
        \P
        \left[ 
          Y(1)
          \in
    [z_{l-1},z_l]
        \,
    |
        \,
    X
        \right]
      \right)
      \,
   \right]^{1/2}
\ 
\lesssim
\ 
\varepsilon^{1/2}
\,.
\end{align*}
Since $m\le 2/\varepsilon$ it holds
  \begin{align*}
    N_{[\,]}
    \left(
\varepsilon^{1/2}
    ,
    \,
    \mathcal{F}\,,\, L^2(\P)
    \right)
    &
    \ 
    \lesssim
    \ 
    \frac{1}{\varepsilon}
    \intertext{and thus}
    N_{[\,]}
    (
    \varepsilon
    ,
    \mathcal{F}, L^2(\P))
    &
    \ 
    \lesssim
    \ 
    \left( 
      \frac{1}{\varepsilon}
    \right)^2
    \,.
  \end{align*}
  Next, we look at $\mathcal{G}$. To this end, we define 
  $m$ brackets by
 \begin{align*}
    \overline{g_l}
    (t,x,y)
    \ 
    :=
    \ 
    \frac{t}{\pi(x)}
    \left( 
      \mathbf{1}_{\left\{  y\,\le\,z_{l} \right\}}
      -
      F_{Y(1)}(z_{l-1}|x)
    \right)
    \ 
    +
    \ 
    F_{Y(1)}(z_{l}|x)
-
F_{Y(1)}(z_{l-1})
\,,
\\
    \underline{g_l}
    (t,x,y)
    \ 
    :=
    \ 
    \frac{t}{\pi(x)}
    \left( 
      \mathbf{1}_{\left\{  y\,\le\,z_{l-1} \right\}}
      -
      F_{Y(1)}(z_l|x)
    \right)
    \ 
    +
    \ 
    F_{Y(1)}(z_{l-1}|x)
-
      F_{Y(1)}(z_l)
\,,
  \end{align*}
  for $l\in \left\{ 1,\ldots,m \right\}$.
  With the same arguments as before, we see that these brackets cover $\mathcal{G}$.
  Let's calculate the size.
  It holds
  \begin{align*}
    &
    \norm{
      \frac{T}{\pi(X)}
      \left( 
      \mathbf{1}_{
      \left\{ 
      Y(T)\in [z_{l-1},z_l] 
    \right\}
    }
      +
      \P
      \left[ 
      Y(1)\in [z_{l-1},z_l] 
      \,
      |
      \,
      X
      \right]
      \right)
    }_{ L^2(\P)}
    \\
    &
    \
    \lesssim
    \
    \left( 
      \E
      \left[ 
        \frac{1}{\pi(X)}
        \frac{T}{\pi(X)}
      \left( 
      \mathbf{1}_{
      \left\{ 
      Y(T)\in [z_{l-1},z_l] 
    \right\}
    }
      +
      \P
      \left[ 
      Y(1)\in [z_{l-1},z_l] 
      \,
      |
      \,
      X
      \right]
      \right)
      \right]
    \right)
    ^{1/2}
    \\
    &
    \
    \lesssim
    \
    \left( 
      \E
      \left[ 
        \frac{1}{\pi(X)}
      \P
      \left[ 
      Y(1)\in [z_{l-1},z_l] 
      \,
      |
      \,
      X
      \right]
      \right]
    \right)
    ^{1/2}
    \\
    &
    \ 
    \lesssim
    \ 
    \left( 
      \norm{1/\pi(X)}_{ L^2(\P)}
      \sqrt{\varepsilon}
    \right)
    ^{1/2}
    \ 
    =
    \ 
    \varepsilon^{1/4}
    \norm{1/\pi(X)}_{ L^2(\P)}^{1/2}
  \end{align*}
  and
  \begin{align*}
     \norm{
      \P
      \left[ 
      Y(1)\in [z_{l-1},z_l] 
      \,
      |
      \,
      X
      \right]
      \ 
     + 
      \ 
      \P
      \left[ 
      Y(1)\in [z_{l-1},z_l] \,
      \right]
    }_{ L^2(\P)}
    \ 
    \lesssim
    \ 
    \varepsilon^{1/2}
    \,.
  \end{align*}
  Thus
  \begin{align*}
  \norm{
    \left( 
  \overline{g_l}
-
  \underline{g_l}
    \right)
  (T,X,Y(T))
}_
{ L^2(\P)}
    &
\ 
\lesssim
\ 
\varepsilon^{1/4}
\left( 
  1
  +
    \norm{1/\pi(X)}_{ L^2(\P)}^{1/2}
\right)
\\
    &
\ 
\lesssim
\ 
\varepsilon^{1/4}
\left( 
  1
  +
    \norm{1/\pi(X)}_{ L^2(\P)}
\right)
\,.
  \end{align*}
As before, it follows
\begin{gather*}
    N_{[\,]}
    (
    \varepsilon
    ,
    \mathcal{G}, L^2(\P))
    \ 
    \lesssim
    \ 
    \left( 
    \frac{
      1+
    \norm{1/\pi(X)}_{ L^2(\P)}
    }
    {\varepsilon}
    \right)^4
    \,.
\end{gather*}
\end{proof}

\begin{lemma}
  \label{markov_max_lemma}
  Let $(\mathcal{H}_N)$ be a sequence of measurable function classes with envelope functions $(H_N)$.
  If
  \begin{gather}
    J_{[\, ]}
    \left( 
    \norm{F_N}_{ L^2(\P)}
    ,
    \mathcal{H}_N
    ,
     L^2(\P)
    \right)
    \ 
    \to
    \ 
    0
    \qquad
    \text{for}
    \ 
    N
    \to
    \infty
  \end{gather}
  it holds 
  $
  \norm{\G_N}^*_{\mathcal{H}_N}\overset{\P}{\to}0
  $.
\end{lemma}
\begin{proof}
  By Markov's inequality and Theorem~\ref{th:max_ineq} it holds for all $\varepsilon>0$
  \begin{align*}
    &
    \P
    [
  \norm{\G_N}^*_{\mathcal{H}_N}
  \ge
  \varepsilon
    ]
    \\
    &
    \ 
    \le
    \ 
    \varepsilon^{-1}
    \E
    [
  \norm{\G_N}^*_{\mathcal{H}_N}
    ]
    \ 
    =
    \ 
    \varepsilon^{-1}
    \E^*
    [
  \norm{\G_N}_{\mathcal{H}_N}
    ]
    \ 
    \lesssim
    \ 
    \varepsilon^{-1}
    J_{[\, ]}
    \left( 
    \norm{H_N}_{ L^2(\P)}
    ,
    \mathcal{H}_N
    ,
     L^2(\P)
    \right)
    \ 
    \to
    \ 
    0
  \end{align*}
  for $N\to\infty$. 
\end{proof}
A technical lemma for products of function classes.

Define the product of two function classes
\begin{gather*}
  \mathcal{F}\cdot \mathcal{G}
  :=
  \left\{ 
    f\cdot g
    \colon
    f\in\mathcal{F},
    g\in\mathcal{G}
  \right\}\,.
\end{gather*}
\begin{lemma}
  \label{lem_prod_br}
  Let
  $\mathcal{F}$ and $\mathcal{G}$ be two function classes 
  with envelope functions $F$ and $G$ satisfying
  $\norm{F}_\infty,\norm{G}_\infty\le 1$.
  For all $\varepsilon>0$ and all $r\in [1,\infty)$ it holds
  \begin{gather*}
    N_{[\,]}(2\varepsilon,\mathcal{F}\cdot\mathcal{G},\mathrm{L}_r(\P))
    \
    \le
    \ 
    N_{[\,]}(\varepsilon,\mathcal{F},\mathrm{L}_r(\P))
    \cdot
    N_{[\,]}(\varepsilon,\mathcal{G},\mathrm{L}_r(\P))
    \,.
  \end{gather*}
\end{lemma}
\begin{proof}
  Let $f\in\mathcal{F}$ and $g\in\mathcal{G}$.
  We can choose two 
  $(\varepsilon,L^r(\P))$
  brackets
  $[\underline{f},\overline{f}]$
  and
  $[\underline{g},\overline{g}]$
  containing $f$ and $g$ with 
  $\norm{\underline{f}}_\infty,\norm{\overline{f}}_\infty\le\norm{F}_\infty\le 1$
  and
  $\norm{\underline{g}}_\infty,\norm{\overline{g}}_\infty\le\norm{G}_\infty\le 1$.
  We the get an 
  $(2\varepsilon,L^r(\P))$
  $[\underline{h},\overline{h}]$
  bracket, containing $f\cdot g$, by
\end{proof}

\begin{lemma}
  \label{aa:r3:lemma:1}
  Let
$(\varepsilon_N)\subset(0,1]$
be 
a decreasing sequence
with $\varepsilon_N\to 0$ for $N\to\infty$ 
and let Assumption~\ref{aa:assumption:1} hold true
for a sequence of function classes $\mathcal{F}_N$.
Then
\begin{gather*}
  J_{[\,]}(
\norm{F_N}_{L^2(\P)}
,\mathcal{F}_N\cdot\mathcal{F},\mathrm{L}_2(\P))
  \to 0
  \quad
  \text{and}
  \quad
  \norm{\G_N}^*_{\mathcal{F}_N\cdot\mathcal{F}}\overset{\P}{\to}0
  \qquad
  \text{for}\ 
  N\to\infty
  \,.
\end{gather*}
\end{lemma}
\begin{proof}
  By Assumption~\ref{aa:assumption:1}
  and Lemma~\ref{aa:mean:l:br} it holds
for some $k<2$
\begin{gather*}
\norm{F_N}_{L^2(\P)}
\ 
\le
\ 
\varepsilon_N
\quad
\text{and}
\quad
  \log
  N_{[\,]}(\varepsilon,\mathcal{F}_N,\mathrm{L}_2(\P))
  \ 
  \lesssim
  \ 
  \left( 
  \frac{1}{\varepsilon}
  \right)^k
  \quad
  \text{for all}
  \ 
  N\in\mathbb{N}
  \,,
\end{gather*}
and
  \begin{gather*}
    N_{[\,]}
    (
    \varepsilon
    ,
    \mathcal{F}, L^2(\P))
    \ 
    \lesssim
    \ 
    \left( 
      \frac{1}{\varepsilon}
    \right)^2
    \qquad
    \text{for all}
    \ 
    \varepsilon>0
    \,.
  \end{gather*}
  Since $\mathcal{F}_N$ and $\mathcal{F}$ have envelope function smaller 1, we can apply Lemma~\ref{lem_prod_br} to get
  \begin{gather*}
  \log
  N_{[\,]}(\varepsilon,\mathcal{F}_N\cdot\mathcal{F},\mathrm{L}_2(\P))
  \ 
  \lesssim
  \ 
  \left( 
  \frac{1}{\varepsilon}
  \right)^k
  +
  \log
  (1/\varepsilon)
  \ 
  \lesssim
  \ 
  \left( 
  \frac{1}{\varepsilon}
  \right)^k
  \quad
  \text{for all}\ 
  \varepsilon>0
  \,.
  \end{gather*}
  Since 
  $k/2\in(0,1)$
  it holds
\begin{align*}
  J_{[\,]}(
\norm{F_N}_{L^2(\P)}
,\mathcal{F}_N\cdot\mathcal{F},\mathrm{L}_2(\P))
  &
  \ 
=
  \ 
\int_0^{
\norm{F_N}_{L^2(\P)}
}
\sqrt{
  \log
  N_{[\,]}(\varepsilon,\mathcal{F}_N\cdot\mathcal{F},\mathrm{L}_2(\P))
}
\,d\varepsilon
\\
&
\ 
\lesssim
\ 
\int_0^{
  \varepsilon_N
}
  \left( 
  \frac{1}{\varepsilon}
\right)^{k/2}
\,d\varepsilon
\\
&
\ 
=
\ 
\frac{
\varepsilon_N^{1-k/2}
}{1-k/2}
\ 
\to 0
\ 
\qquad
\text{for}
\ 
N\to\infty
\,.
\end{align*}
The second statement follows from Lemma~\ref{markov_max_lemma}
for 
$\mathcal{H}_N:=\mathcal{F}_N\cdot\mathcal{F}$ and
$H_N:=F_N$.
\end{proof}

\begin{lemma}
  \label{ps_weights_lemma}
  Let
  $
  g_1\colon
  \mathcal{X}\to\R
  $
  and
  $
  g_2\colon
  \mathcal{Y}\to\R
  $
  be a measurable functions.
  It holds
  \begin{gather*}
    \E
    \left[
    \frac{T}{\pi(X)}
    g_1(X)
    \right]
    \ 
    =
    \ 
    \E
    \left[
    g_1(X)
    \right]
    \,.
  \end{gather*}
  If Assumption~\ref{aa:assumption:treatment_str_ign} holds true, then
  \begin{gather*}
    \E
    \left[
    \frac{T}{\pi(X)}
    g_2(Y(T))
    \right]
    \ 
    =
    \ 
    \E
    \left[
    f(Y(1))
    \right]
    \,.
  \end{gather*}
\end{lemma}
\begin{lemma}
  \label{aa:mean:r3:lem:conv}
  Under conditions it holds
  $\sup_{z\in\R}\left| R_3(z) \right|\overset{\P}{\to}0$.
\end{lemma}
\begin{proof}
  Let $N\ge\underline{N}$,
  $z\in\R$,
  and 
  let
  $g^\dagger$ 
  denote the function \eqref{ghost_function} with 
  $
(\lambda^\dagger,\lambda_0^\dagger)
  $.
  If
  \begin{gather*}
    \left| 
    w^\dagger(X)- \frac{1}{\pi(X)} 
    \right|
    \ 
    \le
    \ 
    \varepsilon_N
  \end{gather*}
  it holds
  \begin{gather*}
    g^\dagger
    (X)
    \cdot
    f_z
    (T,X,Y(T))
    \ 
    =
    \ 
  \left( 
      w^\dagger(X)
      -
      \frac{1}{\pi(X)}
  \right)
      T
      \left( 
        \mathbf{1}
        _{\left\{  Y(T)\,\le\,z \right\}}
        -
        F_{Y(1)}(z|X)
      \right)
      \,.
  \end{gather*}
  By Lemma~\ref{aa:mean:r3:lem:fz_E}
  it holds 
  \begin{align*}
    &
  f_z(T,X,Y(T))
  \ 
  \in
  \ 
  L^1(\P)\,, 
  \\
  &
  f_z(T,X,Y(T))
  \ 
  \perp
  \ 
  D_N
  \,,
  \\
  \E
  [
  &
  f_z(T,X,Y(T))
    |X
  ]
  \ 
  =
  \ 
  0
  \,.
  \end{align*}
  Thus, 
  it follows from Lemma~\ref{w.Z=0}
  \begin{gather*}
    \E
    \left[
      w^\dagger(X)
      \cdot
      f_z(T,X,Y(T))
    \right]
    \ 
    =
    \ 
    0
    \,.
  \end{gather*}
  Since
  \begin{gather*}
    \E
    \left[
      \frac{T}{\pi(X)}
      f_z(T,X,Y(T))
    \right]
    \ 
    =
    \ 
    \E
    \left[
      \frac{T}{\pi(X)}
      \left( 
        \mathbf{1}
        _{\left\{  Y(T)\,\le\,z \right\}}
        -
        F_{Y(1)}(z|X)
      \right)
    \right]
    \  
    =
    \ 
    0
  \end{gather*}
  by Lemma~\ref{ps_weights_lemma},
  it follows
  \begin{align*}
    &
  \E
  \left[
    g^\dagger(X)
    f_z(T,X,Y(T))
  \right]
  \\
    &
  \ 
  =
  \ 
    \E
    \left[
      w^\dagger(X)
      \cdot
      f_z(T,X,Y(T))
    \right]
    \ 
    -
    \ 
    \E
    \left[
      \frac{T}{\pi(X)}
      f_z(T,X,Y(T))
    \right]
    \ 
    =
    \ 
  0
  \,.
  \end{align*}
  But then 
  \begin{align*}
    R_3(z)
    &
    \ 
    =
    \ 
  \frac{1}
  {
\sqrt{N}
  }
    \sum_{i=1}^{N} 
    \left[ 
    \left( 
    w(X_i) 
    -
    \frac{1}{\pi(X_i)}
    \right)
    T_i
    \left( 
    \mathbf{1}_{\left\{ Y_i \le z \right\}}
    -
  F_{Y(1)}(z|X_i)
    \right)
    \right]
    \ 
    =
    \ 
    \G_N
    \left( 
      g^\dagger
      \cdot
      f_z
    \right)
    \,.
  \end{align*}
  It follows
\begin{align*}
    \P
    \left[ 
      \sup_{z\in\R}
     | 
    R_3(z)
    |
      \ge
      \varepsilon
    \right]
    &
    \ 
    \le
    \ 
    \P
    \left[ 
      \sup_{z\in\R}
     | 
    R_3(z)
    |
      \ge
      \varepsilon
      \ 
      \text{and}
      \ 
    | 
    w^\dagger(X)- 1/\pi(X)
    |
    \ 
    \le
    \ 
    \varepsilon_N
    \right]
    \\
    &
    \quad
    +
    \ 
    \P
    \left[ 
    | 
    w^\dagger(X)- 1/\pi(X)
    |
    \ 
    >
    \ 
    \varepsilon_N
    \right]
    \\
    &
    \ 
    \le
    \ 
    \P
    \left[ 
      \norm{\G_N}^*_{\mathcal{F}_N\cdot\mathcal{F}}
      \ge
      \varepsilon
    \right]
    \ 
    +
    \ 
    \ 
    \P
    \left[ 
    | 
    w^\dagger(X)- 1/\pi(X)
    |
    \ 
    >
    \ 
    \varepsilon_N
    \right]
    \\
    &
    \ 
    \to
    \ 
    0
    \,.
\end{align*}
The convergence of the first term follows from
Lemma~\ref{aa:r3:lemma:1}.
The convergence of the second term follows from
Theorem~\ref{aa:weights:th}.

\end{proof}


Until now, all parts of the error decomposition converge to 0.
The last term $R_4$ will decide the profile of the limiting process.
To this end we need the following concept.
\begin{definition}
  We call a class 
  $\mathcal{F}$ of measurable functions 
$\P$-Donsker
if the sequence of processes 
$\left\{ G_N f \colon f\in\mathcal{F}\right\}$
converges in
$l^\infty(\mathcal{F})$
to a tight limit process.
\end{definition}

\begin{theorem}
  Every class $\mathcal{F}$ of measurable functions 
  with
  $
    J
    _{[]}
    (
    1
    ,
    \mathcal{F}
    ,
    L_2(\P)
    )
    <\infty
  $
  is

  $\P$-Donsker, that is,
  the sequence of processes 
$\left\{ G_N f \colon f\in\mathcal{F}\right\}$
  converges 
  in
$l^\infty(\mathcal{F})$
to a Gaussian process with mean 0 and covariance function given by
\begin{gather}
  \mathrm{Cov}(f,g)
  :=
  \E[fg]-\E[f]\E[g]
  \,.
\end{gather}
\end{theorem}
\begin{proof}
  \cite[Theorem~19.5]{Vaart2000}
\end{proof}



\begin{lemma}
  \label{aa:mean:l:r4}
  Let
  $1/\pi(X)\in L^2(\P)$.
  $R_4$ converges
  converges in
  $l^\infty(\R)$
  to a Gaussian process with mean 0 and covariance
\begin{align*}
  &
  \mathbf{Cov}
  (z_1,z_2)
  \\
  &
  =\ 
  \E
  \left[ 
 \frac{
 F_{Y(1)}(z_1 \land z_2\,|\,X)
}{\pi(X)}
\ 
-
\ 
 \frac{1-\pi(X)}{\pi(X)}
 F_{Y(1)}(z_1|X)
 \cdot
 F_{Y(1)}(z_2|X)
  \right]
  \ 
 -
 \ 
 F_{Y(1)}(z_1)
 \cdot
 F_{Y(1)}(z_2)
\end{align*}

\end{lemma}
\begin{proof}
  By Lemma~\ref{aa:mean:r3:lem:fz_E} it follows
  \begin{align*}
    \E
    \left[
      \frac{f_z(T,X,Y(T))}{\pi(X)}
      +
      F_{Y(1)}(z|X)
      -
      F_{Y(1)}(z)
      \right]
      \ 
      =
      \ 
      \E
      \left[
      \frac{1}{\pi(X)}
      \E
      \left[
        f_z(T,X,Y(T))
        |X
      \right]
      \right]
      \ 
      =
      \ 
      0
      \,.
  \end{align*}
  Thus
  \begin{align*}
    R_4(z)
    &
  \
  =
  \ 
  \frac{1}{
  \sqrt{N}
  }
    \sum_{i=1}^{N} 
    \frac{T_i}{\pi(X_i)}
    \left( 
    \mathbf{1}_{\left\{ Y_i \le z \right\}}
    -
  F_{Y(1)}(z|X_i)
    \right)
    \ 
    +
    \ 
    \left( 
  F_{Y(1)}(z|X_i)
    -
  F_{Y(1)}(z)
    \right)
    \\
    &
    \ 
  =
    \ 
  \frac{1}{
  \sqrt{N}
  }
    \sum_{i=1}^{N} 
      \frac{f_z(T_i,X_i,Y_i)}{\pi(X_i)}
      +
      \left( 
      F_{Y(1)}(z|X_i)
      -
      F_{Y(1)}(z)
      \right)
      \\
      &
      \ 
      =
      \ 
      \G_N 
      \left(
       \frac{f_z}{\pi(\cdot)}
      +
      F_{Y(1)}(z|\cdot)
      -
      F_{Y(1)}(z)
      \right)
      \,.
  \end{align*}
  By Lemma~\ref{aa:mean:l:br}
  it holds
  \begin{align*}
    &
  \log
  N_{[\,]}
    (
    \varepsilon
    ,
    \mathcal{G}, L^2(\P))
    \\
    &
    \ 
    \lesssim
    \ 
    \log
    \left(
      \frac
      {
      1+
    \norm{1/\pi(X)}_{ L^2(\P)}
      }
      {\varepsilon}
    \right)
    \ 
    \lesssim
    \ 
      \frac
      {
      1+
    \norm{1/\pi(X)}_{ L^2(\P)}
      }
      {\varepsilon}
    \qquad
    \text{for all}
    \ 
    \varepsilon\in (0,1)
    \,.
  \end{align*}
  Thus
  \begin{gather*}
    J_{[\,]}(1,\mathcal{G},L^2(\P))
    \ 
    \lesssim
    \ 
    \int_0^1
    \sqrt
    {
      \frac
      {
      1+
    \norm{1/\pi(X)}_{ L^2(\P)}
      }
      {\varepsilon}
    }
    \,
    d\varepsilon
    \ 
    \lesssim
    \ 
      1+
    \norm{1/\pi(X)}_{ L^2(\P)}
    \ 
    <
    \ 
    \infty
    \,.
  \end{gather*}
But then $\mathcal{G}$ is $\P$-Donsker.
By the Donsker Theorem \cite[Theorem~19.5]{Vaart2000}
the process $R_4$ converges in $l^\infty(\R)$ to a Gaussian process, called $\P$-Brownian bridge, with mean 0.
We now calculate the covariance of the limiting process.
\subsubsection*{Covariance}
\begin{align*}
  &
  \E
  \left[
  \left( 
  f^{z_1}_{1/\pi}
  +
  F_{Y(1)}(z_1|X)
  -
F_{Y(1)}(z_1)
  \right)
  \left( 
  f^{z_2}_{1/\pi}
  +
  F_{Y(1)}(z_2|X)
  -
F_{Y(1)}(z_2)
  \right)
  \right]
  \\
  &
  \ 
  =
  \ 
\E
\left[
  f^{z_1}_{1/\pi}
  \cdot
  f^{z_2}_{1/\pi}
\right]
\\
  &
  \quad
  +
  \ 
  \E
  \left[
  f^{z_1}_{1/\pi}
  \left( 
  F_{Y(1)}(z_2|X)
  -
F_{Y(1)}(z_2)
  \right)
  \right]
  \ 
  +
  \ 
  \E
  \left[
  f^{z_2}_{1/\pi}
  \left( 
  F_{Y(1)}(z_1|X)
  -
F_{Y(1)}(z_1)
  \right)
  \right]
  \\
  &
  \quad
  +
  \ 
  \E
  \left[
  \left( 
  F_{Y(1)}(z_1|X)
  -
F_{Y(1)}(z_1)
  \right)
  \left( 
  F_{Y(1)}(z_2|X)
  -
F_{Y(1)}(z_2)
  \right)
  \right]
  \\
  &
  \ 
  =:
  \ 
  C_0
  \quad 
  +
  \quad 
  C_1
  +
  C_2
  \quad 
  +
  \quad 
  C_3
  \,.
\end{align*}
It holds
\begin{align*}
  C_0 
  &
  \ 
  =
  \ 
\E
\left[
  f^{z_1}_{1/\pi}
  \cdot
  f^{z_2}_{1/\pi}
\right]
\\
&
\ 
=
\ 
\E
\left[
\frac{1}{\pi(X)}
\frac{T}{\pi(X)}
\left( 
\mathbf{1}_{\left\{ Y(T)\,\le\, z_1 \right\}}
-
F_{Y(1)}(z_1|X)
\right)
\left( 
\mathbf{1}_{\left\{ Y(T)\,\le\, z_2 \right\}}
-
F_{Y(1)}(z_2|X)
\right)
\right]
\\
&
\ 
=
\ 
\E
\left[
\frac{1}{\pi(X)}
\left( 
\mathbf{1}_{\left\{ Y(1)\,\le\, z_1 \right\}}
-
F_{Y(1)}(z_1|X)
\right)
\left( 
\mathbf{1}_{\left\{ Y(1)\,\le\, z_2 \right\}}
-
F_{Y(1)}(z_2|X)
\right)
\right]
\\
&
\ 
=
\ 
\E
\left[
\frac{1}{\pi(X)}
\left( 
F_{Y(1)}(z_1\land z_2|X)
\ 
-
\ 
F_{Y(1)}(z_1|X)
\cdot
F_{Y(1)}(z_2|X)
\right)
\right]
\,.
\end{align*}
\begin{align*}
  C_1
  &
  \ 
  =
  \ 
 \E
  \left[
  f^{z_1}_{1/\pi}
  \left( 
  F_{Y(1)}(z_2|X)
  -
F_{Y(1)}(z_2)
  \right)
  \right]
  \\
  &
  \ 
  =
  \ 
 \E
  \left[
\frac{T}{\pi(X)}
\left( 
\mathbf{1}_{\left\{ Y(T)\,\le\, z_1 \right\}}
-
F_{Y(1)}(z_1|X)
\right)
  \left( 
  F_{Y(1)}(z_2|X)
  -
F_{Y(1)}(z_2)
  \right)
  \right]
  \\
  &
  \ 
  =
  \ 
 \E
  \left[
\left( 
\mathbf{1}_{\left\{ Y(1)\,\le\, z_1 \right\}}
-
F_{Y(1)}(z_1|X)
\right)
  \left( 
  F_{Y(1)}(z_2|X)
  -
F_{Y(1)}(z_2)
  \right)
  \right]
  \\
  &
  \ 
  =
  \ 
  0
  \,.
\end{align*}
In the same way we see $C_2=0$.
\begin{align*}
  C_3
  &
  \ 
  =
  \ 
  \E
  \left[
  \left( 
  F_{Y(1)}(z_1|X)
  -
F_{Y(1)}(z_1)
  \right)
  \left( 
  F_{Y(1)}(z_2|X)
  -
F_{Y(1)}(z_2)
  \right)
  \right]
  \\
  &
  \ 
  =
  \ 
  \E
  \left[
  F_{Y(1)}(z_1|X)
  \cdot
  F_{Y(1)}(z_2|X)
  \right]
  \ 
  -
  \ 
  F_{Y(1)}(z_1)
  \cdot
  F_{Y(1)}(z_2)
  \,.
\end{align*}
Adding up the results gives us \eqref{cov:lp}.
\end{proof}
We have gathered all the results to prove Theorem~\ref{aa:mean:th}.
\begin{proof}
  \emph{(Theorem~\ref{aa:mean:th})}
  We connect the statement of the theorem to the error decomposition by Lemma~\ref{aa:mean:lemma_decomp}.
  By Lemma~\ref{aa:mean:l:r1}, Lemma~\ref{aa:mean:l:r2},
  Lemma~\ref{aa:mean:r3:lem:conv}
   it follows 
   $\sup_{z\in\R}|R_i(z)|\overset{\P}{\to}0$ for $i=1,2,3$.
   Thus, by Slutzky's theorem (cf.\cite[Theorem~13.18]{Klenke2020})
   the behaviour of the limiting process is the one of Lemma~\ref{aa:mean:l:r4}.
\end{proof}

  %\begin{theorem}
  The estimate 
  \begin{gather}
    \sqrt{N}
    \left( 
  \frac{1}{N}
    \sum_{i=1}^{n} 
    w_i
    \mathbf{1}
    _{\left\{ Y_i\le t \right\}}
    -
    \P[Y(1)\le t]
    \right)
  \end{gather}
  is Asymptoticaly normal.
\end{theorem}
\begin{proof}
\begin{align*}
  &
    \frac{1}{N}
    \sum_{i=1}^{n} 
    w_i
    \mathbf{1}_{\left\{ Y_i \le t \right\}}
    \ 
    -
    \ 
    \P[Y(1)\le t]
    \\
  &
  \ 
  =
  \ 
  \frac{1}{N}
  \left( 
    \sum_{i=1}^{n} 
    w_i
    B(X_i)
    -
    \sum_{i=1}^{N} 
    B(X_i)
  \right)
  ^\top
     \mathbf{Y}(1) 
  %%%% 1 %%%%
  \\
  &
  \qquad
  +
  \ 
  \frac{1}{N}
    \sum_{i=1}^{N} 
    \left( T_i\cdot w_i -1 \right)
    \left( 
    \P[Y(1)\le t| X_i]
    -
    \inner
    {B(X_i)}
    { \mathbf{Y}(1) }
    \right)
  %%%% 2 %%%%
  \\
  &
  \qquad
  +
  \ 
  \frac{1}{N}
    \sum_{i=1}^{N} 
    T_i
    (
    w_i 
    -
    1/\pi_i
    )
    \left( 
    \mathbf{1}_{\left\{ Y_i \le t \right\}}
    -
    \P[Y(1)\le t| X_i]
    \right)
  %%%% 3 %%%%
  \\
  &
  \qquad
  +
  \ 
  \frac{1}{N}
    \sum_{i=1}^{N} 
    T_i
    /\pi_i
    \left( 
    \mathbf{1}_{\left\{ Y_i \le t \right\}}
    -
    \E[Y(1)\le t| X_i]
    \right)
    +
    \left( 
    \P[Y(1)\le t| X_i]
    -
    \P[Y(1)\le t]
    \right)
  \\
  %%%%%%%%%%%%%%%
  &
  \ 
  =:
  \ 
  R_1
  +
  R_2
  +
  R_3
  +
  R_4
  \end{align*}

The term 
\begin{gather}
    \left|  
  \frac{1}{N}
    \sum_{i=1}^{N} 
    \left( 
    T_i
    w_i
    -
    1
    \right)
    \left( 
    \P[Y(1)\le t | X_i]
    -
    \inner
    {B(X_i))}
    {\lambda^*}
    \right)
  \right|
\end{gather}
gives back the expectation. We control its rate 
with the factor of basis rates and primal weights.
Choose $d$ to be faster than $\sqrt{n}$.
The rest follows with standard empirical theory.
Expectations are 0.
$\norm{\delta}$ has to converge fast enough.
\begin{align*}
  &
  \sqrt{N}
  \cdot
  \E
  \left[ 
    \left( 
    T
    \cdot
    w(X)
    -
    1
    \right)
    \left( 
      \P
      [
      Y(1)\le t
      |X
      ]
      -
      \inner
      {B(X)}
      {\lambda^*}
    \right)
  \right]
  \\
  &
  \ 
  =
  \ 
  \sqrt{N}
  \cdot
  \E
  \left[ 
    \pi(X)
    \cdot
    \left( 
    w(X)
    -
    1
    /
    \pi(X)
    \right)
    \left( 
      \P
      [
      Y(1)\le t
      |X
      ]
      -
      \inner
      {B(X)}
      {\lambda^*}
    \right)
  \right]
  \\
  &
  \ 
  \le
  \ 
  \sqrt{N}
  \cdot
  \E
  \left[ 
    \left| 
    w(X)
    -
    1
    /
    \pi(X)
    \right|
    ^2
    \right]
    ^{1/2}
  \E
  \left[ 
    \left| 
      \P
      [
      Y(1)\le t
      |X
      ]
      -
      \inner
      {B(X)}
      {\lambda^*}
    \right|
    ^2
    \right]
    ^{1/2}
    \\
  &
  \ 
  \le
  \ 
  C
  \sqrt{N}
  \sqrt{\log(N)}
  N^{-1/(d+2)}
  \cdot
  N^{-1/(d+2)}
    \\
  &
  \ 
  \le
  \ 
  C
  \sqrt{\log(N)}
  N^{-1/6}
  \to 0
  \qquad
  \text{for}\ 
  N\to\infty
\end{align*}
Note, that we get no convergence for $d>1$. 
Also note that
\begin{gather*}
  \E
  \left[ 
    T
    \cdot
    w(X)
    -
    1
    |
    X,
    X_1,
    \ldots,
    X_N
    \right]
    =
  \E
  \left[ 
    T
    |
    X
    \right]
    w(X)
    -
    1
    =
    \pi(X)
    \left( 
      w(X)
      -
      1/\pi(X)
    \right)
\end{gather*}
\end{proof}
Let 
\begin{gather}
  B_{\mathcal{F}}
  :=
  \left\{ 
    f\in \mathcal{F}
    \colon
    \norm{f}_\infty
    \le
    1
  \right\}
\end{gather}
be the intersection of the unit ball with $\mathcal{F}$.
\begin{lemma}
  Let $(\varepsilon_N)$ be a sequence converging to 0
  and let 
  $
  \left( 
    \mathcal{F}_N
  \right)
    :=
    \left( 
    \varepsilon_N\cdot B_\mathcal{F}
    \right)
  $
  denote the sequence of scaled unit balls in $\mathcal{F}$.
  Assume that 
  there exists
  $k<2$ such that 
  the covering number of the unit ball in $\mathcal{F}$
  satisfies
  \begin{gather}
        \log 
      N_{[\,]}
\left( \varepsilon, B_\mathcal{F},\mathrm{L}_2(\P) \right)
\ 
\le
\ 
\left( \frac{1}{\varepsilon} \right)^k
\qquad
\text{for all}\ 
\varepsilon>0
\,.
  \end{gather}
  Then it holds
  $
      \norm{G_N}^*_{\mathcal{F}_N}
    \ 
    \overset{\P}
    {\to}
    \ 
    0
  $
  for $N\to \infty$. 
\end{lemma}
\begin{proof}
  By maximal inequalities it holds
  \begin{align*}
    \E^*
    \left[ 
      \norm{G_N}_{\mathcal{F}_N}
    \right]
    &
      \ 
      \lesssim
      \ 
      J_{[\,]}\left( \varepsilon_N, \mathcal{F}_N,\mathrm{L}_2(\P) \right)
      \\
    &
      \ 
      =
      \ 
      \int_0^{\varepsilon_N}
      \sqrt{
        \log 
      N_{[\,]}
\left( \varepsilon, \mathcal{F}_N,\mathrm{L}_2(\P) \right)
    }
    \,
    d\varepsilon
    \\
    &
    \ 
    =
    \ 
      \int_0^{\varepsilon_N}
      \sqrt{\log 
      N_{[\,]}
\left( \varepsilon/\varepsilon_N, B_\mathcal{F},\mathrm{L}_2(\P) \right)
    }
    \,
    d\varepsilon
    \\
    &
    \ 
    \le
    \ 
      \int_0^{\varepsilon_N}
      \left( 
      \frac{\varepsilon_N}{\varepsilon}
    \right)^{k/2}
    \,
    d\varepsilon
    \\
    &
    \ 
    \lesssim
    \ 
  \varepsilon_N^{k/2}
  \frac{1}{1-k/2}
  \varepsilon_N^{1-k/2}
  \\
    &
    \ 
  \lesssim
    \ 
  \varepsilon_N
  \\
    &
    \ 
  \to
  \ 
  0
  \qquad
  \text{for}\ 
  N\to
  \infty
  \,.
  \end{align*}
  Note, that $k<2$.
  By the boundedness of $\E^*$ there is no measurability problem.
  By Markov's Inequality it holds
  \begin{align*}
    \P
    \left[ 
      \norm{G_N}^*_{\mathcal{F}_N}\ge \varepsilon
    \right]
    \le
    \varepsilon^{-1}
    \,
    \E^*
    \left[ 
      \norm{G_N}_{\mathcal{F}_N}
    \right]
    \ 
    \to
    \ 
    0
    \qquad
    \text{for}\ 
    N\to\infty
    \,.
  \end{align*}
\end{proof}
We define $f^\dagger$ by
\begin{gather}
  f^\dagger(T,X,Y)
  :=
  T
  \left( 
    (f^{'})^{-1}
    \left( 
      \inner
      {B(X)}
      {\lambda^\dagger}
    +
    \lambda_0^\dagger
    \right)
    -
    \frac{1}{\pi(X)}
  \right)
  \left( 
    Y(T)
    -
    \E
    [
    Y(1)|X
    ]
  \right)
\end{gather}
\begin{lemma}
  Assume that 
  $f^\dagger \in \mathcal{F}$ almost surely 
  for some function class $\mathcal{F}$ satisfying the requirements 
  of the previous lemma. It then holds
  $
  \left| G_Nf^\dagger \right|
  \overset{\P}{\to}0$ for 
  $N\to\infty$.
\end{lemma}
\begin{proof}
  Let $(\varepsilon_N)$ be the lr of the primal
  then
  \begin{gather}
    \P
    \left[ 
      f^\dagger
      \in
      \mathcal{F}_N
    \right]
    \to 1
    \qquad
    \text{for}\ 
    N\to\infty
    \,.
  \end{gather}
  For all $\varepsilon>0$ it holds
  \begin{align*}
    \P
    \left[ 
  \left| G_Nf^\dagger \right|
  \le
  \varepsilon
    \right]
    &
    \ 
    \ge
    \ 
    \P
    \left[ 
  \left| G_Nf^\dagger \right|
  \le
  \norm{G_N}^*_{\mathcal{F}_N}
  \le
  \varepsilon
    \right]
    \\
    &
    \ 
    \ge
    \ 
    \P
    \left[ 
      f^\dagger\in\mathcal{F}_N
      \ 
      \text{and}\ 
  \norm{G_N}^*_{\mathcal{F}_N}
  \le
  \varepsilon
    \right]
    \\
    &
    \ 
    \ge
    \ 
    \P
    \left[ 
      f^\dagger\in\mathcal{F}_N
    \right]
    -
    \P
    \left[ 
  \norm{G_N}^*_{\mathcal{F}_N}
    \ 
  \ge
    \ 
  \varepsilon
    \right]
    \\
    &
    \ 
    \to
    \ 
    1
    \,.
  \end{align*}
  The last step is due to Lemma.
\end{proof}


  %We can even view
$
\frac{1}{\sqrt{n}}
\sum_{i=1}^{n}S_i 
$
as an empirical process 
$
\mathbb{G}_n f
$
indexed over 

\begin{gather}
  f_\Phi(T,X,Y)
  =
  \frac{T}{\pi(X)}
 \left( 
   \Phi(Y) - \E [ \Phi(Y) | X ]
 \right)
 +
   \E [ \Phi(Y) | X ] 
   .
\end{gather}
If $\mathcal{F}=\left\{ f_\Phi \colon \Phi \in \ \text{some set}\right\}$
is $\P$-Donsker, the empirical process converges to a tight gaussian process.
Then the functional delta Method is applicable.



  \section{Application to Plug In Estimators}
  A plethora of applications of the delta method to estimates of the distribution function are to be found in \cite{Vaart2000} and \cite{vaart2013}.
This includes Quantile estimation \cite[§21]{Vaart2000}\cite[§3.9.21/24]{vaart2013},
survival analysis via Nelson-Aalen and Kaplan-Meier estimator\cite[§3.9.19/31]{vaart2013},
Wilcoxon Test~\cite[§3.9.4.1]{vaart2013},
and much more.
Maybe Boostrapping from the weighted distribution is also sensible .

  %We leverage (weak) universal consistency
of partitioning estimates~\cite{Gyorfi2002}.
To this end, we choose the basis functions as
\begin{gather}
  B_k(x)
  :=
  \frac{
  \mathbf{1}_{X_k \in A_n(x)}
  }{
  \sum_{j=1}^{n} 
  \mathbf{1}_{X_j \in A_n(x)}
  }
  \,,
  \qquad
  k=
  1,\ldots,n
  \,.
\end{gather}
The euclidian norm of 
$
  B(x)
$
is bounded by one.
\begin{gather}
  \norm{B(x)}^2
  =
  \sum_{k=1}^{n} 
  \left( 
  \frac{
  \mathbf{1}_{X_k \in A_n(x)}
  }{
  \sum_{j=1}^{n} 
  \mathbf{1}_{X_j \in A_n(x)}
  }
  \right)
  ^2
  \le
  \sum_{k=1}^{n} 
  \frac{
  \mathbf{1}_{X_k \in A_n(x)}
  }{
  \sum_{j=1}^{n} 
  \mathbf{1}_{X_j \in A_n(x)}
  }
  =1
  \,.
\end{gather}
The oracle parameters are
$
  \left[ 
    f^{'}
    \left( 
      \frac{1}{\pi_i}
    \right)
  \right] _ { i = 1,\ldots,n }
$
and
$
  \left[ 
    Y_i(1)
  \right] _ { i = 1,\ldots,n }
$
.
It even holds
\begin{gather}
  \lambda ^ \dagger _ i \to 
    f^{'}
    \left( 
      \frac{1}{\pi_i}
    \right)
\end{gather}
in probability.

 \begin{gather}
   G(\lambda)
   :=
      \frac{1}{n}
      \sum_{j = 1}^{n} 
        f^* 
        \left( 
          \inner
          {B(X_i)}
          {\lambda}
        \right)
      -
          \inner
          {B(X_i)}
          {\lambda}
      +
      \inner
      { | \lambda | }
      { \delta }
 \end{gather} 
 The consistency of 
 $
  \lambda ^ \dagger
 $
 relates to the difference being nonnegative.
 For all 
 $
  \varepsilon > 0
 $
 it holds
 \begin{gather}
   \P
   \left[ 
     \norm
     {
      \lambda ^ \dagger
      -
      f ^ { ' }
      \left( 
        1 / \pi 
      \right)
     }
     \le
     \varepsilon
   \right]
   =
   \P
   \left[ 
     \inf _ { \norm{\Delta} = \varepsilon }
     G
     (
      f ^ { ' }
      \left( 
        1 / \pi 
      \right)
      +
      \Delta
     )
     -
     G
     (
      f ^ { ' }
      \left( 
        1 / \pi 
      \right)
     )
     \ge 
     0
   \right]
 \end{gather}

\begin{takeaways}
  To ensure double robustness, we leverage primal and dual optimization
  problem. The dual problem solves propensity score estimation and the 
  primal problem the bias of the final estimate.
  Partitioning estimates as in \cite{Gyorfi2002}. 
\end{takeaways}


  %We leverage (weak) universal consistency
of partitioning estimates~\cite{Gyorfi2002}.
To this end, we choose the basis functions as
\begin{gather}
  B_k(x)
  :=
  \frac{
  \mathbf{1}_{X_k \in A_n(x)}
  }{
  \sum_{j=1}^{n} 
  \mathbf{1}_{X_j \in A_n(x)}
  }
  \,,
  \qquad
  k=
  1,\ldots,n
  \,.
\end{gather}
The euclidian norm of 
$
  B(x)
$
is bounded by one.
\begin{gather}
  \norm{B(x)}^2
  =
  \sum_{k=1}^{n} 
  \left( 
  \frac{
  \mathbf{1}_{X_k \in A_n(x)}
  }{
  \sum_{j=1}^{n} 
  \mathbf{1}_{X_j \in A_n(x)}
  }
  \right)
  ^2
  \le
  \sum_{k=1}^{n} 
  \frac{
  \mathbf{1}_{X_k \in A_n(x)}
  }{
  \sum_{j=1}^{n} 
  \mathbf{1}_{X_j \in A_n(x)}
  }
  =1
  \,.
\end{gather}
The oracle parameters are
$
  \left[ 
    f^{'}
    \left( 
      \frac{1}{\pi_i}
    \right)
  \right] _ { i = 1,\ldots,n }
$
and
$
  \left[ 
    Y_i(1)
  \right] _ { i = 1,\ldots,n }
$
.
It even holds
\begin{gather}
  \lambda ^ \dagger _ i \to 
    f^{'}
    \left( 
      \frac{1}{\pi_i}
    \right)
\end{gather}
in probability.

 \begin{gather}
   G(\lambda)
   :=
      \frac{1}{n}
      \sum_{j = 1}^{n} 
        f^* 
        \left( 
          \inner
          {B(X_i)}
          {\lambda}
        \right)
      -
          \inner
          {B(X_i)}
          {\lambda}
      +
      \inner
      { | \lambda | }
      { \delta }
 \end{gather} 
 The consistency of 
 $
  \lambda ^ \dagger
 $
 relates to the difference being nonnegative.
 For all 
 $
  \varepsilon > 0
 $
 it holds
 \begin{gather}
   \P
   \left[ 
     \norm
     {
      \lambda ^ \dagger
      -
      f ^ { ' }
      \left( 
        1 / \pi 
      \right)
     }
     \le
     \varepsilon
   \right]
   =
   \P
   \left[ 
     \inf _ { \norm{\Delta} = \varepsilon }
     G
     (
      f ^ { ' }
      \left( 
        1 / \pi 
      \right)
      +
      \Delta
     )
     -
     G
     (
      f ^ { ' }
      \left( 
        1 / \pi 
      \right)
     )
     \ge 
     0
   \right]
 \end{gather}

\begin{takeaways}
  To ensure double robustness, we leverage primal and dual optimization
  problem. The dual problem solves propensity score estimation and the 
  primal problem the bias of the final estimate.
  Partitioning estimates as in \cite{Gyorfi2002}. 
\end{takeaways}


 % \section{Introduction}
  %We work in the Rubin Causal Model.

We assume a sample of $n$ units which is drawn from a population distribution.

In i.i.d. fashion.

We observe $ (\mathbf{X}_i, T_i, Y_i) $,
where $\mathbf{X}$ are covariates, 
$T$ is the indicator if treatment has been received
and $Y$ is the observed outcome.

In the Rubin Causal Model we assume that for each unit the potential outcome exist, i.e. $(Y_i^0, Y_i^1)$ where $Y^1$ stands for the potential outcome had the unit received treatment and $Y^0$ for the potential outcome had the unit received \textbf{no} treatment.

It is clear that $Y_i = Y_i^{T_i}$ i.e. we can observe only one of the potential outcomes.

Thus there is a connection to missing data problems.

This is the dilemma of causal inference.
 
On the population level it is possible to estimate both.

Usually the means of the potential outcomes are compared against each other.

In randomized trials this is a valid approach to causal inference.

In observational studies however the treatment assignment is not known and direct comparison can lead to systematically wrong results.

This phenomenon is called \textbf{confounding}.
 
To address the issue of confounding many methods have been proposed.
%%%%%%%%%%%%%%%%%%%%%%%%%%%%%%%%%%%%%%%%%

An intuitive way to think about potential outcomes is to think of a stochastic process $Y(\cdot)$ indexed over $\left\{ 0,1 \right\}$.
By observing $Y_i$ we in fact sample from this process at random index $T$, i.e. from $Y(T)$.
We have 
\begin{gather}
  \E[Y(T)]
  =
  \E[Y(1) \vert T=1] \P[T=1]
    +
  \E[Y(0) \vert T=0] \P[T=0]
  .
\end{gather}
Suppose we observe $T=1$. Clearly we have
\begin{gather}
  \E[Y(T)\vert T=1]
  =
  \E[Y(1) \vert T=1] 
\end{gather}

%\section{Double Robustness}
%%read the paper \cite{Zhao2017a} for discussion of double robustness for balancing weights and \cite{Hahn1998} for semiparametric efficiency bounds.


\subsection{Learning Rates of the weighted mean}

We begin the discussion with an
uncomplicated question:
What is the speed of convergence in the weak law of large numbers?
The next statement gives a clear-cut answer:
The arithmetic mean of independent, identically distributed, square-integrable random variables 
learns with rate $n^{-1/2}$.
Furthermore, the statement is easy to prove using
Bienaymé's formula and Chebyshev's inequality (cf. \cite[Theorem~5.14]{Klenke2020}).


\begin{theorem*}
  Let 
  $
    X_1,X_2,\ldots
  $
  be i.i.d, square-integrable random variables with 
  $
    V:=
    \mathbf{Var}[X_1]
    <\infty
  $.
  Then, for any $\tau \in (0,1]$ and all $n\in\mathbb{N}$, we have
  \begin{gather}
   \P
   \left[
     \left| 
   \frac{1}{n}
   \sum_{i=1}^{n}
   (X_i - \E[X_i])
     \right|
     \le
     \sqrt{V}
     \frac{1}{\sqrt{\tau}}
     \frac{1}{\sqrt{n}}
   \right]
   \ge
   1-\tau\,.
  \end{gather}
\end{theorem*}

Deriving learning rates in this way, we make an implicit assumption.
We assume that
observed outcomes of the treated follow the same distribution as
marginal potential outcomes under treatment.
In other words, we require
\begin{gather}
  Y(1)\,|\,T=1
  \ 
  \sim 
  \ 
  \P_
  {
  Y(1)
  }
  \,.
\end{gather}
Virtually every scenario of interest violates this assumption.
Indeed, any external influence on both $T$ and $Y(1)$ can ruin the 
above assessment simply by imposing 
\begin{gather}
  \E[Y(T)]\ \neq\  \E[Y(1)]
  \,.
\end{gather}
Statisticians have wrestled with this issue for nearly a century.

In experimental studies we usually specify treatment assignment as opposed to merely observing a unit receiving treatment. 

The next statement makes use of the propensity score.

\begin{theorem*}
  Consider the weighted mean estimator with weights
  \begin{gather}
    w_i
    \ 
    =
    \ 
    \frac{1}{n}
    \frac{T_i}{\pi(X_i)}
    \,.
  \end{gather}
  Denote
  $
    V:=
    \E[
    (
      Y(1)
    )^2
    \,/\,
    \pi(X)
    ]
    -
    \E[Y(1)]^2
  $.  
  Assume that \textit{weak unconfoundedness} holds.
  Then, for any $\tau \in (0,1]$ and all $n\in\mathbb{N}$, we have
  \begin{gather}
   \P
   \left[
     \left| 
   \sum_{i=1}^{n}
   w_iY_i - \E[Y(1)]
     \right|
     \le
     \sqrt{V}
     \frac{1}{\sqrt{\tau}}
     \frac{1}{\sqrt{n}}
   \right]
   \ge
   1-\tau\,.
  \end{gather}

\end{theorem*}
\begin{proof}
We want to reinforce the coherent use of the weak law of large numbers.
To this end, we verify 
\begin{align*}
  n\,\E[w(T,X)\,Y(T)]
  \ 
  &=
  \ 
  \E[Y(1)]
  \,,
  \\
  n^2\,
  \mathbf{Var}[w(T,X)\,Y(T)]
  \ 
  &=
  \ 
    \E[
    (
      Y(1)
    )^2
    \,/\,
    \pi(X)
    ]
    -
    \E[Y(1)]^2
    \,.
\end{align*}
Essentially, the random weight $w(T,X)$ acts on $Y(T)$ through $T\,/\,\pi(X)$.
It does so by inducing independence of observed outcome $Y(T)$ and treatment $T$.
This effect can only manifest itself if weak unconfoundedness holds, i.e., 
\begin{gather}
  (Y(0),Y(1))\perp\!\!\!\perp T\, |\, X\,.
\end{gather}
To showcase this detail we added the $n$ and $n^2$ factors in the above display.
The calculations go as follows.
\begin{align}
  \begin{split}
  n\,\E[w(T,X)\,Y(T)]
  &\ 
  =
  \ 
  \E
  \left[ 
    Y(T)
    \cdot
    (T\,/\,\pi(X))
  \right]
  \\
  &\ =\ 
  \E
  \left[ 
    Y(1)
    \,
    /
    \,
    \pi(X)
    \,
    \vert
    \,
    T=1
  \right]
  \cdot
  \P[T=1]
  \\
  &\ =\ 
  \int_\mathcal{X}
  \E
  \left[ 
    Y(1)
    \,
    \vert
    \,
    X=x
    ,
    T=1
  \right]
  \cdot
  (
  \P[T=1]
  \,
  /
  \,
  \pi(x)
  )
  \,
  \P_{X|T}(dx\,|\,1)
  \\
  &\ =\ 
  \int_\mathcal{X}
  \left[ 
    Y(1)
    \vert
    X=x
  \right]
  \P_X(dx)
  =
  \E[Y(1)]
  .
\end{split}
\end{align}
The first equality comes from the definition of the weights.
The second, third and fifth equality stem from 
$
  T\in \left\{ 0,1 \right\}
$,
and the law of total expectation, applied with $T$ and $X$.
The fourth equality is due to the assumption of weak unconfoundedness.
The density transformation is due to Bayes's Theorem.
With slight modifications in the above argument, it follows
\begin{gather}
  n^2
  \E
  \left[ 
  \Big( 
    Y(T)
    \cdot
    (T\,/\,\pi(X))
  \Big)
  ^2
  \right]
  \ 
    =
    \ 
    \E
    \Big[
    (Y(1))^2
    \,
    /
    \,
    \pi(X)
    \Big]
    \,.
\end{gather}
We omit the details.
Invoking the weak law of large numbers finishes the proof.
\end{proof}

We started by asking an easy question, so it is time for a more challenging one: How do we proceed in deriving learning rates if the propensity score is unknown.
How do we generally procede?
[what has been done in the past. why are some methods obsolete]
A naive answer would be: We hope to select a proper model and try to estimate the propensity score.
Stunningly, a lot of practicioners still settle for obsolete methods when it comes to propensity score analysis.




%  \section{Error Decompositions}
%  The following decomposition is flexible in $\Phi.$
We get different causal estimaands $\E[\Phi(Y(1))],$
e.g.
the popolation average of $Y(1)$
for 
$
  \Phi(Y)=Y,
$
i.e.
$
\E[Y(1)]
,
$
ot the distribution function of $Y(1)$ at $t$
for 
$
  \Phi(Y)=
  \mathbf{1}_{(-\infty, t]}
  (Y)
  ,
$
i.e.
$
\P[
Y(1)\le t
]
.
$
\begin{gather}
 \sum_{i=1}^{n}  
 w_i
 T_i
 \Phi(Y_i)
  - \E [ \Phi(Y(1)) ]
  =
  \frac{1}{n}
  \sum_{ i = 1 }^{n} S_i
    + R_0
    + R_1
    + R_2
    ,
\end{gather}
where
\begin{align*}
  S_i 
  &:= 
  \frac{T_i}{\pi_i}
 \left( 
   \Phi(Y_i) - \E [ \Phi(Y_i(1)) | X_i ]
 \right)
 +
 \left( 
   \E [ \Phi(Y_i(1)) | X_i ] - \E [ \Phi(Y(1))]
 \right)
 \quad
 \text{for}\ 
 i \in \left\{ 1, \ldots, n, \right\}
 ,
 \\
  R_0
  &:=
  \sum_{ i = 1 }^{n}
  T_i
    \left(  
      w_i - \frac{1}{n \pi_i}
    \right)
 \left( 
   \Phi(Y_i) - \E [ \Phi(Y_i(1)) | X_i ]
 \right)
 ,\\
  R_1
  &:=
  \sum_{ i = 1 }^{n}
    \left(  
      T_i
      w_i - \frac{1}{n}
    \right)
 \left( 
    \E [ \Phi(Y_i(1)) | X_i ] - B(X_i)^\top \lambda
 \right)
 ,
 \\
  R_2
  &:=
  \sum_{ i = 1 }^{n}
    \left(  
      T_i
      w_i - \frac{1}{n}
    \right)
 B(X_i)^\top \lambda
 \quad 
 \text{for}
 \ 
 \lambda \in \R^K
 .
\end{align*}

We can even view
$
\frac{1}{\sqrt{n}}
\sum_{i=1}^{n}S_i 
$
as an empirical process 
$
\mathbb{G}_n f
$
indexed over 

\begin{gather}
  f_\Phi(T,X,Y)
  =
  \frac{T}{\pi(X)}
 \left( 
   \Phi(Y) - \E [ \Phi(Y) | X ]
 \right)
 +
   \E [ \Phi(Y) | X ] 
   .
\end{gather}
If $\mathcal{F}=\left\{ f_\Phi \colon \Phi \in \ \text{some set}\right\}$
is $\P$-Donsker, the empirical process converges to a tight gaussian process.
Then the functional delta Method is applicable.

%  %\section{Estimating the Population Mean of Potential Outcomes}
%  %We want to estimate the populition mean of the outcome under treatment,
i.e. $\E \left[ Y^1 \right]$.

Since $Y_i^1$ is only observed for the treated units, i.e. if $T_i = 1$
we will consider a weighted mean of the observed outcomes as an estimator, i.e. $\hat{Y}_w^1 = \sum_{i = 1}^{n} w_i T_i Y_i $
where we use convex optimization to compute the weights.

We consider the following decomposition

\begin{gather}
  \hat{Y}_w^1 - \E [ Y^1 ]
  =
  \frac{1}{n}
  \sum_{ i = 1 }^{n} S_i
    + R_0
    + R_1
    + R_2
    ,
\end{gather}
where
\begin{gather*}
  S_i 
  := 
  \frac{T_i}{\pi_i}
 \left( 
    Y_i - \E [ Y^1_i | X_i ]
 \right)
 +
 \left( 
    \E [ Y^1_i | X_i ] - \E [ Y^1]
 \right)
 \quad
 \text{for}\ 
 i \in \left\{ 1, \ldots, n, \right\}
 ,
 \\
  R_0
  :=
  \sum_{ i = 1 }^{n}
  T_i
    \left(  
      w_i - \frac{1}{n \pi_i}
    \right)
 \left( 
    Y_i - \E [ Y^1_i | X_i ]
 \right)
 ,\\
  R_1
  :=
  \sum_{ i = 1 }^{n}
    \left(  
      T_i
      w_i - \frac{1}{n}
    \right)
 \left( 
    \E [ Y^1_i | X_i ] - B(X_i)^\top \lambda
 \right)
 \quad
 \text{and}
 \quad
  R_2
  :=
  \sum_{ i = 1 }^{n}
    \left(  
      T_i
      w_i - \frac{1}{n}
    \right)
 B(X_i)^\top \lambda
 \quad 
 \text{for}
 \ 
 \lambda \in \R^K
 .
\end{gather*}

We want to prove asyptotic normality

\begin{ftheorem}
  Suppose that conditions hold. Then
  \begin{gather*}
    \sqrt{n} 
    \left(
      \hat{Y}_{w^*}^1 - \E [ Y^1 ]
    \right)
    \overset{\mathcal{D}}{\longrightarrow}
    Z \sim \mathcal{N}(0,\sigma_*^2)
    .
  \end{gather*}
\end{ftheorem}
To accomplish this we need

\begin{ftheorem}
  If $T_i = 1$ then $w^*(X_i)$ is a consistent estimator of $\frac{1}{n \pi(X_i)}$.
\end{ftheorem}


We study the following problem:
%%%%%%%%%%%%%%%%%%
% PROBLEM PRIMAL %
%%%%%%%%%%%%%%%%%%
\begin{problem}
  \begin{gather*}
    \underset{w_1, \ldots, w_n \in \R}{\mathrm{minimize}}
    \qquad
    \sum_{i = 1}^{n} T_i f(w_i)
  \end{gather*}
subject to the constraints
\begin{gather*}
    w_i T_i \ge 0,
    \qquad 
    i = 1, \ldots, n,
\end{gather*}
\begin{gather*}
  \sum_{ i = 1 }^{n}
    w_i T_i
  =
  1
\end{gather*}
\begin{gather*}
    \left| 
      \sum_{i = 1}^{n} w_i T_i B_k(X_i)
      - 
      \frac{1}{n} \sum_{i = 1}^{n} B_k(X_i)
    \right|
    \le 
    \delta_k,
    \qquad
    k = 1, \ldots, K
\end{gather*}
\end{problem}
%%%%%%%%%%%%%%%%%%%%%%%%%%%%%%%%%%%%%%%%%%%%%%%%%%%%%%%%%%%%%%%%%

We aim to prove that the solution to 
Problem \eqref{primal_weighting_binary}
is asymptotical consistent with the propensity score, i.e.

\begin{theorem}
  \label{ch_theorem_1}
  Under some (non-optimal) Assumptions,
  there exist constants
  $
    c_{1}
    , 
    c_{2}
    >
    0
  $
  and decreasing sequences
  $
    (\varepsilon_n^1)
    ,
    (\varepsilon_n^2)
    \subset
    (0,1]
  $
  that converge to 0
  such that
  for all
  $
    \tau
    \in
    (0,1]
  $
  there exists a constant
  $
    c_{\tau}
    \in
    [0,\infty)
  $
  only depending on 
  $ \tau $
  such that for all 
  $ n \ge 1 $
  and
  $
    \tau
    \in
    (0,1]
  $
  it holds
  \begin{align}
    \label{ch_theorem_1_alternative}
    \begin{split}
   & \P
    \left( 
  \norm{
    w^*_i
    -
    \frac{1}{n \pi(X_i)}
  }_\infty
  \le 
  c_{1}
  c_{\tau}
  \varepsilon_{n}^1 
   \right)
   \ge
  1 - \tau
  ,
\\
   &\quad \
   \norm{
    w^*_i
    -
    \frac{1}{n \pi(X_i)}
  }_{\P, 2}
  \le 
  c_{2}
  \varepsilon_n^{2},
      \end{split}
  \end{align}
  where 
  $w^*$
  is the solution to 
  Problem \eqref{primal_weighting_binary}.
\end{theorem}



%%%%%%%%%%%%%%%%%%%ASSUMPTION 1%%%%%%%%%%%%%%%%%%%%%%%%%%%%%%%%%%%%%%%


\begin{assumption}
  Assume, the following conditions hold:
  %
  \begin{subassumption}
    \label{assumption_1_i} 
    The minimizer 
    $
    \lambda_0 
    =
    \arg \min_{\lambda \in \Theta}
    \E
    \left[ 
      -T n 
      \rho 
      \left( 
      B(X)^T \lambda
      \right)
      +
      B(X)^T \lambda
    \right]
    $
    is unique,
    where 
    $\Theta \subseteq \R^n$ is the parameter space for $\lambda$.
  \end{subassumption}
  %
  \begin{subassumption}
    \label{assumption_1_ii} 
    The parameter space 
    $\Theta \subseteq \R^n$
    is compact.
  \end{subassumption}
  %
  \begin{subassumption}
    \label{assumption_1_iii}
    $\lambda_0 \in \text{int}(\Theta)$,
    where
    $\text{int}(\cdot)$
    stands for the interior of a set.
  \end{subassumption}
  %
  \begin{subassumption}
    \label{assumption_1_iv}
    There exists
    $\lambda^*_1 \in \Theta$
    such that
    $
      \norm{
        m^*(\cdot)
        -
        B(\cdot)^T \lambda^*_1
      }_\infty
      \le 
      \varphi_{m^*}
    $,
    where
    $
      m^*(\cdot)
      :=
      \left( \rho^{'} \right)^{-1}
      \left( \frac{1}{n \pi(\cdot)} \right).
    $
  \end{subassumption}
  %
  \begin{subassumption}
    \label{assumption_1_v}
    There exists a constant 
    $
      \varphi_{\pi} 
      \in 
      \left(0, \frac{1}{2} \right)
    $
    such that
    $
      \pi(x)
      \in 
      (
      \varphi_{ \pi },
      1 - \varphi_{ \pi }
      )
    $
    for all
    $
      x
      \in
      \mathcal{X}
    $
  \end{subassumption}
  %
  \begin{subassumption}
    \label{assumption_1_vi}
    There exists 
    $ \varphi_{\rho^{''}} > 0 $
    such that
    $ -\rho^{''} \ge \varphi_{\rho^{''}} > 0 $
  \end{subassumption}
  %
  \begin{subassumption}
    \label{assumption_1_vii}
    There exists 
    $ \varphi_{B(x) B(x)^T} > 0 $ 
    such that
    $
      B(x) B(x)^T 
      \succcurlyeq 
      \varphi_{B(x) B(x)^T} I 
    $
  \end{subassumption}
  %
  \begin{subassumption}
    \label{assumption_1_viii}
    There exists
    $ \varphi_{\norm{B}} > 0 $
    such that
    $
      \sup_{x \in \mathcal{X}} \norm{B(x)}_2
      \le 
      \varphi_{\norm{B}}
    $.
  \end{subassumption}
  %
  \begin{subassumption}
    \label{assumption_1_ix}
    The number of basis functions
    satisfies
    $
      K = o(n)
    $.
  \end{subassumption}
\end{assumption}


\section*{Plan of Proof}

It is easier to study the dual of 
Problem \eqref{primal_weighting_binary}.
Thus we employ results from convex analysis
\cite{Mordukhovich2022}
to establish
\begin{proposition}
  \label{ch_1_dual}
  The dual of Problem \eqref{primal_weighting_binary}
  is equivalent to the unconstrained optimization problem
  \begin{gather}
    \label{dual_weighting_binary}
    \underset{\lambda \in \R^K}{\text{minimize}}
    \quad
    \frac{1}{n}
    \sum_{j = 1}^{n} 
      \left[ 
        -T_j n 
        \rho 
        \left( 
          B(X_j)^T \lambda
        \right)
        +
        B(X_j)^T \lambda
      \right]
    +
    |\lambda|^T \delta,
  \end{gather}
  where 
  $
    B(X_j)=(B_k(X_j))_{1 \le k \le K}
  $
  denotes the $K$ basis functions of the covariates,
  $
    \rho(t) 
    :=
    \frac{t}{n}
    -
    t
    (h^{'})^{-1}(t)
    +
    h((h^{'})^{-1}(t))
  $
  with
  $
    h(x)
    :=
    f%
    \left( 
      \frac{1}{n}
      -
      x
    \right)
  $
  and
  $
    | \lambda | 
    :=
    \left(
    | \lambda_k | 
    \right)_{1 \le k \le K}
  $.
  Moreover,
  the primal solution
  $
    w^*_j
  $
  satisfies
  \begin{gather}
    \label{ch_1_w*}
    w^*_j
    =
    \rho^{'}
    \left(
      B(X_j)^T
      \lambda^\dagger
    \right)
  \end{gather}
  for 
  $
    j
    =
    1, \ldots , n
  $,
  where 
  $
    \lambda^\dagger
  $
  is the solution to the dual optimization problem.
\end{proposition}
%%%%%%%%%%%%%%%%%%%%%%%%%%%%%%%%%%%%%%%%%%%%%%%%%%%%%%%%%%%%%%%%%%%%%
%%%%%%%%%%%%%%%%%%%%%%%%%%%%%%%%%%%%%%%%%%%%%%%%%%%%%%%%%%%%%%%%%%

The core of the subsequent analysis is based on
Assumption~\ref{assumption_1_iv}, 
i.e.
the existence of an oracle parameter
$\lambda_1^{*}$
in a sieve estimate of the true propensity score (or a transformation). 
It is then natural to enquire about the convergence of the dual solution
$
  \lambda^\dagger
$
to 
$
  \lambda^*_1
$.
Making certain assumptions and employing
matrix concentration inequalitys \cite{Tropp2015}
we can establish
\begin{proposition}
  \label{ch_1_near_oracle}
  Under some (non-optimal) Assumptions,
  there exists a constant
  $
    c_{3}
    >
    0
  $
  and a decreasing sequence
  $
    (\varepsilon_n^3)
    \subset
    (0,1]
  $
  that converges to 0
  such that
  for all
  $
    \tau
    \in
    (0,1]
  $
  there exists a constant
  $
    \tilde{
    c_{\tau}
    }
    \in
    [0,\infty)
  $
  only depending on 
  $ \tau $
  such that for all 
  $ n \ge 1 $
  and
  $
    \tau
    \in
    (0,1]
  $
  it holds
  \begin{gather}
    \P
    \left( 
      \norm{
        \lambda^\dagger
        -
        \lambda^*_1
      }_2
      \le
      c^3
      \tilde{
        c_{\tau}
      }
      (
        \varepsilon_n^3
      )
    \right)
    \ge 
    1 - \tau.
  \end{gather}
\end{proposition}%
It is then straightforward to prove a more general result
than
Theorem~\ref{ch_theorem_1}%
.
%
\begin{theorem}
  Under some (non-optimal) Assumptions,
  there exist constants
  $
    c_{1}
    , 
    c_{2}
    >
    0
  $
  and decreasing sequences
  $
    (\varepsilon_n^1)
    ,
    (\varepsilon_n^2)
    \subset
    (0,1]
  $
  that converge to 0
  such that
  for all
  $
    \tau
    \in
    (0,1]
  $
  there exists a constant
  $
    c_{\tau}
    \in
    [0,\infty)
  $
  only depending on 
  $ \tau $
  such that for all 
  $ n \ge 1 $
  and
  $
    \tau
    \in
    (0,1]
  $
  it holds
  \begin{align*}
   & \P
    \left( 
  \norm{
    w^*(\cdot)
    -
    \frac{1}{n \pi(\cdot)}
  }_\infty
  \le 
  c_{1}
  c_{\tau}
  \varepsilon_{n}^1 
   \right)
   \ge
  1 - \tau
  ,
\\
   &\quad \
   \norm{
    w^*(X)
    -
    \frac{1}{n \pi(X)}
  }_{\P, 2}
  \le 
  c_{2}
  \varepsilon_n^{2},
  \end{align*}
  where 
  $w^*(X)$
  is as in \eqref{ch_1_w*}
  without the index.
\end{theorem}

\section*{Proof of theorem 2.2}


\begin{proof}
  Motivated by Proposition~\ref{syu_1_result}
  we consider
 \begin{gather}
   \label{ch_1_G}
   G(\lambda)
   :=
      \frac{1}{n}
      \sum_{j = 1}^{n} 
      \left[ 
        -T_j n 
        \rho 
        \left( 
          B(X_j)^T \lambda
        \right)
      +
      B(X_j)^T \lambda
      \right]
      +
      |\lambda|^T \delta.
 \end{gather} 
 Since 
 $\rho \in C^2(\R)$
 we can employ 
 \eqref{ch_1_G},
 Corollary~\ref{syu_taylor_corollary}
 and
 Proposition~\ref{syu_triangle}
 to get
 \begin{align}
   \label{ch_1_proof1_1}
  \begin{split}
  & 
  G(\lambda^*_1 + \Delta) 
  -
  G(\lambda^*_1)
  \\
  &
  \quad
  \ge
      \frac{1}{n}
      \sum_{j = 1}^{n} 
      \left[ 
        -T_j n 
        \rho^{'} 
        \left( 
          B(X_j)^T \lambda^*_1
        \right)
      +
      1
      \right]
      \Delta^T B(X_j)
      \\
  & 
  \qquad
  +
      \frac{1}{2}
      \sum_{j = 1}^{n} 
        -T_j  
        \rho^{''} 
        \left( 
          B(X_j)^T (\lambda^*_1 + \xi \Delta)
        \right)
        \Delta^T
        \left( 
          B(X_j)
          B(X_j)^T
        \right)
        \Delta
        \\
  &
  \qquad
  -
      |\Delta|^T \delta
  \\
  &  
  \quad
  \ge
    - \norm{\Delta}_2
    \left( 
    \norm{
      \frac{1}{n}
      \sum_{j = 1}^{n} 
      \left[ 
        -T_j n 
        \rho^{'} 
        \left( 
          B(X_j)^T \lambda^*_1
        \right)
      +
      1
      \right]
      B(X_j)
    }_2
    +
    \norm{\delta}_2
    \right)
    \\
  & 
  \qquad
  +
  n
  \norm{\Delta}^2_2
   \varphi_{\rho^{''}}
  \underline{\varphi_{BB^T}}
  \\
  &
  \quad
  :=
  -\norm{\Delta}_2
  (I_1 + \norm{\delta}_2)
  +
  \norm{\Delta}^2_2
  I_2
  .
  \end{split}
\end{align}
The second inequality is due to 
the Cauchy-Schwarz-Inequality 
and
Assumptions~\ref{assumption_1_vi} and \ref{assumption_1_vii}
.
We want to establish probabilistic upper bounds of the factor associated with 
$
  -\norm{\Delta}_2
$. 
This will be done with appropriate assumptions on
$
  \norm{\delta}_2
$
and a thorough analysis of $I_1$.
If we then restrict lower bounds of
$
  I_2
$
to approprietely slow convergence to 0,
e.g. by assumptions on
$
\varphi_{\rho^{''}}
$
and
$
\underline{
  \varphi_{BB^T}
}
$,
we can choose
$
  \norm{\Delta}_2
$
large enough,
such that 
\eqref{ch_1_proof1_1}
yields
$
  G(\lambda^*_1 + \Delta) 
  -
  G(\lambda^*_1)
  >
  0
$
with arbitrarily large probabitity for $n$ large enough.
With Proposition~\ref{syu_1_result} it follows then immediately
Proposition~\ref{ch_1_near_oracle}
.
\subsection*{Analysis of $I_1$}
We want to use Assumption~\ref{assumption_1_iii}.
Thus we perform the following split:
\begin{align}
  \label{ch_1_i1}
  \begin{split}
  I_1 
  &\le
    \norm{
      \sum_{j = 1}^{n} 
        T_j  
      \left[ 
        \rho^{'} 
        \left( 
          B(X_j)^T \lambda^*_1
        \right)
      -
      \frac{1}{n \pi(X_j)}
      \right]
      B(X_j)
    }_2
  \\
  &+
    \norm{
      \frac{1}{n}
      \sum_{j = 1}^{n} 
      \left[ 
        \frac{T_j}{\pi(X_j)}
      -
      1
      \right]
      B(X_j)
    }_2
    \\
  &=:
  J_1 + J_2
  \end{split}
\end{align}

\subsubsection*{Analysis of $J_1$}

By the Lipschitz-continuity of 
$\rho^{'}$,
Assumption~\ref{assumption_1_viii}
and
Assumption~\ref{assumption_1_iv},
$T \in \{0, 1\}$
and 
the triangle inequality 
we have
\begin{gather}
  \label{ch_1_j1_result}
  J_1 
  \le
  n L_{\rho^{'}}\varphi_{\norm{B(x)}} \varphi_{m^*}
\end{gather}

\subsubsection*{Analysis of $J_2$}
We want to employ 
Theorem~\ref{rmineq_bernstein}%
. %
To this end we define
the independent random matrices
\begin{align}
  \label{ch_1_j2_1}
  \begin{split}
  A_j
  &:=
      \frac{1}{n}
      \left[ 
        \frac{T_j}{\pi(X_j)}
      -
      1
      \right]
      B(X_j)
      \,,
      \quad 
      j=1,\ldots,n,
\\
  S
  &:=
  \sum_{j = 1}^{n} 
  A_j
  \end{split}
\end{align}
and check conditions 
\eqref{rmineq_bernstein_cond_1} 
and
\eqref{rmineq_bernstein_cond_2}.
Note that 
$
  \norm{S}_2
  =
  J_2
$
.
By the properties of conditional expectation 
it holds
\begin{gather}
  \label{ch_1_j2_2}
  \E
  \left[  
    \frac{T_j}{\pi(X_j)}
    B(X_j)
  \right]
  =
  \E
  \left[  
    \E
    \left[  
      T_j
      \, | \,
      X_j
    \right]
    \frac{1}{\pi(X_j)}
    B(X_j)
  \right]
  =
  \E[B(X_j)]. 
\end{gather}
Taking the expectation in 
\eqref{ch_1_j2_1} 
and using
\eqref{ch_1_j2_2}
we get
$\E[A_j] = 0$
for all 
$
  j=1,\ldots,n
$
.
Since
\begin{gather}
  \label{ch_1_j2_3}
  \left| 
      \frac{T_j}{\pi(X_j)}
      -
      1
      \right|
  \le
  1 + \frac{1 - \varphi_{\pi}}{\varphi_{\pi}}
  =
  \frac{1}{\varphi_{\pi}}
\end{gather}
by Assumption~\ref{assumption_1_v},
we can employ Assumption~\ref{assumption_1_viii}
together with 
\eqref{ch_1_j2_3}
and
\eqref{ch_1_j2_1}
to get
\begin{gather}
  \label{ch_1_j2_4}
  \norm{A_j}_2
  \le
  \frac{\varphi_{\norm{B}}}{n \varphi_{\pi}}
  =:
  L.
\end{gather}
Thus,
condition \eqref{rmineq_bernstein_cond_1}
is satisfied.
Next we turn to the matrix variance statistic
$ v(S) $
\eqref{rmineq_bernstein_cond_2}.
By 
\eqref{ch_1_j2_1}
and
\eqref{ch_1_j2_3}
we have
\begin{gather}
  \label{ch_1_j2_5}
  \E
  \left[ 
    A_j A_j^T
  \right]
  \le
  \left( 
    \frac{1}{n \varphi_{\pi}}
  \right)^2
  \E
  \left[ 
    B(X) B(X)^T
  \right]
\end{gather}
and by \eqref{ch_1_j2_4} 
\begin{gather}
  \label{ch_1_j2_6}
  \E
  \left[ 
    A_j^T A_j
  \right]
  \le
  L^2.
\end{gather}
Since 
$
  \max\{a,b\}
  \le
  | a | + | b |
$
we can use 
\eqref{ch_1_j2_5}
and
\eqref{ch_1_j2_6}
to get
\begin{gather}
  \label{ch_1_j2_7}
  v(S)
  \le
  \frac{1}{n}
  \frac{\lambda_{\max}}{\varphi_{\pi}^2}
  +
  n L^2,
\end{gather}
where 
$\lambda_{\max}$ 
is the maximal eigenvalue of
the symmetric (non-random) matrix
$
  \E
  \left[ 
    B(X) B(X)^T
  \right]
$.
Having dealt with 
\eqref{rmineq_bernstein_cond_1} 
and
\eqref{rmineq_bernstein_cond_2}
we can establish 
the
expectation bound
\eqref{rmineq_bernstein_expectation_bound}
of Theorem~\ref{rmineq_bernstein}.
Together with 
\eqref{ch_1_j2_4}
and
\eqref{ch_1_j2_7}
we get
\begin{align}
  \label{ch_1_j2_8}
  \begin{split}
    &\E[J_2]
\\
  &\le
  \sqrt{
    \frac{
    2 \log (K + 1)
    \left( 
      \lambda_{\max}  + \varphi_{\norm{B}}^2
    \right)
    }
    {
      n \varphi_{\pi}^2
    }
  }
  +
  \frac{
    \log (K + 1)
    \varphi_{\norm{B}}
  }
  {
    3 n \varphi_{\pi}
  }
\\
  &\le
  \frac{1}{\varphi_{\pi}}
  \sqrt{
    \frac{
      \log(K + 1)
    }
    {
      n
    }
  }
  \left[
    \varphi_{\norm{B}}
    \left(
      \sqrt{2}
      +
      \frac{1}{3}
      \sqrt{
        \frac{\log (K + 1)}{n}
      }
    \right)
    +
    \sqrt{
      2
      \lambda_{\max}
    }
  \right].
\end{split}
\end{align}
Since
$
  K = o(n)
$
by Assumption~\ref{assumption_1_ix}
we can discuss the other influences on the quality of the bound
\eqref{ch_1_j2_8}.
On a high-level it is readily clear
that appropriate bounds on
$
\varphi_{\pi},
\varphi_{\norm{B}}
$
and
$
  \lambda_{\max}
$
will shrink
$
  \E[J_2]
$
to 0 
and will assist in establishing learning rates.

We could also have invoked the probabitity bound 
\eqref{rmineq_bernstein_probability_bound} 
of Theorem~\ref{rmineq_bernstein}.
But for the sake of simplicity we prefer the combination 
of the expectation bound 
\eqref{ch_1_j2_8}
and the Markov inequality.
With the latter we get
\begin{align}
  \label{ch_1_j2_result}
  J_2
  \le
  \frac{1}{\tau}
  \frac{1}{\varphi_{\pi}}
  \sqrt{
    \frac{
      \log(K + 1)
    }
    {
      n
    }
  }
  \left[
    \varphi_{\norm{B}}
    \left(
      \sqrt{2}
      +
      \frac{1}{3}
      \sqrt{
        \frac{\log (K + 1)}{n}
      }
    \right)
    +
    \sqrt{
      2
      \lambda_{\max}
    }
  \right]
\end{align}
with probabitity 
$
  \ge
  1 - \tau
$.

%%%%%%%%%%%%%%%%%%%%%%%%%%%%%%%%%%%%%


If we choose
$
\norm{\Delta}_2
$
to be
\begin{align}
  \label{ch_1_choice_norm_Delta}
  \begin{split} 
    &\left(
     \sqrt{2}
    \frac{1}{\tau}
    \frac{1}{\varphi_{\pi}}
    \sqrt{
      \frac{
        \log(K + 1)
      }
      {
        n^3
      }
    }
    \left[
      \varphi_{\norm{B}}
      \left(
        1
        +
        \sqrt{
          \frac{\log (K + 1)}{n}
        }
      \right)
      +
      \sqrt{
        \lambda_{\max}
      }
  \right]
  \right.
\\
  &
  \quad
  +
     L_{\rho^{'}}
     \varphi_{\norm{B}}
     \varphi_{m^*}
  +
  \left.
  \frac{
    \norm{\delta}_2
  }
  {n}
  \right)
  \frac{1}
  {
    \varphi_{\rho^{''}}
    \underline{
      \varphi_{BB^T}
    }
  }
\end{split}
\end{align}
we get by
\eqref{ch_1_proof1_1},
\eqref{ch_1_i1},
\eqref{ch_1_j1_result},
\eqref{ch_1_j2_result}
and 
Proposition~\ref{syu_1_result}
\begin{align}
  \label{ch_1_near_oracle_probabilistic_bound}
  \begin{split}
  \P
  \left( 
      \norm{
        \lambda^\dagger
        -
        \lambda^*_1
      }_2
      \le
      C 
  \right)
  &=
  \P
  \left( 
    \inf_{\norm{\Delta}_2 = C}
    G(\lambda^*_1 + \Delta) 
    -
    G(\lambda^*_1)
    >
    0
  \right)
  \\
  &\ge
  1 - \tau,
  \end{split}
\end{align}
where 
$
  C
$
is as in 
\eqref{ch_1_choice_norm_Delta}.
With appropriate Assumptions (as discussed before) we can then establish
Proposition~\ref{ch_1_near_oracle}.

We can invoke 
\eqref{ch_1_near_oracle_probabilistic_bound}
to derive bounds as in 
Theorem~\ref{ch_theorem_1_alternative}:
\begin{align*}
  \norm{
    w^*(X)
    -
    \frac{1}{n \pi(X)}
  }_{\P, 2}
  &\le
  L_{\rho^{'}}
  \left[ 
    \norm{
      B(X)^T 
      \left(
        \lambda^\dagger
        -
        \lambda^*_1
      \right)
    }_{\P,2}
    \right.
    \\
    &+
    \left. 
    \norm{
      m^*(X)
      -
      B(X)^T 
      \lambda^*_1
    }_{\P,2}
  \right]
  \\
  &\le
  L_{\rho^{'}}
  \left(
    \varphi_{\norm{B}}
    \sqrt{
      C^2(1 - \tau)
      +
      \text{diam}(\Theta)^2
      \tau
    } 
    +
    \varphi_{m^*}
  \right)
\end{align*}
%%%%%%%%%%%%%%%%%%%
\begin{align*}
  \norm{
    w^*(\cdot)
    -
    \frac{1}{n \pi(\cdot)}
  }_{\infty}
  &\le
  L_{\rho^{'}}
  \left[ 
    \norm{
      B(\cdot)^T 
      \left(
        \lambda^\dagger
        -
        \lambda^*_1
      \right)
    }_{\infty}
    \right.
    \\
    &+
    \left. 
    \norm{
      m^*(\cdot)
      -
      B(\cdot)^T 
      \lambda^*_1
    }_{\infty}
  \right]
  \\
  &\le
  L_{\rho^{'}}
  \left(
    \varphi_{\norm{B}}
    C
    +
    \varphi_{m^*}
  \right)
\end{align*}
with probabitity greater than 
$1 - \tau$.
\end{proof}
\begin{remark}
  By Corollary~\ref{rmineq_intrinsic_bernstein_vector_corollary}
  we can get rid of the $\log(K)$ term in
\eqref{ch_1_choice_norm_Delta}.
\end{remark}
\begin{remark}
  By the matrix Rosenthal-Pinelis Inequality
  \cite{Chen2012}[Thm.A.1]
  we can weaken Assumption~\ref{assumption_1_v} to a lower bound on the expectition of 
  $\pi(X)$
\end{remark}
The next step consists of strenghtening the Assumptions to get concrete learning rates.
This can be done in a series of examples.

%  \section{Application of Matrix Concentration Inequalities}
%  We extend the analysis to covariates with unbounded support. To the best of our knowledge this has not been done in this setting.
Indeed, \cite[Assumption 1]{Wang2019} only allows for covariates with bounded support.
Our idea is to apply Matrix Moment Inequalities instead of Bernstein's inequality. We learned from this possibility from \cite[(6.1.6)]{Tropp2015}
\subsubsection*{
  Analysis of 
  $
  \E[\max_{i \le n}\norm{A_i}^2]
  $
}
We start from the premise that the fourth moment of 
the random quantities $B_k(X_i)$ and $1/\pi_i$ is uniformly bounded 
in $k$ and $i$.
\begin{assumptions*}
  \begin{enumerate}[label={(\roman*)}]
    \item
  There exists 
  a constant 
  $
  \,
    C_{\!\scriptscriptstyle B}
  \ge
  1
  \,
  $ such that
  \begin{gather*}
  \E\left[
    B_k(X_i)^{4}
    \,
  \right]
  \ 
  \le
  \ 
    C_{\!\scriptscriptstyle B}
    \quad
    \text{
  for all $\,(k,i)\ \in\  \left\{ 1, \ldots, K \right\}\times \left\{ 1, \ldots, n \right\}$.
    }
  \end{gather*}
  \item
  There exists a constant $\,C_\pi \!\ge 1\,$ such that
  \begin{gather*}
  \E \left[ 
    \left(
      \frac{1}{\pi_{i}}
    \right)^{\!4}
    \,
  \right]
  \ 
  \le
  \ 
  C_\pi
  \quad
  \text{
  for all $\,i\ \in\  \left\{ 1, \ldots, n \right\}$.
  }
  \end{gather*}
  \end{enumerate}
\end{assumptions*}
Note, that these assumptions allow for covariate distributions with unbounded support. The coming example ought to reinforce this observation.
\begin{example*}
  Let us assume a logistic regression model for the propensity score.
  Then there exist coefficients $\,\vartheta \,\in\,  \R^N\,$ and $\,\vartheta_0\,\in\, \R\,$ ($N$ is the number of covariates and $\vartheta_0$ is the intercept of the model) such that
  \begin{align}
    \label{1:5:3a}
    1\,/\,\pi(X)
    &
    \ 
    =
    \ 
    1
    \ 
    +
    \ 
    \exp(
    \,
    \vartheta_0
    \, 
    +
    \,
    \inner{\vartheta}{X}
    \,
    )
    \,,
    \\
    \label{1:5:3b}
  \E \left[ 
    \left(
      \frac{1}{\pi(X)}
    \right)^{\!4}
    \,
  \right]
    &
    \ 
    =
    \ 
  \sum_{j=1}^{4} 
  \,
    \binom{\,4\,}{j}
    \,
    e^{\,j \vartheta_0}
    \,
    M_X(\,j\,\vartheta\,)
    ,
  \end{align}
  where $M_X$ is the momement-generating function of the random vector $X$ (we assume it exists.).
   While the quantity in \eqref{1:5:3a}
   may be unbounded when $X$ has unbounded support, that in \eqref{1:5:3b} remains bounded for moderate decay rates of the
   underlying distribution.

   The multivariate normal distribution with location parameter $\,\mu\,\in\,\R^N\,$ and covariance matrix 
   $\,\Sigma\,\in\, \mathbb{M}_N\,$ has momement-generating function
   \begin{gather*}
     M_{\mathcal{N}(\mu,\,\Sigma)}
     (t)
     \ 
     =
     \ 
     \exp
     \left( 
       \,
       \inner{\,t}{\mu\,}
       +
       \frac{
       \inner{\,t}{\Sigma \,t\,}
       }{2}
       \,
     \right)
     \quad
     \text{for all}\ 
     \,t\ \in\ \R^N
     \,.
   \end{gather*}
   In particular, the expression in~\eqref{1:5:3b} is finite if $X$ follows a multivariate normal distribution.

   Likewise we give a negative example.
   Let $\,N\,=\,1\,$ and $\,X\,\sim\, \mathrm{Exp}(\lambda)$, that is, only one exponentially distributed covariate is taken into account.
   The momement-generating function is then confined to take arguments $\,t\,<\,\lambda\,$, and thus~\eqref{1:5:3b} becomes pointless if $\,4\,\vartheta\,\ge\, \lambda\,$.
\end{example*}

Next, we recall the entity we want to examin.
\begin{gather*}
  A_i
  \ 
  =
  \ 
  \frac{1}{n}
  \,
  \left( 
    \,
    1
    \ 
    -
    \ 
    \frac{T_i}{\pi_i}
    \,
  \right)
  \,
  B(X_i)
  \qquad
  \text{for}
  \ 
  i\in \{1, \ldots, n\}\,.
\end{gather*}
For all
$
  i\in \{1, \ldots, n\}
$
we get the bound
\begin{gather}
  \label{1:5:1}
  \left| 
    \,
    1
    \ 
    -
    \ 
    \frac{T_i}{\pi_i}
    \,
  \right|
  \ 
  \le
  \ 
  \left( 
    \,
  1
  \ 
  \lor
  \ 
  \frac{1-\pi_i}{\pi_i}
  \,
  \right)
  \ 
  \le
  \ 
  1
  \ 
  +
  \ 
  \frac{1-\pi_i}{\pi_i}
  \ 
  =
  \ 
  \frac{1}{\pi_i}
  \,.
\end{gather}
Let
$i^*\in \left\{ 1, \ldots, n \right\}$
be the index where 
$
\norm{A_i}
$
attains its maximum.
\begin{align}
  \label{1:5:2}
  \begin{split}
  \E\left[\max_{i \le n}\norm{A_i}^2\right]
  &
  \ 
  =
  \ 
  \E\left[\norm{A_{i^*}}^2\right]
  \ 
  \le
  \ 
  \E \left[ 
    \left(
      \,
      \frac{
    \norm{B(X_{i^*})}
      }{\pi_{i^*}}
      \,
    \right)
    ^{\!2}
    \,
  \right]
  \ 
  /
  \ 
  n^2
  \\
  &
    \ 
  \le
  \ 
  \E \left[ 
    \left(
      \frac{1}{\pi_{i^*}}
    \right)^{\!4}
    \,
  \right]^{1/2}
  \ 
  \cdot
  \ 
  \E\left[
    \norm{B(X_{i^*})}^4
  \right]^{1/2}
  \,
  /
  \ 
  n^2
  \\
  &
  \ 
  \le
  \ 
  \,
  K
  \,
  /
  \,
  n^2
  \ 
  \cdot
  \ 
  \sqrt{
    \,
    C_\pi
    C_{\!\scriptscriptstyle B}
  }
  \,.
\end{split}
\end{align}
The first inequality comes from the bound \eqref{1:5:1}.
The Cauchy-Schwarz inequality provides the second inequality.
In the last step we use the assumptions made at the start of the section. 
Paying the price of an extra $n$ factor, the 
maximal inequality~\eqref{1:5:2}
yields a bound of the sum, that is,
\begin{gather*}
  \sum_{i=1}^{n}
  \,
  \E\left[\norm{A_i}^2\right]
  \ 
  \le
  \ 
  \frac{K}{n}
  \,
  \sqrt{\,C_\pi 
    C_{\!\scriptscriptstyle B}
  }
\end{gather*}
\begin{assumption}
  .
\end{assumption}
 
\begin{assumption}
  .
\end{assumption}
\begin{remark}
With Assumption we also get a bound on the fourth moment of 
  $
  \norm{B(X_{i})}
  $. Indeed, by the convexity of 
  $x\mapsto x^2$, the monotonicity and linearity of the expectation it holds   
  \begin{align*}
    \begin{split}
  \E[
  \norm{B(X_{i})}^4
  ] 
  &=
  \E
  \left[ 
    \left( 
      \sum_{k=1}^{K}
      B_k^2(X_i)
    \right)^2
  \right]
  =
  K^2
  \E
  \left[ 
    \left( 
      \sum_{k=1}^{K}
      \frac{1}{K}
      B_k^2(X_i)
    \right)^2
  \right]
  \le
  K^2
  \E
  \left[ 
      \sum_{k=1}^{K}
      \frac{1}{K}
      B_k^4(X_i)
  \right]
  \\
  &=
  K
  \sum_{k=1}^{K}
  \E
  \left[ 
      B_k^4(X_i)
  \right]
  \le
  K^2 C_B
  \end{split}
  \end{align*}
\end{remark}
\subsubsection*{Analysis of $v(\mathbf{S})$}
We use the fact that 
$
  \norm{A}_2 
  \le
  \norm{A}_F
  =
  \sqrt{
    \sum_{i,j}^{}
    a_{ij}^2
  }
$
It holds
\begin{gather*}
  \sum_{i=1}^{n}
  \E[A_iA_i^\top]
  =
  \frac{1}{n^2}
  \sum_{i=1}^{n}
  \E
  \left[ 
    \left( 
      \frac{1-\pi_i}{\pi_i}
    \right)^2
    B(X_i)B(X_i)^\top
  \right]
  =
  \frac{1}{n^2}
  \left( 
    \sum_{i=1}^{n}
    \E
    \left[ 
    \left( 
      \frac{1-\pi_i}{\pi_i}
    \right)^2
    B_k(X_i)B_l(X_i)
    \right]
  \right)
  _{1\le k,l \le K}
  .
\end{gather*}
Thus
\begin{align*}
  \begin{split}
  &\norm{
  \sum_{i=1}^{n}
  \E[A_iA_i^\top]
  }_2^2
  \\
  &\le
  \norm{
  \sum_{i=1}^{n}
  \E[A_iA_i^\top]
  }_F^2
  =
  \frac{1}{n^4}
  \sum_{k,l=1}^{K}
  \left( 
    \sum_{i=1}^{n}
    \E
    \left[ 
    \left( 
      \frac{1-\pi_i}{\pi_i}
    \right)^2
    B_k(X_i)B_l(X_i)
    \right]
  \right)^2
  \\
  &\le
  \frac{1}{n^4}
  \sum_{k,l=1}^{K}
  \left( 
    \sum_{i=1}^{n}
    \E
    \left[ 
    \left( 
      \frac{1-\pi_i}{\pi_i}
    \right)^4
    \right]
    ^\frac{1}{2}
    \E[
    B_k(X_i)^4
    ]^\frac{1}{4}
    \E[
    B_l(X_i)^4
    ]^\frac{1}{4}
  \right)^2
  \le
  \left(
  \frac{K}{n}
  \right)^2
  C_\pi C_B
  \end{split}
\end{align*}
On the other hand
\begin{align*}
  \begin{split}
  \norm{
  \sum_{i=1}^{n}
  \E[A^\top_iA_i]
  }_2
  &=
  \sum_{i=1}^{n}
  \E[A^\top_iA_i]
  =
  \frac{1}{n^2}
    \sum_{i=1}^{n}
    \E
    \left[ 
    \left( 
      \frac{1-\pi_i}{\pi_i}
    \right)^2
    \norm{B(X_i)}_2^2
    \right]
    \\
    &\le
  \frac{1}{n^2}
    \sum_{i=1}^{n}
    \E
    \left[ 
    \left( 
      \frac{1-\pi_i}{\pi_i}
    \right)^4
    \right]^\frac{1}{2}
    \norm{B(X_i)}_2^4
    ]^\frac{1}{2}
    \le
    \frac{K}{n}\sqrt{
  C_\pi C_B
    }
  \end{split}
\end{align*}
It follows
\begin{gather*}
  v(\mathbf{S})
  \le
    \frac{K}{n}\sqrt{
  C_\pi C_B
  }
\end{gather*}
Thus we can apply Theorem~\ref{rmineq_rosenthal_pinelis}
to get
\begin{gather*}
  \E[\norm{\mathbf{S}}_2]
  \le
  \sqrt{
    2e 
    \frac{K}{n}\sqrt{
  C_\pi C_B
  }
  \log
  (K+1)
  }
  +
  4e
  \frac{\sqrt{K}}{n}
  \sqrt[4]{
  C_\pi C_B
  }
  \log
  (K+1)
  \le
  14
  C_\pi C_B
  \sqrt{
    \frac{K \log(K+1)}{n}
  }
\end{gather*}
\todo[color=green!40, inline]{Give explanation why simple estimates are insufficient. See notes.}

%
%  \section{Continuous Treatment}
%  We introduce the measure of proximity
\begin{gather}
  d_n(t,s)
  :=
  \frac{
  \mathbf{1} _ 
  { s \in N_n(t) }
}{
  \mathbf{\lambda}[N_n(t)]
}
\end{gather}
where $ N_n(t) $ is a neighborhood of $t$ with 
$  
  \mathbf{\lambda}[N_n(t)]
  \to
  0
  $
  for $ n \to \infty $.
  If we can apply the dominated convergence theorem we get
  \begin{align}
    \E
    \left[ 
      \frac{
        d_n(t,T_i)
      }{
        h_{T|X}(t,X_i)
      }
    \right]
    =
    \E
    \left[ 
      \frac{
        \P
        [ T_i \in N_n(t) | X_i]
      }
      {
        \mathbf{\lambda}[N_n(t)]
      }
      \cdot
      \frac{1}
      {
        h_{T|X}(t,X_i)
      }
    \right]
    \to
    \E
    \left[ 
      \frac
      {
        h_{T|X}(t,X_i)
      }
      {
        h_{T|X}(t,X_i)
      }
    \right]
    =1
    \,.
  \end{align}
  Furthermore, if $ Z_i \sim Y(t)|T_i $
  we have
  \begin{gather}
    \P
    \left[ 
      d_n(t,T_i)
      \left| 
      Y(T_i) - Z_i
      \right|
      \ge
      \varepsilon
    \right]
    =
    \P
    [
      T_i \in N_n(t)
    ]
    \cdot
    \P
    \left[ 
      \left| 
      Y(T_i) - Z_i
      \right|
      \ge
      \varepsilon
      \cdot
      \lambda(N_n(t))
    \right]
    \,.
  \end{gather}
  If $T$ is continuously distributed, 
  $
    \P
    [
      T_i \in N_n(t)
    ]
    \to 0
  $
  which implies 
  $
      d_n(t,T_i)
      \left( 
      Y(T_i) - Z_i
      \right)
      \to
      0
  $
  in probability. We shall employ concentration inequalities to 
  also derive learning rates.
  The optimization problem is for fixed
  $ t\in \mathcal{T} $
  \begin{gather*}
    \underset{w_1, \ldots, w_n \in \R}{\mathrm{minimize}}
    \qquad
    \sum_{i = 1}^{n}
    d_n(t,T_i)
    f(w_i)
  \end{gather*}
subject to the constraints
\begin{gather*}
%    w_i T_i \ge 0,
%   \qquad 
%    i = 1, \ldots, n,
%   \\
%  \sum_{ i = 1 }^{n}
%    w_i T_i
%  =
%  1
%  \\
    \left| 
      \frac{1}{n} 
      \sum_{i = 1}^{n} 
      (
      w_i  
      \cdot
    d_n(t,T_i)
      - 
      1
      )
      \cdot
      B_k(X_i)
    \right|
    \ 
    \le 
    \ 
    \delta_k,
    \qquad
    k = 1, \ldots, K
\end{gather*}
\pagebreak
We shall derive
\begin{gather}
  d_n(t,T_i)
  \cdot
  \left| 
  w_i
  -
  \frac{1}{h_{T|X}(t,X_i)}
  \right|
  \to 0
\end{gather}
in probability for all 
.
From this we control the error
\begin{align*}
  &\left| 
  \frac{1}{n}
  \sum_{i=1}^{n} 
  d_n(t,T_i)
  w_i
  Y_i
  -
  \E[Y(t)]
  \right|
  \\
  &
  \ 
  \le
  \ 
  \left|  
  \frac{1}{n}
    \sum_{i=1}^{n} 
    (
    w_i 
  d_n(t,T_i)
    -
    1
    )
    \inner
    {B(X_i))}
    { \mathbf{Y}(t) }
  \right|
  %%%% 1 %%%%
  \\
  &
  \qquad
  +
  \ 
  \left|  
  \frac{1}{n}
    \sum_{i=1}^{n} 
    (
    w_i 
  d_n(t,T_i)
    -
    1
    )
    \left( 
    \E[Y(t)| X_i]
    -
    \inner
    {B(X_i))}
    { \mathbf{Y}(t) }
    \right)
  \right|
  %%%% 2 %%%%
  \\
  &
  \qquad
  +
  \ 
  \left|  
  \frac{1}{n}
    \sum_{i=1}^{n} 
  d_n(t,T_i)
    \cdot
    (
    w_i 
    -
    1/
h_{T|X}(t,X_i)
    )
    \left( 
      Z_i
    -
    \E[Y(t)| X_i]
    \right)
  \right|
  %%%% 3 %%%%
  \\
  &
  \qquad
  +
  \ 
  \left|  
  \frac{1}{n}
    \sum_{i=1}^{n} 
h_{T}(t)
    /
h_{T|X}(t,X_i)
    \left( 
      Z_i
    -
    \E[Y(t)| X_i]
    \right)
  \right|
  %%%% 3a %%%%
  \\
  &
  \qquad
  +
  \ 
  \left|  
  \frac{1}{n}
    \sum_{i=1}^{n} 
    \left( 
h_{T}(t)
-
  d_n(t,T_i)
    \right)
    /
h_{T|X}(t,X_i)
    \left( 
      Z_i
    -
    \E[Y(t)| X_i]
    \right)
  \right|
  %%%% 3 b %%%%
  \\
  &
  \qquad
  +
  \ 
  \left|  
  \frac{1}{n}
    \sum_{i=1}^{n} 
    \left( 
    \E[Y(t)| X_i]
    -
    \E[Y(t)]
    \right)
  \right|
  %%%% 3 c %%%%
  \\
  &
  \qquad
  +
  \ 
  \left|  
  \frac{1}{n}
    \sum_{i=1}^{n} 
    w_i 
d_n(t,T_i)
    \left( 
      Y_i
    -
    Z_i
    \right)
  \right|
\end{align*}
%%%%%%%%%%%   w_i %%%%%%%%%%
\begin{align*}
&
d_n(t,T_i)
   \cdot
   \left| 
   w_i
   -
   1 / 
h_{T|X}(t,X_i)
   \right|
   \\
   & 
   \ 
   =
   \ 
d_n(t,T_i)
   \cdot
   \left| 
   (
      f ^ { ' }
   )
   ^ {-1}
   \inner
   {B(X_i)}
   { \lambda ^ \dagger }
   -
     (f ^ { ' }) ^ { -1 }
     \left( 
      f ^ { ' }
      \left( 
   1 / 
h_{T|X}(t,X_i)
      \right)
     \right)
   \right|
   %%%% 1 %%%%
   \\
   & 
   \ 
   \le
   \ 
d_n(t,T_i)
   \left| 
   (
      f ^ { ' }
   )
   ^ {-1}
   \inner
   {B(X_i)}
   { \lambda ^ \dagger }
   -
   (
      f ^ { ' }
   )
   ^ {-1}
          \inner
          {B(X_i)}
          {
      f ^ { ' }
      \left( 
   1 / 
h_{T|X}(t,X)
      \right)
          }
   \right|
   %%%% 1 a %%%%
   \\
   &
   \qquad
   +
   \ 
d_n(t,T_i)
   \left| 
   (
      f ^ { ' }
   )
   ^ {-1}
          \inner
          {B(X_i)}
          {
      f ^ { ' }
      \left( 
   1 / 
h_{T|X}(t,X)
      \right)
          }
   -
     (f ^ { ' }) ^ { -1 }
     \left( 
      f ^ { ' }
      \left( 
   1 / 
h_{T|X}(t,X_i)
      \right)
     \right)
   \right|
   %%%% 2 %%%%
   \\
   &
   \ 
   \le
   \ 
   \omega
   \left( 
     (f ^ { ' }) ^ { -1 }
     ,
     \norm{ \lambda ^ \dagger 
     - 
      f ^ { ' }
      \left( 
   1 / 
h_{T|X}(t,X)
      \right)
 }
   \right)
   \ 
   +
   \ 
   \omega
   \left( 
     (f ^ { ' }) ^ { -1 }
     ,
     \varepsilon _ m
   \right)
   %%%% 3 %%%%
   \\
   &
   \ 
   \le
   \ 
   \omega
   \left( 
     (f ^ { ' }) ^ { -1 }
     ,
     \varepsilon _ \dagger
   \right)
   \ 
   +
   \ 
   \omega
   \left( 
     (f ^ { ' }) ^ { -1 }
     ,
     \varepsilon _ m
   \right)
   %%%% 4 %%%%
   \\
   &
   \ 
   \le
   \ 
   \varepsilon
 \end{align*}
 with probability tending to 1.
%%%%%%%%%%%%%%%%%%%%%%%
\begin{align*}
   &
     G
     (
      f ^ { ' }
      \left( 
   1 / 
h_{T|X}(t,X)
      \right)
      +
      \Delta
     )
     \ 
     -
     \ 
     G
     (
      f ^ { ' }
      \left( 
        1 / \pi 
      f ^ { ' }
      \left( 
   1 / 
h_{T|X}(t,X)
      \right)
      \right)
     )
          %%%% 2 %%%%
     \\
     &
     \quad
     \ge
     \ 
     -
     \norm{\Delta}
     \left( 
     \norm{\delta}
     \ 
     +
     \ 
     \norm{B(X_i)}
     \cdot
     \frac{1}{n}
     \sum_{i=1}^{n} 
     \left| 
     \,
      1
     \ 
      -
     \ 
     d_n(t,T_i)
        \cdot
     (f ^ { ' }) ^ { -1 }
          \inner
          {B(X_i)}
          {
      f ^ { ' }
      \left( 
   1 / 
h_{T|X}(t,X)
      \right)
          }
          \,
     \right|
     \right)
 \end{align*}


\chapter{Convex Analysis}
In our application we want to analyse a convex optimization problem by its dual problem.
In particular we want to obtain primal optimal solutions from dual solutions.
To accomplish the task we need technical tools from convex analysis, 
mainly conjugate calculus and some KKT related results.

Our starting point is 
the support function intersection rule~\cite[Theorem 4.23]{Mordukhovich2022}.
We give the details in the case of finite dimensions and refer for the rest of the proof to the book.
The support function intersection rule is applied to give first conjugate sum and then chain rule,
which are vital to calculating convex conjugates. The proofs are omited, since the book is thorough enough. 
%The well known Fenche-Rockafellar Duality theorem is a corollary of conjugate sum and chain rule. It gives general conditions under which dual and primal values coincide.
The material we present is very well known.
As an introduction, we recommend the recent book \cite{Mordukhovich2022} and classical reference \cite{Rockafellar1970}.
We finish the chapter with ideas from \cite{Tseng1991}. 
They provide the high-level ideas to obtain for strictly convex
functions a dual relationship between optimal solutions.
We will deliver the details that are omited in the paper.
  \section{A Convex Analysis Primer}
  \subsection*{My Contribution}
I present the relevant facts from Convex analysis.
I prove some results that I did not find in the literature, but likely are folklore.

Throughout this section let $n\in\mathbb{N}$.
\subsection*{Sets}
A subset $C\subseteq \R^n$ is called \textbf{convex set}, 
\index{convex set}
if for all $x,y\in C$ and all $\theta\in [0,1]$,
we have 
$
  \theta x + (1-\theta)y 
  \in
  C
$.
Many set operations preserve convexity. Among them
forming the 
\textbf{Cartesian product} of two convex sets, 
\textbf{intersection} of a collection of convex sets and 
taking the \textbf{inverse image under linear functions}.

The classical theory evolves around the question 
if convex sets can be separated.
\begin{definition*}
  Let 
  $C_1$ and $C_2$
  be two non-empty convex sets in $\R^n$. 
  A hyperplane $H$ is said to \textbf{separate}
  $C_1$ and $C_2$
  if $C_1$ is contained in one of the closed half-spaces associated with
  $H$ and $C_2$ lies in the opposite closed half-space. It is said to separate 
  $C_1$ and $C_2$
  \textbf{properly} if 
  $C_1$ and $C_2$
  are not both contained in $H$.
\end{definition*}

We need a refined concept of interiors, since some convex sets have empty interior. To this end, 
  we call a set
  \index{affine set}
  $A\subseteq \R^n$ 
  \textbf{affine set}, if
  $
    \alpha x + (1-\alpha)y \in A
    \quad
    \text{for all}
    \ 
    x,y \in A
    \ 
    \text{and all}
    \ 
    \alpha \in \R
  $.
  The \textbf{affine hull} 
  \index{$\mathrm{aff(\cdot)}$, affine hull}
  $\mathrm{aff}(\Omega)$
  of a set 
  $\Omega\subseteq \R^n$
  is the smallest affine set that includes $\Omega$.
  We define the \textbf{relative interior}
  \index{$\mathrm{ri}(\cdot)$, relative interior}
  $\mathrm{ri}\,\Omega$ 
  of a set 
  $\Omega\subseteq \R^n$
  to be the interior relative to the affine hull, that is,
    \begin{gather}
    \mathrm{ri}(\Omega)
    \ 
    :=
    \ 
    \left\{ 
      x \in \Omega 
      \ 
      |
      \ 
      \exists
      \,
      \varepsilon > 0\ 
      \colon
      (
      x+\varepsilon B_{\R^n}
      )
      \cap
      \mathrm{aff}(\Omega)
      \ 
      \subset
      \ 
      \Omega
      \,
    \right\}
    \,.
  \end{gather}

\begin{ftheorem}
  \label{cv:primer:sep}
  \emph{(Convex separation in finite dimension)}
  Let $C_1$ and $C_2$ be two non-empty convex sets in $\R^n$. 
  Then $C_1$ and $C_2$ can be properly separated if and only if 
  $\mathrm{ri}(C_1)\cap\mathrm{ri}(C_2)=\emptyset.$
\end{ftheorem}
\begin{proof}
  \cite[Theorem~11.3]{Rockafellar1970}
\end{proof}
We collect some useful 
properties of relative interiors
before we get on to convex functions.
\begin{proposition}
  \label{cv:primer:prop}
  Let $C$ be a non-empty convex set in $\R^n.$ The following holds:
\begin{enumerate}[label={(\roman*)}]
  \item
    $
      \mathrm{ri}(C)
      \ 
      \neq
      \ 
      \emptyset
      $
if and only if
      $
      C
      \ 
      \neq
      \ 
      \emptyset
    $
  \item
    $
      \mathrm{cl}(\mathrm{ri}\,C)
      \ 
      =
      \ 
      \mathrm{cl}\,C
      $
      and
      $
      \mathrm{ri}(\mathrm{cl}\,C)
      \ 
      =
      \ 
      \mathrm{ri}(C)
    $
  \item
    $
    \mathrm{ri}(C)
      \ 
    =
      \ 
    \left\{ 
      z \in C
      \colon
      \text{for all}\ 
      x \in C \ 
      \text{there exists}\ 
      t > 0 \ 
      \text{such that}\ 
      z + t (z-x)
      \in C
    \right\}
    $
  \item
    Suppose
    $
      \bigcap_{i\in I} C_i
      \ 
      \neq
      \ 
      \emptyset
    $
    for a finite index set $I$.
    Then
    $
      \mathrm{ri}
      \left( 
        \bigcap_{i\in I} C_i
      \right)
      \ 
      =
      \ 
      \bigcap_{i\in I}  
      \mathrm{ri}(C_i)
    $.
    \item
      Let 
      $
        L\,:\,\R^n \to\  \R^m
      $
      be a linear function. Then
      $
        \mathrm{ri}\,L(C)
        \ 
        =
        \ 
        L(\mathrm{ri}\,C)
      $.
      If it also holds
      $
        L^{\!-1}(\mathrm{ri}\,C)
        \ 
        \neq
        \ 
        \emptyset
      $,
      we have
      $
      \mathrm{ri}\,L^{\!-1}(C)
      \ 
        =
      \ 
        L^{\!-1}(\mathrm{ri}\,C)
      $.
      \item
        $
          \mathrm{ri}(C_1\!\times C_2)
          \ 
          = 
          \ 
          \mathrm{ri}\,C_1
          \! 
          \times
          \mathrm{ri}\,C_2
        $
\end{enumerate}

\end{proposition}


\begin{proof}
  For a proof of (i)-(v) we refer to~\cite[Theorem 6.2 - 6.7]{Rockafellar1970}.

To prove (vi) we use (iii).
Let
  $
  (z_1, z_2)
  \in 
  \mathrm{ri}(C_1\!\times C_2).
  $
  Then for all 
  $
  (x_1, x_2)
  \in 
  C_1\!\times C_2
  $
  there exists
  $t>0$
  such that
  \begin{gather}
    \label{cv:primer:prop:1}
      z_i + t (z_i-x_i)
      \in C_i
      \qquad
      \text{for all}\ 
      i\in \left\{ 1,2 \right\}.
  \end{gather}
  Using (iii) again, we get
  $
  \mathrm{ri}(C_1\!\times C_2)
  \, 
  \subseteq
  \,
          \mathrm{ri}\,C_1
          \! 
          \times
          \mathrm{ri}\,C_2
  $.
  Suppose 
  $
  (z_1,z_2)
    \in
    \mathrm{ri}\,C_1
    \!
    \times
    \mathrm{ri}\,C_2
  $.
  By (iii), for all
  $
    (x_1,x_2)\in C_1\times C_2
  $
  there exist
  $
    (t_1,t_2)>0
  $
  such that
  \begin{gather}
    \label{cv:primer:prop:2}
      z_i + t_i (z_i-x_i)
      \in C_i
      \qquad
      \text{for all}\ 
      i\in \left\{ 1,2 \right\}.
  \end{gather}
  If $t_1=t_2$
  we recover
  \eqref{cv:primer:prop:1}
  from
  \eqref{cv:primer:prop:2}.
  By (iii) it holds
  $
  (z_1,z_2)
    \in
    \mathrm{ri}
    (C_1
    \!
    \times
    C_2)
  $.
  If
  $t_1<t_2$
  we
  define $\theta:=\frac{t_1}{t_2}\in (0,1).$
  Consider  
  \eqref{cv:primer:prop:2} with $i=2$,
  together with $z_2 \in C_2$
  and
  the convexity of $C_2$.
  It follows
  \begin{gather}
    \label{cv:primer:prop:3}
    z_2 + t_1 (z_2 - x_2)
    \ 
    =
    \ 
    \theta
    \cdot
    (
    z_2 + t_2 (z_2 - x_2)
    )
    \ 
    +
    \ 
    (1-\theta)
    \cdot
    z_2
    \in C_2
    \,.
  \end{gather}
  Now we consider
  \eqref{cv:primer:prop:3} and
  \eqref{cv:primer:prop:2} with $i=1$.
  This gives \eqref{cv:primer:prop:1} with $t=t_1$.
  As before, it follows
  $
  (z_1,z_2)\in\mathrm{ri}(C_1\!\times C_2)
  $.
  If 
  $t_1>t_2$
  similar arguments lead to the same result.
  We have proven 
  $
  \mathrm{ri}(C_1\!\times C_2)
  \, 
  \supseteq
  \,
          \mathrm{ri}\,C_1
          \! 
          \times
          \mathrm{ri}\,C_2
  $
  and equality.
\end{proof}
\subsection*{Functions}
A function 
$
f
\colon
\R^n
\to
\overline{\R}
$
is called \textbf{convex function},
\index{convex function}
if the area above its graph, that is, its epigraph(cf.\cite[§2.4.1]{Mordukhovich2022}), is convex. We shall often use an equivalent definition.
To this end, 
a function $f$ is convex if and only if 
\begin{gather}
  \label{cv:cf}
  f(\theta x + (1-\theta)y)
  \ 
  \le
  \ 
  \theta f(x)
  +
  (1-\theta)f(y)
  \qquad
  \text{for all}\ 
  x,y\in \R^n
  \ 
  \text{and all}\ 
  \theta\in[0,1]
  \,.
\end{gather}
This definition extends to convex cominbinations
$
  \theta_1,\ldots,\theta_m\in[0,1]
$
with
$
  \sum_{i=1}^{m} 
  \theta_i
  =1
$, that is, 
a function $f$ is convex if and only if 
\begin{gather}
  f
  \left( 
    \sum_{i=1}^{m} 
    \theta_i
    x_i
  \right)
  \ 
  \le
  \ 
    \sum_{i=1}^{m} 
    \theta_i
    f(x_i)
  \qquad
  \text{for all}\ 
  x_1,\ldots,x_m\in \R^n
  \,.
\end{gather}
We call a function \textbf{strictly convex} if the inequality in
\index{strictly convex}
\eqref{cv:cf} is strict.

We define the \textbf{domain} $\mathrm{dom}\,f$
\index{$\mathrm{dom}\,f$, domain of a convex function}
of a convex function $f$ to be the set where $f$ is finite, that is,
\begin{gather}
  \mathrm{dom}\,f
  \ 
  :=
  \ 
  \left\{ 
x\in\R^n
:
f(x)<\infty
  \right\}
  \,.
\end{gather}
The domain of a convex function is convex. 
We say that $f$ is a \textbf{proper function} if  $\mathrm{dom}\,f\neq\emptyset$. 
\index{proper (convex) function}

For any $\overline{x}\in\mathrm{dom}\,f$ we call $x^*\in\R^n$ a 
\textbf{subgradient} of $f$ at $\overline{x}$ if for all 
\index{subgradient}
$x\in\R^n$ it holds
\begin{gather}
  \inner{x^*}{x-\overline{x}}
  \le
  f(x)
  -
  f(\overline{x})
  \,.
\end{gather}
We denote the collection of all subgradients at $\overline{x}$, that is, the \textbf{subdifferential} of $f$ at $\overline{x}$, as
\index{
  $
\partial f(x)
  $,
  subdifferential of $f$ at $x$
}
$
\partial f(\overline{x})
$.
If $f$ is differentiable at $\overline{x}$ it holds
$
\partial f(\overline{x})
=
\left\{ 
  \nabla
  f(\overline{x})
\right\}
$
and thus
\begin{gather}
  \label{cv:primer:mvthe}
  \inner{
  \nabla
  f(\overline{x})
}{x-\overline{x}}
  \le
  f(x)
  -
  f(\overline{x})
  \,.
\end{gather}
\begin{definition}
  \index{$\sigma_\omega$, support function}
  Given a nonempty subset 
  $\Omega \subseteq \R^n$,
  we define
  the \textbf{support function} 
  $
  $
  of $\Omega$
  to be
  \begin{gather*}
  \sigma_\Omega 
  \,
  :
  \,
  \R^n \to\  \overline{\R}
  \,,
  \qquad
  x^*
  \ 
  \mapsto
  \ 
    \sup_{x \in \Omega}
    \ 
    \inner{x^*\!}{x}
    \,.
  \end{gather*}
\end{definition}


\begin{definition}
  \index{$f\square g$, infimal convolution of $f$ and $g$}
  Given functions
  $
    f_i\,:\,
    \R^n \to\  \overline{\R}
  $
  for $ i = 1, \ldots, m $,
  we define the \textbf{infimal convolution} of these functions to be
  \begin{gather*}
    f_1 \square \cdots \square f_m
    \ 
    \colon
    \ 
    \R^n
    \to
    \ 
    \overline{\R}
    \,,
    \quad
    x
    \ 
    \mapsto
    \ 
    \inf
    \left\{ 
    \sum_{i = 1}^{m}
      f_i(x_i)
      \ 
      \colon
      \ 
      x_i \in \R^n 
      \ 
      \mathrm{and}\ 
      \sum_{i = 1}^{m} 
        x_i
      =
      x
    \right\}
    \,.
  \end{gather*}
\end{definition}
 
The next result establishes a connection between the support function of the intersection of two convex sets and the infimal convolution of the support functions of the sets taken by themselfes.
The proof translates the geometric concept of convex separation to the world of convex functions.

\begin{lemma}
  \label{cv:primer:lem}
  Let $C_1$ and $C_2$ be two non-empty convex sets in $\R^n$.
  For any
  $ x^* \in \mathrm{dom}\, \sigma_{C_1\cap C_2} $
  the sets
  \begin{align*}
    \Theta_1
    &
    \ :=\ 
    C_1 \times [\,0,\infty)
    \,,
    \\
    \Theta_2
    (x^*)
    &
    \ :=\ 
    \left\{ 
      (x,\lambda)\in \R^n
      \ 
      \colon
      \ 
      x \in C_2
      \ 
      \text{and}
      \ 
      \lambda
      \,
      \le
      \,
      \inner{x^*\!}{x} 
      \ 
      -
      \ 
      \sigma_{C_1\cap C_2}(x^*)
    \right\}
  \end{align*}
  can by properly separated.
\end{lemma}
\begin{proof}
  We fix 
  $ x^* \in \mathrm{dom}\, \sigma_{C_1\cap C_2} $
  and write
  $ 
  \alpha
  \ 
  :=
  \ 
  \sigma_{C_1\cap C_2}(x^*)
  $.
  In order to apply convex separation in finite dimension 
  (Theorem~\ref{cv:primer:sep})
  to
  the sets
  $ \Theta_1 $ and $ \Theta_2(x^*) $,
  it suffics to show
  their convexity and
  $
    \mathrm{ri}\, 
    \Theta_1
    \cap
    \mathrm{ri}\, 
    \Theta_2(x^*)
    =
    \emptyset
  $.
  \subsubsection*{Convexity of 
  $ \Theta_1 $ and $ \Theta_2(x^*) $
  }
  Clearly, 
  $ \Theta_1 $ is convex by the convexity of 
  $ C_1 $ and $ [0,\infty) $.
 To see that $\Theta_2(x^*)$ is convex consider the linear function
 \begin{gather*}
    L
    \,
    :
    \,
    \R^n\times\,  \R 
    \ 
    \to
    \ 
    \R
    \,,
    \qquad 
    (x,\lambda)
    \ 
    \mapsto
    \ 
    \inner{x^*\!}{x} - \lambda
    \,.
 \end{gather*}
 From the definitions of $L$ and $\Theta_2(x^*)$ we get 
  \begin{gather*}
 \Theta_2
 (x^*)
    \ 
    =
    \ 
    (
    C_2\!\times\R
    )
    \ 
    \cap
    \ 
    L^{\!-1}
    [\,\alpha,\infty)
    \,
    .
  \end{gather*}
  Thus,
  by
  Proposition~\ref{cv:primer:prop}~(v)
  and the convexity of $C_2$ we get the convexity of
  $
    L^{\!-1}
    [\,\alpha,\infty)
  $ and with it that of $\Theta_2(x^*)$.

  \subsubsection*{Relative interiors of
  $ \Theta_1 $ and $ \Theta_2(x^*) $
  are disjoint}
  We start by calculating the relative interiors. It holds
  \begin{alignat*}{3}
    \mathrm{ri}\,
    \Theta_1
    &
    \ 
    =
    \ 
    \mathrm{ri}
    ( C_1\times [0,\infty) )
    &&
    \ 
    =
    \ 
    \mathrm{ri}\,
    C_1
    \!
    \times
    \mathrm{ri}\,
    [0,\infty)
    \ 
    =
    \ 
    \mathrm{ri}\,
    C_1
    \!
    \times
    (0,\infty)
    \,,
    %%%%%%%%%%%%%%%%%
    \\
    \mathrm{ri}\,
    \Theta_2(x^*)
    & 
    \ 
    =
    \ 
    \mathrm{ri}
    (
    L^{\!-1}
    [\,\alpha,\infty)
    )
    &&
    \ 
    =
    \ 
    L^{\!-1}
    (
    \mathrm{ri}\,
    [\,\alpha,\infty)
    )
    \ 
    \ 
    =
    \ 
    L^{\!-1}
    (\alpha,\infty)
    \,.
  \end{alignat*}
  Suppose there exists
  $ (\lambda,x) \in
    \mathrm{ri}\, 
    \Theta_1
    \cap
    \,
    \mathrm{ri}\, 
    \Theta_2(x^*)
  $.
  Then it holds 
  $ x \in C_1\!\times C_2 $
  and 
  $ \lambda >0 $.
  We also note, that
  \begin{gather*}
  \alpha
  \ 
  =
  \ 
  \sigma_{C_1\cap\, C_2}(x^*)
    \ 
  =
    \ 
  \sup_{z \in C_1\cap\, C_2}
  \inner
  {x^*}
  {z}
  \ 
  \ge
  \ 
  \inner
  {x^*}
  {x}
  \,.
  \end{gather*}
  Then it follows
  \begin{gather*}
    \alpha
    \ 
    <
    \ 
  \inner
  {x^*}
  {x}
  - \lambda
    \ 
  \le
    \ 
  \alpha\,,
  \end{gather*}
  a contradiction.
  Thus, the relative interiors of
  $ \Theta_1 $ and $ \Theta_2(x^*) $
  are disjoint.

  Applying Theorem~\ref{cv:primer:sep} finishes the proof.
\end{proof}


\begin{theorem*}
  Let $C_1$ and $C_2$ be two non-empty convex sets in $\R^n$ with
  $\mathrm{ri}\,C_1\cap\mathrm{ri}\,C_2\neq\emptyset.$
  Then the support function of the intersection 
  $
    C_1\! \cap C_2
  $
  is represented as
  \begin{gather}
    (\sigma_{
    C_1 \cap\, C_2
    })
    (x^*)
    =
    (\sigma_{C_1}\square \,\sigma_{C_2})
    (x^*)
    \qquad
    \text{for all}\ 
    x^* \in \R^n.
  \end{gather}
  Furthermore, for any
  $
  x^*\in \mathrm{dom}
    (\sigma_{
    C_1 \cap\, C_2
    })
  $
  there exist dual elements 
  $
    x_1^*
    ,
    x_2^*
    \in \R^n
  $ 
  such that 
  $
    x^*
    =
    x_1^*
    +
    x_2^*.
  $
  and
  \begin{gather}
    (\sigma_{
    C_1 \cap\, C_2
    })
    (x^*)
    =
    \sigma_{C_1}(x_1^*)
    +
    \sigma_{C_2}(x_2^*).
  \end{gather}
\end{theorem*}
\begin{proof}
  Using Lemma~\ref{cv:primer:lem}
  the rest of the proof is as that of
  \emph{\cite[Theorem~4.23(b)]{Mordukhovich2022}}.
\end{proof}

\begin{takeaways}
  The support function intersection rule connects the geometric 
  property of convex separation to an identity of support functions
  This result is central to the analysis of convex conjugates.
\end{takeaways}
One important application of convex functions is in optimization.
There we often analyse a dual problem instead, which relies on the notion of \textbf{convex conjugate} 
\index{$f^*$, convex conjugate of f}
$
    f^*:
    \R^n \to \overline{\R}
  $
  of $f$ defined by
  \begin{gather}
    \label{def:convex_conjugate}
    f^*(x^*)
    :=
    \sup_{ x \in \R^n }
    \inner
    {x^*}{x}
    - f(x)
    \,.
  \end{gather}
  Even for arbitrary functions, the convex conjugate is convex(cf.
  \cite[Proposition~4.2]{Mordukhovich2022}
  ).
  Like in differential calculus, there exist sum and chain rule for computing the convex conjugate.
\begin{theorem}
  \index{conjugate sum rule}
  Let
  $
    f,g:
    \R^n \to (-\infty, \infty]
  $
  be proper convex functions 
  and
  \begin{align*}
  \text{ri}\left( \text{dom}(f) \right)
  \cap
  \text{ri}\left( \text{dom}(g) \right)
  \neq 
  \emptyset
  \,.
  \end{align*}
  Then we have the 
  \textbf{
  conjugate sum rule
  }
  \begin{gather}
    ( f + g )^*(x^*)
    =
    ( f^* \square g^*)(x^*)
  \end{gather}
  for all $x^* \in \R^n$.
  Moreover, the infimum in 
  $
    ( f^* \square g^*)(x^*)
  $
  is attained, i.e., for any
  $
    x^* \in \text{dom}(f+g)^*
  $
  there exists vectors $x_1^*, x_2^*$
  for which
  \begin{gather}
    (f+g)^*(x^*)
    =
    f^*(x_1^*)
    +
    g^*(x_2^*),
    \quad
    x^* = x_1^* + x_2^*.
  \end{gather}
\end{theorem}
\begin{proof}
  \cite[Theorem~4.27(c)]{Mordukhovich2022}
\end{proof}



% conjugate chain rule %
 %%%%%%%%%%%%%%%%%%%%%%
\begin{theorem}
  \index{conjugate chain rule}
  \label{cvxa_conjugate_chain_rule}
  Let 
  $
    A:
      \R^m \to \R^n
  $
  be a linear map (matrix)
  and
  $
    g:
      \R^n \to (-\infty, \infty]
  $
  a proper convex function. If
  $
    \text{Im}(A) \cap \text{ri}(\text{dom}(g))
    \neq
    \emptyset
  $
  it follows
  the 
  \textbf{conjugate chain rule}
  \begin{gather}
    ( g \circ A )^* ( x^* )
    =
    \inf_
          { y^* \in ( A^* )^{ -1 } ( x^* )}
                                          g^*( y^* )
                                          .
  \end{gather}
  Furthermore, 
    for any 
      $
        x^* \in \text{dom}( g \circ A)^*
      $
        there exists
          $
            y^* \in ( A^* )^{ -1 } ( x^* )
          $
            such that
              $
                ( g \circ A)^* ( x^* )
                =
                g^*( y^* )
              $.
\end{theorem}
\begin{proof}
  \cite[Theorem~4.28(c)]{Mordukhovich2022}
\end{proof}
%%%%%%%%%%%%%%%%%
%%%% EXAMPLE %%%%
%%%%%%%%%%%%%%%%%
\begin{example}
  \label{cv:cc:ex}
  Let 
  $
    f:\R\to \overline{\R}
  $
  be a proper convex function, that is, 
  $
    \mathrm{dom}\,f
    \neq
    \emptyset
  $
  and $f$ is convex.
  In steps we apply the conjugate chain and sum rule, together with mathematical induction,
  to prove the conjugate relationship 
  \begin{align*}
    &S_{f,n}:\R^n \to \overline{\R},
    \qquad
    (x_1,\ldots,x_n)
    \mapsto
    \sum_{i=1}^{n} 
    f(x_i)
    ,
    \\
    &S_{f,n}^*:\R^n \to \overline{\R},
    \qquad
    (x^*_1,\ldots,x^*_n)
    \mapsto
    \sum_{i=1}^{n} 
    f^*(x^*_i)
    \,.
  \end{align*}
  This relationship is very natural and the ensuing calculations serve to confirm our intuition.

  First, we work in the projections on the coordinates. 
  For the $i$-th coordinate, where $i=1,\ldots,n$, this is 
  \begin{gather}
    p_i:\R^n\to \R
    ,
    \quad
    (x_1,\ldots,x_n)
    \mapsto
    x_i\,.
  \end{gather}
  All projections 
  $p_i$
  are linear function with matrix representation
  $
    e_i^\top
  $,
  where $e_i$ is $i$-the coordinate vector.
  The adjoint of $p_i$ is therefore
  \begin{gather}
    p^*_i:\R\to \R^n
    ,
    \quad
    x
    \mapsto
    e_i\cdot x
    \,.
  \end{gather}
  For the inverse image of the adjoint of $p_i$ it holds
  \begin{gather}
    (p_i^*)^{-1}
    \left\{ 
    (x_1^*,\ldots,x_n^*)
    \right\}
    \ 
    =
    \ 
    \begin{cases}
      \left\{ x_i^* \right\},
      \quad
      &\text{if}\ 
      x_j^*=0\ \text{for all}\ j\neq i\,,
      \\
      \ \ \emptyset
      \quad
      &\text{else.}
    \end{cases}
  \end{gather}
  Throughout this example we use the asterisk character $^*$ somewhat inconsistently. 
  Note that $f^*$ is the convex conjugate 
  of the function $f$ and $p_i^*$ is the adjoint linear function of the projection on the $i$-th coordinate. Likewise, we denote dual variables, that is, the arguments of convex conjugates, as $x^*$.

  Next, we employ the conjugate chain rule to establish the conjugate relationship 
  \begin{align*}
    f_i&:\R^n\to \overline{\R}
    ,
    \quad
    (x_1,\ldots,x_n)
    \mapsto x_i \mapsto f(x_i)
    \,,
    \\
    f^*_i&:\R^n\to \overline{\R}
    ,
    \quad
    (x^*_1,\ldots,x^*_n)\mapsto 
    \begin{cases}
      f^*(x_i^*),
      \quad
      &\text{if}\ 
      x_j^*=0\ \text{for all}\ j\neq i\,,
      \\
      \infty
      \quad
      &\text{else.}
    \end{cases}
  \end{align*}
  Note, that 
  $
    f_i
    =
    (f\circ p_i)
  $
  and
  $
    f^*_i
    =
    (f\circ p_i)^*
  $.
  Since 
  $
    \mathrm{Im}\,p_i=\R
  $
  and 
  $
    \mathrm{dom}\, f
    \neq
    \emptyset
  $,
  it holds
  $
    \mathrm{Im}\, p_i
    \cap
    \mathrm{ri}(
    \mathrm{dom}\, f
    )
    \neq
    \emptyset
  $.
  Then $f$ and $p_i$ conform with the demands of the conjugate chain rule.
  It follows
  \begin{align*}
    &f_i^*
    (x^*_1,\ldots,x^*_n) 
    \ 
    =
    \ 
    (f\circ p_i)^*
    (x^*_1,\ldots,x^*_n) 
    \ =
    \ 
    \inf
    \left\{ 
    f^*(y)
    \ 
    |
    \ 
    y\in 
    (p_i^*)^{-1}
    \left\{ 
    (x_1^*,\ldots,x_n^*)
    \right\}
    \right\}
    \\
    &\quad=
    \ 
    \begin{cases}
      f^*(x_i^*),
      \quad
      &\text{if}\ 
      x_j^*=0\ \text{for all}\ j\neq i\,,
      \\
      \infty
      \quad
      &\text{else,}
    \end{cases}
  \end{align*}
  where we keep to the convention $\inf\emptyset=\infty$.
  In the same way it follows
  \begin{gather}
    \left( 
      S_{f,n}
      \circ
      p_{\left\{ 1,\ldots,n \right\}}
    \right)^*
    (x^*_1,\ldots,x^*_{n+1})
    =
    \begin{cases}
      S_{f,n}^*
    (x^*_1,\ldots,x^*_{n})
      \quad
      &\text{if}\ 
      x_{n+1}^*=0\,,
      \\
      \infty
      \quad
      &\text{else,}
    \end{cases}
  \end{gather}

  Next, note that for $n=1$ we arrive at the result. Thus, for some $n\in \mathbb{N}$ it holds
  $
  \left( 
    S_{f,n}
  \right)
  ^*
  =
    S_{f,n}^*
  $.
  In order to apply the conjugate sum rule to 
  $
    S_{f,n}
  $
  and
  $
    f_{n+1}
  $
  we note that
  \begin{align*}
    \mathrm{dom}\, f_i
    &
    \ =\ 
    \left\{ 
      (x_1,\ldots,x_{n+1})
      \in \R^{n+1}
      :
      x_i\in \mathrm{dom}\,f
    \right\}
    \ 
    \neq 
    \ 
    \emptyset
    \qquad
    \text{for all}
    \ 
    i=1,\ldots,n+1
    \,,
    \\
    \bigcap_{i=1}^{n+1}
    \mathrm{dom}\, f_i
    &
    \ =\ 
    \left\{ 
      (x_1,\ldots,x_{n+1})
      \in \R^{n+1}
      :
      x_i\in \mathrm{dom}\,f
      \ 
    \text{for all}
    \ 
    i=1,\ldots,n+1
    \right\}
    \ 
    \neq 
    \ 
    \emptyset
    \,,
  \end{align*}
  and
\begin{align*}
  &
  \mathrm{ri}\left( 
    \mathrm{dom}
    \left( 
    S_{f,n}
    \circ
    p_{\left\{ 1,\ldots,n \right\}}
    \right)
  \right)
  \ 
  \cap
  \ 
  \mathrm{ri}\left( 
    \mathrm{dom}\,f_{n+1}
  \right)
  \\
  &
  \hspace{25mm}
  =
  \ 
  \mathrm{ri}\left( 
    \mathrm{dom}
    \left( 
    S_{f,n}
    \circ
    p_{\left\{ 1,\ldots,n \right\}}
    \right)
  \ 
  \cap
  \ 
    \mathrm{dom}\,f_{n+1}
  \right)
  \ 
  =
  \ 
  \mathrm{ri}
  \left( 
    \bigcap_{i=1}^{n+1}
    \mathrm{dom}\, f_i
  \right)
  \ 
  \neq
  \ 
  \emptyset
  \,.
\end{align*}
By the conjugate sum rule it follows
\begin{align*}
  (
  S_{f,n+1}
  )^*
  =
  (
    S_{f,n}
    \circ
    p_{\left\{ 1,\ldots,n \right\}}
  +
  f_{n+1}
  )^*
  =
  (
    S_{f,n}
    \circ
    p_{\left\{ 1,\ldots,n \right\}}
  )
  ^*
  \square
  f_{n+1}^*
  \\
  =
    S_{f,n}^*
    \circ
    p_{\left\{ 1,\ldots,n \right\}}
    +
  f_{n+1}^*
  =
  S^*_{f,n+1}
  \,.
\end{align*}
\end{example}






  %\section{Conjugate Calculus}
  %The goal of this section is to establish the tools to calculate convex conjugates. 
We prove the conjugate sum and chain rule.
After some examples, we will derive the Fenchel-Rockafellar Theorem.
\begin{definition}
  \label{ def_convex_conjugate }
  \emph{(Convex conjugate)}
  Given a function
  $
    f:
    \R^n \to \overline{\R}
  $
  ,
  the 
  \textbf{convex conjugate}
  $
    f^*:
    \R^n \to \overline{\R}
  $
  of $f$ is defined as
  \begin{gather}
    f^*(x^*)
    :=
    \sup_{ x \in \R^n }
    (x^*)^T x - f(x)
  \end{gather}
\end{definition}

Note that $f$ in Definition~\ref{ def_convex_conjugate }
does not have to be convex. On the other hand, the convex conjugate is always convex:

\begin{proposition}
  Let  
  $
    f:
    \R^n \to ( - \infty, \infty ]
  $
  be a proper function. 
  Then its convex conjugate
  $
    f^*:
    \R^n \to ( - \infty, \infty ]
  $
  is convex.
\end{proposition}


\begin{lemma}
  For any proper function
  $
    f:\R^n\to\overline{\R}
  $
  we have
  \begin{gather}
    f^*(x^*) 
    =
    \sigma_{\mathrm{epi}(f)}
    (x^*,-1)
    \qquad
    \text{for}
    \ 
    x^* \in \R^n.
  \end{gather}
\end{lemma}
\begin{proof}
  Let $x^*\in\R^n$
  and
  $
    (x,\lambda)\in \mathrm{epi}(f).
  $
  Then
  $
    x \in \mathrm{dom}(f)
  $
  and
  $
    f(x)\le \lambda.
  $
  Thus
  \begin{gather}
    \inner{x^*}{x} - f(x)
    \ge
    \inner{x^*}{x} - \lambda
    \qquad
    \text{for all}\ 
    (x,\lambda)\in \mathrm{epi}(f).
  \end{gather}
  On the other hand 
  $
    (x,f(x))\in \mathrm{epi}(f)
  $
  for all
  $
    x \in \mathrm{dom}(f).
  $
  It follows
  \begin{gather}
    \inner{x^*}{x} - f(x)
    \le
    \sup_{(x,\lambda)\in\mathrm{epi}(f)}
    \inner{x^*}{x} - \lambda
    \qquad
    \text{for all}\ 
    x \in \mathrm{dom}(f).
  \end{gather}
  Taking the supremum in the last two displays yields
  \begin{align}
    f^*(x^*)
    =
    \sup_{x\in\mathrm{dom}(f)}
    \inner{x^*}{x} - f(x)
    &=
    \sup_{(x,\lambda)\in\mathrm{epi}(f)}
    \inner{x^*}{x} - \lambda
    \\
    &=
    \sup_{(x,\lambda)\in\mathrm{epi}(f)}
    \inner{(x^*,-1)}{(x,\lambda)} 
    =
    \sigma_{\mathrm{epi}(f)}
    (x^*,-1).
  \end{align}
\end{proof}
% conjugate chain rule %
 %%%%%%%%%%%%%%%%%%%%%%
\begin{theorem}
  \emph{(Conjugate Chain Rule)}
  \label{cvxa_conjugate_chain_rule}
  Let 
  $
    A:
      \R^m \to \R^n
  $
  be a linear map (matrix)
  and
  $
    g:
      \R^n \to (-\infty, \infty]
  $
  a proper convex function. If
  $
    \text{Im}(A) \cap \text{ri}(\text{dom}(g))
    \neq
    \emptyset
  $
  it follows
  \begin{gather}
    ( g \circ A )^* ( x^* )
    =
    \inf_
          { y^* \in ( A^* )^{ -1 } ( x^* )}
                                          g^*( y^* )
                                          .
  \end{gather}
  Furthermore, 
    for any 
      $
        x^* \in \text{dom}( g \circ A)^*
      $
        there exists
          $
            y^* \in ( A^* )^{ -1 } ( x^* )
          $
            such that
              $
                ( g \circ A)^* ( x^* )
                =
                g^*( y^* )
              $.
\end{theorem}

% conjugate sum rule %
 %%%%%%%%%%%%%%%%%%%%



\begin{theorem}
  Let
  $
    f,g:
    \R^n \to (-\infty, \infty]
  $
  be proper convex functions 
  and
  $
  \text{ri}\left( \text{dom}(f) \right)
  \cap
  \text{ri}\left( \text{dom}(g) \right)
  \neq 
  \emptyset
  .
  $
  Then we have the conjugate sum rule
  \begin{gather}
    ( f + g )^*(x^*)
    =
    ( f^* \square g^*)(x^*)
  \end{gather}
  for all $x^* \in \R^n$.
  Moreover, the infimum in 
  $
    ( f^* \square g^*)(x^*)
  $
  is attained, i.e., for any
  $
    x^* \in \text{dom}(f+g)^*
  $
  there exists vectors $x_1^*, x_2^*$
  for which
  \begin{gather}
    (f+g)^*(x^*)
    =
    f^*(x_1^*)
    +
    g^*(x_2^*),
    \quad
    x^* = x_1^* + x_2^*.
  \end{gather}
\end{theorem}
\begin{proof}
  Let $x^*\in\R^n$ and fix $x_1^*,x_2^*\in\R^n$ such that
  $x^*=x^*_1+x^*_2$.
  We get
  \begin{align*}
    f^*(x^*_1)+g^*(x^*_2)
    &=
    \sup_{x\in\R^n}
    \inner{x^*_1}{x}-f(x)
    +
    \sup_{x\in\R^n}
    \inner{x^*_2}{x}-g(x)
    \\
    &\ge
    \sup_{x\in\R^n}
    \inner{x^*_1}{x}-f(x)
    +
    \inner{x^*_2}{x}-g(x)
    =
    \sup_{x\in\R^n}
    \inner{x^*_1+x^*_2}{x}-(f(x)+g(x))
    \\
    &=
    \sup_{x\in\R^n}
    \inner{x^*}{x}-(f+g)(x)
    =(f+g)^*(x^*)
  \end{align*}
  Taking the infimum over $x_1^*,x_2^*\in\R^n$ in the above display gives 
  $
  (f^*\square g^*)(x^*)
  \ge
  (f+g)^*(x^*).
  $
  Let us prove now $\le$ under the condition
  $
  \text{ri}\left( \text{dom}(f) \right)
  \cap
  \text{ri}\left( \text{dom}(g) \right)
  \neq 
  \emptyset
  .
  $
  The only case we need to consider is
  $
    (f+g)^*(x^*)<\infty.
  $
  Define two convex sets by
  \begin{align}
    \Omega_1
    &:=
    \left\{ 
      (x,\alpha,\beta)\in\R^{n+2}
      \colon
      \alpha\ge f(x)
    \right\}
    =
    \mathrm{epi}(f)\times \R,
    \\
    \Omega_2
    &:=
    \left\{ 
      (x,\alpha,\beta)\in\R^{n+2}
      \colon
      \beta\ge g(x)
    \right\}.
  \end{align}
  Similar to Lemma we get the representation
  \begin{gather}
    (f+g)^*(x^*)
    =
    \sigma_{\Omega_1\cap\Omega_2}
    (x^*,-1,-1).
  \end{gather}
  Indeed, the only thing we need to verify is
  $
    \mathrm{dom}(f)\cap\mathrm{dom}(g)
    =
    \mathrm{dom}(f+g).
  $
  The inclusion $\subseteq$ is clear.
  Assume towards a contradiction that
  $
    (f+g)(x)<\infty
  $
  and
  $
    f(x)=\infty.
  $
  Since $g(x)>-\infty$ it holds
  \begin{gather}
    \infty
    =
    \infty+g(x)
    =f(x)+g(x)
    =(f+g)(x)
    <
    \infty.
  \end{gather}
  This is a contradiction. The same holds for $f$ and $g$ reversed. It follows the inclusion $\supseteq$ and equality.
  By the support function intersection rule there exist triples
  \begin{gather}
    (x^*_1,-\alpha_1,-\beta_1),
    (x^*_2,-\alpha_2,-\beta_2)
    \in \R^{n+2}
    \quad
    \text{such that}
    \quad
    (x^*,-1,-1)
    =
    (x^*_1+x^*_2,-(\alpha_1+\alpha_2),-(\beta_1+\beta_2))
  \end{gather}
  and
  \begin{gather}
    (f+g)^*(x^*)
    =
    \sigma_{\Omega_1\cap\Omega_2}
    (x^*,-1,-1)
    =
    \sigma_{\Omega_1}
    (x^*_1,-\alpha_1,-\beta_1)
    +
    \sigma_{\Omega_2}
    (x^*_2,-\alpha_2,-\beta_2).
  \end{gather}
  Next we show
  $\beta_1=\alpha_2=0.$
  Suppose towards a contradiction that 
  $\beta_1\neq 0.$ 
  We fix 
  $(\overline{x},\overline{\alpha})\in\mathrm{epi}(f).$
  Then
  \begin{gather}
    \sigma_{\Omega_1}
    (x^*_1,-\alpha_1,-\beta_1)
    =
    \sup_{(x,\alpha,\beta)\in \mathrm{epi}(f)\times \R}
    \inner{x^*}{x}-\alpha \alpha_1 -\beta \beta_1
    \ge
    \sup_{\beta\in \R}
    \inner{x^*}{\overline{x}}-\overline{\alpha} \alpha_1 -\beta \beta_1
    =\infty.
  \end{gather}
  This contradicts
  $
    (f+g)^*(x^*)<\infty.
  $
  In a similar fashion we can derive a contradiction for $\alpha_2\neq0.$
  Employing Lemma and taking into account the structures of the sets 
  $\Omega_1$ and $\Omega_2$ this implies
  \begin{align}
    (f+g)^*(x^*)
    &=
    \sigma_{\Omega_1\cap\Omega_2}
    (x^*,-1,-1)
    =
    \sigma_{\Omega_1}
    (x^*_1,-1,0)
    +
    \sigma_{\Omega_2}
    (x^*_2,0,-1)
    \\
    &=
    \sigma_{\mathrm{epi}(f)}(x^*_1,-1)
    +
    \sigma_{\mathrm{epi}(g)}(x^*_2,-1)
    =
    f^*(x^*_1)
    +
    g^*(x^*_2)
    \ge
    (f^*\square g^*)(x^*).
  \end{align}
  This finishes the proof.
\end{proof}




Given 
proper convex functions $f, g : \R^n \to \overline{\R}$ 
and
a matrix $A \in \R^{n \times n}$,
we define 
the primal minimization problem as follows:
\begin{problem}
  \emph{(Primal)}
  Given proper convex functions 
  $
  f:\R^n \to \overline{\R} 
  $,
  $
  g:\R^m \to \overline{\R} 
  $
  and a matrix 
  $
    A\in \R^{m \times n}
  $
  we define the \textbf{primal optimization problem} to be
  \label{cvxa_primal_problem}
  \begin{gather*}
    \underset{x \in \R^n}{\mathrm{minimize}}
    \qquad
    f(x) + g(Ax)
  \end{gather*}
\end{problem}
\begin{remark}
  Problem \autoref{cvxa_primal_problem}
  appears in the unconstrained form. We can impose constraints by controling for the domains of $f$ and $g$.
  To incorporate linear constraints $Ax \le 0$
  or more general constraints $x\in\Omega$, where $\Omega$ is a convex set,
  we can choose
  \begin{gather}
    g(x)=\delta_\Omega(x):=
  \end{gather}
  where $x\notin\Omega$
  leads to 
  $f(x) + g(x)=\infty$
  and the optimization problem (if feasible) will exclude $x$ from the solutions.
\end{remark}
\begin{problem}
  \emph{(Dual)}
  \label{cvxa_dual_problem}
  Consider the same setting as in Problem~\autoref{cvxa_primal_problem}.
  Using the convex conjugates of $f,g$ and the transpose of $A$
  we define the \textbf{dual problem} of Problem~\autoref{cvxa_primal_problem} to be
  \begin{gather*}
    \underset{y^* \in \R^m}{\mathrm{maximize}}
    \qquad
   - f^*(A^\top y^*) - g^*(y^*).
  \end{gather*}
\end{problem}
\begin{proposition}
  Consider the optimization problem~\autoref{cvxa_primal_problem} and its dual~\autoref{cvxa_dual_problem}, where the functions $f$ and $g$ are not assumed to be convex. Define the \textbf{optimal values} of these problems by
  \begin{gather*}
    \widehat{p}:= \inf_{x\in\R^n}f(x)+g(Ax)
    \quad
    \text{and}
    \quad
    \widehat{d}:= \sup_{y\in\R^m} - f^*(A^\top y) - g^*(y).
  \end{gather*}
  Then we have the relationship
  $\widehat{d}\le \widehat{p}$.
\end{proposition}
\begin{proof}
  It holds
  \begin{align*}
    - f^*(A^\top y^*) - g^*(y^*) 
    &=
      -\sup_{x\in \R^n}\inner{A^\top y^*}{x}-f(x)
      -\sup_{y\in \R^m}\inner{-y^*}{y}-g(y)
      \\
    &=
    \inf_{x\in \R^n}f(x)-\inner{y^*}{Ax}
      +\inf_{y\in \R^m}g(y)+\inner{y^*}{y}
      \\
    &\le
      \inf_{x\in \R^n}f(x)-\inner{y^*}{Ax}
      +\inf_{x\in \R^n}g(Ax)+\inner{y^*}{Ax}
      \\
    &\le
      \inf_{x\in \R^n}f(x)-\inner{y^*}{Ax} + g(Ax) + \inner{y^*}{Ax}
      \\
    &=  
      \inf_{x\in \R^n}f(x)+g(Ax)
    =
      \widehat{p}
  \end{align*}
  The first equality is due to the definition of convex conjugates, the second equality due to $\inner{A^\top y}{x}=\inner{y}{Ax}$ and $\inf \left\{ -B \right\}=-\sup \left\{ B \right\}$ for all $B \subseteq \overline{\R}$ and the first inequality due to $\mathrm{Im}(A)\subseteq \R^m$.
  Taking the supremum with respect to all 
  $y^*\in\R^m$
  yields the result.
\end{proof}
\begin{theorem}
  \label{cvxa_fenchel_theorem}
  Let 
  $f, g : \R^n \to \overline{\R}$ 
  be proper convex functions
  and
  $0 \in \text{ri}(\text{dom}(g) - A (\text{dom}(f)) )$
  .
  Then the optimal values of \eqref{cvxa_primal_problem} and \eqref{cvxa_dual_problem} are equal, 
  i.e.
  \begin{gather}
    \inf_{x \in \R^n} 
    \left\{ f(x) + g(Ax) \right\}
    =
    \sup_{y \in \R^n} \left\{   -f^* \left( A^T y \right) - g^*(-y) \right\}
    .
  \end{gather}
\end{theorem}
%%%%%%%%%%%%%%%%%%%%%%%%%%%%%%%%%%%%%%%%%%%%%%%
\begin{lemma}
  \label{syu_1_result}
  Let 
  $f : \R^n \to (-\infty, \infty]$ 
  be convex.
  Then 
  for all $y \in \R^n$ and $C>0$ 
    \begin{gather}
      \label{7060_0}
      \inf_{\norm{\Delta}=C} f(y+\Delta) - f(y) \ge 0 \quad
      \Longrightarrow
      \quad
    \exists y^* \in \R^n
    \colon
    y^* \,\text{is global minimum of $f$ and}\,
      \norm{y^* - y} \le C.
    \end{gather}
\end{lemma}
\begin{proof}
  Since 
  $\mathcal{C}:=\left\{ \norm{\Delta}\le C \right\}$
  is convex
  $f$ has a local minimum in 
  $
    y + \mathcal{C}
    :=
    \left\{ 
      y + \Delta \,
      \mid \,
      \norm{\Delta}\le C
    \right\}
    .
  $
  Suppose towards a contradiction that
  $
    y^* \in 
            y + \mathcal{C}
  $
  is a local minimum, but not a global minimum 
  and
  the left-hand side of 
  \eqref{7060_0} is true.
  Then it holds
  \begin{gather}
    \label{7060_3}
    f(x) < f(y^*)
    \quad
    \text{for some}\ 
    x 
    \in 
    \R^n 
    \setminus 
      y + \mathcal{C}
    .
  \end{gather}
  Furthermore since $y + \mathcal{C}$ is compact and contains $y^*$,
  the line segment $\mathcal{L}[y^*,x]$ contains a point on the boundary of 
  $y + \mathcal{C}$, i.e.
  \begin{gather}
    \label{7060_4}
    \theta x + (1 - \theta) y^* = y + \Delta_x
    \quad
    \text{for some}\ 
    \theta \in (0,1)\ 
    \text{and}\ 
    \Delta_x \ 
    \text{with}\ 
    \norm{\Delta_x}=C
    .
  \end{gather}
    It follows
    \begin{align}
      \label{7060_5}
      \begin{split}
      f(y^*)
      \le
      f(y)
      \le
      f(y + \Delta_x)
      &=
      f(
        \theta x + (1 - \theta) y^*
      )
      \\
      &\le
      \theta f(x)
      + 
      (1 - \theta)
      f(y^*)
      <
      f(y^*)
      ,
      \end{split}
    \end{align}
    which is a contradiction.
    Thus every local minimum of $f$ in $y + \mathcal{C}$ is also a global minimum.
    The first inequality is due to
    $y^*$ being a local minimum of $f$ in
    $
      y + \mathcal{C},
    $
    the second inequality is due to the left-hand side of 
    \eqref{7060_0} being true,
    the equality is due to \eqref{7060_4},
    the third inequality is due to the convexity of $f$
    and the strict inequality is due to \eqref{7060_3}.
    %It follows the right-hand side of \eqref{7060_0}.
\end{proof}


\begin{takeaways}
  Almost there 
  \lipsum[1]
\end{takeaways}

%%%%%%%%%%%%%%%%%
% INTRODUCTION %
%%%%%%%%%%%%%%%%
In the following we do not expect the reader to be familiar with convex analysis. However, some very well known results will be stated without proof. The interested reader can study \cite{Mordukhovich2022} for the bedrock analysis.


We begin by defining convex sets
%

\begin{definition}
  A subset $\Omega\subseteq \R^n$ is called CONVEX if we have $\lambda x+(1-\lambda)y\in \Omega$ for all $x,y\in \Omega$ and $\lambda\in (0,1)$. 
\end{definition}

Clearly, the line segment 
$[a,b]:=\left\{ \lambda a+(1-\lambda)b\,\mid \, \lambda\in [0,1] \right\}$ is contained in $\Omega$ for all $a,b\in \Omega$ if and only if $\Omega$ is a convex set.
%

Next we define convex functions. 
%

The concept of convex functions is closely related to convex sets.
%  
 
The line segment between two points on the graph of a convex function lies on or above and does not intersect the graph.
%

In other words: The area above the graph of a convex function $f$ is a convex set, i.e. the \textit{epigraph}
$\text{epi}(f):=\left\{ (x,\alpha)\in \R^n\times\R\,\mid\, f(x)\le \alpha\right\}$ is a convex set in $\R^{n+1}$.
%

Often an equivalent characterisation of convex functions is more useful.
%

\begin{theorem}
  The convexity of a function $f:\R^n\to \overline{\R}$ on $\R^n$ is equivalent to the following statement:

  For all $x,y\in \R^n$ and $\lambda\in(0,1)$ we have 
    \begin{align}
      f(\lambda x + (1-\lambda)y)\le \lambda f(x)+(1-\lambda)f(y).
    \end{align}
\end{theorem}


\begin{definition}
  \emph{proper convex function}
\end{definition}

\begin{definition}
  \label{ def_convex_conjugate }
  \emph{(Convex Conjugate)}
  Given a function
  $
    f:
    \R^n \to ( - \infty, \infty ]
  $
  ,
  the convex conjugate
  $
    f^*:
    \R^n \to ( - \infty, \infty ]
  $
  of $f$ is defined as
  \begin{gather}
    f^*(x^*)
    :=
    \sup_{ x \in \R^n }
    (x^*)^T x - f(x)
  \end{gather}
\end{definition}

Note that $f$ in Definition~\ref{ def_convex_conjugate }
does not have to be convex. On the other hand, the convex conjugate is always convex:

\begin{proposition}
  Let  
  $
    f:
    \R^n \to ( - \infty, \infty ]
  $
  be a proper function. 
  Then its convex conjugate
  $
    f^*:
    \R^n \to ( - \infty, \infty ]
  $
  is convex.
\end{proposition}

%%%%%%%%%%%%%%%%%%%%%%
% CONJUGATE CALCULUS %
%%%%%%%%%%%%%%%%%%%%%%
\section{Conjugate Calculus}
When studying different primal problems such as \eqref{primal_weighting_binary} we often turn to the dual instead.
Therefore we need some reliable tools.
Beign able to compute specific convex conjugates is one tool required.



%%%%%%%%%%%%%%%%%%%
% FENCHEL DUALITY %
%%%%%%%%%%%%%%%%%%%
\section{Fenchel Duality}


Given 
proper convex functions $f, g : \R^n \to \overline{\R}$ 
and
a matrix $A \in \R^{n \times n}$,
we define 
the primal minimization problem as follows:

\begin{gather}
  \label{cvxa_primal_problem}
  \text{minimize}
  \quad
  f(x) + g(Ax) 
  \quad
  \text{subject to}
  \quad
  x \in \R^n.
\end{gather}

The Fenchel dual problem is then

\begin{gather}
  \label{cvxa_dual_problem}
  \text{maximize}
  \quad
  -f^* \left( A^T y \right) - g^*(-y) 
  \quad
  \text{subject to}
  \quad
  y \in \R^n.
\end{gather}

\begin{theorem}
  \label{cvxa_fenchel_theorem}
  Let 
  $f, g : \R^n \to \overline{\R}$ 
  be proper convex functions
  and
  $0 \in \text{ri}(\text{dom}(g) - A (\text{dom}(f)) )$
  .
  Then the optimal values of \eqref{cvxa_primal_problem} and \eqref{cvxa_dual_problem} are equal, 
  i.e.
  \begin{gather}
    \inf_{x \in \R^n} 
    \left\{ f(x) + g(Ax) \right\}
    =
    \sup_{y \in \R^n} \left\{   -f^* \left( A^T y \right) - g^*(-y) \right\}
    .
  \end{gather}
\end{theorem}

  \section{Duality of Optimal Solutions}
  \subsection*{My Contribution}
I adapt ideas from \cite{Tseng1991} to take also equality constraints. For this, I had to understand the connection to my version of the primal optimization problem. 
I filled in many details that were omitted in the paper:
I derived the Karush-Kuhn-Tucker conditions for the problem
from the general result \cite[Theorem~28.3]{Rockafellar1970}.
I prove in detail, that they hold for the adapted problem.


We consider a general convex optimization problem 
with matrix equality and inequality constraints.
For this problem there exists a related problem,
which we call its dual.
With ideas from \cite{Tseng1991} we establish 
a functional relationship
between the optimal solution of the original problem 
and
optimal solutions of the dual.
The main assumption is that in the original problem we have a strictly convex objective function 
with continuously differentiable 
convex conjugate(cf. Definition~\ref{cv:cc:d:cc}). 
\begin{assumption}
  \label{cv:ts:asu}
  The objective function $f\colon \R^n\to \overline{\R}$ is strictly convex and its
convex conjugate $f^*$ is continuously differentiable.
\end{assumption}
\begin{ftheorem}
  \label{cv:ts:th}
  Consider the optimization problem
\begin{align}
  \label{cv:ts:primal}
  %%%% objective %%%%
    &\underset{w \in \R^n}
    {\mathrm{minimize}}
    &&\qquad\qquad
    f(w)
    &&&
    \\
    %%%% Ax >= b %%%%
    \nonumber
    &\mathrm{subject}\ \mathrm{to} 
    &&\qquad\qquad
    \mathbf{U}w
    \ 
    \ge
    \ 
    d
    \,.
    \\
    \nonumber
    &
    &&\qquad\qquad
    \mathbf{A}w
    \ 
    =
    \ 
    a
    \,,
\end{align}
and its dual problem
  \begin{alignat}{2}
    \label{cv:ts:dual}
  %%%% objective %%%%
    &\underset{
    \lambda_d \in \R^r
,
    \lambda_a \in \R^s
  }
    {\mathrm{maximize}}
    &&\qquad\qquad
    \inner
    {\lambda_d}
    {d}
    \ 
    +
    \ 
    \inner
    {\lambda_a}
    {a}
    \ 
    -
    \ 
    f^*
    \!
    \left( 
      \mathbf{U}^\top \! \lambda_d
      +
      \mathbf{A}^\top \! \lambda_a
    \right)
    \\
    %%%% Ax >= b %%%%
    \nonumber
    &\mathrm{subject}\ \mathrm{to} 
    &&\qquad\qquad
    \lambda_d
    \ 
    \ge
    \ 
    0
    \,.
\end{alignat}
  Let 
$
(\lambda_d^\dagger,\lambda_a^\dagger)
$
be an optimal solution to \eqref{cv:ts:dual}.
If the objective function $f$ of 
\eqref{cv:ts:primal} is strictly convex and its
convex conjugate $f^*$ is continuously differentiable,
then the unique optimal solution to 
\eqref{cv:ts:primal}
is given by
\begin{gather}
  w^\dagger
  =
  \nabla
    f^*
    \!
    \left( 
      \mathbf{U}^\top  \lambda_d^\dagger
      +
      \mathbf{A}^\top  \lambda_a^\dagger
    \right)
    \,.
\end{gather}
\end{ftheorem}

\subsubsection*{Plan of Proof}
We show that 
$w^\dagger$ and 
$
(\lambda_d^\dagger,\lambda_a^\dagger)
$
meet the 
Karush-Kuhn-Tucker conditions for \ref{cv:ts:primal},
that is,
\textbf{complementary slackness}
\begin{gather}
  \label{cv:ts:comps}
\inner
{\lambda_d^\dagger\,}{d-\mathbf{U} w^\dagger}
\ 
=
\ 
0
\,,
\end{gather}
\textbf{primal} and \textbf{dual feasibility}
\begin{align}
  \label{cv:ts:pfea}
    \mathbf{U}w^\dagger
    &
    \ 
    \ge
    \ 
    d
    \,,
    \\
    \nonumber
    \mathbf{A}w^\dagger
    &
    \ 
    =
    \ 
    a
    \,,
  \\
  \label{cv:ts:dfea}
  \lambda_d^\dagger
    &
    \ 
  \ge
  \ 
  0
    \,,
\end{align}
and 
\textbf{stationarity}
\begin{gather}
  \label{cv:ts:st}
  \mathrm{0}_n
  \ 
  \in
  \ 
  [
  \partial
  f(w^\dagger)
  \ 
  +
  \ 
    \partial
    \left( 
      w
      \mapsto
      d
      -
      \mathbf{U}w
    \right)
    (w^\dagger)
    \cdot
    \lambda_d^\dagger
    \ 
    +
    \ 
    \partial
    \left( 
      w
      \mapsto
      a
      -
      \mathbf{A}w
    \right)
    (w^\dagger)
    \cdot
    \lambda_a^\dagger
    \,
  ]
  \,.
\end{gather}
Applying the well know result\cite[Theorem~28.3]{Rockafellar1970}
finishes the proof.
Apart from elementary calculations, our main tools are the 
strict convexity of $f$, the smoothness of $f^*$ and 
\begin{proposition}
  \emph{
\cite[Theorem~23.5(a)-(b)]{Rockafellar1970}.
  }
  \label{cv:ts:prop}
   For any proper convex function $g$ and any vector $w$, 
   it holds $t\in \partial f(w)$ 
   if and only if 
   $
   x
   \mapsto
   \inner
   {x}{t}
   -
   f(x)
   $
   achieves its supremum at $w$.
\end{proposition}

\begin{proof}
  Let 
$
(\lambda_d^\dagger,\lambda_a^\dagger)
$
be an optimal solution to \eqref{cv:ts:dual}. 

\subsubsection*{Complementary Slackness}
  We fix 
  $
  \lambda_a^\dagger
  $
  and
  work with the objective function $G$ of the dual problem, that is,
  \begin{gather*}
    G
(\lambda_d)
\
:
=
\
    \inner
    {\lambda_d}
    {d}
    \ 
    +
    \ 
    \inner
    {\lambda_a^\dagger}
    {a}
    \ 
    -
    \ 
    f^*
    \left( 
      \mathbf{U}^\top \!\lambda_d
      +
      \mathbf{A}^\top \! \lambda_a^\dagger
    \right)
    \,.
  \end{gather*}
  Since $f^*$ is continuously differentiable, so is $G$.
  Thus
  \begin{gather*}
    \nabla
    G
(\lambda_d^\dagger)
\
:
=
\
d
\ 
    -
    \ 
    \mathbf{U}
    \cdot
    \nabla
    f^*
    \!
    \left( 
      \mathbf{U}^\top \lambda_d^\dagger
      +
      \mathbf{A}^\top  \lambda_a^\dagger
    \right)
    \ 
    =
    \ 
d
\ 
    -
    \ 
    \mathbf{U}
    w^\dagger
    \,.
  \end{gather*}
Let
$\lambda_{d,i}^\dagger$ be the $i$-th coordinate of $\lambda_d^\dagger$ 
and
$
\nabla
G_i
(\lambda_d^\dagger)
$
be the $i$-th coordinate of 
$
\nabla
G
(\lambda_d^\dagger)
$.
  To establish \eqref{cv:ts:comps} we will show
  for all coordinates 
\begin{alignat*}{2}
  \text{either}
  &
  &&
  \qquad
  \lambda_{d,i}^\dagger
  = 0
  \quad
  \text{and}
  \quad
  \nabla
  G
  _i(
  \lambda_{d}^\dagger
  ) \le 0
  \\
  \text{or}
  &
  &&
  \qquad
  \lambda_{d,i}^\dagger
  > 0
  \quad
  \text{and}
  \quad
  \nabla
  G
  _i(
  \lambda_{d}^\dagger
  ) = 0
  \,.
\end{alignat*}
It is well know that a concave functions $g$ satisfies
  \begin{gather}
    \label{cv:ts:concD}
    g(x)-g(y)
    \ge
    \nabla
    g(x)^\top
    (x-y)
    \qquad 
    \text{for all}\ 
    x,y\,.
  \end{gather}
  But $G$ is concave 
  by the convexity of $f^*$(cf. Proposition~\ref{cv:cc:pr:ccc}).

First, we show 
\begin{gather}
  \label{cv:ts:comps:d}
\nabla G_i(\lambda_d^\dagger)
\ 
\le
\ 
0
\qquad
\text{for all}\ 
  i\in \left\{ 1,\ldots, s \right\}
  \,.
\end{gather}
Assume towards a contradiction that 
$
\nabla G_i(\lambda_d^\dagger)>0
$
for some 
$
  i\in \left\{ 1,\ldots, s \right\}
$.
By the continuity of $\nabla G$ there exists $\varepsilon>0$ such that 
$
\nabla G_i(
\lambda_d^\dagger
+
e_i\cdot \varepsilon
)
>
0
$.
It follows from \eqref{cv:ts:concD}
\begin{gather*}
  G
  (
\lambda_d^\dagger
+
e_i\cdot \varepsilon
  )
  \ 
  -
  \ 
  G
  (
\lambda_d^\dagger
  )
  \ 
  \ge
  \ 
\nabla G_i(
\lambda_d^\dagger
+
e_i\cdot \varepsilon
)
\cdot
\varepsilon
\ 
>
\ 
0
\,,
\end{gather*}
which contradicts the optimality of 
$
\lambda_d^\dagger
$
for \eqref{cv:ts:dual}.
It follows \eqref{cv:ts:comps:d}.

Next, 
we assume that
$ \lambda_{d,i}^\dagger>0 $ 
and 
$
  \nabla G_i(\lambda_d^\dagger)< 0
$
for some
$
  i\in \left\{ 1,\ldots, s \right\}
$.
Again, by the 
continuity of $\nabla G$ there exists $\varepsilon>0$ such that
$
  \nabla G_i(\lambda_d^\dagger-e_i\cdot \varepsilon)< 0
$
and
$
\varepsilon
-
\lambda_{d,i}^\dagger
<0
$.
Thus
\begin{gather*}
  G
  (
\lambda_d^\dagger
-
e_i\cdot \varepsilon
  )
  -
  G
  (
\lambda_d^\dagger
  )
  \ge
\nabla G_i(
\lambda_d^\dagger
-
e_i\cdot \varepsilon
)
\cdot
\left( 
  -
\varepsilon
\right)
>0
\,,
\end{gather*}
which contradicts the optimality of 
$
\lambda_d^\dagger
$.
It follows \eqref{cv:ts:comps}, that is,
we proved complementary slackness.
\subsubsection*{Primal Feasibility}
Since $f^*$ is continuously differentiable it holds
\begin{gather*}
  \nabla G(\lambda_d^\dagger)
  \ 
  =
  \ 
  d
  \ 
  -
  \ 
  \mathbf{U}
  \cdot
  \nabla
    f^*
    \left( 
      \mathbf{U}^\top  \lambda_d^\dagger
      +
      \mathbf{A}^\top  \lambda_a^\dagger
    \right)
    \ 
  =
  \ 
  d
  -
  \mathbf{U}w^\dagger
  \,.
\end{gather*}
Thus, by \eqref{cv:ts:comps:d},
$w^\dagger$ satisfies the inequality constraints in \eqref{cv:ts:primal}. 
To prove this for the equality constraints,
we view $G$ from a different angel. Let for fixed
$\lambda^\dagger_d$
\begin{gather*}
  G(\lambda_a)
  \ 
  :=
  \ 
  \inner
  {\lambda_a}{a}
  \ 
  -
  \ 
  \left( 
    f^*
    \left( 
      \mathbf{U}^\top  \lambda_d^\dagger
      +
      \mathbf{A}^\top  \lambda_a
    \right)
    -
  \inner
  {\lambda_d^\dagger}{d}
  \right)
  \ 
  =:
  \ 
  \inner
  {\lambda_a}{a}
  \ 
  -
  \ 
  g(\lambda_a)
  \,
  .
\end{gather*}
The function $g$ inherits convexity and differentiability from 
$f^*$.
From the optimality of $\lambda_a^\dagger$ we know that
$G$ takes its maximum there. But then by Proposition~\ref{cv:ts:prop}
and the differentiability of $g$ it holds
\begin{gather}
  a
  \ 
  \in
  \ 
  \partial
  g(\lambda_a^\dagger)
  \ 
  =
  \ 
  \left\{ 
    \mathbf{A}
    \cdot
    \nabla
    f^*
    \left( 
      \mathbf{U}^\top  \lambda_d^\dagger
      +
      \mathbf{A}^\top  \lambda_a^\dagger
    \right)
  \right\}
  \ 
  =
  \ 
  \left\{ 
    \mathbf{A}
    w^\dagger
  \right\}
  \,.
\end{gather}
Thus $a=
    \mathbf{A}
    w^\dagger
$. But then $w^\dagger$ satisfies also the equality constraints.
We proved \eqref{cv:ts:pfea}.
\subsubsection*{Stationarity}
First we show
\begin{gather}
  \label{cv:ts:st:1}
      \mathbf{U}^\top  \lambda_d^\dagger
      \ 
      +
      \ 
      \mathbf{A}^\top  \lambda_a^\dagger
      \ 
      \in
      \ 
\partial f (w^\dagger)
\,.
\end{gather}
By Proposition~\ref{cv:ts:prop}
it suffices to show
that
\begin{gather*}
  w
  \ 
  \mapsto
  \ 
\inner
{w}
{
      \mathbf{U}^\top  \lambda_d^\dagger
      +
      \mathbf{A}^\top  \lambda_a^\dagger
}
  \ 
-
  \ 
f(w)
\end{gather*}
achieves its supremum at
$w^\dagger$.
Since $f$ is strictly convex 
there exists a unique vector $x^\dagger$
where
the above expression achieves its maximum.
Since
$f^*$ is differentiable it holds
\begin{gather*}
  w^\dagger
  \ 
  =
  \ 
  \nabla
    f^*
    \left( 
      \mathbf{U}^\top  \lambda_d^\dagger
      +
      \mathbf{A}^\top  \lambda_a^\dagger
    \right)
  \ 
    =
  \ 
    \nabla
    \left( 
      \lambda
      \mapsto
\inner
{x^\dagger}
{
  \lambda
}
  \ 
-
  \ 
f(x^\dagger)
    \right)
    \left( 
      \mathbf{U}^\top  \lambda_d^\dagger
      +
      \mathbf{A}^\top  \lambda_a^\dagger
    \right)
  \ 
    =
  \ 
x^\dagger
\,.
\end{gather*}
It follows 
\eqref{cv:ts:st:1}.
Next we show
\begin{gather}
  \label{cv:ts:st:2}
-
\mathbf{U}^\top
\in
\ 
    \partial
    \left( 
      w
      \mapsto
      d
      -
      \mathbf{U}w
    \right)
    (w^\dagger)
    \qquad
    \text{and}
\qquad
-
\mathbf{A}^\top
\in
\ 
    \partial
    \left( 
      w
      \mapsto
      d
      -
      \mathbf{A}w
    \right)
    (w^\dagger)
    \,.
\end{gather}
To this end, note that
\begin{gather*}
  \inner
  {
  -
\mathbf{U}^\top
\!
 e_i
}
  {w-w^\dagger}
  \ 
  =
  \ 
  \left( 
    d-
\mathbf{U}
w
  \right)_i
  \ 
  -
  \ 
  (
    d-
\mathbf{U} w^\dagger
  )
  _i
  \qquad
  \text{for all}\ 
  i\in\left\{ 1,\ldots,r \right\}
  \,.
\end{gather*}
Thus
$
-
\mathbf{U}^\top
\in
    \partial
    \left( 
      w
      \mapsto
      d
      -
      \mathbf{U}w
    \right)
    (w^\dagger)
$.
In the same way it follows
$
-
\mathbf{A}^\top
\in
    \partial
    \left( 
      w
      \mapsto
      d
      -
      \mathbf{A}w
    \right)
    (w^\dagger)
$.
From \eqref{cv:ts:st:1} and \eqref{cv:ts:st:2} we conclude
\begin{align*}
  \mathrm{0}_n
  &
      \ 
  =
      \ 
      \left( 
      \mathbf{U}^\top  \lambda_d^\dagger
      \ 
      +
      \ 
      \mathbf{A}^\top  \lambda_a^\dagger
      \right)
      -
      \mathbf{U}^\top  \lambda_d^\dagger
      \ 
      -
      \ 
      \mathbf{A}^\top  \lambda_a^\dagger
      \\
      &
      \ 
  \in
      \ 
  [
  \partial
  f(w^\dagger)
  \ 
  +
  \ 
    \partial
    \left( 
      w
      \mapsto
      d
      -
      \mathbf{U}w
    \right)
    (w^\dagger)
    \cdot
    \lambda_d^\dagger
  \ 
    +
  \ 
    \partial
    \left( 
      w
      \mapsto
      a
      -
      \mathbf{A}w
    \right)
    (w^\dagger)
    \cdot
    \lambda_a^\dagger
    \, 
  ]
  \,.
\end{align*}
We have proved \eqref{cv:ts:st}, that is, stationarity.
\subsubsection*{Dual Feasibility and Conclusion}
Dual feasibility \eqref{cv:ts:dfea} follows immediately from the optimality of $\lambda_d^\dagger$ for \eqref{cv:ts:dual}.
Thus, $(\lambda_d^\dagger,\lambda_a^\dagger)$ and $w^\dagger$ satisfy the Karush-Kuhn-Tucker conditions for \eqref{cv:ts:primal}.
Applying 
\cite[Theorem~28.3]{Rockafellar1970} finishes the proof.
\end{proof}
\begin{takeaways}
  For strictly convexity objective functions with continuously
  differentiable convex conjugate we get a functional relationship
  of primal and dual solutions via the Karush-Kuhn-Tucker conditions.
\end{takeaways}



%\chapter{Random Matrix Inequalities}
%  In our application we want to bound moments of vector-valued random variables.
%  For this we choose the theory of random matrix inequalities
%  which lately received a lot of attention.
%  In particular an approach via the method of exchangable pairs \cite{Mackey2014}
%  has been fruitful in simplifying the proofs of long standing results such as the matrix Khintchin inequality.
%  The paper offers a comprehensive introduction to this method.
%
%  We will cite the matrix Khintchin inequality and 
%  inequalities for moments of matrices that follow from it\cite{Chen2012}. 
%  As a novelty, we will apply intrinsic dimension results and Hermition Dilitation from \cite{Tropp2015}  
%  to matrix moments inequalities. Even though it is straightforward, to the best of our knowledge the calculations have not been carried out in any publication so far.
%  \section{A Matrix Analysis Primer}
%    The \textbf{trace} of a square matrix, denoted by $\mathrm{tr},$
  is the sum of its diagonal entries, i.e. 
  $
    \mathrm{tr}(\mathbf{B})
    =
    \sum_{j=1}^{d}b_{jj}
    \quad 
    \text{for}\ 
    \mathbf{B} \in \mathbb{M}_d.
  $
  The trace is unitarily invariant, i.e.
  $
    \mathrm{tr}(\mathbf{B})
    =
    \mathrm{tr}(\mathbf{Q}\mathbf{B}\mathbf{Q}^*)
    \quad 
    \text{for all}
    \ 
    \mathbf{B}\in \mathbb{M}_d
    \ 
    \text{for all unitary}\ 
    \mathbf{Q} \in \mathbb{M}_d.
  $
  In particular, the existence of an eigenvalue value decomposition shows 
  that the trace of a Hermitian matrix equals the sum of its  eigenvalues.
  Let
  $
  f: I\to \R
  $
  where 
  $I\subseteq\R$ 
  is an interval.
  Consider a matrix 
  $\mathbf{A}\in \mathbb{H}_d$
  whose eigenvalues are contained in $I.$
  We define the matrix 
  $
    f(\mathbf{A})\in \mathbb{H}_d
  $
  using an eigenvalue decomposition of $\mathbf{A}:$
  \begin{gather}
    f(\mathbf{A})
    =
    \mathbf{Q}
    \begin{bmatrix}
      f(\lambda_1) &&\\
                   &\ddots&\\
                   && f(\lambda_d)
    \end{bmatrix}
    \mathbf{Q}^*
    \qquad
    \text{where}
    \qquad
    \mathbf{A}
    =
    \mathbf{Q}
    \begin{bmatrix}
      \lambda_1 &&\\
                   &\ddots&\\
                   && \lambda_d
    \end{bmatrix}
    \mathbf{Q}^*
    =
    \sum_{i=1}^{d} 
    \lambda_i
    \mathbf{Q}_{\bullet i}
    \mathbf{Q}_{\bullet i}^*
    .
  \end{gather}
  The definition of $f(\mathbf{A})$ does not depend on which 
  eigenvalue decomposition we choose.
  Any matrix function that arises in this fashion is called a \textbf{standard matrix function}.

  For each $p\ge 1$
  the \textbf{Schatten $p$-norm} is defined as 
  $
    \norm{\mathbf{B}}_p
    :=
    (
    \mathrm{tr}(\left| \mathbf{B} \right|^p)
    )
    ^{1/p}
    \ 
    \text{for}
    \ 
    \mathbf{B} \in \mathbb{M}_d.
  $
  In this setting,
  $
  \left| \mathbf{B} \right|
  :=
  (
    \mathbf{B}^*
    \mathbf{B}
    )^{1/2}
    .
  $
  The \textbf{spectral norm}
  of an Hermitian matrix $\mathbf{A}$
  is defined by the relation
  $
    \norm{\mathbf{A}}
    :=
    \lambda_{\mathrm{max}}(\mathbf{A})
    \lor
    (
    -\lambda_{\mathrm{min}}(\mathbf{B})
    )
    .
  $
  For a general matrix $\mathbf{B},$
  the spectral norm is defined to be the largest singular value:
  $
    \norm{\mathbf{B}}
    :=
    \sigma_1(\mathbf{B})
    .
  $
  The Schatten $p$-norm dominates the spectral norm for all $p\ge 1.$



\begin{proposition}
  Let
  $
  f,g: I\to \R
  $
  be real-valued functions on an interval $I\subseteq\R,$ 
  and let
  $\mathbf{A}\in \mathbb{H}_d$
  be a Hermitian matrix
  whose eigenvalues are contained in $I.$

  \begin{enumerate}[label={(\roman*)}]
    \item
      If $\lambda$ is an eigenvalue of of $\mathbf{A},$
      then $f(\lambda)$ is an eigenvalue of $f(\mathbf{A}).$
    \item
      $
        f(a)
        \le
        g(a)
        \quad
        \text{for all}\ 
        a\in I
        \quad
        \text{implies}
        \quad
        f(\mathbf{A})
        \preccurlyeq
        g(\mathbf{A})
        .
      $
  \end{enumerate}
\end{proposition}


\begin{takeaways}
  This Primer is not a prim number.
  \lipsum[5]
\end{takeaways}

%  %\section{The Method of Exchangeable Pairs}
%  %\subsection{Matrix Stein pairs}
We first define an exchangable pair.
\begin{definition}
  Let 
  $Z$ and $Z^{'}$
  random variables taking values
  in a Polish space $\mathcal{Z}.$
  We say that
  $
  (
    Z
    ,
    Z^{'}
  )
  $
  is an \textbf{exchangable pair}
  if it has the same distribution as 
  $
  (
  Z^{'}
    ,
    Z
  )
  .
  $
  In particular, 
  $Z$ and $Z^{'}$
  must share the same distribution.
\end{definition}

The following approach originates in the work of Charles Stein \cite{Stein1972} 
on normal approximation for a sum of dependent random variable.
We will explain how some central ideas of this theory extends to matrices.

We can obtain a lot of information about the fluctuation of a random matrix
$\mathbf{X}$
if we can construct a good exchangable pair 
$
(
  \mathbf{X}
  ,
  \mathbf{X}^{'}
)
.
$
With this motivation in mind, let us introduce a special class of exchangable pairs.
\begin{definition}
  Let
  $
  (
    Z
    ,
    Z^{'}
  )
  $
  be an exchangable pair of random variables taking values
  in a Polish space $\mathcal{Z},$
  and let 
  $
    \mathbf{\Psi}
    : 
    \mathcal{Z}
    \to 
  \mathbb{H}_d
  $
  be a measurable function.
  Define the random Hermitian matrices
  \begin{gather}
    \mathbf{X}
    :=
    \mathbf{\Psi}
    (Z)
    \quad
    \text{and}
    \quad
    \mathbf{X}^{'}
    :=
    \mathbf{\Psi}
    (Z^{'})
    .
  \end{gather}
  We say that 
  $
  (
    \mathbf{X}
    ,
    \mathbf{X}^{'}
  )
  $
  is a \textbf{matrix Stein pair}
  if there is a constant 
  $\alpha\in (0,1]$
  for which
  \begin{gather}
    \E[
    \mathbf{X}
    -
    \mathbf{X}^{'}
    \vert
    Z
    ]
    =
    \alpha 
    \mathbf{X}
    \qquad
    \text{almost surely.}
  \end{gather}
  The constant 
  $\alpha$
  is called the \textbf{scale factor} of the pair.
  We always assume 
  $
    \E
    \left[ 
      \norm{\mathbf{X}}^2
    \right]
    <
    \infty
    .
  $
\end{definition}

A matrix Stein pair 
$
(
\mathbf{X}
,
\mathbf{X}^{'}
)
$
has several useful propreties. First,
$
(
\mathbf{X}
,
\mathbf{X}^{'}
)
$
always forms an exchangable pair. Second, it must be the case that
$\E[\mathbf{X}]=\mathbf{0}.$
Indeed,
\begin{gather*}
  \E[\mathbf{X}]
  =
  \frac{1}{\alpha}
  \E[
  \E
  [
    \mathbf{X}
    -
    \mathbf{X}^{'}
    \vert
    Z
  ]
  ]
  =
  \frac{1}{\alpha}
  \E[
    \mathbf{X}
    -
    \mathbf{X}^{'}
    \vert
  ]
  =
  \mathbf{0}
  .
\end{gather*}
\subsection{The method of exchangable pairs}
A well-chosen matrix Stein pair 
$
(
\mathbf{X}
,
\mathbf{X}^{'}
)
$
provides a surprisingly powerful tool for studying 
the random matrix $\mathbf{X}.$
The technique depends on a fundamental technical lemma.
\begin{lemma}
  Suppose that
  $
  (
    \mathbf{X}
    ,
    \mathbf{X}^{'}
  )
  $
  is a matrix Stein pair with scale factor $\alpha.$
  Let 
  $
    \mathbf{F}
    :
    \mathbb{H}_d
    \to
    \mathbb{H}_d
  $
  be a measurable function that satisfies the regularity condition
  $
  \E
  \left[
  \norm{
    (
    \mathbf{X}
    -
    \mathbf{X}^{'}
  )
  \mathbf{F}(\mathbf{X})
  }
  \right]
  <
  \infty
  .
  $
  Then
\begin{gather}
  \E
  [
    \mathbf{X}
    \cdot
    \mathbf{F}
    (\mathbf{X})
  ]
  =
  \frac{1}{2\alpha}
  \E
  [
    (
    \mathbf{X}
    -
    \mathbf{X}^{'}
    )
    (
    \mathbf{F}
    (
    \mathbf{X}
    )
    -
    \mathbf{F}
    (
    \mathbf{X}^{'}
    )
  )
  ]
  .
\end{gather}
\end{lemma}
In short, the randomness in the Stein pair furnishes an alternative expression for the expected product of $\mathbf{X}$
and a function $\mathbf{F}.$
It allows us to estimate the expectation using the smoothness properties of the function $\mathbf{F}$ and the discrepancy between $\mathbf{X}$
and $\mathbf{X}^{'}.$
\begin{proof}
  \emph{\cite[Lemma~2.4]{Mackey2014}}
  Suppose that
  $
  (
    \mathbf{X}
    ,
    \mathbf{X}^{'}
  )
  $
  constructed from an auxiliary exchangable pair 
  $
  (
  Z,
  Z^{'}
  )
  .
  $
  The defining property implies
 \begin{gather}
   \alpha \cdot 
   \E[
   \mathbf{X}
   \cdot
   \mathbf{F}
   (\mathbf{X})
   ]
   =
   \E[
    \E[
    \mathbf{X}
    -
    \mathbf{X}^{'}
    \vert
    Z
    ]
    \cdot
   \mathbf{F}
   (\mathbf{X})
   ]
   =
    \E[
    (
    \mathbf{X}
    -
    \mathbf{X}^{'}
    )
   \mathbf{F}
   (\mathbf{X})
   ]
 \end{gather} 
\end{proof}
\subsection{The conditional variance}
To each matrix Stein pair 
  $
  (
    \mathbf{X}
    ,
    \mathbf{X}^{'}
  ),
  $
  we may associate a random matrix called the conditional variance of $\mathbf{X}.$
  The purpose of this section is to argue that the spectral norm of $\mathbf{X}$
  is unlikely to be large, when the conditional variance is small.
\begin{definition}
  Suppose that
  $
  (
    \mathbf{X}
    ,
    \mathbf{X}^{'}
  ),
  $
  is a matrix Stein pair, constructed from an auxiliary exchangeable pair
  $
  (
  Z
    ,
  Z^{'}
  ).
  $
  The \textbf{conditional variance}
  is the random matrix
  \begin{gather}
    \mathbf{\Delta_X}
    :=
    \mathbf{\Delta_X}
    (Z)
    :=
    \frac{1}{2\alpha}
    \E
    [
    (
    \mathbf{X}
    -
    \mathbf{X}^{'}
    )^2
    \vert
    Z
    ]
    ,
  \end{gather}
  where $\alpha$ is the scale factor of the pair. We may take any version of the conditional expectation in this definition.
\end{definition}

The conditional variance
$
    \mathbf{\Delta_X}
$
can be regarded as a stochastic estimate for the variance 
of the random matrix $\mathbf{X}.$
To see this, assume
the second moment of $\mathbf{X}$ exists. Then it follows from Lemma 
with $\mathbf{F}(\mathbf{X})=\mathbf{X}$
\begin{gather}
  \E
  [
    \mathbf{\Delta_X}
  ]
  =
  \E
  [
  \mathbf{X}^2
  ]
  .
\end{gather}
To verify the regularity condition, note that
\begin{gather}
  \E
  [
  \lVert
    (
    \mathbf{X}
    -
    \mathbf{X}^{'}
    )
    \mathbf{X}
    \rVert
  ]
  \le
  \E[
  \norm{
    \mathbf{X}
  }
  ^2
  ]
  +
  \E
  [
  \norm{
    \mathbf{X}
  }
  \cdot
  \lVert
    \mathbf{X}^{'}
    \rVert
  ]
  \le
  2
  \E
  [
  \norm{
    \mathbf{X}
  }
  ^2
  ]
  <
  \infty
  .
\end{gather}
\begin{example}
  \emph{\cite[Example~2.4]{Mackey2014}}
\end{example}
nrt\\
is not necessary, as, for example, the normal distribution does not fulfill it. For

% \section{Matrix Khintchin Inequality and Applications}
% In this section we state the matrix Khintchin inequality and matrix moment inequalities as an applications.
% We provide the proof of auxiliary theorems which are cited without proof in \cite{Mackey2014}. They are needed to prove the matrix Khintchin inequality. 
%  \begin{theorem}
  \emph{(Matrix BDG inequality)}
  Let
  $
    p = 1
    \ 
    \text{or}\ 
    p \ge 3/2
    .
  $
  Suppose that 
  $
  (
    \mathbf{X}
   , 
   \mathbf{X}^{'}
  )
  $
  is a matrix Stein pair where
  $
   \E
   [ 
    \norm{\mathbf{X}}
    _{2p}^{2p}
   ]
   <
   \intfy
   .
  $
\end{theorem}




\begin{ftheorem}
  \emph{\cite[Corollary~7.3]{Mackey2014}}
  Suppose that
  $
    p = 1
    \ 
    \text{or}\ 
    p \ge 3/2
    .
  $
  Consider a finite sequence
  $
    (\mathbf{Y}_k)_{k\ge 1}
  $
  of independent, random, Hermitian matrices 
  and a deterministic sequence
  $
    (\mathbf{A}_k)_{k\ge 1}
  $
  for which
  \begin{gather}
    \E[\mathbf{Y}_k]
    =
    0
    \quad 
    \text{and}
    \quad
    \mathbf{Y}_k^2
    \preccurlyeq
    \mathbf{A}_k^2
    \qquad
    \text{almost surely for all}\ 
    k \ge 1.
  \end{gather}
  Then
  \begin{gather}
      \E
      \left[
        \norm{
          \sum_{k\ge 1}
            \mathbf{Y}_k
        }
        _{2p}
        ^{2p}
      \right]
      ^{1/(2p)}
      \le
      \sqrt{
        p - \frac{1}{2}
      }
      \,
      \norm{
        \left( 
          \sum_{k\ge 1}
          (
            \mathbf{A}_k^2
            + 
            \E[
              \mathbf{Y}_k^2
            ]
          )
        \right)
        ^{1/2}
        }
      _{2p}
      .
  \end{gather}
  In particular, when 
  $
    (\xi_k)_{k\ge 1}
  $
  is an independent sequence of Rademacher random variables,
  \begin{gather}
      \E
      \left[
        \norm{
          \sum_{k\ge 1}
            \xi_k
            \mathbf{A}_k
        }
        _{2p}
        ^{2p}
      \right]
      ^{1/(2p)}
      \le
      \sqrt{
        2p - 1
        }
      \,
      \norm{
        \left( 
          \sum_{k\ge 1}
            \mathbf{A}_k^2
        \right)
        ^{1/2}
        }
      _{2p}
      .
  \end{gather}
\end{ftheorem}

%  \section{Generalzed Inequalities by Hermitian Dilation}
%  For an introduction to Hermitian Dilation see \cite[§2.1.16]{Tropp2015}
\begin{definition}
  \emph{(Hermitian Dilation)}
  The Hermitian dilation
  \begin{gather*}
    \mathfrak{H} : \C^{d_1 \times d_2} \to \mathbb{H}_{d_1 \times d_2}
  \end{gather*}
  is a map from a general matrix to an Hermitian matrix defined by
  \begin{gather}
    \label{ rmineq_hermitian_dilation } 
    \mathfrak{H}(B)
    :=
    \begin{bmatrix}
      0   & B \\
      B^* & 0 \\
    \end{bmatrix}
  \end{gather}
\end{definition}

Properties:

$
\norm{\mathbf{B}}
=
\norm{\mathfrak{H}\mathbf{B}}
$
Linearity
confer \cite[§7.7.3]{Tropp2015}
for combination with intrinsic dimension argument.

%  \begin{ftheorem}
  %\label{rmineq_bernstein}
  Let $\mathbf{A}_1, \ldots, \mathbf{A}_n$ be independent, random matrices with dimension 
  $d_1 \times d_2$.
    Introduce the random matrix
      \begin{gather*}
        \mathbf{S}:=\sum_{k=1}^n \mathbf{A}_k.
      \end{gather*}
    Let $v(\mathbf{S})$ be the matrix variance statistic of the sum:
      \begin{align}
        \label{rmineq_bernstein_cond_2}
        v(\mathbf{S}):= \norm{\E[\mathbf{S}\mathbf{S}^\top ]} \lor \norm{\E[\mathbf{S}^\top \mathbf{S}]} 
             = \norm{\sum_{k=1}^n\E[\mathbf{A}_k\mathbf{A}_k^\top]} \lor \norm{\sum_{k=1}^n\E[\mathbf{A}_k^T \mathbf{A}_k]} .
      \end{align}
    Then
      \begin{align}
        \label{rmineq_bernstein_expectation_bound}
        \left(
          \E \left[ \norm{\mathbf{S}}^2 \right]
        \right)^{\frac{1}{2}}
        \le
        \sqrt{
          2ev(\mathbf{S})\log(d_1+d_2)
        } 
        + 
        4e \left( 
          \E[\max_{k \le n}\norm{\mathbf{A}_k}^2]
        \right)^\frac{1}{2}
        \log(d_1+d_2).
      \end{align}
\end{ftheorem}
  

%  
%  \newpage
%  \section{Intrinsic Dimension}
%  \begin{definition*}
  \label{rmineq_intrinsic_bernstein}
  For a positive-semidefinite matrix $\,\mathbf{S}$,
  the \textbf{intrinic dimension} is the quantity
  \begin{gather*}
    \mathrm{intdim}
    \,
    \mathbf{A}
    \ 
    :=
    \ 
\mathrm{tr}\,\mathbf{A}
\,
/
\,
\norm{\mathbf{A}}
    \,
    .
  \end{gather*}
\end{definition*}
\begin{lemma*}
  \emph{(Intrinsic dimenision)}
  Let 
  $
  \,
    \varphi\,:\, [\,0,\infty) \,\to\, \R
    \,
  $
  be a convex function with
  $
    \varphi(0)\,=\,0
  $.
  For any positive-semidefinite matrix $\,\mathbf{S}$ it holds 
  \begin{gather*}
    \mathrm{tr}\,\varphi(\mathbf{S})
    \ 
    \le
    \ 
    \mathrm{intdim}\,\mathbf{S}
    \ 
    \cdot
    \ 
    \varphi\!\norm{\mathbf{S}}
    \,
    .
  \end{gather*}
\end{lemma*}
\begin{proof}
  \emph{\cite[Lemma~7.5.1]{Tropp2015}}
  Since $\varphi$ is convex on any interval $[\,0,L\,]$ with $L\,>\,0$, and $\varphi(0)\,=\,0$, it holds
  \begin{gather*}
    \varphi(a)
    \ 
    \le
    \ 
    \left( 
      1 - a/L
    \right)
    \cdot
    \varphi(0)
    \ 
    +
    \ 
    a/L
    \cdot
    \varphi(L)
    \ 
    =
    \ 
    a/L
    \cdot
    \varphi(L)
    \qquad
    \text{for all}\ 
    a \,\in\, [\,0,L\,]
    \,
    .
  \end{gather*}
  Since $\,\mathbf{S}\,$ is positive-semidefinite, the eigenvalues of $\,\mathbf{S}\,$ 
  fall in the interval $\,[\,0,L\,]\,$, where $\,L\,=\,\norm{\mathbf{S}}\,$.
  It follows 
  \begin{align*}
    \mathrm{tr}\,\varphi(\mathbf{S})
    &
    \ 
    =
    \ 
    \sum_{i=1}^{d}
    \varphi(\lambda_i)
    \ 
    \le
    \ 
    \sum_{i=1}^{d}
    \lambda_i
    /
    \norm{\mathbf{S}}
    \ 
    \cdot
    \ 
    \varphi(\norm{\mathbf{S}})
    \\
    &
    \ 
    =
    \ 
    \mathrm{tr}(\mathbf{S})
    /
    \norm{\mathbf{S}}
    \ 
    \cdot
    \ 
    \varphi(\norm{\mathbf{S}})
    \ 
    =
    \ 
    \mathrm{intdim}\,\mathbf{S}
    \ 
    \cdot
    \ 
    \varphi\!\norm{\mathbf{S}}
    \,
    .
  \end{align*}
\end{proof}
The next example applies the preceding lemma to bound 
the $p$-Schatten-norm, when $\,p\,\ge\, 2\,$, by the spectral norm and the intrinsic dimension.
  \begin{example*}
    Let 
    $
    \,
      \mathbf{B} \in \mathbb{C}^{m\times n}
      \,
    $
    be any rectangular matrix and let $\,p\ge 2\,$.
    Then 
    $
    \,
      \varphi(x)
      \,
      :=
      \,
      \left| x \right|
      \,
    $
  defines a convex function with $\,\varphi(0)\,=\,0\,$.
  The intrinsic dimension lemma yields
  \begin{gather*}
    \norm{\mathbf{B}}^{\,p}_p
    \ 
    =
    \ 
    \mathrm{tr}
      \left|
      \mathbf{B}^* \mathbf{B}
      \right|
      ^{\,p/2}
      \ 
    \le
    \ 
    \mathrm{intdim}
    \,
      \mathbf{B}^* \mathbf{B}
      \ 
      \cdot
      \ 
      \norm{
      \mathbf{B}^* \mathbf{B}
    }^{\,p/2}
    \ 
    =
    \ 
    \mathrm{intdim}
    \,
      \mathbf{B}^* \mathbf{B}
      \ 
      \cdot
      \ 
      \norm{
        \mathbf{B}
    }^{\,p}
    \,.
  \end{gather*}
  If, additionally, $\mathbf{B}$ is self-adjoint and positive-semidefinite
  then it holds 
  \begin{gather*}
    \mathrm{tr}
    \,
      \mathbf{B}^* \mathbf{B}
      \ 
    =
    \ 
    \mathrm{tr}
    \,
    \mathbf{B}^{\,2}
    \ 
    =
    \ 
    \sum_{i=1}^{n} 
    \,
    \lambda_i^{\,2}
    \ 
    \le 
    \ 
    \left( 
      \,
    \sum_{i=1}^{n} 
    \,
    \lambda_i
    \,
    \right)
    ^{\!2}
    \ 
    =
    \ 
    \left( 
    \mathrm{tr}
    \,
    \mathbf{B}
    \right)
    ^{\,2}
    \,,
  \end{gather*}
  and consequently
  \begin{gather*}
    \norm{\mathbf{B}}^{\,p}_{p}
    \ 
    \le
    \ 
    \left( 
    \mathrm{intdim}
    \,
    \mathbf{B}
    \right)
    ^{\,2}
    \ 
      \cdot
      \ 
      \norm{
        \mathbf{B}
    }^{\,p}
    \,.
  \end{gather*}
  \end{example*}


  \todo[color=red!40, inline]{Apply to estimate \cite[above (A.4)]{Chen2012} to get intrinsic dimension version of Rosenthal-Pinelis inequality.}
\begin{takeaways}
  The notion of intrinsic dimension is useful when bounding convex
  functions of a positive-semidefinite matrix by its spectral norm.
  We saw how to derive bounds on the $p\,$-Schatten-norm when $\,p\,\ge\, 2$.
\end{takeaways}

%
%
%\chapter{Empirical Processes}
%Classical references are \cite{Vaart2000} and \cite{vaart2013}.
%For maximal inequalities see \cite[§19]{Vaart2000}
%For Functional Delta-Method see \cite[§20]{Vaart2000}
%For an introduction to empirical processes and outer expectation 
%see the beginning of \cite{vaart2013}.
%\section{A Primer on Empirical Processes}
%\todo[color=red!40,inline]{Add outer probability calculus. \cite{vaart2013} p.6}

Let 
$
(
\mathbb{D}
,
d
)
$
be a metric space, and let 
$
(\P_n)_{n\in \mathbb{N}}
  \P
$
be (Borel) probability measures
on
$
(
\mathbb{D}
,
\mathcal{D}
)
,
$
where 
$\mathcal{D}$
is the Borel $\sigma$-algebra on $\mathbb{D},$
the smallest $\sigma$-algebra
containing all open sets.
Then the sequence 
$\P_n$
\textbf{converges weakly}
to 
$\P,$
which we denote as $\P_n \rightsquigarrow \P,$
if and only if 
\begin{gather}
  \int_\mathbb{D}
  f
  \text{d}
  \P_n
  \to
  \int_\mathbb{D}
  f
  \text{d}
  \P
  \qquad
  \text{for all}
  \ 
  f
  \in
  C_b(\mathbb{D})
  .
\end{gather}
Here 
$
  C_b(\mathbb{D})
$
denotes the set of all bounded, continuous, real functions on $\mathbb{D}.$
Equivalently, if 
$X_n$ and $X$
are 
$\mathbb{D}$-valued
random variables with distribution 
$\P_n$ and $\P$
respectively, then 
$X_n \to X$
if and only if 
\begin{gather}
  \E[f(X_n)]
  \to
  \E[f(X)]
  \qquad
  \text{for all}
  \ 
  f
  \in
  C_b(\mathbb{D})
  .
\end{gather}

This definitions yield the classical theory of weak convergence.
For a modern treatment see \cite{Klenke2020}.

The classical theory requires that 
$\P_n$
is defined, for each $n\in \mathbb{N},$
on the Borel $\sigma$-algebra $\mathcal{D},$
or, equivalently, that $X_n$ is a Borel measurable map for each $n\in \mathbb{N}.$
If 
$
(
\Omega_n,
\mathcal{A}_n,
\P_n
)
$
are the underlying probability spaces on which the maps 
$X_n$
are defined, this means that
$X_n^{-1}(D)\in \mathcal{A}_n$
for every Borel set $D \in \mathcal{D}.$
This required measurability usually holds when $\mathbb{D}$
is a separable metric space such as $\R^k$
or $C([0,1])$ with the supremum metric.

However, this apparently modest requirement can and does easily fail when the metric space $\mathbb{D}$ is not separable.

\begin{example}
  \emph{\cite[Problem~1.7.3]{vaart2013}}
  Let $\mathbb{D}=D([0,1])$ be the \textbf{Skorohod space} of all right-continuous functions on $[0,1]$
  with left limits endowed with the metric induced by the supremum norm.
  Define 
  $
    X:
    [0,1]
    \to
    \mathbb{D}
    ,\ 
    \omega
    \mapsto
    \mathbf{1}_{[\omega,1]}
    .
  $
  If we equip $[0,1]$ with the Borel $\sigma$-algebra 
  $\mathcal{B}([0,1])$, then 
  $X$ is not measurable. To see this, let $B_s$ be the open ball of radius $1/2$ in $\mathbb{D}$ around the function $\mathbf{1}_{[s,1]}.$
  Now $X(\omega)\in B_s$
  if and only if $\omega=s.$ Indeed, if $\omega\neq s$ there exists an $x$ between $\omega$ and $s$ such that the difference of the indicator functions is 1 at $x$. Conversely, if the distance is greater than
  $1/2$ at a point $x\in [0,1]$, it is because $x$ lies between $\omega$ and $s$ and the indicator functions have difference 1.
  Since arbitrary (even uncountable) unions of open sets are open,
  we get for every $S\subseteq [0,1]$ the open set 
  $
  G
  :=
  \bigcup_{s\in S}
    B_s
    \in \mathcal{D}
    .
  $
  It follows
  $
  X^{-1}(G)=S 
  \ 
  \text{for all}
  \ 
  S \subseteq [0,1]
  .
  $
  Since not all subsets of $[0,1]$ are measurable, we have
  $
  X^{-1}(\mathcal{D})\nsubseteq \mathcal{B}([0,1])
.
$
But then $X$ is not measurable. The $\sigma$-algebra $\mathcal{D}$ is to large.
\end{example}



Let 
$
  \left( 
    \Omega,
    \mathcal{A},
    \P
  \right)
$
be a probability space
and
$
  \left( 
    \mathcal{X},
    \Sigma
  \right)
$
a measurable space.
Let 
$
  X_j
  :
  \left( 
    \Omega,
    \mathcal{A},
    \P
  \right)
  \to
  \left( 
    \mathcal{X},
    \Sigma
  \right)
  ,
  j=1,\ldots, n
$ be independent and identically-distributed (i.i.d.)
random variables
with probability distribution $\P_X$ 
and
$\mathcal{F}$ a family of measurable functions
$
  f:
  \left( 
    \mathcal{X},
    \Sigma
  \right)
    \to
  \left( 
    \R,
    \mathcal{B}(\R)
  \right)
$.
Consider the map
\begin{gather}
  f
  \mapsto
  G_n f
  :=
  \sqrt{n}
  \left( 
    \frac{1}{n}
    \sum_{i = 1}^{n}
      f(X_i)
    -
    \P_X f
  \right),
\end{gather}
where
$
  \P_X f 
  :=
  \int_\mathcal{X} f \text{d}\P_X.
$
We call 
$
  \left( 
    G_n f
  \right)_{f \in \mathcal{F}}
$
the empirical process indexed by $\mathcal{F}$.
Furthermore
\begin{gather}
  \norm{G_n f}_\mathcal{F}
  :=
  \sup_
        { f \in \mathcal{F}}
        \left|
          G_n f
        \right|
        .
\end{gather}



%
%\section{Maximal Inequalities}
%%%%%%%%%%%%%%%%%
% INTRODUCTION %
%%%%%%%%%%%%%%%%

Let 
$
  \left( 
    \Omega,
    \mathcal{A},
    \P
  \right)
$
be a probability space
and
$
  \left( 
    \mathcal{X},
    \Sigma
  \right)
$
a measurable space.
Let 
$
  X_j
  :
  \left( 
    \Omega,
    \mathcal{A},
    \P
  \right)
  \to
  \left( 
    \mathcal{X},
    \Sigma
  \right)
  ,
  j=1,\ldots, n
$ be independent and identically-distributed (i.i.d.)
random variables
with probability distribution $\P_X$ 
and
$\mathcal{F}$ a family of measurable functions
$
  f:
  \left( 
    \mathcal{X},
    \Sigma
  \right)
    \to
  \left( 
    \R,
    \mathcal{B}(\R)
  \right)
$.
Consider the map
\begin{gather}
  f
  \mapsto
  G_n f
  :=
  \sqrt{n}
  \left( 
    \frac{1}{n}
    \sum_{i = 1}^{n}
      f(X_i)
    -
    \P_X f
  \right),
\end{gather}
where
$
  \P_X f 
  :=
  \int_\mathcal{X} f \text{d}\P_X.
$
We call 
$
  \left( 
    G_n f
  \right)_{f \in \mathcal{F}}
$
the empirical process indexed by $\mathcal{F}$.
Furthermore
\begin{gather}
  \norm{G_n f}_\mathcal{F}
  :=
  \sup_
        { f \in \mathcal{F}}
        \left|
          G_n f
        \right|
        .
\end{gather}

\begin{lemma}
  \emph{(Bernstein Inequality for Empirical Processes)}
  For any bounded, measurable function $f$
  it holds for all $t > 0$
  \begin{gather}
    \P 
    \left(
      \left| 
        G_n f
      \right|
      >
      t
    \right)
    \le
    2
    \exp
    \left( 
      - \frac{1}{4}
      \frac{t^2}
      {
        \P_X(f^2)
        +
        t
        \norm{f}_\infty
        /
        \sqrt{n}
      }
    \right)
  \end{gather}
\end{lemma}
\begin{proof}
  By the Markov inequality it holds for all $\lambda > 0$
  \begin{gather}
    \P
    \left( 
      G_n f 
      > 
      t
    \right)
    \le
    e^{-\lambda t}
    \E
    \exp
    \left( 
      \lambda
      G_n f 
    \right)
  \end{gather}
\end{proof}


\begin{lemma}
  For any finite class $\mathcal{F}$ of bounded, measurable, square-integrable functions, with $\left| \mathcal{F} \right|$ elements, it holds
  \begin{gather}
    \E \norm{G_n f}_\mathcal{F}
    \lesssim
    \max_
          {f \in \mathcal{F}}
          \frac{\norm{f}_\infty}{\sqrt{n}}
          \log
          \left( 1 + \left| \mathcal{F} \right| \right)
          +
          \max_
              {f \in \mathcal{F}}
              \norm{f}_{\P,2}
              \sqrt{
                \log \left( 
                  1 + \left| \mathcal{F} \right|
                \right)
              }
              .
  \end{gather}
\end{lemma}

\begin{lemma}
  For any class 
  $\mathcal{F}$
  of measurable functions 
  $
    f:
    \mathcal{X}
    \to \R
  $
  such that 
  $
    \P f^2
    <
    \delta^2
  $
  for every $f,$
  we have,
  with
  $
    a(\delta)
    =
    \delta
    /
    \sqrt{
      \mathrm{Log}
    N
    _{[]}
    (
    \delta
    ,
    \mathcal{F}
    ,
    L_2(\P)
    )
    }
    ,
  $
  and $F$
  an envelope function,

  \begin{gather}
    \E^*
    _\P
    [
    \norm{
      \G
      _n
      }
      _\mathcal{F}
    ]
    \lesssim
    J
    _{[]}
    (
      \delta
    ,
    \mathcal{F}
    ,
    L_2(\P)
    )
    +
    \sqrt{n}
    \P^*
    F
    \left\{ 
      F > \sqrt{n}
      a(\delta)
    \right\}
    .
  \end{gather}
\end{lemma}


\begin{corollary}
  For any class $\mathcal{F}$ of measurable functions with envelope function $F,$
  \begin{gather}
    \E^*
    _\P
    [
    \norm{
      \G
      _n
      }
      _\mathcal{F}
    ]
    \lesssim
    J
    _{[]}
    (
    \norm{
      F
    }
    _{\P,2}
    ,
    \mathcal{F}
    ,
    L_2(\P)
    )
    .
  \end{gather}
\end{corollary}





%\section{Functional Delta Method}
%\begin{definition}
  A map 
  $
  \phi:
  \mathbb{D}_\phi
  \to 
  \mathbb{E}
  ,
  $
  defined on a subset 
  $
  \mathbb{D}_\phi
  $
  of a normed space
  $\mathbb{D}$
  that contains 
  $\theta,$
  is called 
  \textbf{Hadamard diffenertiable}
  at $\theta$
  if there exists a continuous,
  linear map
  $
  \phi_\theta^{'}
    :
    \mathbb{D}
    \to 
    \mathbb{E}
  $
  such that
  \begin{gather}
    \norm{
      \frac{
        \phi(\theta + t h_t)
        -
        \phi(\theta)
      }{
        t
      }
      -
      \phi^{'}_\theta
      (h)
    }_\mathbb{E}
    \to
    0
    \quad
    \text{as}
    \ 
    t\searrow 0
    \ 
    \text{for all}
    \ 
    h_t \to h
  \end{gather}
  $
    \text{such that $\theta + th_t$ is contained in $\mathbb{D}_\phi$ for all small $t>0.$}
  $
\end{definition}


\begin{ftheorem}
  \emph{(Delta Method)}
  Let 
  $
    \mathbb{D}
    \ \text{and}
    \ 
    \mathbb{E}
  $
  be normed linear spaces.
  Let
  $
    \phi
    :
    \mathbb{D}_\phi
    \subseteq
    \mathbb{D}
    \to
    \mathbb{E}
  $
  be Hadamard differentiable it $\theta$
  tangentially to 
  $\mathbb{D}_0.$
  Let
  $
    T_n
    :
    \Omega_n
    \to
    \mathbb{D}_\phi
  $
  be maps such that 
  $
    r_n
    (T_n - \theta)
    \rightsquigarrow
    T
  $
  for some sequence of numbers $r_n \to \infty$
  and a random element $T$
  that takes its values in $\mathbb{D}_0.$
  Then 
  $
    r_n(\phi(T_n)-\phi(\theta))
    \rightsquigarrow
    \phi^{'}
    _\theta
    (T)
    .
  $
  If 
  $
    \phi^{'}
    _\theta
  $
  is defined and continuous on the whole space $\mathbb{D},$
  then we also have 
  $
    r_n
    (
      \phi(T_n)
      -
      \phi(\theta)
    )
    =
    \phi^{'}
    _\theta
    (
    r_n
    (
    T_n
    -
    \theta
    )
    )
    +
    o_\P
    (1)
    .
  $
\end{ftheorem}
\begin{proof}
  \cite[Theorem~20.8]{Vaart1998}
\end{proof}


%\chapter{Simple yet useful Calculations} 
%\begin{proposition}
  Let 
  $f : \R^n \to \R$ 
  be continuous such that 
  a minimum $x^*$ exists and is unique.
  Then 
  for all $y \in \R^n$ and $C>0$ 
  it follows
    \begin{gather}
      \inf_{\norm{\Delta}=C} f(y+\Delta) - f(y) > 0 \qquad
      \Rightarrow \qquad 
      \norm{x^* - y} \le C.
    \end{gather}
\end{proposition}


\begin{proof}
Since 
$\mathcal{C}:=\left\{ \norm{\Delta}\le C \right\}$
is compact and
\begin{gather*}
  f(x^*) \le f(y) <  \inf_{\norm{\Delta}=C} f(y+\Delta)
\end{gather*}
the continious function $f(y+\,\cdot\,)$ has a minimum in 
$\overset{\circ}{\mathcal{C}}:=\left\{ \norm{\Delta} < C \right\}$. 
Since 
$x^*$ is the unique minimum of $f$
there exists $\Delta^* \in \overset{\circ}{\mathcal{C}}$ 
such that 
$x^* - y = \Delta^*$.
We conclude that
$\norm{x^* - y} \le C$.
\end{proof}

 %%%%%%%%%%%%%%%%%%%%%%%%%%%%%%%%%%%%%%%%%%%%%%%%%%%%%%%%%%%%%%%%


\begin{theorem}
  \emph{(Multivariate Taylor Theorem)}
  Let 
  $f \in C^2(\R^n, \R)$.
  Then 
  for all $x, \Delta \in \R^n$
  there exists $\xi \in [0,1]$
  such that 
  it holds
  \begin{gather}
    \label{syu_taylor}
    f(x + \Delta) - f(x)
    =
    \sum_{i = 1}^{n} \frac{\partial f(x)}{\partial x_i} \Delta_i    
    + \sum_{i,j = 1}^{n}  \frac{\partial^2 f(x + \xi \Delta)}{\partial x_i \partial x_j} \Delta_i\Delta_j
    + \frac{1}{2}\sum_{i = 1}^{n} \frac{\partial^2 f(x + \xi \Delta)}{\partial x_i^2}\Delta_i^2 
  \end{gather}
\end{theorem}


\begin{corollary}
  Let 
  $f\in C^2(\R)$. 
  Then
  for all $a,x,\Delta \in \R^n$ 
  there exist $\xi \in [0,1]$ 
  such that it holds
  \begin{gather}
    \label{syu_2.result}
    f(a^T (x + \Delta)) - f(a^T x) = 
    f^{'}(a^T x)\, \Delta^T a + 
    \frac{1}{2}f^{''}(a^T (x + \xi \Delta))\, \Delta^T A\ \Delta,
  \end{gather}
  where 
  $A:= a a^T \in \R^{n \times n}$ .
\end{corollary}

\begin{proof}
  By the chain rule 
  we have
  for all $a, x, \Delta \in \R^n$ and $\xi \in [0,1]$
  \begin{gather}
    \label{syu_2.0}
   \frac{\partial^2  f(a^T (x + \xi \Delta)) }{\partial x_i \partial x_j}=
    f^{''}(a^T (x + \xi \Delta))\, a_i a_j.
  \end{gather}
  Since 
  $A:= a a^T$ is symmetric
  we have
  \begin{gather}
    \label{syu_2.1}
    \Delta^T A\ \Delta 
    =
    2
    \sum_{
    \begin{smallmatrix}
      i,j = 1 \\ 
      i \neq j
    \end{smallmatrix}
    }
    ^{n}
    a_i a_j \Delta_i \Delta_j
    + 
    \sum_{i=1}^{n} 
    a_i^2 \Delta_i^2.
  \end{gather}
  Plugging \eqref{syu_2.0} and \eqref{syu_2.1} into \eqref{syu_taylor}
  yields
  \eqref{syu_2.result}.
\end{proof}



%
%makeindex main.nlo -s nomencl.ist -o main.nls

\nomenclature{$\P$}{generic probability measure}
\nomenclature{$\overset{\mathcal{D}}{\longrightarrow}$}{convergence of distributions}

%\printnomenclature

\bibliography{literature}{}
\bibliographystyle{alpha}
\end{document}
